% Generated mechanically from frozen input p06/ko/duality.tex.
% Do not edit this generated file; edit the bound locale source instead.

% OK, start here.
%

\setcounter{chapter}{47}
\chapter{스킴의 쌍대성}



\phantomsection
\label{duality-section-phantom}


\section{서론}
\label{duality-section-introduction}

\noindent
이 장에서는 스킴의 사상에 대한 상대 쌍대성과
스킴 위의 쌍대화 복합체를 연구한다. 참고문헌은\  \cite{RD}이다.

\medskip\noindent
Noether 환의 쌍대화 복합체는\ 
Dualizing Complexes, Section \ref{dualizing-section-dualizing} 이하에서
정의하고 연구하였다. 이 장에서는 이를 이어 스킴 위의
쌍대화 복합체를 연구한다. Section \ref{duality-section-dualizing-schemes}를 보라.

\medskip\noindent
이 장의 대부분은 준연접 코호몰로지 층을 갖는 가군층의
유도 범주에서 추이 함자의 오른쪽 수반을 연구하는 데 할애한다.
Sections
\ref{duality-section-twisted-inverse-image},
\ref{duality-section-restriction-to-opens},
\ref{duality-section-base-change-map},
\ref{duality-section-base-change-II},
\ref{duality-section-trace},
\ref{duality-section-compare-with-pullback},
\ref{duality-section-sections-with-exact-support},
\ref{duality-section-duality-finite},
\ref{duality-section-perfect-proper},
\ref{duality-section-dualizing-Cartier}, 그리고\ 
\ref{duality-section-examples}를 보라.여기서는 논문\ 
\cite{Neeman-Grothendieck}, \cite{LN},
\cite{Lipman-notes}, 그리고\  \cite{Neeman-improvement}를 따른다.

\medskip\noindent
Sections \ref{duality-section-upper-shriek},
\ref{duality-section-upper-shriek-properties}, 그리고\ 
\ref{duality-section-base-change-shriek}에서는\  Noether 스킴 사이의
유한형 분리 사상에 대한 중요하고 유용한 상단 느낌표 함자\  $f^!$를 논하며,
마지막으로 개관을 제시하는\  Section
\ref{duality-section-duality}에 이른다.

\medskip\noindent
Section \ref{duality-section-glue}에서는 쌍대화 복합체가 주어진
고정된 국소\  Noether 기저 위에서 작업할 때의 쌍대성과
쌍대화 복합체에 관한 다른 이론을 설명한다
(이 절은\  Hartshorne의 책에 있는 한 주의에 해당한다).

\medskip\noindent
나머지 절들에서는 몇 가지 응용을 제시한다.

\medskip\noindent
이 장의 내용은 대수 공간 위의 쌍대성을 다루는 다음 장으로 이어진다.
Duality for Spaces, Section \ref{spaces-duality-section-introduction}를 보라.






\section{스킴 위의 쌍대화 복합체}
\label{duality-section-dualizing-schemes}

\noindent
국소\  Noether 스킴 위의 쌍대화 복합체를 아핀 국소적으로
대응하는 환 위의 쌍대화 복합체에서 오는 복합체로 정의한다.
이는 완전히 표준적인 정의는 아니지만, 유한 차원\  Noether 스킴에 관한
문헌의 모든 정의와 일치한다.

\begin{lemma}
\label{duality-lemma-equivalent-definitions}
$X$를 국소\  Noether 스킴이라 하고, $K$를\ 
$D(\mathcal{O}_X)$의 대상이라 하자. 다음 조건들은 서로 동치이다.
\begin{enumerate}
\item 모든 아핀 열린집합\  $U = \Spec(A) \subset X$에 대해\ 
$A$의 쌍대화 복합체\  $\omega_A^\bullet$이 존재하여,
$K|_U$는 함자\  $\widetilde{} : D(A) \to D(\mathcal{O}_U)$에 의한\ 
$\omega_A^\bullet$의 상과 동형이다.
\item 아핀 열린 덮개\  $X = \bigcup U_i$, $U_i = \Spec(A_i)$가 존재하여,
각\  $i$에 대해\  $A_i$의 쌍대화 복합체\  $\omega_i^\bullet$이 존재하고\ 
$K|_{U_i}$는 함자\  $\widetilde{} : D(A_i) \to D(\mathcal{O}_{U_i})$에 의한\ 
$\omega_i^\bullet$의 상과 동형이다.
\end{enumerate}
\end{lemma}

\begin{proof}
(2)를 가정하고\  $U = \Spec(A)$를\  $X$의 아핀 열린집합이라 하자.
조건 (2)로부터\  $K$가\  $D_\QCoh(\mathcal{O}_X)$에 속하므로,
연관된 준연접층 복합체가\  $K|_U$와 동형인\  $D(A)$의 대상\ 
$\omega_A^\bullet$을 얻는다. Derived Categories of Schemes, Lemma
\ref{perfect-lemma-affine-compare-bounded}를 보라.
$\omega_A^\bullet$이\  $A$의 쌍대화 복합체임을 보이면 증명이 끝난다.

\medskip\noindent
$X = \bigcup U_i$가 열린 덮개이므로, 각\  $D(f_j)$가 아핀 열린집합\ 
$U_i$ 중 하나의 표준 열린집합이 되도록 하는 표준 열린 덮개\ 
$U = D(f_1) \cup \ldots \cup D(f_m)$을 찾을 수 있다.
Schemes, Lemma \ref{schemes-lemma-standard-open-two-affines}를 보라.
$g_j \in A_{i_j}$에 대해\  $D(f_j) = D(g_j)$라고 하자.
그러면\  $A_{f_j} \cong (A_{i_j})_{g_j}$이고,
Derived Categories of Schemes, Lemma
\ref{perfect-lemma-affine-compare-bounded}에 의해 유도 범주에서
$$
(\omega_A^\bullet)_{f_j} \cong (\omega_i^\bullet)_{g_j}
$$
이다. Dualizing Complexes, Lemma \ref{dualizing-lemma-dualizing-localize}에 의해
복합체\  $(\omega_A^\bullet)_{f_j}$는\  $j = 1, \ldots, m$에 대해\ 
$A_{f_j}$ 위의 쌍대화 복합체이다. 따라서\ 
Dualizing Complexes, Lemma \ref{dualizing-lemma-dualizing-glue}에 의해\ 
$\omega_A^\bullet$은 쌍대화 복합체이다.
\end{proof}

\begin{definition}
\label{duality-definition-dualizing-scheme}
$X$를 국소\  Noether 스킴이라 하자. $D(\mathcal{O}_X)$의 대상\  $K$가\ 
Lemma \ref{duality-lemma-equivalent-definitions}의 동치 조건들을 만족하면\ 
$K$를 {\it 쌍대화 복합체}라 한다.
\end{definition}

\noindent
이 절의 첫머리에 있는 주의들을 보라.
\begin{lemma}
\label{duality-lemma-affine-duality}
$A$를\  Noether 환이라 하고\  $X = \Spec(A)$라 하자. $K, L$을\ 
$D(A)$의 대상이라 하자. $K \in D_{\textit{Coh}}(A)$이고\  $L$의
단사 차원이 유한하면,
$$
R\SheafHom_{\mathcal{O}_X}(\widetilde{K}, \widetilde{L})
=
\widetilde{R\Hom_A(K, L)}
$$
가\  $D(\mathcal{O}_X)$에서 성립한다.
\end{lemma}

\begin{proof}
$L$이 단사\  $A$-가군들로 이루어진 유한 복합체\  $I^\bullet$로
주어졌다고 해도 좋다. $I^\bullet$의 길이에 대한 귀납법과
이 구성들이 완전 삼각형과 양립한다는 사실을 이용하면,
$I$가 단사\  $A$-가군이고\  $L = I[0]$인 경우로 환원된다.
이 경우\  Derived Categories of Schemes, Lemma
\ref{perfect-lemma-quasi-coherence-internal-hom}에 의해\ 
$R\SheafHom_{\mathcal{O}_X}(\widetilde{K}, \widetilde{L})$의\ 
$n$차 코호몰로지 층은 다음 전층에 연관된 층이다.
$$
D(f) \longmapsto \Ext^n_{A_f}(K \otimes_A A_f, I \otimes_A A_f)
$$
$A$가\  Noether 환이므로\  $A_f$-가군\  $I \otimes_A A_f$는 단사이다
(Dualizing Complexes, Lemma
\ref{dualizing-lemma-localization-injective-modules}). 따라서
\begin{align*}
\Ext^n_{A_f}(K \otimes_A A_f, I \otimes_A A_f)
& =
\Hom_{A_f}(H^{-n}(K \otimes_A A_f), I \otimes_A A_f) \\
& =
\Hom_{A_f}(H^{-n}(K) \otimes_A A_f, I \otimes_A A_f) \\
& =
\Hom_A(H^{-n}(K), I) \otimes_A A_f
\end{align*}
마지막 등식은\  $H^{-n}(K)$가 유한\  $A$-가군이기 때문에 성립한다.
Algebra, Lemma \ref{algebra-lemma-hom-from-finitely-presented}를 보라.
이는 표준적 사상
$$
\widetilde{R\Hom_A(K, L)}
\longrightarrow
R\SheafHom_{\mathcal{O}_X}(\widetilde{K}, \widetilde{L})
$$
이 이 경우 준동형임을 보이며, 증명이 끝난다.
\end{proof}
\begin{lemma}
\label{duality-lemma-internal-hom-evaluate-isom}
$X$를\  Noether 스킴이라 하자. $K, L, M \in D_\QCoh(\mathcal{O}_X)$라 하자.
그러면\  Cohomology, Lemma \ref{cohomology-lemma-internal-hom-evaluate}의 사상
$$
R\SheafHom(L, M) \otimes_{\mathcal{O}_X}^\mathbf{L} K
\longrightarrow
R\SheafHom(R\SheafHom(K, L), M)
$$
은 다음 두 경우에 동형이다.
\begin{enumerate}
\item $K \in D^-_{\textit{Coh}}(\mathcal{O}_X)$,
$L \in D^+_{\textit{Coh}}(\mathcal{O}_X)$이고, $M$은 아핀 국소적으로
단사 차원이 유한하다(증명을 보라).
\item $K$와\  $L$은\  $D_{\textit{Coh}}(\mathcal{O}_X)$에 속하고,
대상\  $R\SheafHom(L, M)$의\  tor 차원은 유한하며,
$L$과\  $M$은 아핀 국소적으로 단사 차원이 유한하다
(특히\  $L$과\  $M$은 유계이다).
\end{enumerate}
\end{lemma}

\begin{proof}
(1)의 증명. $M$의 단사 차원이 아핀 국소적으로 유한하다는 것은\ 
$X$가 아핀 열린집합\  $U = \Spec(A)$들로 덮여 있고,
각 아핀 열린집합에서\  $M|_U$에 대응하는\  $D(A)$의 대상
(Derived Categories of Schemes, Lemma
\ref{perfect-lemma-affine-compare-bounded})의 단사 차원이 유한한 것을
뜻한다고 하자\footnote{이 조건은\ 
Noether 스킴\  $X$의 아핀 열린 덮개의 선택과 무관하다.
세부사항은 생략한다.}. 보조정리를 증명하기 위해\  $X$를\  $U$로
바꾸어도 좋다. 즉 어떤\  Noether 환\  $A$에 대해\ 
$X = \Spec(A)$라고 가정해도 좋다.
Derived Categories of Schemes, Lemma
\ref{perfect-lemma-coherent-internal-hom}에 의해\ 
$R\SheafHom(K, L)$은\  $D^+_{\textit{Coh}}(\mathcal{O}_X)$에 속한다.
더 나아가, 화살표의 왼쪽과 오른쪽 변의 구성은\ 
Lemma \ref{duality-lemma-affine-duality}와\  Derived Categories of Schemes, Lemma
\ref{perfect-lemma-quasi-coherence-internal-hom}에 의해
함자\  $D(A) \to D_\QCoh(\mathcal{O}_X)$와 가환한다
(정확히 말하면 여기서는\  $K$, $L$, $M$에 대한 가정과 방금\ 
$R\SheafHom(K, L)$에 대해 증명한 사실을 사용한다).
마지막으로\  More on Algebra, Lemma
\ref{more-algebra-lemma-internal-hom-evaluate-isomorphism} part (2)에 의해
이 화살표는 동형이다.

\medskip\noindent
(2)의 증명. 위와 같이 논증한다. 여기서의 작은 차이는\ 
$R\SheafHom(K, L)$이\  $D_{\textit{Coh}}(\mathcal{O}_X)$에 속한다는 점이다.
실제로 아핀 국소적으로(Lemma \ref{duality-lemma-affine-duality}에 의해 허용된다)
Dualizing Complexes, Lemma
\ref{dualizing-lemma-finite-ext-into-bounded-injective}를 적용할 수 있다.
이제\  More on Algebra, Lemma
\ref{more-algebra-lemma-internal-hom-evaluate-isomorphism-technical}에 의해
결론을 얻는다.
\end{proof}
\begin{lemma}
\label{duality-lemma-dualizing-schemes}
$K$를 국소\  Noether 스킴\  $X$ 위의 쌍대화 복합체라 하자.
그러면\  $K$는\  $D_{\textit{Coh}}(\mathcal{O}_X)$의 대상이고,
$D = R\SheafHom_{\mathcal{O}_X}(-, K)$는 반동치
$$
D :
D_{\textit{Coh}}(\mathcal{O}_X)
\longrightarrow
D_{\textit{Coh}}(\mathcal{O}_X)
$$
를 유도하며, 표준 동형\  $\text{id} \to D \circ D$가 함께 주어진다.
$X$가 준콤팩트이면\  $D$는\ 
$D^+_{\textit{Coh}}(\mathcal{O}_X)$와\ 
$D^-_{\textit{Coh}}(\mathcal{O}_X)$를 서로 바꾸고,
반동치\  $D^b_{\textit{Coh}}(\mathcal{O}_X) \to D^b_{\textit{Coh}}(\mathcal{O}_X)$를 유도한다.
\end{lemma}

\begin{proof}
$U \subset X$를 아핀 열린집합이라 하자. $U = \Spec(A)$라 쓰고,
Lemma \ref{duality-lemma-equivalent-definitions}에서처럼\  $K|_U$에 대응하는\ 
$A$의 쌍대화 복합체를\  $\omega_A^\bullet$이라 하자.
Lemma \ref{duality-lemma-affine-duality}에 의해 도식
$$
\xymatrix{
D_{\textit{Coh}}(A) \ar[r] \ar[d]_{R\Hom_A(-, \omega_A^\bullet)} &
D_{\textit{Coh}}(\mathcal{O}_U) \ar[d]^{R\SheafHom_{\mathcal{O}_X}(-, K|_U)} \\
D_{\textit{Coh}}(A) \ar[r] &
D(\mathcal{O}_U)
}
$$
은 가환한다. 따라서\  $D$는\  $D_{\textit{Coh}}(\mathcal{O}_X)$를\ 
$D_{\textit{Coh}}(\mathcal{O}_X)$ 안으로 보낸다. 또한 표준적 사상
$$
L
\longrightarrow
R\SheafHom_{\mathcal{O}_X}(K, K) \otimes_{\mathcal{O}_X}^\mathbf{L} L
\longrightarrow
R\SheafHom_{\mathcal{O}_X}(R\SheafHom_{\mathcal{O}_X}(L, K), K)
$$
(두 번째 화살표에는\  Cohomology, Lemma
\ref{cohomology-lemma-internal-hom-evaluate}를 사용한다)은 모든\  $L$에 대해
동형이다. 이는\  Dualizing Complexes, Lemma
\ref{dualizing-lemma-dualizing}\footnote{다른 방법으로는 먼저 아핀 국소적으로
계산하여\  $R\SheafHom_{\mathcal{O}_X}(K, K) = \mathcal{O}_X$임을 보인 뒤,
Lemma \ref{duality-lemma-internal-hom-evaluate-isom} part (2)를 사용하여
이 사상이 동형임을 보일 수 있다.}에 의해 아핀 위에서 참이고,
대수에서 일어나는 일을 아핀 위에서 되찾는다는 사실을 이미 보았기 때문이다.
준콤팩트한 경우 함자\  $D$의 유계성에 관한 명제도\ 
Dualizing Complexes, Lemma \ref{dualizing-lemma-dualizing}의
대응하는 명제들로부터 따른다.
\end{proof}
\noindent
$X$를 국소 환 달린 공간이라 하자. $D(\mathcal{O}_X)$의 대상\  $L$이
유도 텐서곱으로 주어지는\  $D(\mathcal{O}_X)$의 대칭 모노이드 구조에 대한
가역 대상이면 이를 {\it 가역}이라 한다. Cohomology, Lemma
\ref{cohomology-lemma-invertible-derived}에서 이는\  $L$이 완전하고,
어떤 정수\  $n_i$들에 대해\  $L|_{U_i} \cong \mathcal{O}_{U_i}[-n_i]$가
되도록 하는 열린 덮개\  $X = \bigcup U_i$가 존재한다는 뜻임을 보았다.
이 경우 함수
$$
x \mapsto n_x,\quad
\text{where }n_x\text{ is the unique integer such that }
H^{n_x}(L_x) \not = 0
$$
는\  $X$ 위에서 국소 상수이다. 특히\ 
$L = \bigoplus H^n(L)[-n]$이고, 이는\  $L$을 나타내는
미분이 영인 잘 정의된\  $\mathcal{O}_X$-가군 복합체를 준다.

\begin{lemma}
\label{duality-lemma-dualizing-unique-schemes}
$X$를 국소\  Noether 스킴이라 하자. $K$와\  $K'$가\  $X$ 위의
쌍대화 복합체이면, $D(\mathcal{O}_X)$의 어떤 가역 대상\  $L$에 대해\ 
$K'$는\  $K \otimes_{\mathcal{O}_X}^\mathbf{L} L$과 동형이다.
\end{lemma}

\begin{proof}
$$
L = R\SheafHom_{\mathcal{O}_X}(K, K')
$$
로 놓자. 이는\  $D(\mathcal{O}_X)$의 가역 대상이다. 아핀 국소적으로
그렇기 때문이다. Dualizing Complexes, Lemma
\ref{dualizing-lemma-dualizing-unique}와 그 증명을 보라.
같은 이유로 평가 사상\ 
$L \otimes_{\mathcal{O}_X}^\mathbf{L} K \to K'$는 동형이다.
\end{proof}
\begin{lemma}
\label{duality-lemma-dimension-function-scheme}
$X$를 국소\  Noether 스킴이라 하고, $\omega_X^\bullet$을\ 
$X$ 위의 쌍대화 복합체라 하자. 그러면\  $X$는 보편 사슬 조건을 만족하고,
다음으로 정의되는 함수\  $X \to \mathbf{Z}$는 차원 함수이다.
$$
x \longmapsto \delta(x)\text{ such that }
\omega_{X, x}^\bullet[-\delta(x)]
\text{ is a normalized dualizing complex over }
\mathcal{O}_{X, x}
$$
\end{lemma}

\begin{proof}
아핀인 경우인\  Dualizing Complexes, Lemma
\ref{dualizing-lemma-dimension-function}와 정의들로부터 즉시 따른다.
\end{proof}

\begin{lemma}
\label{duality-lemma-sitting-in-degrees}
$X$를 국소\  Noether 스킴이라 하고, $\omega_X^\bullet$을 연관된 차원 함수가\ 
$\delta$인\  $X$ 위의 쌍대화 복합체라 하자. $\mathcal{F}$를 연접\ 
$\mathcal{O}_X$-가군이라 하고\ 
$\mathcal{E}^i = \SheafExt^{-i}_{\mathcal{O}_X}(\mathcal{F}, \omega_X^\bullet)$로 놓자.
그러면\  $\mathcal{E}^i$는 연접\  $\mathcal{O}_X$-가군이고,
$x \in X$에 대해 다음이 성립한다.
\begin{enumerate}
\item $\mathcal{E}^i_x$는\ 
$\delta(x) \leq i \leq \delta(x) + \dim(\text{Supp}(\mathcal{F}_x))$일 때만
영이 아닐 수 있다.
\item $\dim(\text{Supp}(\mathcal{E}^{i + \delta(x)}_x)) \leq i$이다.
\item $\text{depth}(\mathcal{F}_x)$는\ 
$\mathcal{E}_x^{i + \delta(x)} \not = 0$을 만족하는 가장 작은 정수\ 
$i \geq 0$이다.
\item 다음이 성립한다.
$x \in \text{Supp}(\bigoplus_{j \leq i} \mathcal{E}^j)
\Leftrightarrow
\text{depth}_{\mathcal{O}_{X, x}}(\mathcal{F}_x) + \delta(x) \leq i$.
\end{enumerate}
\end{lemma}

\begin{proof}
Lemma \ref{duality-lemma-dualizing-schemes}에 의해\  $\mathcal{E}^i$는 연접이다.
$x$의 아핀 근방을 택하고\  Derived Categories of Schemes, Lemma
\ref{perfect-lemma-quasi-coherence-internal-hom}과\ 
More on Algebra, Lemma
\ref{more-algebra-lemma-base-change-RHom} part (3)을 사용하면
$$
\mathcal{E}^i_x =
\SheafExt^{-i}_{\mathcal{O}_X}(\mathcal{F}, \omega_X^\bullet)_x =
\Ext^{-i}_{\mathcal{O}_{X, x}}(\mathcal{F}_x,
\omega_{X, x}^\bullet) =
\Ext^{\delta(x) - i}_{\mathcal{O}_{X, x}}(\mathcal{F}_x,
\omega_{X, x}^\bullet[-\delta(x)])
$$
을 얻는다. Lemma \ref{duality-lemma-dimension-function-scheme}에서의\  $\delta$의 구성에 의해
(1), (2), (3)은\  Dualizing Complexes, Lemma
\ref{dualizing-lemma-sitting-in-degrees}로 환원된다.
(4)는 (3)과 (1)의 형식적 귀결이다.
\end{proof}



\section{추이 함자의 오른쪽 수반}
\label{duality-section-twisted-inverse-image}

\noindent
이 절과 다음 절들의 참고문헌은\ 
\cite{Neeman-Grothendieck}, \cite{LN},
\cite{Lipman-notes}, 그리고\  \cite{Neeman-improvement}이다.

\medskip\noindent
$f : X \to Y$를 스킴의 사상이라 하자.
이 절에서는 함자\ 
$Rf_* : D_\QCoh(\mathcal{O}_X) \to D_\QCoh(\mathcal{O}_Y)$의
오른쪽 수반을 고찰한다. 문헌에서는 이 함자가 존재할 때
때때로\  $f^{\times}$로 표기한다. 이 표기는 보편적으로 받아들여지지 않으므로
여기서는 사용하지 않는다. 또한 그러한 함자에\  $f^!$라는 표기를 쓰지 않는다.
일반 사상\  $f$에 대해 이 표기는\  \cite{RD}의 표기와 충돌하기 때문이다.

\begin{lemma}
\label{duality-lemma-twisted-inverse-image}
\begin{reference}
이는\  \cite[Example 4.2]{Neeman-Grothendieck}와 거의 같다.
\end{reference}
$f : X \to Y$를 준분리 준콤팩트 스킴 사이의 사상이라 하자.
함자\  $Rf_* : D_\QCoh(\mathcal{O}_X) \to D_\QCoh(\mathcal{O}_Y)$는
오른쪽 수반을 갖는다.
\end{lemma}

\begin{proof}
Derived Categories, Proposition \ref{derived-proposition-brown}의 가정들을
확인하여 오른쪽 수반이 존재함을 보이겠다. 먼저 범주\ 
$D_\QCoh(\mathcal{O}_X)$는 직합을 갖는다. Derived Categories of Schemes, Lemma
\ref{perfect-lemma-quasi-coherence-direct-sums}를 보라.
Derived Categories of Schemes, Theorem
\ref{perfect-theorem-bondal-van-den-Bergh}에 의해 범주\ 
$D_\QCoh(\mathcal{O}_X)$는 콤팩트 생성된다.
$X$와\  $Y$가 준콤팩트이고 준분리이므로\  $f$도 그러하다.
Schemes, Lemmas \ref{schemes-lemma-compose-after-separated}와\ 
\ref{schemes-lemma-quasi-compact-permanence}를 보라.
따라서 함자\  $Rf_*$는 직합과 가환한다. Derived Categories of Schemes, Lemma
\ref{perfect-lemma-quasi-coherence-pushforward-direct-sums}를 보라.
이로써 증명이 끝난다.
\end{proof}
\begin{example}
\label{duality-example-affine-twisted-inverse-image}
$A \to B$를 환 사상이라 하자. $Y = \Spec(A)$, $X = \Spec(B)$라 하고,
$f : X \to Y$를\  $A \to B$에 대응하는 사상이라 하자.
그러면\  $Rf_* : D_\QCoh(\mathcal{O}_X) \to D_\QCoh(\mathcal{O}_Y)$는
동치\  $D(B) \to D_\QCoh(\mathcal{O}_X)$와\ 
$D(A) \to D_\QCoh(\mathcal{O}_Y)$ 아래에서 제한 함자\ 
$D(B) \to D(A)$에 대응한다. 따라서 오른쪽 수반은\ 
Dualizing Complexes, Section \ref{dualizing-section-trivial}의 함자\ 
$K \longmapsto R\Hom(B, K)$에 대응한다.
\end{example}

\begin{example}
\label{duality-example-does-not-preserve-coherent}
$f : X \to Y$가\  Noether 스킴의 유한형 분리 사상이면,
$Rf_* : D_\QCoh(\mathcal{O}_X) \to D_\QCoh(\mathcal{O}_Y)$의
오른쪽 수반은\  $D_{\textit{Coh}}(\mathcal{O}_Y)$를\ 
$D_{\textit{Coh}}(\mathcal{O}_X)$ 안으로 보내지 않는다.
실제로\  $k$를 체라 하고 사상\  $f : \mathbf{A}^1_k \to \Spec(k)$를 생각하자.
Example \ref{duality-example-affine-twisted-inverse-image}에 의해 이는\ 
$A = k$, $B = k[x]$일 때\  $R\Hom(B, -)$가\ 
$D_{\textit{Coh}}(A)$를\  $D_{\textit{Coh}}(B)$ 안으로 보내는지를 묻는 문제에
해당한다. 이는 참이 아니다. 왜냐하면
$$
R\Hom(k[x], k) = \left(\prod\nolimits_{n \geq 0} k\right)[0]
$$
이고 이것은 유한\  $k[x]$-가군이 아니기 때문이다. 따라서\ 
$a(\mathcal{O}_Y)$는 연접 코호몰로지 층들을 갖지 않는다.
\end{example}

\begin{example}
\label{duality-example-does-not-preserve-bounded-above}
$f : X \to Y$가\  Noether 스킴의 고유 사상, 나아가 유한 사상이라 해도,
$Rf_* : D_\QCoh(\mathcal{O}_X) \to D_\QCoh(\mathcal{O}_Y)$의
오른쪽 수반은\  $D_\QCoh^-(\mathcal{O}_Y)$를\ 
$D_\QCoh^-(\mathcal{O}_X)$ 안으로 보내지 않는다. 실제로\  $k$를 체라 하고,
$k[\epsilon]$을\  $k$ 위의 쌍대수라 하며,
$X = \Spec(k)$, $Y = \Spec(k[\epsilon])$라 하자.
그러면 모든\  $i \geq 0$에 대해\  $\Ext^i_{k[\epsilon]}(k, k)$는 영이 아니다.
따라서\  Example \ref{duality-example-affine-twisted-inverse-image}에 의해\ 
$a(\mathcal{O}_Y)$는 위로 유계가 아니다.
\end{example}
\begin{lemma}
\label{duality-lemma-twisted-inverse-image-bounded-below}
$f : X \to Y$를 준콤팩트 준분리 스킴의 사상이라 하자.
$a : D_\QCoh(\mathcal{O}_Y) \to D_\QCoh(\mathcal{O}_X)$를\ 
Lemma \ref{duality-lemma-twisted-inverse-image}의\  $Rf_*$의 오른쪽 수반이라 하자.
그러면\  $a$는\  $D^+_\QCoh(\mathcal{O}_Y)$를\ 
$D^+_\QCoh(\mathcal{O}_X)$ 안으로 보낸다. 더 정확히 말하면,
$H^i(K) = 0$ for $i \leq c$이면\ 
$H^i(a(K)) = 0$ for $i \leq c - N$이 되게 하는 정수\  $N$이 존재한다.
\end{lemma}

\begin{proof}
Derived Categories of Schemes, Lemma
\ref{perfect-lemma-quasi-coherence-direct-image}에 의해
함자\  $Rf_*$의 코호몰로지 차원은 유한하다. 다시 말해,
$H^i(L) = 0$ for $i \geq c$이면\ 
$H^i(Rf_*L) = 0$ for $i \geq N + c$가 되게 하는 정수\  $N$이 존재한다.
$K \in D^+_\QCoh(\mathcal{O}_Y)$이고\ 
$H^i(K) = 0$ for $i \leq c$라고 하자. 그러면 위에서 말한 사실에 의해
$$
\Hom_{D(\mathcal{O}_X)}(\tau_{\leq c - N}a(K), a(K)) =
\Hom_{D(\mathcal{O}_Y)}(Rf_*\tau_{\leq c - N}a(K), K) = 0
$$
이다. 이는 분명히\  $H^i(a(K)) = 0$ for $i \leq c - N$을 함의한다.
\end{proof}
\noindent
$f : X \to Y$를 준분리 준콤팩트 스킴의 사상이라 하자.
$a$를\  $Rf_* : D_\QCoh(\mathcal{O}_X) \to D_\QCoh(\mathcal{O}_Y)$의
오른쪽 수반이라 하자. 모든\  $K \in D_\QCoh(\mathcal{O}_Y)$와\ 
$L \in D_\QCoh(\mathcal{O}_X)$에 대해 표준적 사상
\begin{equation}
\label{duality-equation-sheafy-trace}
Rf_*R\SheafHom_{\mathcal{O}_X}(L, a(K))
\longrightarrow
R\SheafHom_{\mathcal{O}_Y}(Rf_*L, K)
\end{equation}
을 얻는다. 이 사상은 합성
$$
Rf_*R\SheafHom_{\mathcal{O}_X}(L, a(K)) \to
R\SheafHom_{\mathcal{O}_Y}(Rf_*L, Rf_*a(K)) \to
R\SheafHom_{\mathcal{O}_Y}(Rf_*L, K)
$$
으로 구성된다. 첫 번째 화살표는\  Cohomology, Remark
\ref{cohomology-remark-projection-formula-for-internal-hom}의 것이고,
두 번째 화살표는 수반의 여단위원\  $Rf_*a(K) \to K$이다.

\begin{lemma}
\label{duality-lemma-iso-on-RSheafHom}
$f : X \to Y$를 준콤팩트 준분리 스킴의 사상이라 하자.
$a$를\  $Rf_* : D_\QCoh(\mathcal{O}_X) \to D_\QCoh(\mathcal{O}_Y)$의
오른쪽 수반이라 하자. $L \in D_\QCoh(\mathcal{O}_X)$이고\ 
$K \in D_\QCoh(\mathcal{O}_Y)$라 하자. 그러면 사상
(\ref{duality-equation-sheafy-trace})
$$
Rf_*R\SheafHom_{\mathcal{O}_X}(L, a(K))
\longrightarrow
R\SheafHom_{\mathcal{O}_Y}(Rf_*L, K)
$$
은\  Derived Categories of Schemes, Section
\ref{perfect-section-better-coherator}에서 논한 함자\ 
$DQ_Y : D(\mathcal{O}_Y) \to D_\QCoh(\mathcal{O}_Y)$를 적용한 뒤 동형이 된다.
\end{lemma}

\begin{proof}
Derived Categories of Schemes, Lemma
\ref{perfect-lemma-better-coherator}에 의해\  $DQ_Y$가 존재하므로
이 명제는 의미가 있다. $DQ_Y$는 포함 함자\ 
$D_\QCoh(\mathcal{O}_Y) \to D(\mathcal{O}_Y)$의 오른쪽 수반이므로,
보조정리를 증명하려면 임의의\  $M \in D_\QCoh(\mathcal{O}_Y)$에 대해
사상 (\ref{duality-equation-sheafy-trace})이 전단사
$$
\Hom_Y(M, Rf_*R\SheafHom_{\mathcal{O}_X}(L, a(K)))
\longrightarrow
\Hom_Y(M, R\SheafHom_{\mathcal{O}_Y}(Rf_*L, K))
$$
를 유도함을 보이면 된다. 이를 위해 다음 등식열을 사용한다.
\begin{align*}
\Hom_Y(M, Rf_*R\SheafHom_{\mathcal{O}_X}(L, a(K)))
& =
\Hom_X(Lf^*M, R\SheafHom_{\mathcal{O}_X}(L, a(K))) \\
& =
\Hom_X(Lf^*M \otimes_{\mathcal{O}_X}^\mathbf{L} L, a(K)) \\
& =
\Hom_Y(Rf_*(Lf^*M \otimes_{\mathcal{O}_X}^\mathbf{L} L), K) \\
& =
\Hom_Y(M \otimes_{\mathcal{O}_Y}^\mathbf{L} Rf_*L, K) \\
& =
\Hom_Y(M, R\SheafHom_{\mathcal{O}_Y}(Rf_*L, K))
\end{align*}
첫 번째 등식은\  Cohomology, Lemma \ref{cohomology-lemma-adjoint}에 의해,
두 번째 등식은\  Cohomology, Lemma \ref{cohomology-lemma-internal-hom}에 의해,
세 번째 등식은\  $a$의 구성에 의해 성립한다. 네 번째 등식은\ 
Derived Categories of Schemes, Lemma
\ref{perfect-lemma-cohomology-base-change}에 의해 성립하며, 이것이 중요한 단계이다.
다섯 번째 등식은\  Cohomology, Lemma
\ref{cohomology-lemma-internal-hom}에 의해 성립한다.
\end{proof}
\begin{example}
\label{duality-example-iso-on-RSheafHom}
Lemma \ref{duality-lemma-iso-on-RSheafHom}의 명제는 ``coherator'' $DQ_Y$를
적용하지 않으면 참이 아니다. 실제로\  $Y = \Spec(R)$이고\ 
$X = \mathbf{A}^1_R$이라 하자. $L = \mathcal{O}_X$와\ 
$K = \mathcal{O}_Y$를 택하자. 화살표의 왼쪽 변은\ 
$D_\QCoh(\mathcal{O}_Y)$에 속하지만, 오른쪽 변은
준연접이 아닌\  $\prod_{n \geq 0} \mathcal{O}_Y$와 동형이다.
\end{example}

\begin{remark}
\label{duality-remark-iso-on-RSheafHom}
Lemma \ref{duality-lemma-iso-on-RSheafHom}의 상황에서\ 
Derived Categories of Schemes, Lemma
\ref{perfect-lemma-pushforward-better-coherator}에 의해
$$
DQ_Y(Rf_*R\SheafHom_{\mathcal{O}_X}(L, a(K))) =
Rf_* DQ_X(R\SheafHom_{\mathcal{O}_X}(L, a(K)))
$$
이다. 따라서\  $R\SheafHom_{\mathcal{O}_X}(L, a(K)) \in D_\QCoh(\mathcal{O}_X)$이면
화살표의 왼쪽 변에서\  $DQ_Y$를 ``지울'' 수 있다. 한편\ 
$R\SheafHom_{\mathcal{O}_Y}(Rf_*L, K) \in D_\QCoh(\mathcal{O}_Y)$임을 알면
화살표의 오른쪽 변에서\  $DQ_Y$를 ``지울'' 수 있다.
둘 다 참이면 (\ref{duality-equation-sheafy-trace})가 동형임을 알 수 있다.
이를\  Derived Categories of Schemes, Lemma
\ref{perfect-lemma-quasi-coherence-internal-hom}과 결합하면,
$Rf_*R\SheafHom_{\mathcal{O}_X}(L, a(K)) \to
R\SheafHom_{\mathcal{O}_Y}(Rf_*L, K)$는 다음 경우에 동형이다.
\begin{enumerate}
\item $L$과\  $Rf_*L$이 완전한 경우.
\item $K$가 아래로 유계이고\  $L$과\  $Rf_*L$이 유사연접인 경우.
\end{enumerate}
(2)에서는\  $K$가 아래로 유계이면\  $a(K)$도 아래로 유계라는 사실을 사용한다.
Lemma \ref{duality-lemma-twisted-inverse-image-bounded-below}를 보라.
\end{remark}

\begin{example}
\label{duality-example-iso-on-RSheafHom-noetherian}
$f : X \to Y$를\  Noether 스킴의 고유 사상이라 하고,
$L \in D^-_{\textit{Coh}}(X)$ 및\  $K \in D^+_{\QCoh}(\mathcal{O}_Y)$라 하자.
그러면 사상\  $Rf_*R\SheafHom_{\mathcal{O}_X}(L, a(K)) \to
R\SheafHom_{\mathcal{O}_Y}(Rf_*L, K)$는 동형이다.
실제로\  Derived Categories of Schemes, Lemmas
\ref{perfect-lemma-identify-pseudo-coherent-noetherian}와\ 
\ref{perfect-lemma-direct-image-coherent}에 의해 복합체\  $L$과\  $Rf_*L$은
유사연접이고, Remark \ref{duality-remark-iso-on-RSheafHom}의 논의를 적용할 수 있다.
\end{example}
\begin{lemma}
\label{duality-lemma-iso-global-hom}
$f : X \to Y$를 준분리 준콤팩트 스킴의 사상이라 하자.
모든\  $L \in D_\QCoh(\mathcal{O}_X)$와\  $K \in D_\QCoh(\mathcal{O}_Y)$에 대해
(\ref{duality-equation-sheafy-trace})는 대역 유도\  hom의 동형\ 
$R\Hom_X(L, a(K)) \to R\Hom_Y(Rf_*L, K)$를 유도한다.
\end{lemma}

\begin{proof}
Cohomology, Section \ref{cohomology-section-global-RHom}의 구성에 의해
$$
R\Hom_X(L, a(K)) =
R\Gamma(X, R\SheafHom_{\mathcal{O}_X}(L, a(K))) =
R\Gamma(Y, Rf_*R\SheafHom_{\mathcal{O}_X}(L, a(K)))
$$
이고
$$
R\Hom_Y(Rf_*L, K) = R\Gamma(Y, R\SheafHom_{\mathcal{O}_Y}(Rf_*L, K))
$$
이다. 따라서 보조정리는\  Lemma \ref{duality-lemma-iso-on-RSheafHom}의 귀결이다.
구체적으로, $D(\mathcal{O}_Y)$의 사상\  $E \to E'$가 동형\ 
$DQ_Y(E) \to DQ_Y(E')$를 유도하면 준동형\ 
$R\Gamma(Y, E) \to R\Gamma(Y, E')$를 유도한다. 실제로\ 
$\mathcal{O}_Y[-i]$는\  $D_\QCoh(\mathcal{O}_Y)$에 속하고\  $DQ_Y$는
포함 함자\  $D_\QCoh(\mathcal{O}_Y) \to D(\mathcal{O}_Y)$의 오른쪽 수반이므로\ 
$H^i(Y, E) = \Ext^i_Y(\mathcal{O}_Y, E) = \Hom(\mathcal{O}_Y[-i], E) =
\Hom(\mathcal{O}_Y[-i], DQ_Y(E))$이다.
\end{proof}






\section{추이 함자의 오른쪽 수반과 열린집합으로의 제한}
\label{duality-section-restriction-to-opens}

\noindent
이 절에서는 추이 함자의 오른쪽 수반이 열린 부분스킴으로의 제한과
어느 정도 가환하는지를 연구한다. 이는 기저변환 문제이므로,
먼저 더 일반적인 상황을 논하자.

\medskip\noindent
우리는 흔히 추이 함자의 오른쪽 수반이 기저변환과 가환하는지를 알고 싶다.
따라서 준콤팩트 준분리 스킴들의\  Cartesian 정사각형
\begin{equation}
\label{duality-equation-base-change}
\vcenter{
\xymatrix{
X' \ar[r]_{g'} \ar[d]_{f'} & X \ar[d]^f \\
Y' \ar[r]^g & Y
}
}
\end{equation}
을 생각한다. $Rf_*$와\  $Rf'_*$의 오른쪽 수반
(Lemma \ref{duality-lemma-twisted-inverse-image})을 각각
$$
a  : D_\QCoh(\mathcal{O}_Y) \to D_\QCoh(\mathcal{O}_X)
\quad\text{and}\quad
a'  : D_\QCoh(\mathcal{O}_{Y'}) \to D_\QCoh(\mathcal{O}_{X'})
$$
로 나타내자. Cohomology, Remark \ref{cohomology-remark-base-change}의
기저변환 사상을 생각하자. 이는 준연접 코호몰로지를 갖는 층들의
유도 범주 위에서 함자들의 변환
$$
Lg^* \circ Rf_* \longrightarrow Rf'_* \circ L(g')^*
$$
을 유도한다. 따라서 반대 방향으로 오른쪽 수반들 사이의 변환
$$
a \circ Rg_* \longleftarrow Rg'_* \circ a'
$$
을 얻는다.

\begin{lemma}
\label{duality-lemma-flat-precompose-pus}
도식 (\ref{duality-equation-base-change})에서\  $g$가 평탄이라고 가정하거나,
더 일반적으로\  $f$와\  $g$가\  Tor 독립이라고 가정하자. 그러면\ 
$a \circ Rg_* \leftarrow Rg'_* \circ a'$는 동형이다.
\end{lemma}

\begin{proof}
이 경우 모든\  $D_\QCoh(\mathcal{O}_X)$의\  $K$에 대해 기저변환 사상\ 
$Lg^* \circ Rf_* K \longrightarrow Rf'_* \circ L(g')^*K$는\ 
Derived Categories of Schemes, Lemma
\ref{perfect-lemma-compare-base-change}에 의해 동형이다.
따라서 대응하는 수반 함자들 사이의 변환도 동형이다.
\end{proof}
\noindent
$f : X \to Y$를 준콤팩트 준분리 스킴의 사상이라 하자.
$V \subset Y$를 준콤팩트 열린 부분스킴이라 하고\ 
$U = f^{-1}(V)$로 놓자. 이는 (\ref{duality-equation-base-change})에서와 같은\ 
Cartesian 정사각형
$$
\xymatrix{
U \ar[r]_{j'} \ar[d]_{f|_U} & X \ar[d]^f \\
V \ar[r]^j & Y
}
$$
을 준다. Lemma \ref{duality-lemma-flat-precompose-pus}에 의해
사상\  $\xi : a \circ Rj_* \leftarrow Rj'_* \circ a'$는 동형이다.
여기서\  $a$와\  $a'$는 각각\  $Rf_*$와\  $R(f|_U)_*$의 오른쪽 수반이다.
함자\  $D_\QCoh(\mathcal{O}_Y) \to D_\QCoh(\mathcal{O}_U)$의 변환
\begin{equation}
\label{duality-equation-sheafy}
(j')^* \circ a \to
(j')^* \circ a \circ Rj_* \circ j^* \xrightarrow{\xi^{-1}}
(j')^* \circ Rj'_* \circ a' \circ j^* \to a' \circ j^*
\end{equation}
을 얻는다. 첫 번째 화살표는\  $\text{id} \to Rj_* \circ j^*$에서 오고,
마지막 화살표는 동형\  $(j')^* \circ Rj'_* \to \text{id}$에서 온다.
특히 (\ref{duality-equation-sheafy})를\  $K$에서 평가한 것이 동형일 필요충분조건은\ 
$a(K)|_U \to a(Rj_*(K|_V))|_U$가 동형인 것이다.

\begin{example}
\label{duality-example-not-supported-on-inverse-image}
어떤\  Noether 스킴의 유한 사상\  $f : X \to Y$와\ 
$K \in D_{\textit{Coh}}(\mathcal{O}_Y)$가 존재하여,
(\ref{duality-equation-sheafy})를 이 대상에서 평가한 것은 동형이 아니다.
실제로\  $k$를 체라 하고\  $A = k[x, \epsilon]$, $\epsilon^2 = 0$이라 하며,
$B = k[x] = A/(\epsilon)$라 하자. $X = \Spec(B) \to Y = \Spec(A)$를
생각한다. $n \in \mathbf{N}$에 대해\  $M_n = A/(\epsilon, x^n)$로 놓자.
$B$는\  $\epsilon$의 곱셈으로 주어지는 사상들을 갖는 자유 주기 분해\ 
$\ldots \to A \to A \to A$를 가지므로
$$
\Ext^i_A(B, M_n) = M_n,\quad i \geq 0
$$
이다. $D_{\textit{Coh}}(A)$의 대상\ 
$K = \bigoplus M_n[n] = \prod M_n[n]$를 생각하자
(Derived Categories, Lemmas \ref{derived-lemma-direct-sums}와\ 
\ref{derived-lemma-products}에 의해\  $D(A)$에서 등식이다).
Example \ref{duality-example-affine-twisted-inverse-image}에 의해\  $a(K)$는\ 
$R\Hom(B, K)$에 대응하고, 위의 계산에 의해
$$
H^0(R\Hom(B, K)) = \Ext^0_A(B, K) =
\prod\nolimits_{n \geq 1} \Ext^n_A(B. M_n) = 
\prod\nolimits_{n \geq 1} M_n
$$
이다. 그러나 이 가군에는\  $x$의 어떤 거듭제곱으로도 소멸하지 않는
원소가 있는 반면, 복합체\  $K$의 코호몰로지의 모든 원소는\ 
$x$의 어떤 거듭제곱으로 소멸한다. 다시 말해\ 
$V = D(x)$, $U = D(x)$ 및 복합체\  $K$에 대한
(\ref{duality-equation-sheafy})의 사상은 동형일 수 없다. 실제로\ 
$(j')^*(a(K))$는 영이 아니고\  $a'(j^*K)$는 영이다.
\end{example}
\begin{lemma}
\label{duality-lemma-when-sheafy}
$f : X \to Y$를 준콤팩트 준분리 스킴의 사상이라 하자.
$a$를\  $Rf_* : D_\QCoh(\mathcal{O}_X) \to D_\QCoh(\mathcal{O}_Y)$의
오른쪽 수반이라 하자. $V \subset Y$를 준콤팩트 열린집합이라 하고,
그 역상을\  $U \subset X$라 하자.
\begin{enumerate}
\item $Y \setminus V$ 위에 지지된 모든\ 
$Q \in D_\QCoh^+(\mathcal{O}_Y)$에 대해\  $a(Q)$가\ 
$X \setminus U$ 위에 지지될 필요충분조건은
(\ref{duality-equation-sheafy})가\  $D_\QCoh^+(\mathcal{O}_Y)$의 모든\  $K$에서
동형인 것이다.
\item $Y \setminus V$ 위에 지지된 모든\ 
$Q \in D_\QCoh(\mathcal{O}_Y)$에 대해\  $a(Q)$가\ 
$X \setminus U$ 위에 지지될 필요충분조건은
(\ref{duality-equation-sheafy})가\  $D_\QCoh(\mathcal{O}_Y)$의 모든\  $K$에서
동형인 것이다.
\item $a$가 직합과 가환하면, (1)의 동치 조건들은 (2)의 동치 조건들을
함의한다.
\end{enumerate}
\end{lemma}

\begin{proof}
(1)의 증명. $K \in D_\QCoh^+(\mathcal{O}_Y)$라 하자.
완전 삼각형
$$
K \to Rj_*K|_V \to Q \to K[1]
$$
을 택하자. $Q$는\  $D_\QCoh^+(\mathcal{O}_Y)$에 속하고
(Derived Categories of Schemes, Lemma
\ref{perfect-lemma-quasi-coherence-direct-image}),
$Y \setminus V$ 위에 지지된다
(Derived Categories of Schemes, Definition
\ref{perfect-definition-supported-on}).
$a$를 적용하면\  $X$ 위의 완전 삼각형
$$
a(K) \to a(Rj_*K|_V) \to a(Q) \to a(K)[1]
$$
을 얻는다. $a(Q)$가\  $X \setminus U$ 위에 지지되면,
$U$로 제한했을 때 사상\  $a(K)|_U \to a(Rj_*K|_V)|_U$는 동형이다.
즉 (\ref{duality-equation-sheafy})는\  $K$에서 동형이다. 역은 자명하다.

\medskip\noindent
(2)의 증명은 (1)의 증명과 정확히 같다.

\medskip\noindent
(3)의 증명. (1)의 동치 조건들이 성립한다고 가정하자.
$T = Y \setminus V$로 놓자. 코호몰로지 층이 각각\ 
$T$와\  $f^{-1}(T)$ 위에 지지된 복합체들의 범주를\ 
$D_{\QCoh, T}(\mathcal{O}_Y)$와\ 
$D_{\QCoh, f^{-1}(T)}(\mathcal{O}_X)$로 나타내겠다.
$a$가 직합과 가환하므로, 대상들이
$$
\{Q \in D_{\QCoh, T}(\mathcal{O}_Y) \mid
a(Q) \in D_{\QCoh, f^{-1}(T)}(\mathcal{O}_X)\}
$$
인 충만하고 포화된 삼각 부분범주\  $\mathcal{D}$는 직합에 의해 보존되며,
따라서 유도 쌍극한에 의해서도 보존된다. 한편 범주\ 
$D_{\QCoh, T}(\mathcal{O}_Y)$는 완전 대상\  $E$로 생성된다
(Derived Categories of Schemes, Lemma
\ref{perfect-lemma-generator-with-support}를 보라).
가정에 의해\  $E \in \mathcal{D}$이다. Derived Categories, Lemma
\ref{derived-lemma-write-as-colimit}에 의해\ 
$D_{\QCoh, T}(\mathcal{O}_Y)$의 모든 대상\  $Q$는 전이 사상의 원뿔들이\ 
$E$의 평행이동들의 직합인 계\ 
$Q_1 \to Q_2 \to Q_3 \to \ldots$의 유도 쌍극한이다.
귀납적으로 모든\  $n$에 대해\  $Q_n \in \mathcal{D}$임을 알 수 있고,
마지막으로\  $Q$가\  $\mathcal{D}$에 속함을 얻는다.
따라서 (2)의 동치 조건들이 성립한다.
\end{proof}
\begin{lemma}
\label{duality-lemma-proper-noetherian}
$Y$를 준콤팩트 준분리 스킴이라 하고, $f : X \to Y$를 고유 사상이라 하자.
다음 중 하나가 성립한다고 하자\footnote{이 증명은 준콤팩트 준분리 스킴의
사상 중에서\  $X$ 위의 모든 완전한\  $P$에 대해\  $Rf_*P$가 유사연접인
사상에도 적용된다. Kiehl의 정리\  \cite{Kiehl}로부터\ 
$f$가 고유하고 유사연접이면 이 조건이 성립함을 쉽게 알 수 있다.
이것이 이 보조정리와 이 장의 몇몇 다른 결과에 대한 올바른 일반성이다.}.
\begin{enumerate}
\item $f$가 평탄하고 유한 표시인 경우.
\item $Y$가\  Noether 스킴인 경우.
\end{enumerate}
그러면\  $Y$의 모든 준콤팩트 열린집합\  $V$에 대해\ 
Lemma \ref{duality-lemma-when-sheafy} part (1)의 동치 조건들이 성립한다.
\end{lemma}

\begin{proof}
$Q \in D^+_\QCoh(\mathcal{O}_Y)$가\  $Y \setminus V$ 위에 지지된다고 하자.
모순을 얻기 위해\  $a(Q)$가\  $X \setminus U$ 위에 지지되지 않는다고
가정하자. 그러면\  $U$ 위의 완전 복합체\  $P_U$와 영이 아닌 사상\ 
$P_U \to a(Q)|_U$를 찾을 수 있다
(Derived Categories of Schemes, Theorem
\ref{perfect-theorem-bondal-van-den-Bergh}로부터 따른다).
이제\  Derived Categories of Schemes, Lemma
\ref{perfect-lemma-lift-map-from-perfect-complex-with-support}를 사용하면,
$X$ 위의 완전 복합체\  $P$와\  $U$로의 제한이 영이 아닌 사상\ 
$P \to a(Q)$가 존재한다고 가정해도 좋다.
$a$의 정의에 의해 이 사상은 사상\  $Rf_*P \to Q$에 수반된다.

\medskip\noindent
복합체\  $Rf_*P$는 유사연접이다. (1)의 경우 이는\ 
Derived Categories of Schemes, Lemma
\ref{perfect-lemma-flat-proper-pseudo-coherent-direct-image-general}로부터 따른다.
(2)의 경우 이는\  Derived Categories of Schemes, Lemmas
\ref{perfect-lemma-direct-image-coherent}와\ 
\ref{perfect-lemma-identify-pseudo-coherent-noetherian}로부터 따른다.
따라서\  Derived Categories of Schemes, Lemma
\ref{perfect-lemma-map-from-pseudo-coherent-to-complex-with-support}를 적용하여,
$V$로의 제한이 동형이고 합성\ 
$I \otimes^\mathbf{L}_{\mathcal{O}_Y} Rf_*P \to Rf_*P \to Q$가 영인
완전 복합체들의 사상\  $I \to \mathcal{O}_Y$를 얻는다.
Derived Categories of Schemes, Lemma
\ref{perfect-lemma-cohomology-base-change}에 의해\ 
$I \otimes^\mathbf{L}_{\mathcal{O}_Y} Rf_*P =
Rf_*(Lf^*I \otimes^\mathbf{L}_{\mathcal{O}_X} P)$이다.
따라서 합성
$$
Lf^*I \otimes^\mathbf{L}_{\mathcal{O}_X} P \to P \to a(Q)
$$
은 영이다. 그러나\  $U$로 제한하면 이는 영이 아니라고 가정한 사상\ 
$P|_U \to a(Q)|_U$이다. 이 모순으로 증명이 끝난다.
\end{proof}











\section{추이 함자의 오른쪽 수반과 기저변환, I}
\label{duality-section-base-change-map}

\noindent
사상 (\ref{duality-equation-sheafy})은 기저변환 사상의 특수한 경우이다.
구체적으로, 준콤팩트 준분리 스킴들의\  Cartesian 정사각형
$$
\xymatrix{
X' \ar[r]_{g'} \ar[d]_{f'} & X \ar[d]^f \\
Y' \ar[r]^g & Y
}
$$
즉 (\ref{duality-equation-base-change})와 같은 도식이 주어졌다고 하자.
$f$와\  $g$가 {\bf Tor 독립}이라고 가정하자. 그러면 합성
\begin{equation}
\label{duality-equation-base-change-map}
L(g')^* \circ a \to
L(g')^* \circ a \circ Rg_* \circ Lg^* \leftarrow
L(g')^* \circ Rg'_* \circ a' \circ Lg^* \to a' \circ Lg^*
\end{equation}
으로 주어지는 함자\ 
$D_\QCoh(\mathcal{O}_Y) \to D_\QCoh(\mathcal{O}_{X'})$의 사상을
생각할 수 있다. 첫 번째 화살표는 수반 사상\ 
$\text{id} \to Rg_* Lg^*$에서 오고, 마지막 화살표는 수반 사상\ 
$L(g')^*Rg'_* \to \text{id}$에서 온다. 가운데 화살표를 역으로
바꾸려면\  Tor 독립성 가정이 필요하다. Lemma
\ref{duality-lemma-flat-precompose-pus}를 보라. 달리 말해,
$L(g')^*$와\  $R(g')_*$의 수반성에 의해
(\ref{duality-equation-base-change-map})를 자연 변환
$$
a \to a \circ Rg_* \circ Lg^* \leftarrow Rg'_* \circ a' \circ Lg^*
$$
으로 생각할 수 있으며, 여기서도 두 번째 화살표는 가역이다.
$M \in D_\QCoh(\mathcal{O}_X)$이고\ 
$K \in D_\QCoh(\mathcal{O}_Y)$이면\  Yoneda 함자에서 이 사상은
\begin{align*}
\Hom_X(M, a(K))
& =
\Hom_Y(Rf_*M, K) \\
& \to
\Hom_Y(Rf_*M, Rg_* Lg^*K) \\
& =
\Hom_{Y'}(Lg^*Rf_*M, Lg^*K) \\
& \leftarrow
\Hom_{Y'}(Rf'_* L(g')^*M, Lg^*K) \\
& =
\Hom_{X'}(L(g')^*M, a'(Lg^*K)) \\
& =
\Hom_X(M, Rg'_*a'(Lg^*K))
\end{align*}
으로 주어진다. 여기서 왼쪽을 가리키는 화살표는\ 
Derived Categories of Schemes, Lemma
\ref{perfect-lemma-compare-base-change}의 기저변환 정리에 의해 가역이다.
이는 구성을 조금 더 명시적으로 나타낸다.

\medskip\noindent
이 절에서는 먼저 이 기저변환 사상이 정사각형들을 쌓을 때
몇 가지 자연스러운 호환성을 만족함을 증명한다. 이는 통상적인
기저변환 사상에 대한\  Cohomology, Remarks
\ref{cohomology-remark-compose-base-change}와\ 
\ref{cohomology-remark-compose-base-change-horizontal}의 호환성과 같다.
처음 읽을 때에는 이 절의 나머지를 건너뛰기를 권한다.
\begin{lemma}
\label{duality-lemma-compose-base-change-maps}
준콤팩트 준분리 스킴들의 가환 도식
$$
\xymatrix{
X' \ar[r]_k \ar[d]_{f'} & X \ar[d]^f \\
Y' \ar[r]^l \ar[d]_{g'} & Y \ar[d]^g \\
Z' \ar[r]^m & Z
}
$$
을 생각하자. 두 정사각형이 모두\  Cartesian이고,
$f$와\  $l$뿐 아니라\  $g$와\  $m$도\  Tor 독립이라고 하자.
그러면 두 정사각형에 대한 사상
(\ref{duality-equation-base-change-map})들을 합성하면 바깥 직사각형의
기저변환 사상을 얻는다(정확한 명제는 증명을 보라).
\end{lemma}

\begin{proof}
가정으로부터\  $g \circ f$와\  $m$이\  Tor 독립임이 따른다
(세부사항은 생략한다). 따라서 명제는 의미가 있다.
이 증명에서는\  $Lk^*$ 대신\  $k^*$를, $Rf_*$ 대신\  $f_*$를 쓴다.
$a$, $b$, $c$를 각각\  $f$, $g$, $g \circ f$에 대한\ 
Lemma \ref{duality-lemma-twisted-inverse-image}의 오른쪽 수반이라 하고,
프라임이 붙은 것들도 마찬가지로 정하자. 위쪽 정사각형에 대응하는
화살표는 합성
$$
\gamma_{top} :
k^* \circ a \to k^* \circ a \circ l_* \circ l^*
\xleftarrow{\xi_{top}} k^* \circ k_* \circ a' \circ l^* \to a' \circ l^*
$$
이다. 여기서\  $\xi_{top} : k_* \circ a' \to a \circ l_*$는 동형이고
(따라서 역으로 바꿀 수 있고), 기저변환 사상\ 
$l^* \circ f_* \to f'_* \circ k^*$에 ``쌍대''인 화살표이다.
바깥쪽 화살표들은 표준 사상\  $1 \to l_* \circ l^*$와\ 
$k^* \circ k_* \to 1$에서 온다. 두 번째 정사각형에 대해서도 마찬가지로
$$
\gamma_{bot} :
l^* \circ b \to l^* \circ b \circ m_* \circ m^*
\xleftarrow{\xi_{bot}} l^* \circ l_* \circ b' \circ m^* \to b' \circ m^*
$$
를 얻는다. 바깥 직사각형에 대해서는
$$
\gamma_{rect} :
k^* \circ c \to k^* \circ c \circ m_* \circ m^*
\xleftarrow{\xi_{rect}} k^* \circ k_* \circ c' \circ m^* \to c' \circ m^*
$$
를 얻는다. $(g \circ f)_* = g_* \circ f_*$이므로\ 
$c = a \circ b$이고, 마찬가지로\  $c' = a' \circ b'$이다.
보조정리의 명제는\  $\gamma_{rect}$가 합성
$$
k^* \circ c = k^* \circ a \circ b \xrightarrow{\gamma_{top}}
a' \circ l^* \circ b \xrightarrow{\gamma_{bot}}
a' \circ b' \circ m^* = c' \circ m^*
$$
과 같다는 것이다. 이를 보기 위해 다음 도식을 살펴보자.
$$
\xymatrix{
& & k^* \circ a \circ b \ar[d] \ar[lldd] \\
& & k^* \circ a \circ l_* \circ l^* \circ b \ar[ld] \\
k^* \circ a \circ b \circ m_* \circ m^* \ar[r] &
k^* \circ a \circ l_* \circ l^* \circ b \circ m_* \circ m^* &
k^* \circ k_* \circ a' \circ l^* \circ b \ar[u]_{\xi_{top}} \ar[d] \ar[ld] \\
& k^*\circ k_* \circ a' \circ l^* \circ b \circ m_* \circ m^*
\ar[u]_{\xi_{top}} \ar[rd] &
a' \circ l^* \circ b \ar[d] \\
k^* \circ k_* \circ a' \circ b' \circ m^* \ar[uu]_{\xi_{rect}} \ar[ddrr] &
k^*\circ k_* \circ a' \circ l^* \circ l_* \circ b' \circ m^*
\ar[u]_{\xi_{bot}} \ar[l] \ar[dr] &
a' \circ l^* \circ b \circ m_* \circ m^* \\
& & a' \circ l^* \circ l_* \circ b' \circ m^* \ar[u]_{\xi_{bot}} \ar[d] \\
& & a' \circ b' \circ m^*
}
$$
오른쪽을 따라 내려가면 위의 합성을 얻고, 왼쪽을 따라 내려가면\ 
$\gamma_{rect}$를 얻는다. 이 도식의 오른쪽에 있는 모든 사각형은\ 
Categories, Lemma \ref{categories-lemma-properties-2-cat-cats},
또는 더 간단하게\  Categories, Definition
\ref{categories-definition-horizontal-composition} 앞의 논의에 의해 가환한다.
따라서 화살표\  $\xi_{top}$, $\xi_{bot}$, $\xi_{rect}$를 역으로 바꾸었을 때
도식
$$
\xymatrix{
a \circ l_* \circ l^* \circ b \circ m_* &
a \circ b \circ m_* \ar[l] \\
k_* \circ a' \circ l^* \circ b \circ m_* \ar[u]_{\xi_{top}} & \\
k_* \circ a' \circ l^* \circ l_* \circ b' \ar[u]_{\xi_{bot}} \ar[r] &
k_* \circ a' \circ b' \ar[uu]_{\xi_{rect}}
}
$$
이 가환함을 보이면 충분하다(이는 원래 도식이 가환한다고 요구하는 것과
다름에 유의하라). 그런데 도식
$$
\xymatrix{
& a \circ l_* \circ l^* \circ b \circ m_* \\
a \circ l_* \circ l^* \circ l_* \circ b'
\ar[ru]^{\xi_{bot}} & &
k_* \circ a' \circ l^* \circ b \circ m_* \ar[ul]_{\xi_{top}} \\
& k_* \circ a' \circ l^* \circ l_* \circ b'
\ar[ul]^{\xi_{top}} \ar[ur]_{\xi_{bot}}
}
$$
은\  Categories, Lemma \ref{categories-lemma-properties-2-cat-cats}에 의해
가환한다. 또한 도식들
$$
\vcenter{
\xymatrix{
a \circ l_* \circ l^* \circ b \circ m_* & a \circ b \circ m \ar[l] \\
a \circ l_* \circ l^* \circ l_* \circ b' \ar[u] &
a \circ l_* \circ b' \ar[l] \ar[u]
}
}
\quad\text{and}\quad
\vcenter{
\xymatrix{
a \circ l_* \circ l^* \circ l_* \circ b' \ar[r] & a \circ l_* \circ b' \\
k_* \circ a' \circ l^* \circ l_* \circ b' \ar[u] \ar[r] &
k_* \circ a' \circ b' \ar[u]
}
}
$$
은 가환하고(인용한 문헌을 보라), 합성\ 
$l_* \to l_* \circ l^* \circ l_* \to l_*$가 항등이므로
$$
k \circ a' \circ b' \xrightarrow{\xi_{bot}} a \circ l_* \circ b
\xrightarrow{\xi_{top}} a \circ b \circ m_*
$$
가\  $\xi_{rect}$와 같음을 보이면 충분하다
($a \circ b = c$와\  $a' \circ b' = c'$라는 동일시를 사용한다).
이는\  Cohomology, Remark
\ref{cohomology-remark-compose-base-change}의 명제에 쌍대인 명제이므로
증명이 끝난다.
\end{proof}
\begin{lemma}
\label{duality-lemma-compose-base-change-maps-horizontal}
준콤팩트 준분리 스킴들의 가환 도식
$$
\xymatrix{
X'' \ar[r]_{g'} \ar[d]_{f''} & X' \ar[r]_g \ar[d]_{f'} & X \ar[d]^f \\
Y'' \ar[r]^{h'} & Y' \ar[r]^h & Y
}
$$
을 생각하자. 두 정사각형이 모두\  Cartesian이고,
$f$와\  $h$뿐 아니라\  $f'$와\  $h'$도\  Tor 독립이라고 하자.
그러면 두 정사각형에 대한 사상
(\ref{duality-equation-base-change-map})들을 합성하면 바깥 직사각형의
기저변환 사상을 얻는다(정확한 명제는 증명을 보라).
\end{lemma}

\begin{proof}
가정으로부터\  $f$와\  $h \circ h'$가\  Tor 독립임이 따른다
(세부사항은 생략한다). 따라서 명제는 의미가 있다.
이 증명에서는\  $Lg^*$ 대신\  $g^*$를, $Rf_*$ 대신\  $f_*$를 쓴다.
$a$, $a'$, $a''$를 각각\  $f$, $f'$, $f''$에 대한\ 
Lemma \ref{duality-lemma-twisted-inverse-image}의 오른쪽 수반이라 하자.
오른쪽 정사각형에 대응하는 화살표는 합성
$$
\gamma_{right} :
g^* \circ a \to g^* \circ a \circ h_* \circ h^*
\xleftarrow{\xi_{right}} g^* \circ g_* \circ a' \circ h^* \to a' \circ h^*
$$
이다. 여기서\  $\xi_{right} : g_* \circ a' \to a \circ h_*$는 동형이고
(따라서 역으로 바꿀 수 있고), 기저변환 사상\ 
$h^* \circ f_* \to f'_* \circ g^*$에 ``쌍대''인 화살표이다.
바깥쪽 화살표들은 표준 사상\  $1 \to h_* \circ h^*$와\ 
$g^* \circ g_* \to 1$에서 온다. 왼쪽 정사각형에 대해서도 마찬가지로
$$
\gamma_{left} :
(g')^* \circ a' \to (g')^* \circ a' \circ (h')_* \circ (h')^*
\xleftarrow{\xi_{left}}
(g')^* \circ (g')_* \circ a'' \circ (h')^* \to a'' \circ (h')^*
$$
를 얻는다. 바깥 직사각형에 대해서는
$$
\gamma_{rect} :
k^* \circ a \to
k^* \circ a \circ m_* \circ m^* \xleftarrow{\xi_{rect}}
k^* \circ k_* \circ a'' \circ m^* \to
a'' \circ m^*
$$
를 얻는다. 여기서\  $k = g \circ g'$이고\  $m = h \circ h'$이다.
$k^* = (g')^* \circ g^*$이고\  $m^* = (h')^* \circ h^*$이다.
보조정리의 명제는\  $\gamma_{rect}$가 합성
$$
k^* \circ a =
(g')^* \circ g^* \circ a \xrightarrow{\gamma_{right}}
(g')^* \circ a' \circ h^* \xrightarrow{\gamma_{left}}
a'' \circ (h')^* \circ h^* = a'' \circ m^*
$$
과 같다는 것이다. 이를 보기 위해 다음 도식을 살펴보자.
$$
\xymatrix{
& (g')^* \circ g^* \circ a \ar[d] \ar[ddl] \\
& (g')^* \circ g^* \circ a \circ h_* \circ h^* \ar[ld] \\
(g')^* \circ g^* \circ a \circ h_* \circ (h')_* \circ (h')^* \circ h^* &
(g')^* \circ g^* \circ g_* \circ a' \circ h^*
\ar[u]_{\xi_{right}} \ar[d] \ar[ld] \\
(g')^* \circ g^* \circ g_* \circ a' \circ (h')_* \circ (h')^* \circ h^*
\ar[u]_{\xi_{right}} \ar[dr] &
(g')^* \circ a' \circ h^* \ar[d] \\
(g')^* \circ g^* \circ g_* \circ (g')_* \circ a'' \circ (h')^* \circ h^*
\ar[u]_{\xi_{left}} \ar[ddr] \ar[dr] &
(g')^* \circ a' \circ (h')_* \circ (h')^* \circ h^* \\
& (g')^*\circ (g')_* \circ a'' \circ (h')^* \circ h^*
\ar[u]_{\xi_{left}} \ar[d] \\
& a'' \circ (h')^* \circ h^*
}
$$
오른쪽을 따라 내려가면 위의 합성을 얻고, 왼쪽을 따라 내려가면\ 
$\gamma_{rect}$를 얻는다. 이 도식의 오른쪽에 있는 모든 사각형은\ 
Categories, Lemma \ref{categories-lemma-properties-2-cat-cats},
또는 더 간단하게\  Categories, Definition
\ref{categories-definition-horizontal-composition} 앞의 논의에 의해 가환한다.
따라서
$$
g_* \circ (g')_* \circ a'' \xrightarrow{\xi_{left}}
g_* \circ a' \circ (h')_* \xrightarrow{\xi_{right}}
a \circ h_* \circ (h')_*
$$
가\  $\xi_{rect}$와 같음을 보이면 충분하다. 이는\  Cohomology, Remark
\ref{cohomology-remark-compose-base-change-horizontal}의 명제에 쌍대인
명제이므로 증명이 끝난다.
\end{proof}
\begin{remark}
\label{duality-remark-going-around}
준콤팩트 준분리 스킴들의 가환 도식
$$
\xymatrix{
X'' \ar[r]_{k'} \ar[d]_{f''} & X' \ar[r]_k \ar[d]_{f'} & X \ar[d]^f \\
Y'' \ar[r]^{l'} \ar[d]_{g''} & Y' \ar[r]^l \ar[d]_{g'} & Y \ar[d]^g \\
Z'' \ar[r]^{m'} & Z' \ar[r]^m & Z
}
$$
을 생각하자. 모든 정사각형이\  Cartesian이고,
$(f, l)$, $(g, m)$, $(f', l')$, $(g', m')$이\ 
Tor 독립인 사상들의 쌍이라고 하자.
$a$, $a'$, $a''$, $b$, $b'$, $b''$를 각각\ 
$f$, $f'$, $f''$, $g$, $g'$, $g''$에 대한\ 
Lemma \ref{duality-lemma-twisted-inverse-image}의 오른쪽 수반이라 하자.
도식의 정사각형들을 다음과 같이\  $A$, $B$, $C$, $D$로 표시하자.
$$
\begin{matrix}
A & B \\
C & D
\end{matrix}
$$
그러면 각 정사각형에 대한 사상
(\ref{duality-equation-base-change-map})들은 다음과 같다
(여기서는\  $k^* = Lk^*$ 등을 쓴다).
$$
\begin{matrix}
\gamma_A : (k')^* \circ a' \to a'' \circ (l')^* &
\gamma_B : k^* \circ a \to a' \circ l^* \\
\gamma_C : (l')^* \circ b' \to b'' \circ (m')^* &
\gamma_D : l^* \circ b \to b' \circ m^*
\end{matrix}
$$
$2 \times 1$ 및\  $1 \times 2$ 직사각형에 대해서는 네 개의
기저변환 사상을 더 얻는다.
$$
\begin{matrix}
\gamma_{A + B} : (k \circ k')^* \circ a \to a'' \circ (l \circ l')^* \\
\gamma_{C + D} : (l \circ l')^* \circ b \to b'' \circ (m \circ m')^* \\
\gamma_{A + C} : (k')^* \circ (a' \circ b') \to (a'' \circ b'') \circ (m')^* \\
\gamma_{B + D} : k^* \circ (a \circ b) \to (a' \circ b') \circ m^*
\end{matrix}
$$
Lemma \ref{duality-lemma-compose-base-change-maps-horizontal}에 의해
$$
\gamma_{A + B} = \gamma_A \circ \gamma_B, \quad
\gamma_{C + D} = \gamma_C \circ \gamma_D
$$
이고, Lemma \ref{duality-lemma-compose-base-change-maps}에 의해
$$
\gamma_{A + C} = \gamma_C \circ \gamma_A, \quad
\gamma_{B + D} = \gamma_D \circ \gamma_B
$$
이다. 더 정확하게는\  Categories, Section
\ref{categories-section-formal-cat-cat}의 표기법으로\ 
$\gamma_{A + B} = (\gamma_A \star \text{id}_{l^*}) \circ
(\text{id}_{(k')^*} \star \gamma_B)$라고 써야 하며, 나머지도 마찬가지이다.
그러나 변환의 출발점과 도착점을 알면 어떤 항등변환을 덧붙여야 하는지
알 수 있으므로, Lemmas \ref{duality-lemma-compose-base-change-maps}와\ 
\ref{duality-lemma-compose-base-change-maps-horizontal}의 증명에서처럼
항등변환과의\  $\star$ 곱을 생략하는 표기의 남용을 계속하겠다.
이 모든 논의를 종합하면 선험적으로 두 변환
$$
(k')^* \circ k^* \circ a \circ b
\longrightarrow
a'' \circ b'' \circ (m')^* \circ m^*
$$
을 얻는다. 즉
$$
\gamma_C \circ \gamma_A \circ \gamma_D \circ \gamma_B =
\gamma_{A + C} \circ \gamma_{B + D}
$$
와
$$
\gamma_C \circ \gamma_D \circ \gamma_A \circ \gamma_B =
\gamma_{C + D} \circ \gamma_{A + B}
$$
이다. 이 두 변환이 같다는 사실을 지적하는 것이 이 주의의 요점이다.
실제로 이를 보이려면 도식
$$
\xymatrix{
(k')^* \circ a' \circ l^* \circ b \ar[r]_{\gamma_D} \ar[d]_{\gamma_A} &
(k')^* \circ a' \circ b' \circ m^* \ar[d]^{\gamma_A} \\
a'' \circ (l')^* \circ l^* \circ b \ar[r]^{\gamma_D} &
a'' \circ (l')^* \circ b' \circ m^*
}
$$
이 가환함을 보이면 충분하다. 이는\  Categories, Lemma
\ref{categories-lemma-properties-2-cat-cats}, 또는 더 간단하게\ 
Categories, Definition
\ref{categories-definition-horizontal-composition} 앞의 논의에 의해 성립한다.
\end{remark}






\section{추이 함자의 오른쪽 수반과 기저변환, II}
\label{duality-section-base-change-II}

\noindent
이 절에서는\  Section \ref{duality-section-base-change-map}의 기저변환 사상이
몇 가지 경우에 동형임을 증명한다. 먼저 추이 함자의 오른쪽 수반을
만드는 것이 열린집합으로의 제한과 가환한다면 아핀 열린집합 위에서
확인하는 것으로 충분함을 관찰한다.

\begin{remark}
\label{duality-remark-check-over-affines}
준콤팩트 준분리 스킴들의\  Cartesian 도식
$$
\xymatrix{
X' \ar[r]_{g'} \ar[d]_{f'} & X \ar[d]^f \\
Y' \ar[r]^g & Y
}
$$
을 생각하자. 여기서\  $(g, f)$는\  Tor 독립이다.
$V \subset Y$와\  $V' \subset Y'$를\  $g(V') \subset V$를 만족하는
아핀 열린집합이라 하자. Cartesian 도식들
$$
\vcenter{
\xymatrix{
U \ar[r] \ar[d] & X \ar[d] \\
V \ar[r] & Y
}
}
\quad\text{and}\quad
\vcenter{
\xymatrix{
U' \ar[r] \ar[d] & X' \ar[d] \\
V' \ar[r] & Y'
}
}
$$
을 만들자. $K$와 첫 번째 도식에 대한 (\ref{duality-equation-sheafy}) 및\ 
$Lg^*K$와 두 번째 도식에 대한 (\ref{duality-equation-sheafy})가
동형이라고 가정하자. 그러면 기저변환 사상
(\ref{duality-equation-base-change-map})
$$
L(g')^*a(K) \longrightarrow a'(Lg^*K)
$$
을\  $U'$에 제한한 것은\  $K|_V$와\  Cartesian 도식
$$
\xymatrix{
U' \ar[r] \ar[d] & U \ar[d] \\
V' \ar[r] & V
}
$$
에 대한 기저변환 사상 (\ref{duality-equation-base-change-map})과 동형이다.
이는 (\ref{duality-equation-sheafy})가 기저변환 사상
(\ref{duality-equation-base-change-map})의 특수한 경우이고, 정사각형들을
수평으로 쌓을 때 기저변환 사상들이 올바르게 합성되기 때문이다.
Lemma \ref{duality-lemma-compose-base-change-maps-horizontal}을 보라.
따라서\  $U'$에 제한한 기저변환 사상이 동형인지 확인하려면
마지막 도식만 다루면 충분하다.
\end{remark}
\begin{lemma}
\label{duality-lemma-more-base-change}
도식 (\ref{duality-equation-base-change})에서 다음을 가정하자.
\begin{enumerate}
\item $g : Y' \to Y$는 아핀 스킴의 사상이고,
\item $f : X \to Y$는 고유하며,
\item $f$와\  $g$는\  Tor 독립이다.
\end{enumerate}
그러면 다음 경우에 기저변환 사상 (\ref{duality-equation-base-change-map})은
동형
$$
L(g')^*a(K) \longrightarrow a'(Lg^*K)
$$
을 유도한다.
\begin{enumerate}
\item $f$가 평탄하고 유한 표시이면 모든\ 
$K \in D_\QCoh(\mathcal{O}_Y)$에 대해,
\item $f$가 완전하고\  $Y$가\  Noether이면 모든\ 
$K \in D_\QCoh(\mathcal{O}_Y)$에 대해,
\item $g$가 유한\  Tor 차원을 갖고\  $Y$가\  Noether이면\ 
$K \in D_\QCoh^+(\mathcal{O}_Y)$에 대해.
\end{enumerate}
\end{lemma}
\begin{proof}
$Y = \Spec(A)$와\  $Y' = \Spec(A')$로 쓰자. 아핀 사상의 기저변환이므로
사상\  $g'$는 아핀이다. $M$을\  $D_\QCoh(\mathcal{O}_X)$의 완전 생성자라 하자.
Derived Categories of Schemes, Theorem
\ref{perfect-theorem-bondal-van-den-Bergh}를 보라. 그러면\ 
$L(g')^*M$은\  $D_\QCoh(\mathcal{O}_{X'})$의 생성자이다.
Derived Categories of Schemes, Remark
\ref{perfect-remark-pullback-generator}를 보라. 따라서
(\ref{duality-equation-base-change-map})이 대역\  Hom 복합체들의 동형
\begin{equation}
\label{duality-equation-iso}
R\Hom_{X'}(L(g')^*M, L(g')^*a(K))
\longrightarrow
R\Hom_{X'}(L(g')^*M, a'(Lg^*K))
\end{equation}
을 유도함을 보이면 충분하다. Cohomology, Section
\ref{cohomology-section-global-RHom}을 보라. 이는\ 
$L(g')^*a(K) \to a'(Lg^*K)$의 원뿔이 영임을 함의하기 때문이다.
증명의 구조는 다음과 같다. 먼저 이\  Hom 복합체들이 동형임을 보이고,
증명의 마지막 부분에서 그 동형이 (\ref{duality-equation-iso})에 의해
유도됨을 보이겠다.

\medskip\noindent
왼쪽 변. $M$이 완전하므로 표준 사상
$$
R\Hom_X(M, a(K)) \otimes^\mathbf{L}_A A'
\longrightarrow
R\Hom_{X'}(L(g')^*M, L(g')^*a(K))
$$
은\  Derived Categories of Schemes, Lemma
\ref{perfect-lemma-affine-morphism-and-hom-out-of-perfect}에 의해 동형이다.
이를\  Lemma \ref{duality-lemma-iso-global-hom}의 동형\ 
$R\Hom_Y(Rf_*M, K) = R\Hom_X(M, a(K))$과 결합하면 왼쪽 변은\ 
$R\Hom_Y(Rf_*M, K) \otimes^\mathbf{L}_A A'$와 같다.

\medskip\noindent
오른쪽 변. 먼저\  Lemma \ref{duality-lemma-iso-global-hom}의 동형
$$
R\Hom_{X'}(L(g')^*M, a'(Lg^*K)) = R\Hom_{Y'}(Rf'_*L(g')^*M, Lg^*K)
$$
을 사용한다. 그다음\  Derived Categories of Schemes, Lemma
\ref{perfect-lemma-compare-base-change}에 의해 기저변환 사상\ 
$Lg^*Rf_*M \to Rf'_*L(g')^*M$이 동형임을 사용한다. 따라서 이를\ 
$R\Hom_{Y'}(Lg^*Rf_*M, Lg^*K)$로 다시 쓸 수 있다.
$Y$, $Y'$가 아핀이고\  $K$, $Rf_*M$이\ 
$D_\QCoh(\mathcal{O}_Y)$에 속하므로(Derived Categories of Schemes,
Lemma \ref{perfect-lemma-quasi-coherence-direct-image}), $D(A')$에서
표준 사상
$$
\beta :
R\Hom_Y(Rf_*M, K) \otimes^\mathbf{L}_A A'
\longrightarrow
R\Hom_{Y'}(Lg^*Rf_*M, Lg^*K)
$$
을 얻는다. 이는\  More on Algebra, Equation
(\ref{more-algebra-equation-base-change-RHom})의 화살표이다. 여기서는\ 
Derived Categories of Schemes, Lemmas
\ref{perfect-lemma-affine-compare-bounded}와\ 
\ref{perfect-lemma-quasi-coherence-internal-hom}을 사용하여 대수적 표현과
층의 표현 사이를 오갔다.
\begin{enumerate}
\item $f$가 평탄하고 유한 표시이면 복합체\  $Rf_*M$은\ 
Derived Categories of Schemes, Lemma
\ref{perfect-lemma-flat-proper-perfect-direct-image-general}에 의해\ 
$Y$ 위에서 완전하고, $\beta$는\  More on Algebra, Lemma
\ref{more-algebra-lemma-base-change-RHom} part (1)에 의해 동형이다.
\item $f$가 완전하고\  $Y$가\  Noether이면 복합체\  $Rf_*M$은\ 
More on Morphisms, Lemma
\ref{more-morphisms-lemma-perfect-proper-perfect-direct-image}에 의해\ 
$Y$ 위에서 완전하고, 앞에서와 같이\  $\beta$는 동형이다.
\item $g$가 유한\  tor 차원을 갖고\  $Y$가\  Noether이면 복합체\  $Rf_*M$은\ 
$Y$ 위에서 유사연접이다(Derived Categories of Schemes, Lemmas
\ref{perfect-lemma-direct-image-coherent}와\ 
\ref{perfect-lemma-identify-pseudo-coherent-noetherian}). 그리고\ 
$\beta$는\  More on Algebra, Lemma
\ref{more-algebra-lemma-base-change-RHom} part (4)에 의해 동형이다.
\end{enumerate}
따라서 앞 단락과 같은 답을 얻는다.

\medskip\noindent
증명의 나머지 부분에서는 둘째 및 셋째 단락에서 얻은
(\ref{duality-equation-iso})의 왼쪽 변과 오른쪽 변의 동일시가 실제로
(\ref{duality-equation-iso})에 의해 주어짐을 보인다. 식을 다루기 쉽게 하기 위해\ 
$(-, -)_X = R\Hom_X(-, -)$를 쓰고, $- \otimes_A^\mathbf{L} A'$ 대신\ 
$- \otimes A'$를 쓰며, $g^* = Lg^*$와\  $f_* = Rf_*$로 줄여 쓰겠다.
다음 가환 도식을 생각하자.
$$
\xymatrix{
((g')^*M, (g')^*a(K))_{X'} \ar[d] &
(M, a(K))_X \otimes A' \ar[l]^-\alpha \ar[d] &
(f_*M, K)_Y \otimes A' \ar@{=}[l] \ar[d] \\
((g')^*M, (g')^*a(g_*g^*K))_{X'} &
(M, a(g_*g^*K))_X \otimes A' \ar[l]^-\alpha &
(f_*M, g_*g^*K)_Y \otimes A' \ar@{=}[l] \ar@/_4pc/[dd]_{\mu'} \\
((g')^*M, (g')^*g'_*a'(g^*K))_{X'} \ar[u] \ar[d] &
(M, g'_*a'(g^*K))_X \otimes A' \ar[u] \ar[l]^-\alpha \ar[ld]^\mu &
(f_*M, K) \otimes A' \ar[d]^\beta \\
((g')^*M, a'(g^*K))_{X'} &
(f'_*(g')^*M, g^*K)_{Y'} \ar@{=}[l] \ar[r] &
(g^*f_*M, g^*K)_{Y'}
}
$$
$\alpha$로 표시한 화살표들은 네 꼭짓점이\  $X', X, Y', Y$인
도식에 대해\  Derived Categories of Schemes, Lemma
\ref{perfect-lemma-affine-morphism-and-hom-out-of-perfect}가 주는 사상들이다.
수평 화살표들이 각 항에 대해 함자적이므로 도식의 위쪽 부분은 가환한다.
가운데의 수직 화살표들은\  Lemma \ref{duality-lemma-flat-precompose-pus}의
가역 변환\  $g'_* \circ a' \to a \circ g_*$에서 오므로 가운데 정사각형은
가환한다. 왼쪽을 따라 내려가는 것이 (\ref{duality-equation-iso})이다.
위쪽 수평 화살표들은 증명의 둘째 단락에서 사용한 동일시들을 주고,
$\beta$를 포함한 아래쪽 수평 화살표들은 셋째 단락에서 사용한
동일시들을 준다. $E \in D(A)$, $E' \in D(A')$ 및\  $D(A)$의 사상\ 
$c : E \to E'$가 주어졌다고 하자. $c$와 제한 및 기저변환의 수반성에
의해 유도되는 사상을\  $\mu_c : E \otimes A' \to E'$로 나타내겠다.
$c$가 명확하면\  $\mu = \mu_c$로 쓰며, 즉 표기에서\  $c$를 생략한다.
도식의 사상\  $\mu$는 동일시\ 
$(M, g'_*a(g^*K))_X = ((g')^*M, a'(g^*K))_{X'}$가 주는\  $c$에 대한
이러한 사상이다. $\mu$를 포함하는 삼각형은\ 
Derived Categories of Schemes, Remark
\ref{perfect-remark-multiplication-map}에 의해 가환한다.

\medskip\noindent
도식
$$
\xymatrix{
(M, a(g_*g^*K))_X &
(f_*M, g_* g^*K)_Y \ar@{=}[l] &
(g^*f_*M, g^*K)_{Y'} \ar@{=}[l] \\
(M, g'_* a'(g^*K))_X \ar[u] &
((g')^*M, a'(g^*K))_{X'} \ar@{=}[l] &
(f'_*(g')^*M, g^*K)_{Y'} \ar@{=}[l] \ar[u]
}
$$
이 변환\  $g'_* \circ a' \to a \circ g_*$의 정의에 의해 가환함을
관찰하라. 위와 같이\  $\mu'$가 동일시\ 
$(f_*M, g_*g^*K)_X = (g^*f_*M, g^*K)_{Y'}$에 대응한다고 하면
육각형도 가환한다. 따라서\  $\beta$가\ 
$(f_*M, K)_Y \otimes A' \to (f_*M, g_*g^*K)_X \otimes A'$와\ 
$\mu'$의 합성과 같음을 보이면 충분하다. 이를 위해서는 유도된 두 사상\ 
$(f_*M, K)_Y \to (g^*f_*M, g^*K)_{Y'}$가 같음을 보이면 충분하다.
다시 말해 모든\  $E, K \in D(A)$에 대해 도식
$$
\xymatrix{
R\Hom_A(E, K) \ar[rr]_{\text{induced by }\beta} \ar[rd] & &
R\Hom_{A'}(E \otimes_A^\mathbf{L} A', K \otimes_A^\mathbf{L} A') \\
& R\Hom_A(E, K \otimes_A^\mathbf{L} A') \ar[ru]
}
$$
이 가환함을 보이면 충분하다. More on Algebra, Section
\ref{more-algebra-section-base-change-RHom}에서\  $\beta$를 바로 이런 방식으로
구성하므로 증명이 끝난다.
\end{proof}







\section{추이 함자의 오른쪽 수반과 대각합 사상}
\label{duality-section-trace}

\noindent
$f : X \to Y$를 준콤팩트 준분리 스킴의 사상이라 하자.
$a : D_\QCoh(\mathcal{O}_Y) \to D_\QCoh(\mathcal{O}_X)$를\ 
Lemma \ref{duality-lemma-twisted-inverse-image}의 오른쪽 수반이라 하자.
Categories, Section \ref{categories-section-adjoint}에 의해 함자의 변환
$$
\text{Tr}_f : Rf_* \circ a \longrightarrow \text{id}
$$
을 얻는다. $K \in D_\QCoh(\mathcal{O}_Y)$에 대응하는 사상\ 
$\text{Tr}_{f, K} : Rf_*a(K) \longrightarrow K$를 때때로
{\it 대각합 사상}이라 부른다. 이 사상은 오른쪽 수반을 특징짓는\ 
$L \in D_\QCoh(\mathcal{O}_X)$에 대한 전단사
$$
\Hom_X(L, a(K)) \longrightarrow \Hom_Y(Rf_*L, K)
$$
가
$$
\varphi \longmapsto \text{Tr}_{f, K} \circ Rf_*\varphi
$$
로 주어진다는 성질을 갖는다. 사상 (\ref{duality-equation-sheafy-trace})
$$
Rf_*R\SheafHom_{\mathcal{O}_X}(L, a(K))
\longrightarrow
R\SheafHom_{\mathcal{O}_Y}(Rf_*L, K)
$$
은\  $\text{Tr}_{f, K}$와의 합성으로 얻어진다.
이 절에서 생각할 모든 대각합 사상은 이 대각합 사상의 특수한 경우이다.
몇 가지 특수한 경우를 논하기 전에 대각합 사상을 만드는 것이
기저변환과 가환함을 보이겠다.
\begin{lemma}[대각합 사상과 기저변환]
\label{duality-lemma-trace-map-and-base-change}
$f$와\  $g$가\  tor 독립인 도식 (\ref{duality-equation-base-change})가 주어졌다고 하자.
그러면 사상들\ 
$1 \star \text{Tr}_f : Lg^* \circ Rf_* \circ a \to Lg^*$와\ 
$\text{Tr}_{f'} \star 1 : Rf'_* \circ a' \circ Lg^* \to Lg^*$는
기저변환 사상들\ 
$\beta : Lg^* \circ Rf_* \to Rf'_* \circ L(g')^*$
(Cohomology, Remark \ref{cohomology-remark-base-change})와\ 
$\alpha : L(g')^* \circ a \to a' \circ Lg^*$
((\ref{duality-equation-base-change-map}))를 통해 일치한다.
더 정확히 말해 함자의 변환으로 이루어진 도식
$$
\xymatrix{
Lg^* \circ Rf_* \circ a
\ar[d]_{\beta \star 1} \ar[r]_-{1 \star \text{Tr}_f} &
Lg^* \\
Rf'_* \circ L(g')^* \circ a \ar[r]^{1 \star \alpha} &
Rf'_* \circ a' \circ Lg^* \ar[u]_{\text{Tr}_{f'} \star 1}
}
$$
은 가환한다.
\end{lemma}

\begin{proof}
이 증명에서는\  $Rf_*$ 대신\  $f_*$를, $Lg^*$ 대신\  $g^*$를 쓰고,
변환의 출발점과 도착점을 알면 어떤 항등변환을 덧붙여야 하는지
알 수 있으므로 항등변환과의\  $\star$ 곱을 생략한다.
$\beta : g^* \circ f_* \to f'_* \circ (g')^*$가 동형이고,
$\alpha$는\  $\beta$의 수반인 동형\ 
$\beta^\vee : g'_* \circ a' \to a \circ g_*$를 사용하여 정의함을 상기하자.
Lemma \ref{duality-lemma-flat-precompose-pus}와 그 증명을 보라.
먼저 보조정리의 도식에서 위쪽 수평 화살표는 합성
$$
g^* \circ f_* \circ a \to
g^* \circ f_* \circ a \circ g_* \circ g^* \to
g^* \circ g_* \circ g^* \to g^*
$$
과 같다. 여기서 첫 번째 화살표는\  $(g^*, g_*)$의 단위원이고,
두 번째 화살표는\  $\text{Tr}_f$이며, 세 번째 화살표는\ 
$(g^*, g_*)$의 여단위원이다. 이는 단위원과 여단위원의 합성\ 
$g^* \to g^* \circ g_* \circ g^* \to g^*$가 항등이라는 사실의
간단한 귀결이다. 다음 도식을 생각하자.
$$
\xymatrix{
& g^* \circ f_* \circ a \ar[ld]_\beta \ar[d] \ar[r]_{\text{Tr}_f} & g^* \\
f'_* \circ (g')^* \circ a \ar[dr] &
g^* \circ f_* \circ a \circ g_* \circ g^* \ar[d]_\beta \ar[ru] &
g^* \circ f_* \circ g'_* \circ a' \circ g^* \ar[l]_{\beta^\vee} \ar[d]_\beta &
f'_* \circ a' \circ g^* \ar[lu]_{\text{Tr}_{f'}} \\
& f'_* \circ (g')^* \circ a \circ g_* \circ g^* &
f'_* \circ (g')^* \circ g'_* \circ a' \circ g^* \ar[ru] \ar[l]_{\beta^\vee}
}
$$
이 도식의 두 정사각형은\  Categories, Lemma
\ref{categories-lemma-properties-2-cat-cats}, 또는 더 간단하게\ 
Categories, Definition
\ref{categories-definition-horizontal-composition} 앞의 논의에 의해 가환한다.
위의 논의에 의해 삼각형도 가환한다. Categories, Lemma
\ref{categories-lemma-transformation-between-functors-and-adjoints}에 의해
정사각형
$$
\xymatrix{
g^* \circ f_* \circ g'_* \circ a' \ar[d]_{\beta^\vee} \ar[r]_-\beta &
f'_* \circ (g')^* \circ g'_* \circ a' \ar[d] \\
g^* \circ f_* \circ a \circ g_* \ar[r] &
\text{id}
}
$$
은 가환하며, 따라서 큰 도식의 오각형도 가환한다.
$\beta$와\  $\beta^\vee$가 동형이고, 큰 도식의 바깥쪽을 따라가는 것이
정의에 의해\  $\text{Tr}_f \circ \alpha \circ \beta$와 같으므로
보조정리가 증명된다.
\end{proof}
\noindent
$f : X \to Y$를 준콤팩트 준분리 스킴의 사상이라 하자.
$a : D_\QCoh(\mathcal{O}_Y) \to D_\QCoh(\mathcal{O}_X)$를\ 
Lemma \ref{duality-lemma-twisted-inverse-image}에서와 같이\  $Rf_*$의 오른쪽
수반이라 하자. Categories, Section \ref{categories-section-adjoint}에 의해
함자의 변환
$$
\eta_f : \text{id} \to  a \circ Rf_*
$$
을 얻으며, 이를 수반의 단위원이라 부른다.
\begin{lemma}
\label{duality-lemma-unit-and-base-change}
$f$와\  $g$가\  tor 독립인 도식 (\ref{duality-equation-base-change})가 주어졌다고 하자.
그러면 사상들\ 
$1 \star \eta_f : L(g')^* \to L(g')^* \circ a \circ Rf_*$와\ 
$\eta_{f'} \star 1 : L(g')^* \to a' \circ Rf'_* \circ L(g')^*$는
기저변환 사상들\ 
$\beta : Lg^* \circ Rf_* \to Rf'_* \circ L(g')^*$
(Cohomology, Remark \ref{cohomology-remark-base-change})와\ 
$\alpha : L(g')^* \circ a \to a' \circ Lg^*$
((\ref{duality-equation-base-change-map}))를 통해 일치한다.
더 정확히 말해 함자의 변환으로 이루어진 도식
$$
\xymatrix{
L(g')^* \ar[r]_-{1 \star \eta_f} \ar[d]_{\eta_{f'} \star 1} &
L(g')^* \circ a \circ Rf_* \ar[d]^\alpha \\
a' \circ Rf'_* \circ L(g')^* &
a' \circ Lg^* \circ Rf_* \ar[l]_-\beta
}
$$
은 가환한다.
\end{lemma}

\begin{proof}
이 증명은\  Lemma \ref{duality-lemma-trace-map-and-base-change}의 증명에 쌍대이다.
이 증명에서는\  $Rf_*$ 대신\  $f_*$를, $Lg^*$ 대신\  $g^*$를 쓰고,
변환의 출발점과 도착점을 알면 어떤 항등변환을 덧붙여야 하는지
알 수 있으므로 항등변환과의\  $\star$ 곱을 생략한다.
$\beta : g^* \circ f_* \to f'_* \circ (g')^*$가 동형이고,
$\alpha$는\  $\beta$의 수반인 동형\ 
$\beta^\vee : g'_* \circ a' \to a \circ g_*$를 사용하여 정의함을 상기하자.
Lemma \ref{duality-lemma-flat-precompose-pus}와 그 증명을 보라.
먼저 보조정리의 도식에서 왼쪽 수직 화살표는 합성
$$
(g')^* \to (g')^* \circ g'_* \circ (g')^* \to
(g')^* \circ g'_* \circ a' \circ f'_* \circ (g')^* \to
a' \circ f'_* \circ (g')^*
$$
과 같다. 여기서 첫 번째 화살표는\  $((g')^*, g'_*)$의 단위원이고,
두 번째 화살표는\  $\eta_{f'}$이며, 세 번째 화살표는\ 
$((g')^*, g'_*)$의 여단위원이다. 이는 단위원과 여단위원의 합성\ 
$(g')^* \to (g')^* \circ (g')_* \circ (g')^* \to (g')^*$가 항등이라는
사실의 간단한 귀결이다. 다음 도식을 생각하자.
$$
\xymatrix{
& (g')^* \circ a \circ f_* \ar[r] &
(g')^* \circ a \circ g_* \circ g^* \circ f_*
\ar[ld]_\beta \\
(g')^* \ar[ru]^{\eta_f} \ar[dd]_{\eta_{f'}} \ar[rd] &
(g')^* \circ a \circ g_* \circ f'_* \circ (g')^* &
(g')^* \circ g'_* \circ a' \circ g^* \circ f_*
\ar[u]_{\beta^\vee} \ar[ld]_\beta \ar[d] \\
& (g')^* \circ g'_* \circ a' \circ f'_* \circ (g')^*
\ar[ld] \ar[u]_{\beta^\vee} &
a' \circ g^* \circ f_* \ar[lld]^\beta \\
a' \circ f'_* \circ (g')^*
}
$$
이 도식의 두 정사각형은\  Categories, Lemma
\ref{categories-lemma-properties-2-cat-cats}, 또는 더 간단하게\ 
Categories, Definition
\ref{categories-definition-horizontal-composition} 앞의 논의에 의해 가환한다.
위의 논의에 의해 삼각형도 가환한다. Categories, Lemma
\ref{categories-lemma-transformation-between-functors-and-adjoints}의
쌍대에 의해 정사각형
$$
\xymatrix{
\text{id} \ar[r] \ar[d] &
g'_* \circ a' \circ g^* \circ f_* \ar[d]^\beta \\
g'_* \circ a' \circ g^* \circ f_* \ar[r]^{\beta^\vee} &
a \circ g_* \circ f'_* \circ (g')^*
}
$$
은 가환하며, 따라서 큰 도식의 오각형도 가환한다.
$\beta$와\  $\beta^\vee$가 동형이고, 큰 도식의 바깥쪽을 따라가는 것이
정의에 의해\  $\beta \circ \alpha \circ \eta_f$와 같으므로
보조정리가 증명된다.
\end{proof}
\begin{example}
\label{duality-example-trace-affine}
$A \to B$를 환 사상이라 하자. $Y = \Spec(A)$, $X = \Spec(B)$라 하고,
$f : X \to Y$를\  $A \to B$에 대응하는 사상이라 하자.
Example \ref{duality-example-affine-twisted-inverse-image}에서 본 것처럼\ 
$Rf_* : D_\QCoh(\mathcal{O}_X) \to D_\QCoh(\mathcal{O}_Y)$의 오른쪽 수반은\ 
$D(A) = D_\QCoh(\mathcal{O}_Y)$의 대상\  $K$를\ 
$D(B) = D_\QCoh(\mathcal{O}_X)$의\  $R\Hom(B, K)$로 보낸다.
대각합 사상은\  $A$-가군 사상\  $A \to B$가 유도하는 사상
$$
\text{Tr}_{f, K} : R\Hom(B, K) \longrightarrow R\Hom(A, K) = K
$$
이다.
\end{example}



\section{추이 함자의 오른쪽 수반과 당김}
\label{duality-section-compare-with-pullback}

\noindent
$f : X \to Y$를 준콤팩트 준분리 스킴의 사상이라 하고,
$a$를\  Lemma \ref{duality-lemma-twisted-inverse-image}의 추이 함자의 오른쪽
수반이라 하자. $K, L \in D_\QCoh(\mathcal{O}_Y)$에 대해 표준 사상
$$
Lf^*K \otimes^\mathbf{L}_{\mathcal{O}_X} a(L)
\longrightarrow
a(K \otimes_{\mathcal{O}_Y}^\mathbf{L} L)
$$
이 있다. 구체적으로 이 사상은 사상
$$
Rf_*(Lf^*K \otimes^\mathbf{L}_{\mathcal{O}_X} a(L)) =
K \otimes^\mathbf{L}_{\mathcal{O}_Y} Rf_*(a(L))
\longrightarrow
K \otimes^\mathbf{L}_{\mathcal{O}_Y} L
$$
에 수반된다(등식은\  Derived Categories of Schemes, Lemma
\ref{perfect-lemma-cohomology-base-change}에 의한다). 여기서는 대각합 사상\ 
$Rf_*a(L) \to L$을 사용한다. $L = \mathcal{O}_Y$이면
사상
\begin{equation}
\label{duality-equation-compare-with-pullback}
Lf^*K \otimes^\mathbf{L}_{\mathcal{O}_X} a(\mathcal{O}_Y) \longrightarrow a(K)
\end{equation}
을 얻으며, 이는\  $K$에 대해 함자적이고 완전 삼각형과 양립한다.
\begin{lemma}
\label{duality-lemma-compare-with-pullback-perfect}
$f : X \to Y$를 준콤팩트 준분리 스킴의 사상이라 하자.
위에서\  $K, L \in D_\QCoh(\mathcal{O}_Y)$에 대해 정의한 사상\ 
$Lf^*K \otimes^\mathbf{L}_{\mathcal{O}_X} a(L) \to
a(K \otimes_{\mathcal{O}_Y}^\mathbf{L} L)$은\  $K$가 완전이면 동형이다.
특히\  $K$가 완전이면 (\ref{duality-equation-compare-with-pullback})은 동형이다.
\end{lemma}

\begin{proof}
$K^\vee$를\  $K$의 ``쌍대''라 하자. Cohomology, Lemma
\ref{cohomology-lemma-dual-perfect-complex}를 보라.
$M \in D_\QCoh(\mathcal{O}_X)$에 대해
\begin{align*}
\Hom_{D(\mathcal{O}_Y)}(Rf_*M, K \otimes^\mathbf{L}_{\mathcal{O}_Y} L)
& =
\Hom_{D(\mathcal{O}_Y)}(
Rf_*M \otimes^\mathbf{L}_{\mathcal{O}_Y} K^\vee, L) \\
& =
\Hom_{D(\mathcal{O}_X)}(
M \otimes^\mathbf{L}_{\mathcal{O}_X} Lf^*K^\vee, a(L)) \\
& =
\Hom_{D(\mathcal{O}_X)}(M,
Lf^*K \otimes^\mathbf{L}_{\mathcal{O}_X} a(L))
\end{align*}
이다. 두 번째 등식은\  $a$의 정의와 사영 공식
(Cohomology, Lemma \ref{cohomology-lemma-projection-formula-perfect}),
또는 더 일반적으로\  Derived Categories of Schemes, Lemma
\ref{perfect-lemma-cohomology-base-change}에 의한다.
따라서\  Yoneda 보조정리에 의해 결과를 얻는다.
\end{proof}
\begin{lemma}
\label{duality-lemma-restriction-compare-with-pullback}
$f$와\  $g$가\  tor 독립인 도식 (\ref{duality-equation-base-change})가 주어졌다고 하자.
$K \in D_\QCoh(\mathcal{O}_Y)$라 하자. 도식
$$
\xymatrix{
L(g')^*(Lf^*K \otimes^\mathbf{L}_{\mathcal{O}_X} a(\mathcal{O}_Y))
\ar[r] \ar[d] & L(g')^*a(K) \ar[d] \\
L(f')^*Lg^*K \otimes_{\mathcal{O}_{X'}}^\mathbf{L} a'(\mathcal{O}_{Y'})
\ar[r] & a'(Lg^*K)
}
$$
은 가환한다. 여기서 수평 화살표들은\  $K$와\  $Lg^*K$에 대한 사상
(\ref{duality-equation-compare-with-pullback})이고, 수직 사상들은\ 
Cohomology, Remark \ref{cohomology-remark-base-change}와
(\ref{duality-equation-base-change-map})을 사용하여 만든다.
\end{lemma}

\begin{proof}
이 증명에서는\  $Rf_*$ 대신\  $f_*$를, $Lf^*$ 대신\  $f^*$ 등을 쓰고,
$\otimes^\mathbf{L}_{\mathcal{O}_X}$ 대신\  $\otimes$ 등을 쓰겠다.
(\ref{duality-equation-compare-with-pullback})을 합성
\begin{align*}
f^*K \otimes a(\mathcal{O}_Y)
& \to
a(f_*(f^*K \otimes a(\mathcal{O}_Y))) \\
& \leftarrow
a(K \otimes f_*a(\mathcal{O}_K)) \\
& \to
a(K \otimes \mathcal{O}_Y) \\
& \to
a(K)
\end{align*}
으로 쓰자. 여기서 첫 번째 화살표는 단위원\  $\eta_f$이고,
두 번째 화살표는\  $a$를\  Cohomology, Equation
(\ref{cohomology-equation-projection-formula-map})에 적용한 것이다.
이 화살표는\  Derived Categories of Schemes, Lemma
\ref{perfect-lemma-cohomology-base-change}에 의해 동형이다.
세 번째 화살표는\  $a$를\  $\text{id}_K \otimes \text{Tr}_f$에 적용한 것이고,
네 번째 화살표는\  $a$를 동형\  $K \otimes \mathcal{O}_Y = K$에
적용한 것이다. 보조정리를 증명하려면 이 각각의 사상이 보조정리의
명제에 있는 것과 같은 가환 정사각형을 준다는 것을 보이면 된다.
$\eta_f$와\  $\text{Tr}_f$에 대해서는\  Lemmas
\ref{duality-lemma-unit-and-base-change}와\ 
\ref{duality-lemma-trace-map-and-base-change}가 이를 보인다.
Cohomology, Equation
(\ref{cohomology-equation-projection-formula-map})을 사용하는 화살표에 대해서는\ 
Cohomology, Remark \ref{cohomology-remark-compatible-with-diagram}가 이를 보인다.
곱셈 사상에 대해서는 명백하다. 이로써 증명이 끝난다.
\end{proof}
\begin{lemma}
\label{duality-lemma-compare-on-open}
$f : X \to Y$를\  Noether 스킴의 고유 사상이라 하자.
$V \subset Y$를\  $f^{-1}(V) \to V$가 동형인 열린집합이라 하자.
그러면\  $K \in D_\QCoh^+(\mathcal{O}_Y)$에 대해 사상
(\ref{duality-equation-compare-with-pullback})은\  $f^{-1}(V)$ 위에서 동형으로 제한된다.
\end{lemma}

\begin{proof}
Lemma \ref{duality-lemma-proper-noetherian}에 의해 사상
(\ref{duality-equation-sheafy})는\  $D_\QCoh^+(\mathcal{O}_Y)$의 대상들에 대해
동형이다. 따라서\  Lemma \ref{duality-lemma-restriction-compare-with-pullback}에 의해\ 
$K$에 대한 (\ref{duality-equation-compare-with-pullback})을\  $f^{-1}(V)$에 제한한 것은\ 
$K|_V$와\  $f^{-1}(V) \to V$에 대한 사상
(\ref{duality-equation-compare-with-pullback})이다. 그러므로\  $f$가 항등 사상인 경우에
그 사상이 동형임을 보이면 충분하다. 이는 명백하다.
\end{proof}
\begin{lemma}
\label{duality-lemma-transitivity-compare-with-pullback}
$f : X \to Y$와\  $g : Y \to Z$를 준콤팩트 준분리 스킴의 합성 가능한
사상이라 하고\  $h = g \circ f$로 놓자. $a, b, c$를 각각\  $f, g, h$에 대한\ 
Lemma \ref{duality-lemma-twisted-inverse-image}의 수반이라 하자.
임의의\  $K \in D_\QCoh(\mathcal{O}_Z)$에 대해 도식
$$
\xymatrix{
Lf^*(Lg^*K \otimes_{\mathcal{O}_Y}^\mathbf{L}
b(\mathcal{O}_Z)) \otimes_{\mathcal{O}_X}^\mathbf{L} a(\mathcal{O}_Y)
\ar@{=}[d] \ar[r] &
a(Lg^*K \otimes_{\mathcal{O}_Y}^\mathbf{L} b(\mathcal{O}_Z)) \ar[r] &
a(b(K)) \ar@{=}[d] \\
Lh^*K \otimes_{\mathcal{O}_X}^\mathbf{L} Lf^*b(\mathcal{O}_Z)
\otimes_{\mathcal{O}_X}^\mathbf{L} a(\mathcal{O}_Y) \ar[r] &
Lh^*K \otimes_{\mathcal{O}_X}^\mathbf{L} c(\mathcal{O}_Z) \ar[r] &
c(K)
}
$$
은 가환한다. 여기서 화살표들은 (\ref{duality-equation-compare-with-pullback})이고,
$Lh^* = Lf^* \circ Lg^*$와\  $c = a \circ b$를 사용하였다.
\end{lemma}

\begin{proof}
이 증명에서는\  $Rf_*$ 대신\  $f_*$를, $Lf^*$ 대신\  $f^*$ 등을 쓰고,
$\otimes^\mathbf{L}_{\mathcal{O}_X}$ 대신\  $\otimes$ 등을 쓰겠다.
위쪽 화살표들의 합성은 사상
$$
g_*f_*(f^*(g^*K \otimes b(\mathcal{O}_Z)) \otimes a(\mathcal{O}_Y)) \to K
$$
에 수반된다. 왼쪽 항은\  Derived Categories of Schemes, Lemma
\ref{perfect-lemma-cohomology-base-change}에 의해\ 
$K \otimes g_*f_*(f^*b(\mathcal{O}_Z) \otimes a(\mathcal{O}_Y))$와 같고,
정의들을 살펴보면 이 사상은
$$
g_*f_*(f^*b(\mathcal{O}_Z) \otimes a(\mathcal{O}_Y))
\xleftarrow{g_*\epsilon}
g_*(b(\mathcal{O}_Z) \otimes f_*a(\mathcal{O}_Y)) \xrightarrow{g_*\alpha}
g_*(b(\mathcal{O}_Z)) \xrightarrow{\beta} \mathcal{O}_Z
$$
를\  $\text{id}_K$와 텐서하여 얻는다. 여기서\  $\epsilon$은\  Derived Categories
of Schemes, Lemma \ref{perfect-lemma-cohomology-base-change}의 동형이고,
$\beta$는 여단위원 사상\  $g_*b \to \text{id}$에서 온다. 마찬가지로 아래쪽
수평 화살표들의 합성은\  $\text{id}_K$와 다음 합성을 텐서한 사상에 수반된다.
$$
g_*f_*(f^*b(\mathcal{O}_Z) \otimes a(\mathcal{O}_Y)) \xrightarrow{g_*f_*\delta}
g_*f_*(ab(\mathcal{O}_Z)) \xrightarrow{g_*\gamma}
g_*(b(\mathcal{O}_Z)) \xrightarrow{\beta}
\mathcal{O}_Z
$$
여기서\  $\gamma$는 여단위원 사상\  $f_*a \to \text{id}$에서 오고,
$\delta$는 다음 합성을 수반으로 갖는 사상이다.
$$
f_*(f^*b(\mathcal{O}_Z) \otimes a(\mathcal{O}_Y))
\xleftarrow{\epsilon}
b(\mathcal{O}_Z) \otimes f_*a(\mathcal{O}_Y) \xrightarrow{\alpha}
b(\mathcal{O}_Z)
$$
수반 함자, 수반 사상 및 여단위원의 일반 성질
(Categories, Section \ref{categories-section-adjoint} 참조)에 의해 원하는 대로\ 
$\gamma \circ f_*\delta = \alpha \circ \epsilon^{-1}$이다.
\end{proof}




\section{닫힌 몰입에 대한 추이 함자의 오른쪽 수반}
\label{duality-section-sections-with-exact-support}

\noindent
$i : (Z, \mathcal{O}_Z) \to (X, \mathcal{O}_X)$를 환 달린 공간의 사상이라 하자.
$i$는 닫힌 부분집합 위로의 위상동형이고\ 
$i^\sharp : \mathcal{O}_X \to i_*\mathcal{O}_Z$는 전사라고 가정한다.
(예를 들어 스킴의 닫힌 몰입이다.) $\mathcal{I} = \Ker(i^\sharp)$로 놓자.
$\mathcal{O}_X$-가군의 층\  $\mathcal{F}$에 대해 층
$$
\SheafHom_{\mathcal{O}_X}(i_*\mathcal{O}_Z, \mathcal{F})
$$
은\  $\mathcal{I}$에 의해 소멸되는\  $\mathcal{O}_X$-가군의 층이다.
따라서\  Modules, Lemma \ref{modules-lemma-i-star-equivalence}에 의해\ 
$\mathcal{O}_Z$-가군의 층이 존재하며, 이를\ 
$\SheafHom(\mathcal{O}_Z, \mathcal{F})$로 나타내겠다. 이때
$$
i_*\SheafHom(\mathcal{O}_Z, \mathcal{F}) =
\SheafHom_{\mathcal{O}_X}(i_*\mathcal{O}_Z, \mathcal{F})
$$
가\  $\mathcal{O}_X$-가군으로서 성립한다. 이 말의 뜻을 자세히 서술하겠다.
\begin{lemma}
\label{duality-lemma-compute-sheaf-with-exact-support}
위와 같은 표기법을 사용하자. 함자\  $\SheafHom(\mathcal{O}_Z, -)$는 함자\ 
$i_* : \textit{Mod}(\mathcal{O}_Z) \to \textit{Mod}(\mathcal{O}_X)$의
오른쪽 수반이다. 열린집합\  $V \subset Z$에 대해
$$
\Gamma(V, \SheafHom(\mathcal{O}_Z, \mathcal{F})) =
\{s \in \Gamma(U, \mathcal{F}) \mid \mathcal{I}s = 0\}
$$
이다. 여기서\  $U \subset X$는\  $Z$와의 교집합이\  $V$인 열린집합이다.
\end{lemma}

\begin{proof}
$\mathcal{G}$를\  $\mathcal{O}_Z$-가군의 층이라 하자. 그러면
$$
\Hom_{\mathcal{O}_X}(i_*\mathcal{G}, \mathcal{F}) =
\Hom_{i_*\mathcal{O}_Z}(i_*\mathcal{G},
\SheafHom_{\mathcal{O}_X}(i_*\mathcal{O}_Z, \mathcal{F})) =
\Hom_{\mathcal{O}_Z}(\mathcal{G}, \SheafHom(\mathcal{O}_Z, \mathcal{F}))
$$
이다. 첫 번째 등식은\  Modules, Lemma
\ref{modules-lemma-adjoint-tensor-restrict}에 의하고, 두 번째 등식은\ 
$i_*$의 충실충만성에 의한다. Modules, Lemma
\ref{modules-lemma-i-star-equivalence}를 보라. 단면에 관한 서술은 독자에게 맡긴다.
\end{proof}
\noindent
함자
$$
\textit{Mod}(\mathcal{O}_X)
\longrightarrow
\textit{Mod}(\mathcal{O}_Z),
\quad
\mathcal{F} \longmapsto \SheafHom(\mathcal{O}_Z, \mathcal{F})
$$
는 왼쪽 완전이고 다음과 같은 유도 확장을 갖는다.
$$
R\SheafHom(\mathcal{O}_Z, -) : D(\mathcal{O}_X) \to D(\mathcal{O}_Z).
$$

\begin{lemma}
\label{duality-lemma-sheaf-with-exact-support-adjoint}
위와 같은 표기법을 사용하자. 함자\  $R\SheafHom(\mathcal{O}_Z, -)$는 함자\ 
$Ri_* : D(\mathcal{O}_Z) \to D(\mathcal{O}_X)$의 오른쪽 수반이다.
\end{lemma}

\begin{proof}
이는\  Lemma \ref{duality-lemma-compute-sheaf-with-exact-support}에 의해\  $i_*$와\ 
$\SheafHom(\mathcal{O}_Z, -)$가 수반 함자라는 사실의 귀결이다.
Derived Categories, Lemma \ref{derived-lemma-derived-adjoint-functors}를 보라.
\end{proof}
\begin{lemma}
\label{duality-lemma-sheaf-with-exact-support-ext}
위와 같은 표기법을 사용하자. $D(\mathcal{O}_X)$의 모든\  $K$에 대해
$$
Ri_*R\SheafHom(\mathcal{O}_Z, K) =
R\SheafHom_{\mathcal{O}_X}(i_*\mathcal{O}_Z, K)
$$
가\  $D(\mathcal{O}_X)$에서 성립한다.
\end{lemma}

\begin{proof}
이는 함자\  $R\SheafHom(\mathcal{O}_Z, -)$의 구성에서 즉시 따른다.
\end{proof}
\begin{lemma}
\label{duality-lemma-sheaf-with-exact-support-internal-home}
위와 같은 표기법을 사용하자. $M \in D(\mathcal{O}_Z)$에 대해
$$
R\SheafHom_{\mathcal{O}_X}(Ri_*M, K) =
Ri_*R\SheafHom_{\mathcal{O}_Z}(M, R\SheafHom(\mathcal{O}_Z, K))
$$
가\  $D(\mathcal{O}_Z)$에서 모든\  $D(\mathcal{O}_X)$의\  $K$에 대해 성립한다.
\end{lemma}

\begin{proof}
이는 함자\  $R\SheafHom(\mathcal{O}_Z, -)$의 구성과 다음 사실에서 즉시 따른다.
$\mathcal{K}^\bullet$가\  $\mathcal{O}_X$-가군의\  K-단사 복합체이면\ 
$\SheafHom(\mathcal{O}_Z, \mathcal{K}^\bullet)$는\  $\mathcal{O}_Z$-가군의\ 
K-단사 복합체이다. Derived Categories, Lemma
\ref{derived-lemma-adjoint-preserve-K-injectives}를 보라.
\end{proof}
\begin{lemma}
\label{duality-lemma-sheaf-with-exact-support-quasi-coherent}
$i : Z \to X$를 스킴의 유사코히런트 닫힌 몰입이라 하자
($X$가 국소\  Noether이면 임의의 닫힌 몰입). 그러면
\begin{enumerate}
\item $R\SheafHom(\mathcal{O}_Z, -)$는\  $D^+_\QCoh(\mathcal{O}_X)$를\ 
$D^+_\QCoh(\mathcal{O}_Z)$ 안으로 보내고,
\item $X = \Spec(A)$이고\  $Z = \Spec(B)$이면 도식
$$
\xymatrix{
D^+(B) \ar[r] & D_\QCoh^+(\mathcal{O}_Z) \\
D^+(A) \ar[r] \ar[u]^{R\Hom(B, -)} &
D_\QCoh^+(\mathcal{O}_X) \ar[u]_{R\SheafHom(\mathcal{O}_Z, -)}
}
$$
은 가환한다.
\end{enumerate}
\end{lemma}

\begin{proof}
괄호 안의 주장을 설명하자. $X$가 국소\  Noether이면\  More on Morphisms,
Lemma \ref{more-morphisms-lemma-Noetherian-pseudo-coherent}에 의해\ 
$i$는 유사코히런트이다.

\medskip\noindent
$K$를\  $D^+_\QCoh(\mathcal{O}_X)$의 대상이라 하자. (1)을 증명하기 위해\ 
Morphisms, Lemma \ref{morphisms-lemma-i-star-equivalence}에 의해\ 
$i_*$를\  $H^n(R\SheafHom(\mathcal{O}_Z, K))$에 적용하면\  $X$ 위의
준연접 가군이 된다는 것을 보이면 충분하다. Lemma
\ref{duality-lemma-sheaf-with-exact-support-ext}에 의해 이는\ 
$R\SheafHom_{\mathcal{O}_X}(i_*\mathcal{O}_Z, K)$가\ 
$D_\QCoh(\mathcal{O}_X)$에 속함을 보여야 한다는 뜻이다. $i$가
유사코히런트이므로 층\  $\mathcal{O}_Z$는 유사코히런트\ 
$\mathcal{O}_X$-가군이다. 따라서 결과는\  Derived Categories of Schemes,
Lemma \ref{perfect-lemma-quasi-coherence-internal-hom}에서 따른다.

\medskip\noindent
(2)에서와 같이\  $X = \Spec(A)$이고\  $Z = \Spec(B)$라고 가정하자.
$I^\bullet$을\  $D^+(A)$의 대상\  $K$를 나타내는 아래로 유계인 단사\ 
$A$-가군 복합체라 하자. 그러면\  $R\Hom(B, K) = \Hom_A(B, I^\bullet)$이며,
이를\  $B$-가군의 복합체로 본다. 준동형
$$
\widetilde{I^\bullet} \longrightarrow \mathcal{I}^\bullet
$$
을 택하자. 여기서\  $\mathcal{I}^\bullet$은 아래로 유계인 단사\ 
$\mathcal{O}_X$-가군 복합체이다. Lemma
\ref{duality-lemma-compute-sheaf-with-exact-support}에 있는 함자\ 
$\SheafHom(\mathcal{O}_Z, -)$의 서술로부터 사상
$$
\Hom_A(B, I^\bullet)
\longrightarrow
\Gamma(Z, \SheafHom(\mathcal{O}_Z, \mathcal{I}^\bullet))
$$
이 존재한다. $\SheafHom(\mathcal{O}_Z, \mathcal{I}^\bullet)$가\ 
$R\SheafHom(\mathcal{O}_Z, \widetilde{K})$를 나타냄을 관찰하라.
$\widetilde{\ }$ 함자의 보편성을 적용하면\  $D(\mathcal{O}_Z)$의 사상
$$
\widetilde{\Hom_A(B, I^\bullet)}
\longrightarrow
R\SheafHom(\mathcal{O}_Z, \widetilde{K})
$$
을 얻는다. $i_*$를 적용한 뒤 이 사상이\  $D(\mathcal{O}_Z)$에서
동형인지 확인해도 된다. 그런데\  $i_*$를 적용하면\  Derived Categories
of Schemes, Lemma \ref{perfect-lemma-quasi-coherence-internal-hom}의 동형을\ 
Lemma \ref{duality-lemma-sheaf-with-exact-support-ext}의 동일시를 통해 얻는다.
\end{proof}
\begin{lemma}
\label{duality-lemma-sheaf-with-exact-support-coherent}
$i : Z \to X$를 스킴의 닫힌 몰입이라 하자. $X$가 국소\  Noether이라고
가정하자. 그러면\  $R\SheafHom(\mathcal{O}_Z, -)$는\ 
$D^+_{\textit{Coh}}(\mathcal{O}_X)$를\ 
$D^+_{\textit{Coh}}(\mathcal{O}_Z)$ 안으로 보낸다.
\end{lemma}

\begin{proof}
문제는\  $X$ 위에서 국소적이므로\  $X$가 아핀이라고 가정해도 된다.
$A$가\  Noether이고\  $A \to B$가 전사인\ 
$X = \Spec(A)$ 및\  $Z = \Spec(B)$로 쓰자. 이 경우\  Lemma
\ref{duality-lemma-sheaf-with-exact-support-quasi-coherent}를 적용하여 문제를 대수로
옮길 수 있다. 대응하는 대수적 결과는\  Dualizing Complexes, Lemma
\ref{dualizing-lemma-exact-support-coherent}의 귀결이다.
\end{proof}
\begin{lemma}
\label{duality-lemma-twisted-inverse-image-closed}
$X$를 준콤팩트 준분리 스킴이라 하자. $i : Z \to X$를 유사코히런트
닫힌 몰입이라 하자($X$가\  Noether이면 임의의 닫힌 몰입이 유사코히런트이다).
$a : D_\QCoh(\mathcal{O}_X) \to D_\QCoh(\mathcal{O}_Z)$를\  $Ri_*$의 오른쪽
수반이라 하자. 그러면\  $K \in D_\QCoh^+(\mathcal{O}_X)$에 대해 함자적 동형
$$
a(K) = R\SheafHom(\mathcal{O}_Z, K)
$$
이 존재한다.
\end{lemma}

\begin{proof}
(괄호 안의 주장은\  More on Morphisms, Lemma
\ref{more-morphisms-lemma-Noetherian-pseudo-coherent}에서 따른다.)
Lemma \ref{duality-lemma-sheaf-with-exact-support-adjoint}에 의해 함자\ 
$R\SheafHom(\mathcal{O}_Z, -)$는\ 
$Ri_* : D(\mathcal{O}_Z) \to D(\mathcal{O}_X)$의 오른쪽 수반이다.
더욱이\  Lemma \ref{duality-lemma-sheaf-with-exact-support-quasi-coherent}와\  Lemma
\ref{duality-lemma-twisted-inverse-image-bounded-below}에 의해\ 
$R\SheafHom(\mathcal{O}_Z, -)$와\  $a$는 모두\ 
$D_\QCoh^+(\mathcal{O}_X)$를\  $D_\QCoh^+(\mathcal{O}_Z)$ 안으로 보낸다.
따라서 수반 함자의 유일성에 의해 동형을 얻는다.
\end{proof}
\begin{example}
\label{duality-example-trace-closed-immersion}
$i : Z \to X$가\  Noether 스킴의 닫힌 몰입이면 도식
$$
\xymatrix{
i_*a(K) \ar[rr]_-{\text{Tr}_{i, K}} \ar@{=}[d] & &
K \ar@{=}[d] \\
i_*R\SheafHom(\mathcal{O}_Z, K) \ar@{=}[r] &
R\SheafHom_{\mathcal{O}_X}(i_*\mathcal{O}_Z, K)
\ar[r] & K
}
$$
은\  $K \in D_\QCoh^+(\mathcal{O}_X)$에 대해 가환한다. 여기서 수평 등호는\ 
Lemma \ref{duality-lemma-sheaf-with-exact-support-ext}이고, 아래쪽 수평 화살표는
사상\  $\mathcal{O}_X \to i_*\mathcal{O}_Z$에 의해 유도된다. 도식의 가환성은\ 
Lemma \ref{duality-lemma-twisted-inverse-image-closed}의 귀결이다.
\end{example}



\section{닫힌 몰입에 대한 추이 함자의 오른쪽 수반과 기저변환}
\label{duality-section-sections-with-exact-support-base-change}

\noindent
다음 스킴들의\  Cartesian 도식을 생각하자.
$$
\xymatrix{
Z' \ar[r]_{i'} \ar[d]_g & X' \ar[d]^f \\
Z \ar[r]^i & X
}
$$
여기서\  $i$는 닫힌 몰입이다. $Z$와\  $X'$가\  $X$ 위에서\  tor 독립이면\ 
$D(\mathcal{O}_{Z'})$에는\  $D(\mathcal{O}_X)$의\  $K$에 대해 함자적인 정준
기저변환 사상
\begin{equation}
\label{duality-equation-base-change-exact-support}
Lg^*R\SheafHom(\mathcal{O}_Z, K)
\longrightarrow
R\SheafHom(\mathcal{O}_{Z'}, Lf^*K)
\end{equation}
이 있다. 실제로\  Lemma \ref{duality-lemma-sheaf-with-exact-support-adjoint}의 수반성에
의해 이러한 화살표를 주는 것은\  $D(\mathcal{O}_{X'})$의 사상
$$
Ri'_*Lg^*R\SheafHom(\mathcal{O}_Z, K)
\longrightarrow
Lf^*K
$$
을 주는 것과 같다. tor 독립성에 의해\ 
$Ri'_* \circ Lg^* = Lf^* \circ Ri_*$이다(Derived Categories of Schemes,
Lemma \ref{perfect-lemma-compare-base-change-closed-immersion}를 보라).
따라서 이는 사상
$$
Lf^*Ri_*R\SheafHom(\mathcal{O}_Z, K)
\longrightarrow
Lf^*K
$$
을 주는 것과 같다. 여기에는\  $can : Ri_* R\SheafHom(\mathcal{O}_Z, K) \to K$가
수반의 여단위원일 때\  $Lf^*(can)$을 사용할 수 있다.
\begin{lemma}
\label{duality-lemma-check-base-change-is-iso}
위 상황에서 사상 (\ref{duality-equation-base-change-exact-support})가 동형인 것과\ 
Cohomology, Remark \ref{cohomology-remark-prepare-fancy-base-change}의
기저변환 사상
$$
Lf^*R\SheafHom_{\mathcal{O}_X}(\mathcal{O}_Z, K)
\longrightarrow
R\SheafHom_{\mathcal{O}_{X'}}(\mathcal{O}_{Z'}, Lf^*K)
$$
이 동형인 것은 동치이다.
\end{lemma}

\begin{proof}
가정된\  tor 독립성에 의해\  $\mathcal{O}_{Z'} = Lf^*\mathcal{O}_Z$이므로 이 명제는
의미가 있다. $i'_*$가 정확하고 충실하므로 사상
(\ref{duality-equation-base-change-exact-support})에\  $Ri'_*$를 적용한 뒤 동형임을 보이면
충분하다. 가정된\  tor 독립성과\  Derived Categories of Schemes, Lemma
\ref{perfect-lemma-compare-base-change-closed-immersion}에 의해\ 
$Ri'_* \circ Lg^* = Lf^* \circ Ri_*$이므로 사상
$$
Lf^*Ri_*R\SheafHom(\mathcal{O}_Z, K)
\longrightarrow
Ri'_*R\SheafHom(\mathcal{O}_{Z'}, Lf^*K)
$$
을 얻는다. Lemma \ref{duality-lemma-sheaf-with-exact-support-ext}에 의해 그 시작과 끝은
보조정리의 명제에 나타난 것과 같다. 이 사상이\  Cohomology, Remark
\ref{cohomology-remark-prepare-fancy-base-change}에서 구성한 사상과 같다는
확인은 생략한다.
\end{proof}
\begin{lemma}
\label{duality-lemma-flat-bc-sheaf-with-exact-support}
위 상황에서\  $f$가 평탄하고\  $i$가 유사코히런트라고 가정하자. 그러면
(\ref{duality-equation-base-change-exact-support})는\ 
$D^+_\QCoh(\mathcal{O}_X)$의\  $K$에 대해 동형이다.
\end{lemma}

\begin{proof}
첫 번째 증명. 이 사상이 동형임을 증명하기 위해 국소적으로 작업해도 된다.
따라서\  $X$, $X'$, $Z$, $Z'$가 각각 환\  $A$, $A'$, $B$, $B'$에 대응하는
아핀 스킴이라고 가정해도 된다. 그러면\  $B$와\  $A'$는\  $A$ 위에서\  tor 독립이다.
Lemma \ref{duality-lemma-check-base-change-is-iso}에 의해 모든\  $K \in D^+(A)$에 대해\ 
$D(A')$에서
$$
R\Hom_A(B, K) \otimes_A^\mathbf{L} A' =
R\Hom_{A'}(B', K \otimes_A^\mathbf{L} A')
$$
임을 확인하면 충분하다. 여기서 환과 스킴의 수준에서 유도\  Hom을 비교하기 위해\ 
Derived Categories of Schemes, Lemma
\ref{perfect-lemma-quasi-coherence-internal-hom}과\  $B$ 및 각각\  $B'$가
각각\  $A$-가군 및\  $A'$-가군으로서 유사코히런트라는 사실을 사용한다.
표시된 등식은\  More on Algebra, Lemma
\ref{more-algebra-lemma-internal-hom-evaluate-tensor-isomorphism}의 (3)에서 따른다.
Dualizing Complexes, Section
\ref{dualizing-section-base-change-trivial-duality}의 논의도 보라.

\medskip\noindent
두 번째 증명\footnote{이 증명은\  $K$가\  $D^+(\mathcal{O}_X)$에 속한다고만
가정해도 충분함을 보인다.}. $z' \in Z'$이고 그 상이\  $z \in Z$라고 하자.
먼저 (\ref{duality-equation-base-change-exact-support})가\  $z'$에서의 줄기 위에서\ 
Dualizing Complexes, Equation (\ref{dualizing-equation-base-change})의 사상
$$
R\Hom(\mathcal{O}_{Z, z}, K_z)
\otimes_{\mathcal{O}_{Z, x}}^\mathbf{L} \mathcal{O}_{Z', z'}
\longrightarrow
R\Hom(\mathcal{O}_{Z', z'},
K_z \otimes_{\mathcal{O}_{X, z}}^\mathbf{L} \mathcal{O}_{X', z'})
$$
을 유도함을 먼저 보이자. 즉, 이 사상들의 구성은 동일하다.
그런 다음\  Dualizing Complexes, Lemma
\ref{dualizing-lemma-flat-bc-surjection}을 적용한다.
\end{proof}
\begin{lemma}
\label{duality-lemma-sheaf-with-exact-support-tensor}
$i : Z \to X$를 스킴의 유사코히런트 닫힌 몰입이라 하자.
$M \in D_\QCoh(\mathcal{O}_X)$이 국소적으로\  $[a, \infty)$ 안의\  tor-진폭을
가진다고 하자. $K \in D_\QCoh^+(\mathcal{O}_X)$라 하자. 그러면\ 
$D(\mathcal{O}_Z)$에 정준 동형
$$
R\SheafHom(\mathcal{O}_Z, K) \otimes_{\mathcal{O}_Z}^\mathbf{L} Li^*M =
R\SheafHom(\mathcal{O}_Z, K \otimes_{\mathcal{O}_X}^\mathbf{L} M)
$$
이 존재한다.
\end{lemma}

\begin{proof}
왼쪽 변에서 오른쪽 변으로 가는 사상을 주는 것은\  Lemma
\ref{duality-lemma-sheaf-with-exact-support-adjoint}와\ 
\ref{duality-lemma-sheaf-with-exact-support-ext}에 의해 사상
$$
Ri_*R\SheafHom(\mathcal{O}_Z, K) \otimes_{\mathcal{O}_X}^\mathbf{L} M
\longrightarrow
K \otimes_{\mathcal{O}_X}^\mathbf{L} M
$$
을 주는 것과 같다. 이 사상으로 여단위원\ 
$Ri_*R\SheafHom(\mathcal{O}_Z, K) \to K$를\  $\text{id}_M$과 텐서한 사상을
취한다. 주어진 가정 아래 이 사상이 동형임을 보이기 위해\  Lemma
\ref{duality-lemma-sheaf-with-exact-support-quasi-coherent}를 사용하여 대수로 옮긴 뒤,
예를 들어\  More on Algebra, Lemma
\ref{more-algebra-lemma-internal-hom-evaluate-tensor-isomorphism}의 (3)을 사용한다.
Lemma \ref{duality-lemma-sheaf-with-exact-support-quasi-coherent}를 사용하는 대신\  Lemma
\ref{duality-lemma-flat-bc-sheaf-with-exact-support}의 두 번째 증명에서처럼 줄기를
살펴보아도 된다.
\end{proof}


\section{유한 사상에 대한 추이 함자의 오른쪽 수반}
\label{duality-section-duality-finite}

\noindent
$i : Z \to X$가 스킴의 닫힌 몰입이면 함자\ 
$i_* : \textit{Mod}(\mathcal{O}_Z) \to \textit{Mod}(\mathcal{O}_X)$에는 오른쪽
수반\  $\SheafHom(\mathcal{O}_Z, -)$가 있으며, 그 유도 확장\ 
$R\SheafHom(\mathcal{O}_Z, -)$는\ 
$Ri_* : D(\mathcal{O}_Z) \to D(\mathcal{O}_X)$의 오른쪽 수반이다. Section
\ref{duality-section-sections-with-exact-support}를 보라. 유한 사상\ 
$f : Y \to X$의 경우에는 함자\ 
$f_* : \textit{Mod}(\mathcal{O}_Y) \to \textit{Mod}(\mathcal{O}_X)$가 일반적으로
정확하지 않고 따라서 오른쪽 수반을 갖지 않으므로 이 방법은 작동할 수 없다.
그 대신 정확한 함자\ 
$\textit{Mod}(f_*\mathcal{O}_Y) \to \textit{Mod}(\mathcal{O}_X)$를 생각하고,
이에 대응하는 오른쪽 수반과 그 유도 확장을 고찰한다.
\medskip\noindent
$f : Y \to X$를 스킴의 아핀 사상이라 하자. $\mathcal{O}_X$-가군의 층\ 
$\mathcal{F}$에 대해 층
$$
\SheafHom_{\mathcal{O}_X}(f_*\mathcal{O}_Y, \mathcal{F})
$$
은\  $f_*\mathcal{O}_Y$-가군의 층이다. 이로부터 함자\ 
$\textit{Mod}(\mathcal{O}_X) \to \textit{Mod}(f_*\mathcal{O}_Y)$를 얻으며,
이를\  $\SheafHom(f_*\mathcal{O}_Y, -)$로 나타내겠다.

\begin{lemma}
\label{duality-lemma-compute-sheafhom-affine}
위와 같은 표기법을 사용하자. 함자\  $\SheafHom(f_*\mathcal{O}_Y, -)$는 제한 함자\ 
$\textit{Mod}(f_*\mathcal{O}_Y) \to \textit{Mod}(\mathcal{O}_X)$의 오른쪽
수반이다. 아핀 열린집합\  $U \subset X$에 대해
$$
\Gamma(U, \SheafHom(f_*\mathcal{O}_Y, \mathcal{F})) =
\Hom_A(B, \mathcal{F}(U))
$$
가 성립한다. 여기서\  $A = \mathcal{O}_X(U)$이고\ 
$B = \mathcal{O}_Y(f^{-1}(U))$이다.
\end{lemma}

\begin{proof}
수반성은\  Modules, Lemma \ref{modules-lemma-adjoint-tensor-restrict}에서 따른다.
$f$가 아핀이므로\  $f_*\mathcal{O}_Y$는\  $B$를\  $A$-가군으로 보았을 때 이에
대응하는 준연접 층이다. 따라서\  $U$ 위의 절단에 대한 서술은\  Schemes, Lemma
\ref{schemes-lemma-compare-constructions}에서 따른다.
\end{proof}
\noindent
함자\  $\SheafHom(f_*\mathcal{O}_Y, -)$는 왼쪽 정확하다. 그 유도 확장을
$$
R\SheafHom(f_*\mathcal{O}_Y, -) :
D(\mathcal{O}_X)
\longrightarrow
D(f_*\mathcal{O}_Y)
$$
라 하자.

\begin{lemma}
\label{duality-lemma-sheafhom-affine-adjoint}
위와 같은 표기법을 사용하자. 함자\  $R\SheafHom(f_*\mathcal{O}_Y, -)$는 함자\ 
$D(f_*\mathcal{O}_Y) \to D(\mathcal{O}_X)$의 오른쪽 수반이다.
\end{lemma}

\begin{proof}
Lemma \ref{duality-lemma-compute-sheafhom-affine}와\  Derived Categories, Lemma
\ref{derived-lemma-derived-adjoint-functors}에서 따른다.
\end{proof}
\begin{lemma}
\label{duality-lemma-sheafhom-affine-ext}
위와 같은 표기법을 사용하자. 합성
$$
D(\mathcal{O}_X) \xrightarrow{R\SheafHom(f_*\mathcal{O}_Y, -)}
D(f_*\mathcal{O}_Y) \to D(\mathcal{O}_X)
$$
은 함자\  $K \mapsto R\SheafHom_{\mathcal{O}_X}(f_*\mathcal{O}_Y, K)$이다.
\end{lemma}

\begin{proof}
구성에서 즉시 따른다.
\end{proof}
\begin{lemma}
\label{duality-lemma-finite-twisted}
$f : Y \to X$를 스킴의 유한 유사코히런트 사상이라 하자
(Noether 스킴의 유한 사상은 유사코히런트이다). 함자\ 
$R\SheafHom(f_*\mathcal{O}_Y, -)$는\  $D_\QCoh^+(\mathcal{O}_X)$를\ 
$D_\QCoh^+(f_*\mathcal{O}_Y)$ 안으로 보낸다. $X$가 준콤팩트 준분리이면 도식
$$
\xymatrix{
D_\QCoh^+(\mathcal{O}_X) \ar[rr]_a \ar[rd]_{R\SheafHom(f_*\mathcal{O}_Y, -)}
& & D_\QCoh^+(\mathcal{O}_Y) \ar[ld]^\Phi \\
& D_\QCoh^+(f_*\mathcal{O}_Y)
}
$$
은 가환한다. 여기서\  $a$는\  $f$에 대한\  Lemma
\ref{duality-lemma-twisted-inverse-image}의 오른쪽 수반이고, $\Phi$는\  Derived Categories
of Schemes, Lemma \ref{perfect-lemma-affine-morphism-equivalence}의 동치이다.
\end{lemma}

\begin{proof}
(괄호 안의 주장은\  More on Morphisms, Lemma
\ref{more-morphisms-lemma-Noetherian-pseudo-coherent}에서 따른다.)
$f$가 유사코히런트이므로\  $\mathcal{O}_X$-가군\  $f_*\mathcal{O}_Y$는
유사코히런트이다. More on Morphisms, Lemma
\ref{more-morphisms-lemma-finite-pseudo-coherent}를 보라. 따라서\ 
$R\SheafHom(f_*\mathcal{O}_Y, -)$는\  $D_\QCoh^+(\mathcal{O}_X)$를\ 
$D_\QCoh^+(f_*\mathcal{O}_Y)$ 안으로 보낸다. Derived Categories of Schemes,
Lemma \ref{perfect-lemma-quasi-coherence-internal-hom}을 보라.
이제\  $\Phi \circ a$와\  $R\SheafHom(f_*\mathcal{O}_Y, -)$는\ 
$D_\QCoh^+(\mathcal{O}_X)$ 위에서 일치한다. 실제로 두 함자 모두 제한 함자\ 
$D_\QCoh^+(f_*\mathcal{O}_Y) \to D_\QCoh^+(\mathcal{O}_X)$의 오른쪽
수반이다. 이를 보려면\  Lemma
\ref{duality-lemma-twisted-inverse-image-bounded-below}와\ 
\ref{duality-lemma-sheafhom-affine-adjoint}를 사용한다.
\end{proof}
\begin{remark}
\label{duality-remark-trace-map-finite}
$f : Y \to X$가\  Noether 스킴의 유한 사상이면 도식
$$
\xymatrix{
Rf_*a(K) \ar[r]_-{\text{Tr}_{f, K}} \ar@{=}[d] & K \ar@{=}[d] \\
R\SheafHom_{\mathcal{O}_X}(f_*\mathcal{O}_Y, K) \ar[r] & K
}
$$
은\  $K \in D_\QCoh^+(\mathcal{O}_X)$에 대해 가환한다. 이는\  Lemma
\ref{duality-lemma-finite-twisted}에서 따른다. 아래쪽 수평 화살표는 사상\ 
$\mathcal{O}_X \to f_*\mathcal{O}_Y$에 의해 유도되고, 위쪽 수평 화살표는\ 
Section \ref{duality-section-trace}에서 논의한 대각합 사상이다.
\end{remark}





\section{고유 평탄 사상에 대한 추이 함자의 오른쪽 수반}
\label{duality-section-proper-flat}

\noindent
준콤팩트 준분리 스킴 사이의 고유, 평탄, 유한 표시 사상에 대해
추이 함자의 오른쪽 수반은 몇 가지 현저한 성질을 갖는다.

\begin{lemma}
\label{duality-lemma-proper-flat}
$Y$를 준콤팩트 준분리 스킴이라 하자.
$f : X \to Y$를 고유하고 평탄하며 유한 표시인 스킴의 사상이라 하자.
Lemma \ref{duality-lemma-twisted-inverse-image}에서\ 
$Rf_* : D_\QCoh(\mathcal{O}_X) \to D_\QCoh(\mathcal{O}_Y)$의 오른쪽 수반을\ 
$a$라 하자. 그러면\  $a$는 직합과 가환한다.
\end{lemma}

\begin{proof}
$P$를\  $D(\mathcal{O}_X)$의 완전 대상이라 하자. Derived Categories of Schemes,
Lemma \ref{perfect-lemma-flat-proper-perfect-direct-image-general}에 의해
복합체\  $Rf_*P$는\  $Y$ 위에서 완전하다.
$K_i$를\  $D_\QCoh(\mathcal{O}_Y)$의 대상들로 이루어진 족이라 하자. 그러면
\begin{align*}
\Hom_{D(\mathcal{O}_X)}(P, a(\bigoplus K_i))
& =
\Hom_{D(\mathcal{O}_Y)}(Rf_*P, \bigoplus K_i) \\
& =
\bigoplus \Hom_{D(\mathcal{O}_Y)}(Rf_*P, K_i) \\
& =
\bigoplus \Hom_{D(\mathcal{O}_X)}(P, a(K_i))
\end{align*}
이다. 이는 완전 대상이 콤팩트이기 때문이다(Derived Categories of Schemes,
Proposition \ref{perfect-proposition-compact-is-perfect}).
$D_\QCoh(\mathcal{O}_X)$에는 완전 생성자가 있으므로(Derived Categories of Schemes,
Theorem \ref{perfect-theorem-bondal-van-den-Bergh}), 사상\ 
$\bigoplus a(K_i) \to a(\bigoplus K_i)$가 동형임을 얻는다. 즉, $a$는 직합과
가환한다.
\end{proof}
\begin{lemma}
\label{duality-lemma-proper-flat-relative}
$Y$를 준콤팩트 준분리 스킴이라 하자.
$f : X \to Y$를 고유하고 평탄하며 유한 표시인 스킴의 사상이라 하자.
Lemma \ref{duality-lemma-twisted-inverse-image}에서\ 
$Rf_* : D_\QCoh(\mathcal{O}_X) \to D_\QCoh(\mathcal{O}_Y)$의 오른쪽 수반을\ 
$a$라 하자. 그러면
\begin{enumerate}
\item 모든 닫힌집합\  $T \subset Y$에 대해\  $Q \in D_\QCoh(Y)$가\  $T$ 위에
지지되면, $a(Q)$는\  $f^{-1}(T)$ 위에 지지된다.
\item 모든 준콤팩트 열린집합\  $V \subset Y$와 임의의\ 
$K \in D_\QCoh(\mathcal{O}_Y)$에 대해 사상 (\ref{duality-equation-sheafy})는 동형이다.
\end{enumerate}
\end{lemma}

\begin{proof}
이는\  Lemmas \ref{duality-lemma-when-sheafy},
\ref{duality-lemma-proper-noetherian}, \ref{duality-lemma-proper-flat}에서 따른다.
엄밀히 말하면 (1)은 여집합이 준콤팩트 열린집합인\  $T$에 대해서만 직접
따른다. 그러나 모든 닫힌 부분집합은 그러한 닫힌 부분집합들의 교집합이므로
일반적인 경우도 따른다.
\end{proof}
\begin{lemma}
\label{duality-lemma-compare-with-pullback-flat-proper}
$Y$를 준콤팩트 준분리 스킴이라 하자.
$f : X \to Y$를 고유하고 평탄하며 유한 표시인 스킴의 사상이라 하자.
사상 (\ref{duality-equation-compare-with-pullback})은\ 
$D_\QCoh(\mathcal{O}_Y)$의 모든 대상\  $K$에 대해 동형이다.
\end{lemma}

\begin{proof}
Lemma \ref{duality-lemma-proper-flat}에 의해\  $a$는 직합과 가환한다. 따라서
(\ref{duality-equation-compare-with-pullback})가 동형이 되는\ 
$D_\QCoh(\mathcal{O}_Y)$의 대상들은\  $D_\QCoh(\mathcal{O}_Y)$의 엄밀 충만하고
포화된 삼각 부분범주를 이루며, 더욱이 직합을 취하는 연산에 의해 보존된다.
$D_\QCoh(\mathcal{O}_Y)$는 하나의 완전 대상이 생성하는 가군 범주이므로
(Derived Categories of Schemes, Theorems
\ref{perfect-theorem-DQCoh-is-Ddga}와\ 
\ref{perfect-theorem-bondal-van-den-Bergh}), More on Algebra, Remark
\ref{more-algebra-remark-P-resolution}에서와 같이 논증하여
(\ref{duality-equation-compare-with-pullback})가 하나의 완전 대상에 대해 동형임을
보이면 충분하다는 것을 알 수 있다. 그런데 이 결과는 완전 대상들에 대해
성립한다. Lemma \ref{duality-lemma-compare-with-pullback-perfect}를 보라.
\end{proof}
\noindent
다음 보조정리는 유한 표시인 고유 평탄 사상에 대해 기저변환 사상
(\ref{duality-equation-base-change-map})이 동형임을 보인다. Example
\ref{duality-example-base-change-wrong}에서 완전 고유 사상에 대해서는 이 명제가
더 이상 성립하지 않음을 볼 것이다. 그 경우에는\  tor 독립성 조건을 부과해야 한다.

\begin{lemma}
\label{duality-lemma-proper-flat-base-change}
$g : Y' \to Y$를 준콤팩트 준분리 스킴의 사상이라 하자.
$f : X \to Y$를 유한 표시인 고유 평탄 사상이라 하자.
그러면 모든\  $K \in D_\QCoh(\mathcal{O}_Y)$에 대해 기저변환 사상
(\ref{duality-equation-base-change-map})은 동형이다.
\end{lemma}

\begin{proof}
Lemma \ref{duality-lemma-proper-flat-relative}에 의해 함자\  $a$와\  $a'$의 구성은\ 
$Y$와\  $Y'$의 열린집합에 대한 제한과 가환한다. 따라서\  Remark
\ref{duality-remark-check-over-affines}를 사용하면\  $Y' \to Y$가 아핀 스킴의 사상이라고
가정해도 된다. 이 경우 명제는\  Lemma \ref{duality-lemma-more-base-change}에서 따른다.
\end{proof}
\begin{remark}
\label{duality-remark-relative-dualizing-complex}
$Y$를 준콤팩트 준분리 스킴이라 하자.
$f : X \to Y$를 유한 표시인 고유 평탄 사상이라 하자.
$a$를\  $f$에 대한\  Lemma \ref{duality-lemma-twisted-inverse-image}의 수반이라 하자.
이 상황에서\  $\omega_{X/Y}^\bullet = a(\mathcal{O}_Y)$를 때때로
{\it 상대 쌍대화 복합체}라고 부른다. Lemma
\ref{duality-lemma-compare-with-pullback-flat-proper}에 의해\ 
$K \in D_\QCoh(\mathcal{O}_Y)$에 대한 함자적 동형\ 
$a(K) = Lf^*K \otimes_{\mathcal{O}_X}^\mathbf{L} \omega_{X/Y}^\bullet$이 있다.
또한\  Section \ref{duality-section-trace}의 대각합 사상
$$
\text{Tr}_{f, \mathcal{O}_Y} : Rf_*\omega_{X/Y}^\bullet \to \mathcal{O}_Y
$$
은\  $D_\QCoh(\mathcal{O}_Y)$의 모든\  $K$에 대한 대각합 사상을 유도한다.
더 정확히 말하면 도식
$$
\xymatrix{
Rf_*a(K) \ar[rrr]_{\text{Tr}_{f, K}} \ar@{=}[d] & & &
K \ar@{=}[d] \\
Rf_*(Lf^*K \otimes_{\mathcal{O}_X}^\mathbf{L} \omega_{X/Y}^\bullet)
\ar@{=}[r] &
K \otimes_{\mathcal{O}_Y}^\mathbf{L} Rf_*\omega_{X/Y}^\bullet
\ar[rr]^-{\text{id}_K \otimes \text{Tr}_{f, \mathcal{O}_Y}} & & K
}
$$
은 가환한다. 여기서 오른쪽 아래의 등식은\  Derived Categories of Schemes,
Lemma \ref{perfect-lemma-cohomology-base-change}이다.
$g : Y' \to Y$가 준콤팩트 준분리 스킴의 사상이고\ 
$X' = Y' \times_Y X$라고 하자. 그러면\  Lemma
\ref{duality-lemma-proper-flat-base-change}에 의해\ 
$\omega_{X'/Y'}^\bullet = L(g')^*\omega_{X/Y}^\bullet$이다. 여기서\ 
$g' : X' \to X$는 사영이다. 또한\  Lemma
\ref{duality-lemma-trace-map-and-base-change}에 의해\  $f' : X' \to Y'$에 대한 대각합 사상
$$
\text{Tr}_{f', \mathcal{O}_{Y'}} :
Rf'_*\omega_{X'/Y'}^\bullet \to \mathcal{O}_{Y'}
$$
은 기저변환 동형을 통한\  $\text{Tr}_{f, \mathcal{O}_Y}$의 기저변환이다.
\end{remark}
\begin{remark}
\label{duality-remark-relative-dualizing-complex-relative-cup-product}
$f : X \to Y$, $\omega^\bullet_{X/Y}$, $\text{Tr}_{f, \mathcal{O}_Y}$를\ 
Remark \ref{duality-remark-relative-dualizing-complex}에서와 같이 두자.
$K$와\  $M$을\  $D_\QCoh(\mathcal{O}_X)$의 대상으로 두고, $M$이 유사코히런트라고
하자(예를 들어 완전이라고 하자). Cohomology, Equation
(\ref{cohomology-equation-internal-hom})을 통해 동형\ 
$K \to R\SheafHom_{\mathcal{O}_X}(M, \omega^\bullet_{X/Y})$에 대응하는 사상\ 
$K \otimes_{\mathcal{O}_X}^\mathbf{L} M \to \omega^\bullet_{X/Y}$가 주어졌다고 하자.
그러면 상대 컵곱(Cohomology, Remark \ref{cohomology-remark-cup-product})
$$
Rf_*K \otimes_{\mathcal{O}_Y}^\mathbf{L} Rf_*M
\to
Rf_*(K \otimes_{\mathcal{O}_X}^\mathbf{L} M)
\to
Rf_*\omega^\bullet_{X/Y}
\xrightarrow{\text{Tr}_{f, \mathcal{O}_Y}}
\mathcal{O}_Y
$$
은 동형\ 
$Rf_*K \to R\SheafHom_{\mathcal{O}_Y}(Rf_*M, \mathcal{O}_Y)$를 정한다.
실제로\  $\omega^\bullet_{X/Y} = a(\mathcal{O}_Y)$이므로 정준 사상
(\ref{duality-equation-sheafy-trace})
$$
Rf_*R\SheafHom_{\mathcal{O}_X}(M, \omega^\bullet_{X/Y}) \to
R\SheafHom_{\mathcal{O}_Y}(Rf_*M, \mathcal{O}_Y)
$$
은\  Lemma \ref{duality-lemma-iso-on-RSheafHom}, Remark
\ref{duality-remark-iso-on-RSheafHom}, 그리고\  $M$과\  $Rf_*M$이 유사코히런트라는 사실에
의해 동형이다. 마지막 사실에 대해서는\  Derived Categories of Schemes, Lemma
\ref{perfect-lemma-flat-proper-pseudo-coherent-direct-image-general}을 보라.
상대 컵곱이 이 동형을 유도함을 보이려면\  Cohomology, Remark
\ref{cohomology-remark-relative-cup-and-composition}의 도식이 가환함을 사용한다.
\end{remark}
\begin{lemma}
\label{duality-lemma-properties-relative-dualizing}
$Y$를 준콤팩트 준분리 스킴이라 하자.
$f : X \to Y$를 고유하고 평탄하며 유한 표시인 스킴의 사상이라 하고,
그 상대 쌍대화 복합체를\  $\omega_{X/Y}^\bullet$이라 하자
(Remark \ref{duality-remark-relative-dualizing-complex}). 그러면
\begin{enumerate}
\item $\omega_{X/Y}^\bullet$은\  $D(\mathcal{O}_X)$의\  $Y$-완전 대상이다.
\item $Rf_*\omega_{X/Y}^\bullet$의 양의 차수 코호몰로지 층들은 영이다.
\item $\mathcal{O}_X \to
R\SheafHom_{\mathcal{O}_X}(\omega_{X/Y}^\bullet, \omega_{X/Y}^\bullet)$
은 동형이다.
\end{enumerate}
\end{lemma}

\begin{proof}
$\omega_{X/Y}^\bullet$의 구성이 기저변환과 가환하므로(Remark
\ref{duality-remark-relative-dualizing-complex}를 보라), $Y$가 아핀이라고 가정할 수
있고 실제로 그렇게 가정한다. $D(\mathcal{O}_X)$의 완전 대상\  $E$에 대해
\begin{align*}
Rf_*(E \otimes_{\mathcal{O}_X}^\mathbf{L} \omega_{X/Y}^\bullet)
& =
Rf_*R\SheafHom_{\mathcal{O}_X}(E^\vee, \omega_{X/Y}^\bullet) \\
& =
R\SheafHom_{\mathcal{O}_Y}(Rf_*E^\vee, \mathcal{O}_Y) \\
& =
(Rf_*E^\vee)^\vee
\end{align*}
이다. 첫 번째 등식은\  Cohomology, Lemma
\ref{cohomology-lemma-dual-perfect-complex}을 보라. 두 번째 등식은\  Lemma
\ref{duality-lemma-iso-on-RSheafHom}, Remark \ref{duality-remark-iso-on-RSheafHom}, Derived
Categories of Schemes, Lemma
\ref{perfect-lemma-flat-proper-perfect-direct-image-general}을 보라.
세 번째 등식은 쌍대의 정의이다. 특히 이 인용 결과들은 얻어진 대상이\ 
$D(\mathcal{O}_Y)$의 완전 대상임도 보인다. 따라서\  More on Morphisms, Lemma
\ref{more-morphisms-lemma-characterize-relatively-perfect}에 의해\ 
$\omega_{X/Y}^\bullet$이\  $Y$-완전임을 얻는다. 이는 (1)을 증명한다.

\medskip\noindent
$M$을\  $D_\QCoh(\mathcal{O}_Y)$의 대상이라 하자. 그러면
\begin{align*}
\Hom_Y(M, Rf_*\omega_{X/Y}^\bullet) & =
\Hom_X(Lf^*M, \omega_{X/Y}^\bullet) \\
& =
\Hom_Y(Rf_*Lf^*M, \mathcal{O}_Y) \\
& =
\Hom_Y(M \otimes_{\mathcal{O}_Y}^\mathbf{L} Rf_*\mathcal{O}_X, \mathcal{O}_Y)
\end{align*}
이다. 첫 번째 등식은\  Cohomology, Lemma \ref{cohomology-lemma-adjoint}에 의해
성립한다. 두 번째 등식은\  $a$의 구성에 의한다. 세 번째 등식은\  Derived Categories
of Schemes, Lemma \ref{perfect-lemma-cohomology-base-change}에 의한다.
$Rf_*\mathcal{O}_X$가 어떤\  $N$에 대해\  $[0, N]$ 안의\  tor-진폭을 갖는 완전
대상임을 상기하자. Derived Categories of Schemes, Lemma
\ref{perfect-lemma-flat-proper-perfect-direct-image-general}을 보라.
따라서\  $Rf_*\mathcal{O}_X$를\  $[0, N]$ 차수에 놓인 유한 사영 가군들의 복합체로
표현할 수 있다(More on Algebra, Lemma \ref{more-algebra-lemma-perfect}과\ 
$Y$가 아핀이라는 사실을 사용한다). 그러므로 어떤\  $i > 0$에 대해\ 
$M = \mathcal{O}_Y[-i]$이면 마지막 군은 영이다. $Y$가 아핀이므로\ 
$i > 0$에 대해\  $H^i(Rf_*\omega_{X/Y}^\bullet) = 0$임을 얻는다.
이는 (2)를 증명한다.

\medskip\noindent
$E$를\  $D_\QCoh(\mathcal{O}_X)$의 완전 대상이라 하자. 그러면
\begin{align*}
\Hom_X(E, R\SheafHom_{\mathcal{O}_X}(\omega_{X/Y}^\bullet, \omega_{X/Y}^\bullet)
& =
\Hom_X(E \otimes_{\mathcal{O}_X}^\mathbf{L} \omega_{X/Y}^\bullet,
\omega_{X/Y}^\bullet) \\
& =
\Hom_Y(Rf_*(E \otimes_{\mathcal{O}_X}^\mathbf{L} \omega_{X/Y}^\bullet),
\mathcal{O}_Y) \\
& =
\Hom_Y(Rf_*(R\SheafHom_{\mathcal{O}_X}(E^\vee, \omega_{X/Y}^\bullet)),
\mathcal{O}_Y) \\
& =
\Hom_Y(R\SheafHom_{\mathcal{O}_Y}(Rf_*E^\vee, \mathcal{O}_Y),
\mathcal{O}_Y) \\
& =
R\Gamma(Y, Rf_*E^\vee) \\
& =
\Hom_X(E, \mathcal{O}_X)
\end{align*}
이다. 첫 번째 등식은\  Cohomology, Lemma
\ref{cohomology-lemma-internal-hom}에 의해 성립한다. 두 번째 등식은\ 
$\omega_{X/Y}^\bullet$의 정의이다. 세 번째 등식은 쌍대 완전 복합체\ 
$E^\vee$의 구성에서 나온다. Cohomology, Lemma
\ref{cohomology-lemma-dual-perfect-complex}을 보라. 네 번째 등식은 증명의 첫
문단에서 보인 등식\ 
$Rf_*R\SheafHom_{\mathcal{O}_X}(E^\vee, \omega_{X/Y}^\bullet) =
R\SheafHom_{\mathcal{O}_Y}(Rf_*E^\vee, \mathcal{O}_Y)$에서 따른다.
다섯 번째 등식은 완전 복합체의 이중 쌍대성(Cohomology, Lemma
\ref{cohomology-lemma-dual-perfect-complex})과\  $Rf_*E$가 완전하다는 사실에
의해 성립한다. 후자는\  Derived Categories of Schemes, Lemma
\ref{perfect-lemma-flat-proper-perfect-direct-image-general}을 보라.
마지막 등식은\  $f$에 대한\  Leray이다. 이 일련의 등식들은 본질적으로\  Yoneda
보조정리에 의해 (3)이 성립함을 보인다. 구체적으로 대상\ 
$R\SheafHom(\omega_{X/Y}^\bullet, \omega_{X/Y}^\bullet)$은\  Derived Categories
of Schemes, Lemma \ref{perfect-lemma-quasi-coherence-internal-hom}에 의해\ 
$D_\QCoh(\mathcal{O}_X)$에 속한다. 위에서\  $E = \mathcal{O}_X$로 놓으면\ 
$\text{id}_{\mathcal{O}_X} \in \Hom_X(\mathcal{O}_X, \mathcal{O}_X)$에 대응하는
사상\ 
$\alpha : \mathcal{O}_X \to
R\SheafHom_{\mathcal{O}_X}(\omega_{X/Y}^\bullet, \omega_{X/Y}^\bullet)$
를 얻는다. 위의 모든 동형은\  $E$에 함자적이므로\  $\alpha$의 원뿔은
모든 완전 대상\  $E$에 대해\  $\Hom(E, C) = 0$을 만족하는\ 
$D_\QCoh(\mathcal{O}_X)$의 대상\  $C$임을 알 수 있다. 완전 대상들이 생성하므로
(Derived Categories of Schemes, Theorem
\ref{perfect-theorem-bondal-van-den-Bergh}), $\alpha$가 동형임을 얻는다.
\end{proof}
\begin{lemma}[강직성]
\label{duality-lemma-van-den-bergh}
$Y$를 준콤팩트 준분리 스킴이라 하자.
$f : X \to Y$를 유한 표시인 고유 평탄 사상이라 하고, 그 상대 쌍대화 복합체를\ 
$\omega_{X/Y}^\bullet$이라 하자(Remark
\ref{duality-remark-relative-dualizing-complex}). 정준 동형
\begin{equation}
\label{duality-equation-pre-rigid}
\mathcal{O}_X =
c(L\text{pr}_1^*\omega_{X/Y}^\bullet) =
c(L\text{pr}_2^*\omega_{X/Y}^\bullet)
\end{equation}
과 정준 동형
\begin{equation}
\label{duality-equation-rigid}
\omega_{X/Y}^\bullet =
c\left(L\text{pr}_1^*\omega_{X/Y}^\bullet
\otimes_{\mathcal{O}_{X \times_Y X}}^\mathbf{L}
L\text{pr}_2^*\omega_{X/Y}^\bullet\right)
\end{equation}
이 존재한다. 여기서\  $c$는 대각 사상\  $\Delta : X \to X \times_Y X$에 대한\ 
Lemma \ref{duality-lemma-twisted-inverse-image}의 오른쪽 수반이다.
\end{lemma}

\begin{proof}
$a$를\  Lemma \ref{duality-lemma-twisted-inverse-image}에서의\  $Rf_*$의 오른쪽 수반이라
하자. Cartesian 정사각형
$$
\xymatrix{
X \times_Y X \ar[r]_q \ar[d]_p & X \ar[d]_f \\
X \ar[r]^f & Y
}
$$
을 생각하자. $b$를\  Lemma \ref{duality-lemma-twisted-inverse-image}에서의\  $p$의 오른쪽
수반이라 하자. 그러면
\begin{align*}
\omega_{X/Y}^\bullet
& =
c(b(\omega_{X/Y}^\bullet)) \\
& =
c(Lp^*\omega_{X/Y}^\bullet
\otimes_{\mathcal{O}_{X \times_Y X}}^\mathbf{L} b(\mathcal{O}_X)) \\
& =
c(Lp^*\omega_{X/Y}^\bullet
\otimes_{\mathcal{O}_{X \times_Y X}}^\mathbf{L}
Lq^*a(\mathcal{O}_Y)) \\
& =
c(Lp^*\omega_{X/Y}^\bullet
\otimes_{\mathcal{O}_{X \times_Y X}}^\mathbf{L}
Lq^*\omega_{X/Y}^\bullet)
\end{align*}
이며, 이는 (\ref{duality-equation-rigid})과 같다. 설명은 다음과 같다.
\begin{enumerate}
\item $\text{id}_X = p \circ \Delta$이므로\  $\text{id} = c \circ b$이고,
첫 번째 등식이 성립한다.
\item 두 번째 등식은\  Lemma
\ref{duality-lemma-compare-with-pullback-flat-proper}에 의해 성립한다.
\item 세 번째 등식은\  Lemma \ref{duality-lemma-proper-flat-base-change}와\ 
$\mathcal{O}_X = Lf^*\mathcal{O}_Y$라는 사실에 의해 성립한다.
\item 네 번째 등식은\  $\omega_{X/Y}^\bullet = a(\mathcal{O}_Y)$이므로 성립한다.
\end{enumerate}
Equation (\ref{duality-equation-pre-rigid})도 정확히 같은 방식으로 증명된다.
\end{proof}
\begin{remark}
\label{duality-remark-van-den-bergh}
Lemma \ref{duality-lemma-van-den-bergh}는 우리의 상대 쌍대화 복합체가\ 
\cite{vdB-rigid}에서 도입된 개념과 유사한 의미에서 {\it 강직함}을 뜻한다.
실제로 (\ref{duality-equation-rigid})의 오른쪽 함자는\  $\omega_{X/Y}^\bullet$에 대해
``이차적''이고, (\ref{duality-equation-rigid})의 왼쪽 함자는 ``선형적''이므로 이는
복합체\  $\omega_{X/Y}^\bullet$을 어느 정도 ``고정한다''. ``강직한'' (상대)
쌍대화 복합체를 사용하는 쌍대론의 접근법이 있다. 예를 들어\ 
\cite{Neeman-rigid}, \cite{Yekutieli-rigid}, \cite{Yekutieli-Zhang}을 보라.
Section \ref{duality-section-relative-dualizing-complexes}에서 이 문제로 돌아올 것이다.
\end{remark}




\section{완전 고유 사상에 대한 추이 함자의 오른쪽 수반}
\label{duality-section-perfect-proper}

\noindent
이 절에서 적절한 일반성은 준콤팩트 준분리 스킴 사이의 완전 고유 사상을
고찰하는 것이지만, 이에 대해서는\  \cite{LN}을 보라.
\begin{lemma}
\label{duality-lemma-proper-flat-noetherian}
$f : X \to Y$를\  Noether 스킴 사이의 완전 고유 사상이라 하자.
Lemma \ref{duality-lemma-twisted-inverse-image}에서\ 
$Rf_* : D_\QCoh(\mathcal{O}_X) \to D_\QCoh(\mathcal{O}_Y)$의 오른쪽 수반을\ 
$a$라 하자. 그러면\  $a$는 직합과 가환한다.
\end{lemma}

\begin{proof}
$P$를\  $D(\mathcal{O}_X)$의 완전 대상이라 하자. More on Morphisms, Lemma
\ref{more-morphisms-lemma-perfect-proper-perfect-direct-image}에 의해
복합체\  $Rf_*P$는\  $Y$ 위에서 완전하다.
$K_i$를\  $D_\QCoh(\mathcal{O}_Y)$의 대상들로 이루어진 족이라 하자. 그러면
\begin{align*}
\Hom_{D(\mathcal{O}_X)}(P, a(\bigoplus K_i))
& =
\Hom_{D(\mathcal{O}_Y)}(Rf_*P, \bigoplus K_i) \\
& =
\bigoplus \Hom_{D(\mathcal{O}_Y)}(Rf_*P, K_i) \\
& =
\bigoplus \Hom_{D(\mathcal{O}_X)}(P, a(K_i))
\end{align*}
이다. 이는 완전 대상이 콤팩트이기 때문이다(Derived Categories of Schemes,
Proposition \ref{perfect-proposition-compact-is-perfect}).
$D_\QCoh(\mathcal{O}_X)$에는 완전 생성자가 있으므로(Derived Categories of Schemes,
Theorem \ref{perfect-theorem-bondal-van-den-Bergh}), 사상\ 
$\bigoplus a(K_i) \to a(\bigoplus K_i)$가 동형임을 얻는다. 즉, $a$는 직합과
가환한다.
\end{proof}
\begin{lemma}
\label{duality-lemma-proper-flat-noetherian-relative}
$f : X \to Y$를\  Noether 스킴 사이의 완전 고유 사상이라 하자.
Lemma \ref{duality-lemma-twisted-inverse-image}에서\ 
$Rf_* : D_\QCoh(\mathcal{O}_X) \to D_\QCoh(\mathcal{O}_Y)$의 오른쪽 수반을\ 
$a$라 하자. 그러면
\begin{enumerate}
\item 모든 닫힌집합\  $T \subset Y$에 대해\  $Q \in D_\QCoh(Y)$가\  $T$ 위에
지지되면, $a(Q)$는\  $f^{-1}(T)$ 위에 지지된다.
\item 모든 열린집합\  $V \subset Y$와 임의의\  $K \in D_\QCoh(\mathcal{O}_Y)$에
대해 사상 (\ref{duality-equation-sheafy})는 동형이다.
\end{enumerate}
\end{lemma}

\begin{proof}
이는\  Lemmas \ref{duality-lemma-when-sheafy},
\ref{duality-lemma-proper-noetherian}, \ref{duality-lemma-proper-flat-noetherian}에서 따른다.
\end{proof}
\begin{lemma}
\label{duality-lemma-compare-with-pullback-flat-proper-noetherian}
$f : X \to Y$를\  Noether 스킴 사이의 완전 고유 사상이라 하자.
사상 (\ref{duality-equation-compare-with-pullback})은\  $D_\QCoh(\mathcal{O}_Y)$의
모든 대상\  $K$에 대해 동형이다.
\end{lemma}

\begin{proof}
Lemma \ref{duality-lemma-proper-flat-noetherian}에 의해\  $a$는 직합과 가환한다. 따라서
(\ref{duality-equation-compare-with-pullback})가 동형이 되는\ 
$D_\QCoh(\mathcal{O}_Y)$의 대상들은\  $D_\QCoh(\mathcal{O}_Y)$의 엄밀 충만하고
포화된 삼각 부분범주를 이루며, 더욱이 직합을 취하는 연산에 의해 보존된다.
$D_\QCoh(\mathcal{O}_Y)$는 하나의 완전 대상이 생성하는 가군 범주이므로
(Derived Categories of Schemes, Theorems
\ref{perfect-theorem-DQCoh-is-Ddga}와\ 
\ref{perfect-theorem-bondal-van-den-Bergh}), More on Algebra, Remark
\ref{more-algebra-remark-P-resolution}에서와 같이 논증하여
(\ref{duality-equation-compare-with-pullback})가 하나의 완전 대상에 대해 동형임을
보이면 충분하다는 것을 알 수 있다. 그런데 이 결과는 완전 대상들에 대해
성립한다. Lemma \ref{duality-lemma-compare-with-pullback-perfect}를 보라.
\end{proof}
\begin{lemma}
\label{duality-lemma-proper-perfect-base-change}
$f : X \to Y$를\  Noether 스킴 사이의 완전 고유 사상이라 하자.
$g : Y' \to Y$를\  $Y'$이\  Noether인 사상이라 하자. $X$와\  $Y'$이\  $Y$ 위에서\ 
tor 독립이면, 모든\  $K \in D_\QCoh(\mathcal{O}_Y)$에 대해 기저변환 사상
(\ref{duality-equation-base-change-map})은 동형이다.
\end{lemma}

\begin{proof}
Lemma \ref{duality-lemma-proper-flat-noetherian-relative}에 의해 함자\  $a$와\  $a'$의
구성은\  $Y$와\  $Y'$의 열린집합에 대한 제한과 가환한다. 따라서\  Remark
\ref{duality-remark-check-over-affines}를 사용하면\  $Y' \to Y$가 아핀 스킴의 사상이라고
가정해도 된다. 이 경우 명제는\  Lemma \ref{duality-lemma-more-base-change}에서 따른다.
\end{proof}


\section{유효\  Cartier 약수에 대한 추이 함자의 오른쪽 수반}
\label{duality-section-dualizing-Cartier}

\noindent
$X$를 스킴이라 하고, $i : D \to X$를 유효\  Cartier 약수의 포함이라 하자.
$\mathcal{N} = i^*\mathcal{O}_X(D)$를\  $i$의 법선층이라 하자. Morphisms, Section
\ref{morphisms-section-conormal-sheaf}와\  Divisors, Section
\ref{divisors-section-effective-Cartier-divisors}를 보라.
$R\SheafHom(\mathcal{O}_D, -)$는\ 
$i_* : D(\mathcal{O}_D) \to D(\mathcal{O}_X)$의 오른쪽 수반을 나타내며\ 
$i_*R\SheafHom(\mathcal{O}_D, -) =
R\SheafHom_{\mathcal{O}_X}(i_*\mathcal{O}_D, -)$
의 성질을 가짐을 상기하자. Section
\ref{duality-section-sections-with-exact-support}를 보라.
\begin{lemma}
\label{duality-lemma-compute-for-effective-Cartier}
위와 같이\  $X$를 스킴이라 하고\  $D \subset X$를 유효\  Cartier 약수라 하자.
$D(\mathcal{O}_D)$에서 정준 동형\ 
$R\SheafHom(\mathcal{O}_D, \mathcal{O}_X) = \mathcal{N}[-1]$이 존재한다.
\end{lemma}

\begin{proof}
동치로, $R\SheafHom(\mathcal{O}_D, \mathcal{O}_X)$의 유일한 영이 아닌
코호몰로지 층은 차수\  $1$에 있으며 이 층이\  $\mathcal{N}$과 동형이라고 말하는
것이다. $i_*$는 정확하고 충실충만하므로,
$i_*R\SheafHom(\mathcal{O}_D, \mathcal{O}_X)$가\ 
$i_*\mathcal{N}[-1]$과 동형임을 증명하면 충분하다. Lemma
\ref{duality-lemma-sheaf-with-exact-support-ext}에 의해\ 
$i_*R\SheafHom(\mathcal{O}_D, \mathcal{O}_X) =
R\SheafHom_{\mathcal{O}_X}(i_*\mathcal{O}_D, \mathcal{O}_X)$
이다. $D$의 아이디얼 층을\  $\mathcal{I}$라 하면 계산에 사용할 수 있는 해소
$$
0 \to \mathcal{I} \to \mathcal{O}_X \to i_*\mathcal{O}_D \to 0
$$
이 있다. 국소 계산에 의해\ 
$R\SheafHom_{\mathcal{O}_X}(\mathcal{O}_X, \mathcal{O}_X) = \mathcal{O}_X$이고\ 
$R\SheafHom_{\mathcal{O}_X}(\mathcal{I}, \mathcal{O}_X) = \mathcal{O}_X(D)$이므로
$$
R\SheafHom_{\mathcal{O}_X}(i_*\mathcal{O}_D, \mathcal{O}_X) =
(\mathcal{O}_X \to \mathcal{O}_X(D))
$$
임을 알 수 있다. 여기서 오른쪽의\  $\mathcal{O}_X$는 차수\  $0$에 있고\ 
$\mathcal{O}_X(D)$는 차수\  $1$에 있다. 결과는\  $D$가\ 
$\mathcal{O}_X(D)$의 정준 절단의 영점 스킴이라는 사실과\ 
$\mathcal{N} = i^*\mathcal{O}_X(D)$라는 사실에서 나오는 짧은 완전열
$$
0 \to \mathcal{O}_X \to \mathcal{O}_X(D) \to i_*\mathcal{N} \to 0
$$
에서 따른다.
\end{proof}
\noindent
$D(\mathcal{O}_X)$의 모든 대상\  $K$에 대해\  $K$에 함자적이고 완전 삼각형과
양립하는\  $D(\mathcal{O}_D)$의 정준 사상
\begin{equation}
\label{duality-equation-map-effective-Cartier}
Li^*K
\otimes_{\mathcal{O}_D}^\mathbf{L}
R\SheafHom(\mathcal{O}_D, \mathcal{O}_X)
\longrightarrow
R\SheafHom(\mathcal{O}_D, K)
\end{equation}
이 존재한다. 실제로 이 사상은 사상
$$
i_*(Li^*K \otimes^\mathbf{L}_{\mathcal{O}_D}
R\SheafHom(\mathcal{O}_D, \mathcal{O}_X)) =
K \otimes^\mathbf{L}_{\mathcal{O}_X}
R\SheafHom_{\mathcal{O}_X}(i_*\mathcal{O}_D, \mathcal{O}_X)
\longrightarrow K
$$
에 수반된다. 여기서 등식은\  Cohomology, Lemma
\ref{cohomology-lemma-projection-formula-closed-immersion}이고, 화살표는\ 
$\mathcal{O}_X \to i_*\mathcal{O}_D$가 유도하는 정준 사상\ 
$R\SheafHom_{\mathcal{O}_X}(i_*\mathcal{O}_D, \mathcal{O}_X) \to \mathcal{O}_X$
에서 나온다.

\medskip\noindent
$K \in D_\QCoh(\mathcal{O}_X)$이면,
(\ref{duality-equation-map-effective-Cartier})는
(\ref{duality-equation-compare-with-pullback})와 같다. 여기서는\  Lemma
\ref{duality-lemma-twisted-inverse-image-closed}의 식별\ 
$a(K) = R\SheafHom(\mathcal{O}_D, K)$을 사용한다.
$K \in D_\QCoh(\mathcal{O}_X)$이고\  $X$가\  Noether이면, 다음 보조정리는\ 
Lemma \ref{duality-lemma-compare-with-pullback-flat-proper-noetherian}의 특수한
경우이다.
\begin{lemma}
\label{duality-lemma-sheaf-with-exact-support-effective-Cartier}
위와 같이\  $X$를 스킴이라 하고\  $D \subset X$를 유효\  Cartier 약수라 하자.
그러면 (\ref{duality-equation-map-effective-Cartier})를\  Lemma
\ref{duality-lemma-compute-for-effective-Cartier}와 결합하면\  $D(\mathcal{O}_X)$의\ 
$K$에 함자적인 동형
$$
Li^*K \otimes_{\mathcal{O}_D}^\mathbf{L} \mathcal{N}[-1]
\longrightarrow
R\SheafHom(\mathcal{O}_D, K)
$$
이 정의된다.
\end{lemma}

\begin{proof}
$i_*$는 가군들 위에서 정확하고 충실충만하므로, 이 사상이 동형임을 증명하려면\ 
$i_*$를 적용한 뒤 동형임을 보이면 충분하다. Lemma
\ref{duality-lemma-compute-for-effective-Cartier}의 증명에서 사용한 짧은 완전열\ 
$0 \to \mathcal{I} \to \mathcal{O}_X \to i_*\mathcal{O}_D \to 0$과\ 
$0 \to \mathcal{O}_X \to \mathcal{O}_X(D) \to i_*\mathcal{N} \to 0$을
더 언급하지 않고 사용하겠다. Cohomology, Lemma
\ref{cohomology-lemma-projection-formula-closed-immersion}(이는 사상
(\ref{duality-equation-map-effective-Cartier})를 정의할 때 사용되었다)에 의해 왼쪽 변은
$$
K \otimes_{\mathcal{O}_X}^\mathbf{L} i_*\mathcal{N}[-1] =
K \otimes_{\mathcal{O}_X}^\mathbf{L} (\mathcal{O}_X \to \mathcal{O}_X(D))
$$
이 된다. 오른쪽 변은
\begin{align*}
R\SheafHom_{\mathcal{O}_X}(i_*\mathcal{O}_D, K) & =
R\SheafHom_{\mathcal{O}_X}((\mathcal{I} \to \mathcal{O}_X), K) \\
& =
R\SheafHom_{\mathcal{O}_X}((\mathcal{I} \to \mathcal{O}_X), \mathcal{O}_X)
\otimes_{\mathcal{O}_X}^\mathbf{L} K
\end{align*}
이 된다. 마지막 등식은\  Cohomology, Lemma
\ref{cohomology-lemma-dual-perfect-complex}에 의한다. 이 사상은 동형
$$
R\SheafHom_{\mathcal{O}_X}((\mathcal{I} \to \mathcal{O}_X), \mathcal{O}_X)
= (\mathcal{O}_X \to \mathcal{O}_X(D))
$$
에서 나오므로 보조정리는 명백하다.
\end{proof}




\section{예들에서의 추이 함자의 오른쪽 수반}
\label{duality-section-examples}

\noindent
이 절에서는 몇 가지 예에서 추이 함자의 오른쪽 수반을 계산한다. 이 동형들은
정준적이지만 가능한 가장 약한 의미에서만 그러하다. 즉, 이 동형들이 기저변환이나
사상의 합성과 같은 여러 연산과 양립함을 증명하거나 주장하지 않는다. 이런 종류의
문제에 관한 문헌은 방대하다. 독자는\  \cite{RD}, \cite{Conrad-GD}의 내용에서
시작할 수 있고(이 문헌들은 쌍대론의 출발점은 다르지만 상대 쌍대화 복합체의 정준
대표자를 구성하는 문제를 다룬다), 그다음\  Joseph Lipman과 공동연구자들의 저작을
계속 찾아볼 수 있다.
\begin{lemma}
\label{duality-lemma-upper-shriek-P1}
$Y$를 준콤팩트 준분리 스킴이라 하자. $\mathcal{E}$를 계수\  $n + 1$인
유한 국소 자유\  $\mathcal{O}_Y$-가군이라 하고, 그 행렬식을\ 
$\mathcal{L} = \wedge^{n + 1}(\mathcal{E})$이라 하자.
$f : X = \mathbf{P}(\mathcal{E}) \to Y$를 사영사상이라 하자.
$a$를\  Lemma \ref{duality-lemma-twisted-inverse-image}에서의\ 
$Rf_* : D_\QCoh(\mathcal{O}_X) \to D_\QCoh(\mathcal{O}_Y)$의 오른쪽 수반이라
하자. 그러면 동형
$$
c : f^*\mathcal{L}(-n - 1)[n] \longrightarrow a(\mathcal{O}_Y)
$$
이 존재한다. 특히\  $\mathcal{E} = \mathcal{O}_Y^{\oplus n + 1}$이면\ 
$X = \mathbf{P}^n_Y$이고\ 
$a(\mathcal{O}_Y) = \mathcal{O}_X(-n - 1)[n]$을 얻는다.
\end{lemma}

\begin{proof}
Cohomology of Schemes, Lemma
\ref{coherent-lemma-cohomology-projective-bundle}의 증명에서 정준 동형
$$
R^nf_*(f^*\mathcal{L}(-n - 1)) \longrightarrow \mathcal{O}_Y
$$
을 구성하였다. 더욱이\ 
$Rf_*(f^*\mathcal{L}(-n - 1))[n] = R^nf_*(f^*\mathcal{L}(-n - 1))$,
즉 다른 고차 직접상들은 영이다. 따라서 동형
$$
Rf_*(f^*\mathcal{L}(-n - 1)[n]) \longrightarrow \mathcal{O}_Y
$$
을 얻는다. $a$가\  $Rf_*$의 오른쪽 수반이므로 이 동형은 보조정리의 명제에
나오는\  $c$를 결정한다. Lemma \ref{duality-lemma-proper-noetherian}에 의해\  $a$의
구성은 밑에서 국소적이다. 특히\  $c$가 동형인지 확인하기 위해\  $Y$ 위에서
국소적으로 작업할 수 있다. 다시 말해\  $Y$가 아핀이고\ 
$\mathcal{E} = \mathcal{O}_Y^{\oplus n + 1}$이라고 가정해도 된다.
이 경우 층들\  $\mathcal{O}_X, \mathcal{O}_X(-1), \ldots,
\mathcal{O}_X(-n)$은\  $D_\QCoh(X)$를 생성한다. Derived Categories of Schemes,
Lemma \ref{perfect-lemma-generator-P1}를 보라. 따라서\ 
$c : \mathcal{O}_X(-n - 1)[n] \to a(\mathcal{O}_Y)$가\ 
$i \in \{-n, \ldots, 0\}$와\  $p \in \mathbf{Z}$에 대한 함자
$$
F_{i, p}(-) = \Hom_{D(\mathcal{O}_X)}(\mathcal{O}_X(i), (-)[p])
$$
아래에서 동형으로 보내짐을 보이면 충분하다. $F_{0, p}$에 대해서는 화살표\ 
$c$의 구성에 의해 성립한다! $i \in \{-n, \ldots, -1\}$에 대해서는
사영공간의 코호몰로지 계산(Cohomology of Schemes, Lemma
\ref{coherent-lemma-cohomology-projective-space-over-ring})에 의해
$$
\Hom_{D(\mathcal{O}_X)}(\mathcal{O}_X(i), \mathcal{O}_X(-n - 1)[n + p]) =
H^p(X, \mathcal{O}_X(-n - 1 - i)) = 0
$$
이고, 같은 보조정리에 의해\  $Rf_*\mathcal{O}_X(i) = 0$이므로
$$
\Hom_{D(\mathcal{O}_X)}(\mathcal{O}_X(i), a(\mathcal{O}_Y)[p]) =
\Hom_{D(\mathcal{O}_Y)}(Rf_*\mathcal{O}_X(i), \mathcal{O}_Y[p]) = 0
$$
이다. 따라서\  $F_{i, p}(c)$의 원천과 표적은 소멸하고,
$F_{i, p}(c)$는 반드시 동형이다. 이로써 증명이 끝난다.
\end{proof}
\begin{example}
\label{duality-example-base-change-wrong}
$f$가 완전 고유이고\  $g$가 완전이면 기저변환 사상
(\ref{duality-equation-base-change-map})은 동형일 필요가 없다. $k$를 체라 하자.
$Y = \mathbf{A}^2_k$라 하고, $f : X \to Y$를\  $Y$의 원점에서의 부풀리기라
하자. $E \subset X$를 예외 약수라 쓰자. 그러면\  $f$를 다음과 같이
인수분해할 수 있다.
$$
X \xrightarrow{i} \mathbf{P}^1_Y \xrightarrow{p} Y
$$
이에 따라 인수분해\  $a = c \circ b$를 얻는데, 여기서\  $a$, $b$, $c$는 각각\ 
$Rf_*$, $Rp_*$, $Ri_*$에 대한\  Lemma
\ref{duality-lemma-twisted-inverse-image}의 오른쪽 수반이다. $\mathcal{O}(n)$을\ 
$\mathbf{P}^1_Y$의 구조층의\  Serre 꼬임이라 하고, 그\  $X$로의 제한을\ 
$\mathcal{O}_X(n)$이라 하자. $X \subset \mathbf{P}^1_Y$는 차수\  1인 방정식으로
잘려 나오므로\  $\mathcal{O}(X) = \mathcal{O}(1)$임을 유의하자. Lemma
\ref{duality-lemma-upper-shriek-P1}에 의해\ 
$b(\mathcal{O}_Y) = \mathcal{O}(-2)[1]$이다. Lemma
\ref{duality-lemma-twisted-inverse-image-closed}에 의해
$$
a(\mathcal{O}_Y) = c(b(\mathcal{O}_Y)) =
c(\mathcal{O}(-2)[1]) =
R\SheafHom(\mathcal{O}_X, \mathcal{O}(-2)[1]) =
\mathcal{O}_X(-1)
$$
이다. 마지막 등식은\  Lemma
\ref{duality-lemma-sheaf-with-exact-support-effective-Cartier}에 의한다.
$Y' = \Spec(k)$를\  $Y$의 원점이라 하자. $a(\mathcal{O}_Y)$의\ 
$X' = E = \mathbf{P}^1_k$로의 제한은 코호몰로지 차수\  $0$에 놓인 차수\ 
$-1$의 가역층이다. 그러나 다른 한편으로\ 
$a'(\mathcal{O}_{\Spec(k)}) = \mathcal{O}_E(-2)[1]$은 코호몰로지 차수\ 
$-1$에 놓인 차수\  $-2$의 가역층이므로 서로 다르다. 이 예에서는\  Lemma
\ref{duality-lemma-more-base-change}의\  Tor 독립 가정이 위배된다.
\end{example}
\begin{lemma}
\label{duality-lemma-ext}
$Y$를 환 달린 공간이라 하자. $\mathcal{I} \subset \mathcal{O}_Y$를
아이디얼 층이라 하자. $\mathcal{O}_X = \mathcal{O}_Y/\mathcal{I}$라 두고\ 
$\mathcal{N} =
\SheafHom_{\mathcal{O}_Y}(\mathcal{I}/\mathcal{I}^2, \mathcal{O}_X)$라 두자.
정준 동형\ 
$c : \mathcal{N} \to
\SheafExt^1_{\mathcal{O}_Y}(\mathcal{O}_X, \mathcal{O}_X)
$
가 존재한다.
\end{lemma}

\begin{proof}
정준 짧은 완전열
\begin{equation}
\label{duality-equation-second-order-thickening}
0 \to \mathcal{I}/\mathcal{I}^2 \to \mathcal{O}_Y/\mathcal{I}^2 \to
\mathcal{O}_X \to 0
\end{equation}
을 생각하자. $U \subset X$를 열린집합이라 하고\  $s \in \mathcal{N}(U)$라 하자.
그러면 (\ref{duality-equation-second-order-thickening})을\  $s$를 따라 밀어내어\ 
$\mathcal{O}_X|_U$의\  $\mathcal{O}_X|_U$에 의한 확장\  $E_s$를 얻을 수 있다.
이는 다시\  $U$ 위에서\ 
$\SheafExt^1_{\mathcal{O}_Y}(\mathcal{O}_X, \mathcal{O}_X)$의 절단\  $c(s)$를
정의한다. Cohomology, Lemma \ref{cohomology-lemma-section-RHom-over-U}와\ 
Derived Categories, Lemma \ref{derived-lemma-ext-1}을 보라. 역으로\ 
$\mathcal{O}_U$-가군의 확장
$$
0 \to \mathcal{O}_X|_U \to \mathcal{E} \to \mathcal{O}_X|_U \to 0
$$
이 주어졌다고 하자. 열린 덮개\  $U = \bigcup U_i$와\ 
$1 \in \mathcal{O}_X(U_i)$로 가는 절단\  $e_i \in \mathcal{E}(U_i)$를 찾을 수
있다. 그러면\  $e_i$는 사상\  $\mathcal{O}_Y|_{U_i} \to \mathcal{E}|_{U_i}$를
정의하고 그 핵은\  $\mathcal{I}^2$을 포함한다. 이와 같이\ 
$\mathcal{E}|_{U_i}$가 위와 같은 밀어내기에서 나옴을 알 수 있다.
이는\  $c$가 전사임을 보인다. 단사성의 증명은 생략한다.
\end{proof}
\begin{lemma}
\label{duality-lemma-regular-ideal-ext}
$Y$를 환 달린 공간이라 하자. $\mathcal{I} \subset \mathcal{O}_Y$를
아이디얼 층이라 하고\  $\mathcal{O}_X = \mathcal{O}_Y/\mathcal{I}$라 두자.
$\mathcal{I}$가\  Koszul-정칙이면
(Divisors, Definition \ref{divisors-definition-regular-ideal-sheaf}),
$R\SheafHom_{\mathcal{O}_Y}(\mathcal{O}_X, \mathcal{O}_X)$ 위의 합성은 모든\ 
$i$에 대해 동형
$$
\wedge^i(\SheafExt^1_{\mathcal{O}_Y}(\mathcal{O}_X, \mathcal{O}_X))
\longrightarrow
\SheafExt^i_{\mathcal{O}_Y}(\mathcal{O}_X, \mathcal{O}_X)
$$
을 정의한다.
\end{lemma}

\begin{proof}
합성이란\  Cohomology, Lemma
\ref{cohomology-lemma-internal-hom-composition}의 사상
$$
R\SheafHom_{\mathcal{O}_Y}(\mathcal{O}_X, \mathcal{O}_X)
\otimes_{\mathcal{O}_Y}^\mathbf{L}
R\SheafHom_{\mathcal{O}_Y}(\mathcal{O}_X, \mathcal{O}_X)
\longrightarrow
R\SheafHom_{\mathcal{O}_Y}(\mathcal{O}_X, \mathcal{O}_X)
$$
을 뜻한다. 이는 곱셈 사상
$$
\SheafExt^a_{\mathcal{O}_Y}(\mathcal{O}_X, \mathcal{O}_X)
\otimes_{\mathcal{O}_Y}
\SheafExt^b_{\mathcal{O}_Y}(\mathcal{O}_X, \mathcal{O}_X)
\longrightarrow
\SheafExt^{a + b}_{\mathcal{O}_Y}(\mathcal{O}_X, \mathcal{O}_X)
$$
들을 유도한다. More on Algebra, Equation
(\ref{more-algebra-equation-simple-tor-product})과 비교하라. 보조정리의 명제는
유도된 사상
$$
\SheafExt^1_{\mathcal{O}_Y}(\mathcal{O}_X, \mathcal{O}_X)
\otimes \ldots \otimes
\SheafExt^1_{\mathcal{O}_Y}(\mathcal{O}_X, \mathcal{O}_X)
\longrightarrow
\SheafExt^i_{\mathcal{O}_Y}(\mathcal{O}_X, \mathcal{O}_X)
$$
이 외승을 통해 인수분해되고 이어서 동형을 유도한다는 뜻이다. 이를 보이기
위해\  $Y$ 위에서 국소적으로 작업해도 된다. 따라서\  $\mathcal{O}_Y$의 대역절단\ 
$f_1, \ldots, f_r$이\  $\mathcal{I}$를 생성하고\  Koszul 정칙열을 이룬다고
가정해도 된다. 다음 층을 생각하자.
$$
\mathcal{A} = \mathcal{O}_Y\langle \xi_1, \ldots, \xi_r\rangle
$$
이는 엄밀 가환 미분 등급\  $\mathcal{O}_Y$-대수의 층으로서,
$\mathcal{O}_Y$ 위에서 차수\  $-1$인\  $\xi_1, \ldots, \xi_r$에 대한
(분할멱) 다항식 대수이고 미분\  $\text{d}$는 규칙\  $\text{d}\xi_i = f_i$로
주어진다. $\mathcal{A}^\bullet$을 그 기저\  $\mathcal{O}_Y$-가군 복합체라
쓰자. 이는 위에서 말한\  Koszul 복합체이다. 따라서 정준 사상\ 
$\mathcal{A}^\bullet \to \mathcal{O}_X$는 준동형이다. Cohomology, Lemma
\ref{cohomology-lemma-Rhom-strictly-perfect}에 의해 준동형들
$$
R\SheafHom_{\mathcal{O}_Y}(\mathcal{O}_X, \mathcal{O}_X) \to
\SheafHom^\bullet(\mathcal{A}^\bullet, \mathcal{A}^\bullet) \to
\SheafHom^\bullet(\mathcal{A}^\bullet, \mathcal{O}_X)
$$
을 얻는다. 마지막 복합체의 미분들은 영이고, 따라서
$$
\SheafExt^i_{\mathcal{O}_Y}(\mathcal{O}_X, \mathcal{O}_X)
\cong \SheafHom_{\mathcal{O}_Y}(\mathcal{A}^{-i}, \mathcal{O}_X)
$$
이다. $j \in \{1, \ldots, r\}$에 대해\  $\delta_j : \mathcal{A} \to \mathcal{A}$를\ 
$\delta_j(\xi_i) = \delta_{ij}$인 차수\  $1$의 미분법이라 하자(Kronecker 델타).
계산하면\  $\delta_j \circ \text{d} = - \text{d} \circ \delta_j$이고, 따라서
복합체의 사상
$$
\delta_j : \mathcal{A}^\bullet \to \mathcal{A}^\bullet[1].
$$
을 얻는다. 그러므로\  $\delta_j$는 대응하는\  $\SheafExt$-층의 절단을 정의한다.
다른 계산에 의해\  $\delta_1, \ldots, \delta_r$은\  $\mathcal{O}_X$ 위에서\ 
$\SheafHom_{\mathcal{O}_Y}(\mathcal{A}^{-1}, \mathcal{O}_X)$의 기저로 간다.
$\mathcal{A}$의 자기준동형으로서, 따라서\  $\SheafExt$-층들에서도\ 
$\delta_j \circ \delta_j = 0$이고\ 
$\delta_j \circ \delta_{j'} = - \delta_{j'} \circ \delta_j$임이 명백하므로,
위의 사상이 외승을 통해 인수분해된다는 명제를 얻는다. 원하는 동형을 얻는다는
것을 보이기 위해서는 독자가\  $j_1 < \ldots < j_i$에 대한 원소들
$$
\delta_{j_1} \circ \ldots \circ \delta_{j_i}
$$
이\  $\mathcal{O}_X$ 위에서 층\ 
$\SheafHom_{\mathcal{O}_Y}(\mathcal{A}^{-i}, \mathcal{O}_X)$의 기저로 감을
확인하면 된다.
\end{proof}
\begin{lemma}
\label{duality-lemma-regular-immersion-ext}
$Y$를 환 달린 공간이라 하자. $\mathcal{I} \subset \mathcal{O}_Y$를
아이디얼 층이라 하자. $\mathcal{O}_X = \mathcal{O}_Y/\mathcal{I}$라 두고\ 
$\mathcal{N} =
\SheafHom_{\mathcal{O}_Y}(\mathcal{I}/\mathcal{I}^2, \mathcal{O}_X)$라 두자.
$\mathcal{I}$가\  Koszul-정칙이면
(Divisors, Definition \ref{divisors-definition-regular-ideal-sheaf}),
$$
R\SheafHom_{\mathcal{O}_Y}(\mathcal{O}_X, \mathcal{O}_Y) =
\wedge^r \mathcal{N}[-r]
$$
이다. 여기서\  $r : Y \to \{1, 2, 3, \ldots \}$은\  $y$를\  $y$의 근방에서\ 
$\mathcal{I}$를 생성하는 데 필요한 최소 생성원 수로 보내는 함수이다.
\end{lemma}

\begin{proof}
Lemmas \ref{duality-lemma-ext}와\  \ref{duality-lemma-regular-ideal-ext}를 사용하면\  $i \geq 0$에
대해 동형들\  $\wedge^i\mathcal{N} \to
\SheafExt^i_{\mathcal{O}_Y}(\mathcal{O}_X, \mathcal{O}_X)$을 얻는다. 따라서\ 
$\mathcal{O}_Y \to \mathcal{O}_X$가 동형
$$
\SheafExt^r_{\mathcal{O}_Y}(\mathcal{O}_X, \mathcal{O}_Y)
\longrightarrow
\SheafExt^r_{\mathcal{O}_Y}(\mathcal{O}_X, \mathcal{O}_X)
$$
을 유도하고, $i \not = r$이면\ 
$\SheafExt^i_{\mathcal{O}_Y}(\mathcal{O}_X, \mathcal{O}_Y)$가 영임을 보이면
충분하다. 이 명제들은\  $Y$ 위에서 국소적이다. 따라서\  $\mathcal{O}_Y$의
대역절단\  $f_1, \ldots, f_r$이\  $\mathcal{I}$를 생성하고\  Koszul 정칙열을
이룬다고 가정해도 된다. $\mathcal{A}^\bullet$을\  Lemma
\ref{duality-lemma-regular-ideal-ext}의 증명에서 도입한\  $f_1, \ldots, f_r$에 대한\ 
Koszul 복합체라 하자. 그러면\  Cohomology, Lemma
\ref{cohomology-lemma-Rhom-strictly-perfect}에 의해
$$
R\SheafHom_{\mathcal{O}_Y}(\mathcal{O}_X, \mathcal{O}_Y) =
\SheafHom^\bullet(\mathcal{A}^\bullet, \mathcal{O}_Y)
$$
이다. $1 : \mathcal{A}^\bullet \to \mathcal{O}_Y$을 차수\  $0$에서의 항등사상\ 
$\mathcal{A}^0 = \mathcal{O}_Y \to \mathcal{O}_Y$로 주어지는 미분 등급\ 
$\mathcal{O}_Y$-대수의 사상이라 하자. Lemma
\ref{duality-lemma-regular-ideal-ext}의 증명에서와 같은\  $\delta_j$를 사용하면,
차수\  $i$의\  $\delta_{j_1} \ldots \delta_{j_i}$를\ 
$1 \circ \delta_{j_1} \circ \ldots \circ \delta_{j_i}$로 보내어 등급\ 
$\mathcal{O}_Y$-가군의 동형
$$
\mathcal{O}_Y\langle \delta_1, \ldots, \delta_r\rangle
\longrightarrow
\SheafHom^\bullet(\mathcal{A}^\bullet, \mathcal{O}_Y)
$$
을 얻는다. 이 동형을 통해 오른쪽의 미분은 왼쪽의 미분\  $\text{d}$를 유도한다.
우리의 부호 규칙에 의해\  $\text{d}(1) = - \sum f_j \delta_j$이다.
$\delta_j : \mathcal{A}^\bullet \to \mathcal{A}^\bullet[1]$은 복합체의
사상이므로
$$
\text{d}(\delta_{j_1} \ldots \delta_{j_i}) =
(- \sum f_j \delta_j )\delta_{j_1} \ldots \delta_{j_i}
$$
임이 따른다. 미분 등급 대수\  $\mathcal{A}$ 위에서는\ 
$\text{d} = \sum f_j \delta_j$임을 관찰하자. 따라서 규칙
$$
1 \circ \delta_{j_1} \ldots \delta_{j_i} \longmapsto
(\delta_{j_1} \circ \ldots \circ \delta_{j_i})(\xi_1 \ldots \xi_r)
$$
으로 정의한 사상은\  $r$이 홀수이면 복합체의 동형
$$
\SheafHom^\bullet(\mathcal{A}^\bullet, \mathcal{O}_Y)
\longrightarrow \mathcal{A}^\bullet[-r]
$$
을 정의하고, $r$이 짝수이면 부호를 빼고 미분과 가환한다. 어느 경우든 이
복합체들의 코호몰로지는 동형이며, 이는 원하는 소멸을 보인다. 사상
$$
\SheafHom^\bullet(\mathcal{A}^\bullet, \mathcal{O}_Y)
\longrightarrow
\SheafHom^\bullet(\mathcal{A}^\bullet, \mathcal{O}_X)
$$
아래에서 차수\  $r$의 코호몰로지 위의 동형도 여기서 곧바로 따른다.
(Koszul 복합체에 관한 우리의 규약 선택은 동형\ 
$R\SheafHom_{\mathcal{O}_X}(\mathcal{O}_X, \mathcal{O}_Y) =
\wedge^r \mathcal{N}[-r]$의 정의에 실제로 관여함을 유의하자.)
\end{proof}
\begin{lemma}
\label{duality-lemma-regular-immersion}
$Y$를 준콤팩트 준분리 스킴이라 하자. $i : X \to Y$를\  Koszul-정칙 닫힌
몰입이라 하자. $a$를\  Lemma \ref{duality-lemma-twisted-inverse-image}에서의\ 
$Ri_* : D_\QCoh(\mathcal{O}_X) \to D_\QCoh(\mathcal{O}_Y)$의 오른쪽 수반이라
하자. 그러면 동형
$$
\wedge^r\mathcal{N}[-r] \longrightarrow a(\mathcal{O}_Y)
$$
이 존재한다. 여기서\ 
$\mathcal{N} = \SheafHom_{\mathcal{O}_X}(\mathcal{C}_{X/Y}, \mathcal{O}_X)$는\ 
$i$의 법선층이고(Morphisms, Section
\ref{morphisms-section-conormal-sheaf}), $r$은 그 계수를\  $X$ 위의 국소상수
함수로 본 것이다.
\end{lemma}

\begin{proof}
Lemmas \ref{duality-lemma-twisted-inverse-image-closed}와\ 
\ref{duality-lemma-sheaf-with-exact-support-ext}에서\  $a(\mathcal{O}_Y)$는\ 
$D_\QCoh(\mathcal{O}_X)$의 대상이고, 그\  $Y$로의 추이는\ 
$R\SheafHom_{\mathcal{O}_Y}(i_*\mathcal{O}_X, \mathcal{O}_Y)$임을 상기하자.
따라서 결과는\  Lemma \ref{duality-lemma-regular-immersion-ext}에서 따른다.
\end{proof}
\begin{lemma}
\label{duality-lemma-smooth-proper}
$S$를\  Noether 스킴이라 하자. $f : X \to S$를 상대차원\  $d$인 매끄러운 고유
사상이라 하자. $a$를\  Lemma \ref{duality-lemma-twisted-inverse-image}에서와 같은\ 
$Rf_* : D_\QCoh(\mathcal{O}_X) \to D_\QCoh(\mathcal{O}_S)$의 오른쪽 수반이라
하자. 그러면\  $D(\mathcal{O}_X)$에서 동형
$$
\wedge^d \Omega_{X/S}[d] \longrightarrow a(\mathcal{O}_S)
$$
이 존재한다.
\end{lemma}

\begin{proof}
Remark \ref{duality-remark-relative-dualizing-complex}에서와 같이\ 
$\omega_{X/S}^\bullet = a(\mathcal{O}_S)$라 두자. $\Delta : X \to X \times_S X$에
대한\  Lemma \ref{duality-lemma-twisted-inverse-image}의 오른쪽 수반을\  $c$라 하자.
$\Delta$는 매끄러운 사상의 대각사상이므로\  Koszul-정칙 몰입이다. Divisors,
Lemma \ref{divisors-lemma-immersion-smooth-into-smooth-regular-immersion}을 보라.
특히\  $\Delta$는 완전 고유 사상이고(More on Morphisms, Lemma
\ref{more-morphisms-lemma-regular-immersion-perfect}), 다음을 얻는다.
\begin{align*}
\mathcal{O}_X
& =
c(L\text{pr}_1^*\omega_{X/S}^\bullet) \\
& =
L\Delta^*(L\text{pr}_1^*\omega_{X/S}^\bullet)
\otimes_{\mathcal{O}_X}^\mathbf{L}
c(\mathcal{O}_{X \times_S X}) \\
& =
\omega_{X/S}^\bullet \otimes_{\mathcal{O}_X}^\mathbf{L}
c(\mathcal{O}_{X \times_S X}) \\
& =
\omega_{X/S}^\bullet
\otimes_{\mathcal{O}_X}^\mathbf{L}
\wedge^d(\mathcal{N}_\Delta)[-d]
\end{align*}
첫 번째 등식은\  $\omega_{X/S}^\bullet = a(\mathcal{O}_S)$이므로
(\ref{duality-equation-pre-rigid})이다. 두 번째 등식은\  Lemma
\ref{duality-lemma-compare-with-pullback-flat-proper-noetherian}에 의한다. 세 번째
등식은\  $\text{pr}_1 \circ \Delta = \text{id}_X$이기 때문이다. 네 번째 등식은\ 
Lemma \ref{duality-lemma-regular-immersion}에 의한다.
$\wedge^d(\mathcal{N}_\Delta)$가 가역\  $\mathcal{O}_X$-가군임을 관찰하자.
따라서\  $\wedge^d(\mathcal{N}_\Delta)[-d]$는\  $D(\mathcal{O}_X)$의 가역 대상이고,
$a(\mathcal{O}_S) = \omega_{X/S}^\bullet = \wedge^d(\mathcal{C}_\Delta)[d]$라고
결론짓는다. $\Delta$의 여법선층\  $\mathcal{C}_\Delta$는\ 
Morphisms, Lemma \ref{morphisms-lemma-differentials-diagonal}에 의해\ 
$\Omega_{X/S}$이므로 증명이 끝난다.
\end{proof}






\section{상단 느낌표 함자}
\label{duality-section-upper-shriek}

\noindent
이 절에서는 고정된\  Noether 밑 위에서 유한형이고 분리인 스킴들 사이의 사상에
대해 콤팩트화를 사용하여 함자\  $f^!$를 구성한다. 연접 쌍대성에서 흔히 그렇듯이
가환함을 보여야 하는 도식들이 여럿 있다. 독자에게는 구성을 읽은 뒤 보조정리들의
검증은 건너뛰고 상단 느낌표 함자들의 성질을 논하는 다음 절로 넘어가기를 권한다.

\begin{situation}
\label{duality-situation-shriek}
여기서\  $S$는\  Noether 스킴이고\  $\textit{FTS}_S$는 다음 범주이다.
\begin{enumerate}
\item 대상은\  $S$ 위의 스킴\  $X$로서 구조사상\  $X \to S$가 분리이고 유한형인
것들이고,
\item 대상들 사이의 사상\  $f : X \to Y$는\  $S$ 위의 스킴 사상들이다.
\end{enumerate}
\end{situation}

\noindent
Situation \ref{duality-situation-shriek}에서\  $\textit{FTS}_S$의 사상\ 
$f : X \to Y$가 주어지면 삼각 범주의 완전함자
$$
f^! : D_\QCoh^+(\mathcal{O}_Y) \to D_\QCoh^+(\mathcal{O}_X)
$$
를 정의하겠다. 즉, More on Flatness, Theorem \ref{flat-theorem-nagata}와\ 
Lemma \ref{flat-lemma-compactifyable}에 의해 가능한\  $Y$ 위의 콤팩트화\ 
$X \to \overline{X}$를 선택한다. 구조사상을\ 
$\overline{f} : \overline{X} \to Y$라 쓰자. Lemma
\ref{duality-lemma-twisted-inverse-image}에서 구성한\  $R\overline{f}_*$의 오른쪽 수반을\ 
$\overline{a} : D_\QCoh(\mathcal{O}_Y) \to D_\QCoh(\mathcal{O}_{\overline{X}})$라
하자. 그러면\  $K \in D_\QCoh^+(\mathcal{O}_Y)$에 대해
$$
f^!K  = \overline{a}(K)|_X
$$
로 둔다. Lemma \ref{duality-lemma-twisted-inverse-image-bounded-below}에 의해 그 결과는\ 
$D_\QCoh^+(\mathcal{O}_X)$의 대상이다.
\begin{lemma}
\label{duality-lemma-shriek-well-defined}
Situation \ref{duality-situation-shriek}에서\  $f : X \to Y$를\  $\textit{FTS}_S$의 사상이라
하자. 함자\  $f^!$는 정준 동형을 제외하면 콤팩트화의 선택과 무관하다.
\end{lemma}

\begin{proof}
$Y$ 위에서\  $X$의 콤팩트화들이 이루는 범주는\  More on Flatness, Section
\ref{flat-section-compactify}에서 정의된다. More on Flatness, Theorem
\ref{flat-theorem-nagata}와\  Lemma \ref{flat-lemma-compactifyable}에 의해 이
범주는 공집합이 아니다. 콤팩트화의 각 선택
$$
j : X \to \overline{X},\quad \overline{f} : \overline{X} \to Y
$$
에 위 구성은 함자\  $j^* \circ \overline{a} :
D_\QCoh^+(\mathcal{O}_Y) \to D_\QCoh^+(\mathcal{O}_X)$를 결부시킨다. 여기서\ 
$\overline{a}$는\  Lemma \ref{duality-lemma-twisted-inverse-image}에서 구성한\ 
$R\overline{f}_*$의 오른쪽 수반이다.

\medskip\noindent
$Y$ 위의 콤팩트화들\  $j_i : X \to \overline{X}_i$ 사이의 사상\ 
$g : \overline{X}_1 \to \overline{X}_2$가 주어지고\ 
$g^{-1}(j_2(X)) = j_1(X)$라고 하자\footnote{우리의 콤팩트화 정의에서는 이것이
성립하지 않을 수도 있다. More on Flatness, Section
\ref{flat-section-compactify}를 보라.}. $g$에 대한\  Lemma
\ref{duality-lemma-twisted-inverse-image}의 오른쪽 수반을\  $\overline{c}$라 하자. 그러면
이 함자들은\ 
$R\overline{f}_{2, *} \circ Rg_* = R(\overline{f}_2 \circ g)_*$에 수반되므로\ 
$\overline{c} \circ \overline{a}_2 = \overline{a}_1$이다.
(\ref{duality-equation-sheafy})에 의해 함자\ 
$D^+_\QCoh(\mathcal{O}_{\overline{X}_2}) \to D^+_\QCoh(\mathcal{O}_X)$의 정준 변환
$$
j_1^* \circ \overline{c} \longrightarrow j_2^*
$$
이 있고, Lemma \ref{duality-lemma-proper-noetherian}에 의해 이는 동형이다. 합성
$$
j_1^* \circ \overline{a}_1 \longrightarrow
j_1^* \circ \overline{c} \circ \overline{a}_2 \longrightarrow
j_2^* \circ \overline{a}_2
$$
은 함자들의 동형이며 이를\  $\alpha_g$라 쓰겠다.

\medskip\noindent
$Y$ 위에서\  $X$의 두 콤팩트화\  $j_i : X \to \overline{X}_i$, $i = 1, 2$를
생각하자. More on Flatness, Lemma
\ref{flat-lemma-compactifications-cofiltered}의 (b)에 의해 상이 조밀한 콤팩트화\ 
$j : X \to \overline{X}$와 콤팩트화들의 사상\ 
$g_i : \overline{X} \to \overline{X}_i$를 찾을 수 있다. More on Flatness,
Lemma \ref{flat-lemma-compactifications-cofiltered}의 (c)에 의해\ 
$g_i^{-1}(j_i(X)) = j(X)$이다. 따라서 앞 단락으로부터 동형들
$$
\alpha_{g_i} :
j^* \circ \overline{a}
\longrightarrow
j_i^* \circ \overline{a}_i
$$
을 얻는다. 이로부터 동형
$$
\alpha_{g_2} \circ \alpha_{g_1}^{-1} :
j_1^* \circ \overline{a}_1 \to j_2^* \circ \overline{a}_2
$$
을 얻는다. 증명을 끝내려면 이 동형들이 잘 정의됨을 보여야 한다. 앞 단락에서
구성한 동형들의 합성이 다시 그러한 동형임을 보이면 충분하다고 주장한다(정확한
명제는 다음 단락을 보라). 독자는 냅킨 위에서 이것이 참임을 확인해도 되지만,
이 단락의 나머지에서 모두 자세히 적겠다. 즉, 상이 조밀한 콤팩트화의 두 번째 선택\ 
$j' : X \to \overline{X}'$과 콤팩트화들의 사상\ 
$g'_i : \overline{X}' \to \overline{X}_i$를 생각하자. More on Flatness, Lemma
\ref{flat-lemma-compactifications-cofiltered}에 의해 상이 조밀한 콤팩트화\ 
$j'' : X \to \overline{X}''$과 콤팩트화들의 사상\ 
$h : \overline{X}'' \to \overline{X}$ 및\ 
$h' : \overline{X}'' \to \overline{X}'$을 찾을 수 있다. 더 나아가\ 
$g_1 \circ h = g'_1 \circ h'$이고\  $g_2 \circ h = g'_2 \circ h'$이라고 가정해도
된다. 다음 단락의 결과에 의해
$$
\alpha_{g_i} \circ \alpha_h = \alpha_{g_i \circ h} =
\alpha_{g'_i \circ h'} = \alpha_{g'_i} \circ \alpha_{h'}
$$
이다. 여기서\  $i = 1, 2$이다. 이들은 모두 함자들의 동형이므로 원하는 대로\ 
$\alpha_{g_2} \circ \alpha_{g_1}^{-1} =
\alpha_{g'_2} \circ \alpha_{g'_1}^{-1}$라고 결론짓는다.

\medskip\noindent
$i = 1, 2, 3$에 대해 콤팩트화\  $j_i : X \to \overline{X}_i$들이 주어졌다고
하자. 콤팩트화들의 사상\  $g : \overline{X}_1 \to \overline{X}_2$와\ 
$h : \overline{X}_2 \to \overline{X}_3$가 주어지고\ 
$g^{-1}(j_2(X)) = j_1(X)$ 및\  $h^{-1}(j_2(X)) = j_3(X)$라고 하자.
$\overline{a}_i$를 위와 같이 두자. 위 주장은
$$
\alpha_g \circ \alpha_h = \alpha_{g \circ h} :
j_1^* \circ \overline{a}_1 \to j_3^* \circ \overline{a}_3
$$
라는 뜻이다. $g$와\  $h$에 대한\  Lemma
\ref{duality-lemma-twisted-inverse-image}의 오른쪽 수반을 각각\  $\overline{c}$와\ 
$\overline{d}$라 하자. 그러면\ 
$\overline{c} \circ \overline{a}_2 = \overline{a}_1$이고\ 
$\overline{d} \circ \overline{a}_3 = \overline{a}_2$이며, 위와 같은 이유로
함자들의 정준 변환
$$
j_1^* \circ \overline{c} \longrightarrow j_2^*
\quad\text{and}\quad
j_2^* \circ \overline{d} \longrightarrow j_3^*
$$
이 각각\ 
$D^+_\QCoh(\mathcal{O}_{\overline{X}_2}) \to D^+_\QCoh(\mathcal{O}_X)$와\ 
$D^+_\QCoh(\mathcal{O}_{\overline{X}_3}) \to D^+_\QCoh(\mathcal{O}_X)$의 함자로서
존재한다. $h \circ g$에 대한\  Lemma
\ref{duality-lemma-twisted-inverse-image}의 오른쪽 수반을\  $\overline{e}$라 하자.
(\ref{duality-equation-sheafy})로 주어지는 함자\ 
$D^+_\QCoh(\mathcal{O}_{\overline{X}_3}) \to D^+_\QCoh(\mathcal{O}_X)$의 정준 변환
$$
j_1^* \circ \overline{e} \longrightarrow j_3^*
$$
이 있다. 모든 것을 자세히 적으면 합성
$$
\alpha_h \circ \alpha_g :
j_1^* \circ \overline{a}_1 \to
j_1^* \circ \overline{c} \circ \overline{a}_2 \to
j_2^* \circ \overline{a}_2 \to
j_2^* \circ \overline{d} \circ \overline{a}_3 \to
j_3^* \circ \overline{a}_3
$$
이 합성
$$
\alpha_{h \circ g} :
j_1^* \circ \overline{a}_1 \to
j_1^* \circ \overline{e} \circ \overline{a}_3 \to
j_3^* \circ \overline{a}_3
$$
과 같음을 보여야 한다. 이를 두 부분으로 나눈다. 첫 번째는 도식
$$
\xymatrix{
\overline{a}_1 \ar[r] \ar[d] & \overline{c} \circ \overline{a}_2 \ar[d] \\
\overline{e} \circ \overline{a}_3 \ar[r] &
\overline{c} \circ \overline{d} \circ \overline{a}_3
}
$$
이 가환함을 보이는 것이다. 여기서 아래 수평 화살표는 식별\ 
$\overline{e} = \overline{c} \circ \overline{d}$에서 나온다. 이에 대응하는 전체
직접상 함자들의 도식
$$
\xymatrix{
R\overline{f}_{1, *} \ar[r] \ar[d] & Rg_* \circ R\overline{f}_{2, *} \ar[d] \\
R(h \circ g)_* \circ R\overline{f}_{3, *} \ar[r] &
Rg_* \circ Rh_* \circ R\overline{f}_{3, *}
}
$$
이 가환하므로 이는 참이다(향후 참조를 여기에 삽입할 것). 두 번째 부분은 합성
$$
j_1^* \circ \overline{c} \circ \overline{d} \to
j_2^* \circ \overline{d} \to j_3^*
$$
이 식별\  $\overline{e} = \overline{c} \circ \overline{d}$를 통해 사상
$$
j_1^* \circ \overline{e} \to j_3^*
$$
과 같음을 보이는 것이다. 이는\  Lemma
\ref{duality-lemma-compose-base-change-maps}에서 증명하였다(현재의 경우 그 보조정리의
사상\  $f', g'$은\  $\text{id}_X$와 같음을 유의하라).
\end{proof}
\begin{lemma}
\label{duality-lemma-upper-shriek-composition}
Situation \ref{duality-situation-shriek}에서\  $f : X \to Y$와\  $g : Y \to Z$를\ 
$\textit{FTS}_S$의 합성 가능한 사상들이라 하자. 정준 동형\ 
$(g \circ f)^! \to f^! \circ g^!$이 존재한다.
\end{lemma}

\begin{proof}
$Z$ 위에서\  $Y$의 콤팩트화\  $i : Y \to \overline{Y}$를 선택하고,
$\overline{Y}$ 위에서\  $X$의 콤팩트화\  $X \to \overline{X}$를 선택한다. 여기서는\ 
More on Flatness, Theorem \ref{flat-theorem-nagata}와\  Lemma
\ref{flat-lemma-compactifyable}를 두 번 사용한다.
$\overline{X} \to \overline{Y}$에 대한\  Lemma
\ref{duality-lemma-twisted-inverse-image}의 오른쪽 수반을\  $\overline{a}$라 하고,
$\overline{Y} \to Z$에 대한\  Lemma \ref{duality-lemma-twisted-inverse-image}의 오른쪽
수반을\  $\overline{b}$라 하자. 그러면\  $\overline{a} \circ \overline{b}$는 합성\ 
$\overline{X} \to Z$에 대한\  Lemma \ref{duality-lemma-twisted-inverse-image}의 오른쪽
수반이다. 따라서\  $g^! = i^* \circ \overline{b}$이고\ 
$(g \circ f)^! = (X \to \overline{X})^* \circ \overline{a} \circ \overline{b}$이다.
$U$를\  $\overline{X}$ 안에서\  $Y$의 역상이라 하면 가환도식
$$
\xymatrix{
X \ar[r]_j \ar[d] & U \ar[dl] \ar[r]_{j'} & \overline{X} \ar[dl] \\
Y \ar[r]_i \ar[d] & \overline{Y} \ar[dl] \\
Z
}
$$
을 얻는다. $U \to Y$에 대한\  Lemma \ref{duality-lemma-twisted-inverse-image}의 오른쪽
수반을\  $\overline{a}'$라 하자. 그러면\  $f^! = j^* \circ \overline{a}'$이다.
(\ref{duality-equation-sheafy})에 의해
$$
\gamma : (j')^* \circ \overline{a} \to \overline{a}' \circ i^*
$$
를 얻고, 이를 사용하여
$$
(g \circ f)^! =
(j' \circ j)^* \circ \overline{a} \circ \overline{b} =
j^* \circ (j')^* \circ \overline{a} \circ \overline{b}
\to
j^* \circ \overline{a}' \circ i^* \circ \overline{b} =
f^! \circ g^!
$$
를 정의할 수 있다. Lemma \ref{duality-lemma-proper-noetherian}에 의해 이는\ 
$D_\QCoh^+(\mathcal{O}_Z)$의 대상들 위에서 동형이다. 증명을 끝내기 위해 이
동형이 선택들과 무관함을 보이겠다.

\medskip\noindent
두 도식
$$
\vcenter{
\xymatrix{
X \ar[r]_{j_1} \ar[d] & U_1 \ar[dl] \ar[r]_{j'_1} & \overline{X}_1 \ar[dl] \\
Y \ar[r]_{i_1} \ar[d] & \overline{Y}_1 \ar[dl] \\
Z
}
}
\quad\text{and}\quad
\vcenter{
\xymatrix{
X \ar[r]_{j_2} \ar[d] & U_2 \ar[dl] \ar[r]_{j'_2} & \overline{X}_2 \ar[dl] \\
Y \ar[r]_{i_2} \ar[d] & \overline{Y}_2 \ar[dl] \\
Z
}
}
$$
이 있다고 하자. 먼저\  $Z$ 위에서 상이 조밀하고\  $\overline{Y}_1$과\ 
$\overline{Y}_2$를 모두 지배하는\  $Y$의 콤팩트화\ 
$i : Y \to \overline{Y}$를 선택할 수 있다. More on Flatness, Lemma
\ref{flat-lemma-compactifications-cofiltered}를 보라. More on Flatness, Lemma
\ref{flat-lemma-right-multiplicative-system}과\  Categories, Lemmas
\ref{categories-lemma-morphisms-right-localization} 및\ 
\ref{categories-lemma-equality-morphisms-right-localization}에 의해,
$\overline{Y}$ 위에서 상이 조밀한\  $X$의 콤팩트화\  $X \to \overline{X}$를 다음과
같이 선택할 수 있다. 사상들\  $\overline{X} \to \overline{X}_1$과\ 
$\overline{X} \to \overline{X}_2$가 존재하며, 합성\ 
$\overline{X} \to \overline{Y} \to \overline{Y}_1$은 합성\ 
$\overline{X} \to \overline{X}_1 \to \overline{Y}_1$과 같고, 합성\ 
$\overline{X} \to \overline{Y} \to \overline{Y}_2$는 합성\ 
$\overline{X} \to \overline{X}_2 \to \overline{Y}_2$와 같다. 따라서 다음과
같은 가환도식이 있는 경우에 우리 도식들이 정하는 사상들을 비교하면 충분하다.
$$
\xymatrix{
X \ar[rr]_{j_1} \ar@{=}[d] & &
U_1 \ar[d] \ar[ddll] \ar[rr]_{j'_1} & &
\overline{X}_1 \ar[d] \ar[ddll] \\
X \ar'[r][rr]^-{j_2} \ar[d] & &
U_2 \ar'[dl][ddll] \ar'[r][rr]^-{j'_2} & &
\overline{X}_2 \ar[ddll] \\
Y \ar[rr]^{i_1} \ar@{=}[d] & & \overline{Y}_1 \ar[d] \\
Y \ar[rr]^{i_2} \ar[d] & & \overline{Y}_2 \ar[dll] \\
Z
}
$$
더욱이 콤팩트화들\  $X \to \overline{X}_1$과\  $Y \to \overline{Y}_2$의 상이
조밀하다고 하자. $\overline{X}_i \to \overline{Y}_i$,
$U_i \to Y$, $\overline{X}_1 \to \overline{X}_2$, $U_1 \to U_2$에 대한\  Lemma
\ref{duality-lemma-twisted-inverse-image}의 오른쪽 수반을 각각\  $\overline{a}_i$,
$\overline{a}'_i$, $\overline{c}$, $\overline{c}'$라 쓰겠다. 도식의 각 정사각형
$$
\xymatrix{
X \ar[r] \ar[d] \ar@{}[dr]|A & U_1 \ar[d] \\
X \ar[r] & U_2
}
\quad
\xymatrix{
U_2 \ar[r] \ar[d] \ar@{}[dr]|B & \overline{X}_2 \ar[d] \\
Y \ar[r] & \overline{Y}_2
}
\quad
\xymatrix{
U_1 \ar[r] \ar[d] \ar@{}[dr]|C & \overline{X}_1 \ar[d] \\
Y \ar[r] & \overline{Y}_1
}
\quad
\xymatrix{
Y \ar[r] \ar[d] \ar@{}[dr]|D & \overline{Y}_1 \ar[d] \\
Y \ar[r] & \overline{Y}_2
}
\quad
\xymatrix{
X \ar[r] \ar[d] \ar@{}[dr]|E & \overline{X}_1 \ar[d] \\
X \ar[r] & \overline{X}_2
}
$$
은\  Cartesian이다(A, D, E에 대해서는\  More on Flatness, Lemma
\ref{flat-lemma-compactifications-cofiltered}의 (c)를 보고, B, C에 대해서는\ 
$U_i$가\  $\overline{X}_i \to \overline{Y}_i$에 의한\  $Y$의 역상임을 상기하라).
따라서 각 정사각형은 다음과 같이 기저변환 사상
(\ref{duality-equation-sheafy})를 유도한다.
$$
\begin{matrix}
\gamma_A : j_1^* \circ \overline{c}' \to j_2^* &
\gamma_B : (j_2')^* \circ \overline{a}_2 \to \overline{a}'_2 \circ i_2^* &
\gamma_C : (j_1')^* \circ \overline{a}_1 \to \overline{a}'_1 \circ i_1^* \\
\gamma_D : i_1^* \circ \overline{d} \to i_2^* &
\gamma_E : (j'_1 \circ j_1)^* \circ \overline{c} \to (j'_2 \circ j_2)^*
\end{matrix}
$$
$f_1^! = j_1^* \circ \overline{a}'_1$,
$f_2^! = j_2^* \circ \overline{a}'_2$,
$g_1^! = i_1^* \circ \overline{b}_1$,
$g_2^! = i_2^* \circ \overline{b}_2$,
$(g \circ f)_1^! =
(j_1' \circ j_1)^* \circ \overline{a}_1 \circ \overline{b}_1$, 그리고\ 
$(g \circ f)^!_2 =
(j_2' \circ j_2)^* \circ \overline{a}_2 \circ \overline{b}_2$라 쓰자.
증명의 첫 단락과\  Lemma \ref{duality-lemma-shriek-well-defined}에서 주어진 구성은
\begin{enumerate}
\item $(g \circ f)^!_1 \to f_1^! \circ g_1^!$인 사상에\  $\gamma_C$를,
\item $(g \circ f)^!_2 \to f_2^! \circ g_2^!$인 사상에\  $\gamma_B$를,
\item $f_1^! \to f_2^!$인 사상에\  $\gamma_A$를,
\item $g_1^! \to g_2^!$인 사상에\  $\gamma_D$를, 그리고
\item $(g \circ f)^!_1 \to (g \circ f)^!_2$인 사상에\  $\gamma_E$를
\end{enumerate}
사용한다. 도식
$$
\xymatrix{
(g \circ f)^!_1 \ar[r]_{\gamma_E} \ar[d]_{\gamma_C} &
(g \circ f)^!_2 \ar[d]_{\gamma_B} \\
f_1^! \circ g_1^! \ar[r]^{\gamma_A \circ \gamma_D} & f_2^! \circ g_2^!
}
$$
이 가환함을 보여야 한다. Lemmas \ref{duality-lemma-compose-base-change-maps}와\ 
\ref{duality-lemma-compose-base-change-maps-horizontal}를 사용하며, Remark
\ref{duality-remark-going-around}에서와 같은 표기를 약간 남용하겠다(특히 항등 변환과의\ 
$\star$ 곱은 표기에서 생략한다). 다음 정사각형을 두면
$$
\xymatrix{
U_1 \ar[r] \ar[d] \ar@{}[rd]|F & \overline{X}_1 \ar[d] \\
U_2 \ar[r] & \overline{X}_2
}
$$
$\gamma_E = \gamma_A \circ \gamma_F$라고 쓸 수 있다. 따라서
$$
\gamma_B \circ \gamma_E = \gamma_B \circ \gamma_A  \circ \gamma_F
= \gamma_A \circ \gamma_B \circ \gamma_F
$$
임을 알 수 있다. 마지막 등식은 두 정사각형\  $A$와\  $B$가 한 점에서만 만나기
때문이다(Remark \ref{duality-remark-going-around}의 마지막 논증과 유사하다). 그러므로\ 
$\gamma_D \circ \gamma_C = \gamma_B \circ \gamma_F$임을 보이면 충분하다.
이 둘은 모두 정사각형
$$
\xymatrix{
U_1 \ar[r] \ar[d] & \overline{X}_1 \ar[d] \\
Y \ar[r] & \overline{Y}_2
}
$$
에 대한 사상 (\ref{duality-equation-sheafy})과 같으므로 결론을 얻는다.
\end{proof}
\begin{lemma}
\label{duality-lemma-pseudo-functor}
Situation \ref{duality-situation-shriek}에서\  Lemmas
\ref{duality-lemma-shriek-well-defined}와\  \ref{duality-lemma-upper-shriek-composition}의 구성들은
범주\  $\textit{FTS}_S$에서 범주들의\  $2$-범주로 가는 유사함자를 정의한다
(Categories, Definition \ref{categories-definition-functor-into-2-category}을 보라).
\end{lemma}

\begin{proof}
이를 보이려면 사상들\  $f : X \to Y$, $g : Y \to Z$, $h : Z \to T$가 주어졌을
때 도식
$$
\xymatrix{
(h \circ g \circ f)^! \ar[r]_{\gamma_{A + B}} \ar[d]_{\gamma_{B + C}} &
f^! \circ (h \circ g)^! \ar[d]^{\gamma_C} \\
(g \circ f)^! \circ h^! \ar[r]^{\gamma_A} & f^! \circ g^! \circ h^!
}
$$
이 가환함을 증명해야 한다($\gamma$들의 의미는 아래를 보라). 이를 위해 먼저\ 
$T$ 위에서\  $Z$의 콤팩트화\  $\overline{Z}$를 선택하고, 이어서\ 
$\overline{Z}$ 위에서\  $Y$의 콤팩트화\  $\overline{Y}$를 선택한 다음,
$\overline{Y}$ 위에서\  $X$의 콤팩트화\  $\overline{X}$를 선택한다. 여기서는\  More on
Flatness, Theorem \ref{flat-theorem-nagata}와\  Lemma
\ref{flat-lemma-compactifyable}를 사용한다. $W \subset \overline{Y}$를\ 
$\overline{Y} \to \overline{Z}$에 의한\  $Z$의 역상이라 하고,
$U \subset V \subset \overline{X}$를\  $\overline{X} \to \overline{Y}$에 의한\ 
$Y \subset W$의 역상들이라 하자. 그러면 다음 도식을 얻는다.
$$
\xymatrix{
X \ar[d]_f \ar[r] & U \ar[r] \ar[d] \ar@{}[dr]|A &
V \ar[d] \ar[r] \ar@{}[rd]|B & \overline{X} \ar[d] \\
Y \ar[d]_g \ar[r] & Y \ar[r] \ar[d] & W \ar[r] \ar[d] \ar@{}[rd]|C &
\overline{Y} \ar[d] \\
Z \ar[d]_h \ar[r] & Z \ar[d] \ar[r] & Z \ar[d] \ar[r] & \overline{Z} \ar[d] \\
T \ar[r] & T \ar[r] & T \ar[r] & T
}
$$
많은 표기를 도입하지 않고\  Lemma \ref{duality-lemma-upper-shriek-composition}의 증명과
정확히 같은 방식으로 논증하면, 첫 번째 표시 도식의 사상들은 표시된 직사각형\ 
$A + B$, $B + C$, $A$, $C$에 대한 사상 (\ref{duality-equation-sheafy})를 사용함을 알 수
있다. Lemmas \ref{duality-lemma-compose-base-change-maps}와\ 
\ref{duality-lemma-compose-base-change-maps-horizontal}에 의해\ 
$\gamma_{A + B} = \gamma_A \circ \gamma_B$이고\ 
$\gamma_{B + C} = \gamma_C \circ \gamma_B$이므로, 원하는 등식은\ 
$\gamma_A \circ \gamma_C = \gamma_C \circ \gamma_A$이면 성립한다. 이는 두
정사각형\  $A$와\  $C$가 한 점에서만 만나기 때문에 참이다(Remark
\ref{duality-remark-going-around}의 마지막 논증과 유사하다).
\end{proof}
\begin{lemma}
\label{duality-lemma-map-pullback-to-shriek-well-defined}
Situation \ref{duality-situation-shriek}에서\  $f : X \to Y$를\  $\textit{FTS}_S$의 사상이라
하자. $D^+_\QCoh(\mathcal{O}_Y)$의\  $K$에 함자적인 정준 사상들
$$
\mu_{f, K} :
Lf^*K \otimes_{\mathcal{O}_X}^\mathbf{L} f^!\mathcal{O}_Y
\longrightarrow
f^!K
$$
이 존재한다. $g : Y \to Z$가\  $\textit{FTS}_S$의 또 다른 사상이면, 모든\ 
$K \in D^+_\QCoh(\mathcal{O}_Z)$에 대해 도식
$$
\xymatrix{
Lf^*(Lg^*K \otimes_{\mathcal{O}_Y}^\mathbf{L} g^!\mathcal{O}_Z)
\otimes_{\mathcal{O}_X}^\mathbf{L} f^!\mathcal{O}_Y
\ar@{=}[d] \ar[r]_-{\mu_f} &
f^!(Lg^*K \otimes_{\mathcal{O}_Y}^\mathbf{L} g^!\mathcal{O}_Z)
\ar[r]_-{f^!\mu_g} &
f^!g^!K \ar@{=}[d] \\
Lf^*Lg^*K \otimes_{\mathcal{O}_X}^\mathbf{L} Lf^* g^!\mathcal{O}_Z
\otimes_{\mathcal{O}_X}^\mathbf{L} f^!\mathcal{O}_Y \ar[r]^-{\mu_f} &
Lf^*Lg^*K \otimes_{\mathcal{O}_X}^\mathbf{L} f^!g^!\mathcal{O}_Z
\ar[r]^-{\mu_{g \circ f}} & f^!g^!K
}
$$
은 가환한다.
\end{lemma}

\begin{proof}
$f$가 고유이면\  $f^! = a$이고 (\ref{duality-equation-compare-with-pullback})를 사용할 수
있다. $g$도 고유이면\  Lemma \ref{duality-lemma-transitivity-compare-with-pullback}가 위
도식의 가환성을 더 큰 일반성에서 증명한다.

\medskip\noindent
사상\  $\mu_{f, K}$를 정의하자. $Y$ 위에서\  $X$의 콤팩트화\ 
$j : X \to \overline{X}$를 선택하자. $f^!$가\  $j^* \circ \overline{a}$로 정의되므로,
사상 (\ref{duality-equation-compare-with-pullback})
$$
L\overline{f}^*K \otimes_{\mathcal{O}_{\overline{X}}}^\mathbf{L}
\overline{a}(\mathcal{O}_Y)
\longrightarrow
\overline{a}(K)
$$
을\  $X$에 제한하여\  $\mu_{f, K}$를 얻는다. 이것이 콤팩트화의 선택과 무관함을
보이기 위해\  Lemma \ref{duality-lemma-shriek-well-defined}의 증명에서처럼 논증한다.
독자는 먼저 그 보조정리의 증명을 읽기 바란다.

\medskip\noindent
$Y$ 위의 콤팩트화들\  $j_i : X \to \overline{X}_i$ 사이의 사상\ 
$g : \overline{X}_1 \to \overline{X}_2$가 주어지고\ 
$g^{-1}(j_2(X)) = j_1(X)$라고 하자. 사상\  $g$에 대한\  Lemma
\ref{duality-lemma-twisted-inverse-image}의 추이 함자의 오른쪽 수반을\  $\overline{c}$라
쓰자. 사상들
$$
L\overline{f}_1^*K \otimes_{\mathcal{O}_{\overline{X}}}^\mathbf{L}
\overline{a}_1(\mathcal{O}_Y)
\longrightarrow
\overline{a}_1(K)
\quad\text{and}\quad
L\overline{f}_2^*K \otimes_{\mathcal{O}_{\overline{X}}}^\mathbf{L}
\overline{a}_2(\mathcal{O}_Y)
\longrightarrow
\overline{a}_2(K)
$$
은\  Lemma \ref{duality-lemma-transitivity-compare-with-pullback}에 의해 가환도식
$$
\xymatrix{
Lg^*(L\overline{f}_2^*K \otimes^\mathbf{L}
\overline{a}_2(\mathcal{O}_Y))
\otimes^\mathbf{L} \overline{c}(\mathcal{O}_{\overline{X}_2})
\ar@{=}[d] \ar[r]_-\sigma &
\overline{c}(L\overline{f}_2^*K \otimes^\mathbf{L}
\overline{a}_2(\mathcal{O}_Y)) \ar[r] &
\overline{c}(\overline{a}_2(K)) \ar@{=}[d] \\
L\overline{f}_1^*K \otimes^\mathbf{L} Lg^*\overline{a}_2(\mathcal{O}_Y)
\otimes^\mathbf{L} \overline{c}(\mathcal{O}_{\overline{X}_2})
\ar[r]^-{1 \otimes \tau} &
L\overline{f}_1^*K \otimes^\mathbf{L} \overline{a}_1(\mathcal{O}_Y) \ar[r] &
\overline{a}_1(K)
}
$$
에 들어간다. Lemma \ref{duality-lemma-compare-on-open}에 의해 사상들\  $\sigma$와\  $\tau$는\ 
$X$ 위에서 동형으로 제한된다. 실제로 더 말할 수 있다. Lemma
\ref{duality-lemma-shriek-well-defined}의 증명에서 사상 (\ref{duality-equation-sheafy})
$\gamma : j_1^* \circ \overline{c} \to j_2^*$를 사용하여 동형\ 
$\alpha_g : j_1^* \circ \overline{a}_1 \to j_2^* \circ \overline{a}_2$를
구성했음을 상기하자. 사상\  $\sigma$를\  $j_1$로 당기고\ 
$j_1^* \circ \overline{c} \to j_2^*$를 통해\ 
$j_1^*\overline{c}(\mathcal{O}_{\overline{X}_2})$를\  $\mathcal{O}_X$와 식별하면,
$j_2^*\left(L\overline{f}_2^*K \otimes^\mathbf{L}
\overline{a}_2(\mathcal{O}_Y)\right)$ 위의 항등사상을 얻는다. Lemma
\ref{duality-lemma-restriction-compare-with-pullback}을 보라. 마찬가지로 사상\ 
$\tau : Lg^*\overline{a}_2(\mathcal{O}_Y)
\otimes^\mathbf{L} \overline{c}(\mathcal{O}_{\overline{X}_2}) \to
\overline{a}_1(\mathcal{O}_Y) = \overline{c}(\overline{a}_2(\mathcal{O}_Y))$를
당기면\  $j_2^*\overline{a}_2(\mathcal{O}_Y)$ 위의 항등사상이 된다. 따라서\ 
$j_1$로 당기고 가능한 곳마다\  $\gamma$를 적용하면 가환도식
$$
\xymatrix{
j_2^*\left(L\overline{f}_2^*K \otimes^\mathbf{L}
\overline{a}_2(\mathcal{O}_Y)\right) \ar[r] \ar[d] &
j_2^*\overline{a}_2(K) \\
j_1^*L\overline{f}_1^*K \otimes^\mathbf{L} j_2^*\overline{a}_2(\mathcal{O}_Y) &
j_1^*(L\overline{f}_1^*K \otimes^\mathbf{L} \overline{a}_1(\mathcal{O}_Y))
\ar[r] \ar[l]_{1 \otimes \alpha_g} &
j_1^* \overline{a}_1(K) \ar[lu]_{\alpha_g}
}
$$
을 얻는다. 이 도식의 가환성은 콤팩트화\  $\overline{X}_1$을 사용하여 구성한
사상\  $\mu_{f, K}$가\  Lemma \ref{duality-lemma-shriek-well-defined}의 증명에서 사용한 식별\ 
$\alpha_g$를 통해 콤팩트화\  $\overline{X}_2$를 사용하여 구성한 사상\ 
$\mu_{f, K}$와 정확히 같음을 말해 준다. 이제\  Lemma
\ref{duality-lemma-shriek-well-defined}의 증명에서와 정확히 같은 몇 가지 범주론적
논증에 의해\  $\mu_{f, K}$가 잘 정의됨을 알 수 있다(작은 세부사항은 생략한다).

\medskip\noindent
이제 우리 보조정리의 명제에 있는 도식의 가환성은 동형\ 
$(g \circ f)^! \to f^! \circ g^!$의 구성(콤팩트화들\ 
$\overline{X} \to \overline{Y} \to Z$를 사용한\  Lemma
\ref{duality-lemma-upper-shriek-composition}의 증명의 첫 부분)과\ 
$\overline{X} \to \overline{Y} \to Z$에 대한\  Lemma
\ref{duality-lemma-transitivity-compare-with-pullback}의 결과에서 따른다.
\end{proof}


\section{상단 느낌표 함자의 성질}
\label{duality-section-upper-shriek-properties}

\noindent
상단 느낌표 함자의 몇 가지 성질을 살펴보자.

\begin{lemma}
\label{duality-lemma-shriek-open-immersion}
Situation \ref{duality-situation-shriek}에서\  $Y$를\  $\textit{FTS}_S$의 대상이라 하고\ 
$j : X \to Y$를 열린 몰입이라 하자. 그러면 함자들의 정준 동형\ 
$j^! = j^*$가 존재한다.
\end{lemma}

\noindent
$\textit{FTS}_S$의 \'etale 사상\  $f : X \to Y$에 대해서도\ 
$f^* \cong f^!$이다. Lemma \ref{duality-lemma-shriek-etale}를 보라.

\begin{proof}
이 경우에는 콤팩트화로\  $\overline{X} = Y$를 선택할 수 있다. Lemma
\ref{duality-lemma-twisted-inverse-image}에서\  $\text{id} : Y \to Y$에 대한 오른쪽 수반은
항등 함자이므로, 정의에 의해\  $j^! = j^*$이다.
\end{proof}

\begin{lemma}
\label{duality-lemma-restrict-before-or-after}
Situation \ref{duality-situation-shriek}에서
$$
\xymatrix{
U \ar[r]_j \ar[d]_g & X \ar[d]^f \\
V \ar[r]^{j'} & Y
}
$$
를\  $\textit{FTS}_S$의 가환도식이라 하자. 여기서\  $j$와\  $j'$는 열린 몰입이다.
그러면 함자\  $D^+_\QCoh(\mathcal{O}_Y) \to D^+(\mathcal{O}_U)$로서\ 
$j^* \circ f^! = g^! \circ (j')^*$이다.
\end{lemma}

\begin{proof}
$h = f \circ j = j' \circ g$로 놓자. Lemma
\ref{duality-lemma-upper-shriek-composition}에 의해\ 
$h^! = j^! \circ f^! = g^! \circ (j')^!$이다. Lemma
\ref{duality-lemma-shriek-open-immersion}에 의해\  $j^! = j^*$이고\ 
$(j')^! = (j')^*$이다.
\end{proof}

\begin{lemma}
\label{duality-lemma-shriek-affine-line}
Situation \ref{duality-situation-shriek}에서\  $Y$를\  $\textit{FTS}_S$의 대상이라 하고\ 
$f : X = \mathbf{A}^1_Y \to Y$를 사영이라 하자. 그러면 함자들의 (비정준) 동형\ 
$f^!(-) \cong Lf^*(-) [1]$이 존재한다.
\end{lemma}

\begin{proof}
$X = \mathbf{A}^1_Y \subset \mathbf{P}^1_Y$이고\ 
$\mathcal{O}_{\mathbf{P}^1_Y}(-2)|_X \cong \mathcal{O}_X$이므로, 이는\  Lemmas
\ref{duality-lemma-upper-shriek-P1}와\ 
\ref{duality-lemma-compare-with-pullback-flat-proper-noetherian}에서 따른다.
\end{proof}

\begin{lemma}
\label{duality-lemma-shriek-closed-immersion}
Situation \ref{duality-situation-shriek}에서\  $Y$를\  $\textit{FTS}_S$의 대상이라 하고\ 
$i : X \to Y$를 닫힌 몰입이라 하자. 그러면 함자들의 정준 동형\ 
$i^!(-) = R\SheafHom(\mathcal{O}_X, -)$이 존재한다.
\end{lemma}

\begin{proof}
이는\  Lemma \ref{duality-lemma-twisted-inverse-image-closed}를 다시 서술한 것이다.
\end{proof}
\begin{remark}[상단 느낌표의 국소적 기술]
\label{duality-remark-local-calculation-shriek}
Situation \ref{duality-situation-shriek}에서\  $f : X \to Y$를\  $\textit{FTS}_S$의 사상이라 하자.
위의 보조정리들을 사용하면\  $f^!$를 다음과 같이 국소적으로 계산할 수 있다. 아핀
열린집합들의 도식
$$
\xymatrix{
U \ar[r]_j \ar[d]_g & X \ar[d]^f \\
V \ar[r]^i & Y
}
$$
이 주어졌다고 하자. $j^! \circ f^! = g^! \circ i^!$이고
(Lemma \ref{duality-lemma-upper-shriek-composition}), $j^!$와\  $i^!$가 제한으로 주어지므로
(Lemma \ref{duality-lemma-shriek-open-immersion}), 임의의\ 
$E \in D^+_\QCoh(\mathcal{O}_X)$에 대해
$$
(f^!E)|_U = g^!(E|_V)
$$
이다. $U = \Spec(A)$와\  $V = \Spec(R)$로 쓰고, $\varphi : R \to A$를\  $g$에
대응하는 유한형 환 사상이라 하자. $P = R[x_1, \ldots, x_n]$가\  $R$ 위의\ 
$n$변수 다항식 대수가 되도록 표현\  $A = P/I$를 선택하자. $E|_V$에 대응하는
대상\  $K \in D^+(R)$을 선택한다(Derived Categories of Schemes, Lemma
\ref{perfect-lemma-affine-compare-bounded}). 그러면\  $f^!E|_U$는
$$
\varphi^!(K) = R\Hom(A, K \otimes_R^\mathbf{L} P)[n]
$$
에 대응한다고 주장한다. 여기서\  $R\Hom(A, -) : D(P) \to D(A)$는\  Dualizing
Complexes, Section \ref{dualizing-section-trivial}의 함자이고,
$\varphi^! : D(R) \to D(A)$는\  Dualizing Complexes, Section
\ref{dualizing-section-relative-dualizing-complex-algebraic}의 함자이다. 실제로
표현의 선택은 인수분해
$$
U \rightarrow \mathbf{A}^n_V \to \mathbf{A}^{n - 1}_V \to \ldots \to
\mathbf{A}^1_V \to V
$$
를 준다. Lemma \ref{duality-lemma-shriek-affine-line}을 정확히\  $n$번 적용하면\ 
$(\mathbf{A}^n_V \to V)^!(E|_V)$가\  $K \otimes_R^\mathbf{L} P[n]$에 대응함을
알 수 있다. Lemmas \ref{duality-lemma-sheaf-with-exact-support-quasi-coherent}와\ 
\ref{duality-lemma-shriek-closed-immersion}에 의해 마지막 단계는\  $R\Hom(A, -)$의 적용에
대응한다.
\end{remark}
\begin{lemma}
\label{duality-lemma-shriek-coherent}
Situation \ref{duality-situation-shriek}에서\  $f : X \to Y$를\  $\textit{FTS}_S$의 사상이라
하자. 그러면\  $f^!$는\  $D_{\textit{Coh}}^+(\mathcal{O}_Y)$를\ 
$D_{\textit{Coh}}^+(\mathcal{O}_X)$ 안으로 보낸다.
\end{lemma}

\begin{proof}
이 문제는\  $X$ 위에서 국소적이므로\  $X$와\  $Y$가 아핀 스킴이라고 가정해도 된다.
이 경우\  $f : X \to Y$를
$$
X \xrightarrow{i} \mathbf{A}^n_Y \to \mathbf{A}^{n - 1}_Y \to \ldots \to
\mathbf{A}^1_Y \to Y
$$
로 인수분해할 수 있으며, 여기서\  $i$는 닫힌 몰입이다. 이 보조정리는\  Lemmas
\ref{duality-lemma-shriek-affine-line}와\ 
\ref{duality-lemma-sheaf-with-exact-support-coherent}, Dualizing Complexes, Lemma
\ref{dualizing-lemma-dualizing-polynomial-ring}, 그리고 귀납법에서 따른다.
\end{proof}
\begin{lemma}
\label{duality-lemma-shriek-dualizing}
Situation \ref{duality-situation-shriek}에서\  $f : X \to Y$를\  $\textit{FTS}_S$의 사상이라
하자. $K$가\  $Y$의 쌍대화 복합체이면\  $f^!K$는\  $X$의 쌍대화 복합체이다.
\end{lemma}

\begin{proof}
이 문제는\  $X$ 위에서 국소적이므로\  $X$와\  $Y$가 아핀 스킴이라고 가정해도 된다.
이 경우\  $f : X \to Y$를
$$
X \xrightarrow{i} \mathbf{A}^n_Y \to \mathbf{A}^{n - 1}_Y \to \ldots \to
\mathbf{A}^1_Y \to Y
$$
로 인수분해할 수 있으며, 여기서\  $i$는 닫힌 몰입이다. Lemma
\ref{duality-lemma-shriek-affine-line}, Dualizing Complexes, Lemma
\ref{dualizing-lemma-dualizing-polynomial-ring}, 그리고 귀납법에 의해\ 
$p : \mathbf{A}^n_Y \to Y$가 사영일 때\  $p^!K$는\  $\mathbf{A}^n_Y$ 위의 쌍대화
복합체임을 알 수 있다. 마찬가지로\  Dualizing Complexes, Lemma
\ref{dualizing-lemma-dualizing-quotient}와\  Lemmas
\ref{duality-lemma-sheaf-with-exact-support-quasi-coherent} 및\ 
\ref{duality-lemma-shriek-closed-immersion}에 의해\  $i^!$는 쌍대화 복합체를 쌍대화
복합체로 보낸다.
\end{proof}
\begin{lemma}
\label{duality-lemma-shriek-via-duality}
Situation \ref{duality-situation-shriek}에서\  $f : X \to Y$를\  $\textit{FTS}_S$의 사상이라
하자. $K$를\  $Y$ 위의 쌍대화 복합체라 하자.
$M \in D_{\textit{Coh}}(\mathcal{O}_Y)$에 대해\ 
$D_Y(M) = R\SheafHom_{\mathcal{O}_Y}(M, K)$로 놓고,
$E \in D_{\textit{Coh}}(\mathcal{O}_X)$에 대해\ 
$D_X(E) = R\SheafHom_{\mathcal{O}_X}(E, f^!K)$로 놓자. 그러면\ 
$M \in D_{\textit{Coh}}^+(\mathcal{O}_Y)$에 대해 정준 동형
$$
f^!M \longrightarrow D_X(Lf^*D_Y(M))
$$
이 존재한다.
\end{lemma}

\begin{proof}
$Y$ 위에서\  $X$의 콤팩트화\  $j : X \subset \overline{X}$를 선택하자(More on
Flatness, Theorem \ref{flat-theorem-nagata}와\  Lemma
\ref{flat-lemma-compactifyable}). $a$를\  $\overline{X} \to Y$에 대한\  Lemma
\ref{duality-lemma-twisted-inverse-image}의 오른쪽 수반이라 하자.
$E \in D_{\textit{Coh}}(\mathcal{O}_{\overline{X}})$에 대해\ 
$D_{\overline{X}}(E) = R\SheafHom_{\mathcal{O}_{\overline{X}}}(E, a(K))$로 놓는다.
$R\SheafHom$의 형성은 열린집합으로의 제한과 가환하고\ 
$f^! = j^* \circ a$이므로, $M \in D_{\textit{Coh}}(\mathcal{O}_Y)$에 대해 정준 동형
$$
a(M) \longrightarrow D_{\overline{X}}(L\overline{f}^*D_Y(M))
$$
이 존재함을 보이면 충분하다. $F \in D_\QCoh(\mathcal{O}_X)$에 대해
\begin{align*}
\Hom_{\overline{X}}(
F, D_{\overline{X}}(L\overline{f}^*D_Y(M)))
& =
\Hom_{\overline{X}}(
F \otimes_{\mathcal{O}_X}^\mathbf{L} L\overline{f}^*D_Y(M), a(K)) \\
& =
\Hom_Y(
R\overline{f}_*(F \otimes_{\mathcal{O}_X}^\mathbf{L} L\overline{f}^*D_Y(M)),
K) \\
& =
\Hom_Y(
R\overline{f}_*(F) \otimes_{\mathcal{O}_Y}^\mathbf{L} D_Y(M),
K) \\
& =
\Hom_Y(
R\overline{f}_*(F), D_Y(D_Y(M))) \\
& =
\Hom_Y(R\overline{f}_*(F), M) \\
& = \Hom_{\overline{X}}(F, a(M))
\end{align*}
이다. 첫 번째 등식은\  Cohomology, Lemma \ref{cohomology-lemma-internal-hom}에
의한 것이다. 두 번째 등식은\  $a$의 정의에 의한 것이다. 세 번째 등식은\  Derived
Categories of Schemes, Lemma \ref{perfect-lemma-cohomology-base-change}에 의한
것이다. 네 번째 등식은\  Cohomology, Lemma
\ref{cohomology-lemma-internal-hom}과\  $D_Y$의 정의에 의한 것이다. 다섯 번째
등식은\  Lemma \ref{duality-lemma-dualizing-schemes}에 의한 것이다. 마지막 등식은\  $a$의
정의에 의한 것이다. 따라서\  Yoneda 보조정리에 의해\ 
$a(M) = D_{\overline{X}}(L\overline{f}^*D_Y(M))$임을 알 수 있다.
\end{proof}
\begin{lemma}
\label{duality-lemma-perfect-comparison-shriek}
Situation \ref{duality-situation-shriek}에서\  $f : X \to Y$를\  $\textit{FTS}_S$의 사상이라
하자. $f$가 완전하다고 가정하자(예를 들어 평탄). 그러면
\begin{enumerate}
\item[(a)] $f^!\mathcal{O}_Y$는\  $D_{\textit{Coh}}^b(\mathcal{O}_X)$에 속하고,
\item[(b)] $f^!\mathcal{O}_Y$는\  $D(f^{-1}\mathcal{O}_Y)$에서 유한\  Tor 차원을 가지며,
\item[(c)] $\mathcal{O}_X \to
R\SheafHom_{\mathcal{O}_X}(f^!\mathcal{O}_Y, f^!\mathcal{O}_Y)$
는 동형이고,
\item[(d)] $f^!$는\  $D_{\textit{Coh}}^b(\mathcal{O}_Y)$를\ 
$D_{\textit{Coh}}^b(\mathcal{O}_X)$ 안으로 보내며,
\item[(e)] Lemma \ref{duality-lemma-map-pullback-to-shriek-well-defined}의 사상\ 
$\mu_{f,  K} :
Lf^*K \otimes_{\mathcal{O}_X}^\mathbf{L} f^!\mathcal{O}_Y
\to
f^!K$
는 모든\  $K \in D_\QCoh^+(\mathcal{O}_Y)$에 대해 동형이다.
\end{enumerate}
\end{lemma}

\begin{proof}
(유한 표시인 평탄 사상은 완전하다. More on Morphisms, Lemma
\ref{more-morphisms-lemma-flat-finite-presentation-perfect}를 보라.)
명제 (a), (b), (c)는\  $X$ 위에서 국소적이다. 따라서\  $X$와\  $Y$가 아핀이라고
가정해도 된다. 그러면\  Remark \ref{duality-remark-local-calculation-shriek}는 (a), (b),
(c)를\  Dualizing Complexes, Lemma
\ref{dualizing-lemma-relative-dualizing-algebraic}의 (1)(a), (1)(b), (1)(c)로
바꾼다. (명제들을 서로 맞추려면\  Derived Categories of Schemes, Lemmas
\ref{perfect-lemma-tor-dimension-rel-affine},
\ref{perfect-lemma-pseudo-coherent},
\ref{perfect-lemma-identify-pseudo-coherent-noetherian}
\ref{perfect-lemma-quasi-coherence-internal-hom}을 사용한다.)

\medskip\noindent
(d)와 (e)를 증명하기 위해 몇 가지 예비 관찰로 시작한다.
\begin{enumerate}
\item 이미\  $f^!$가\  $D_{\textit{Coh}}^+(\mathcal{O}_Y)$를\ 
$D_{\textit{Coh}}^+(\mathcal{O}_X)$ 안으로 보낸다는 것을 알고 있다. Lemma
\ref{duality-lemma-shriek-coherent}를 보라.
\item $f$가 열린 몰입이면\  $f^! = f^*$이고, $f^!$와\  $\mu_f$의 구성에서\ 
$\overline{X} = Y$를 취할 수 있으므로 (d)와 (e)는 참이다. Lemma
\ref{duality-lemma-shriek-open-immersion}을 보라.
\item $f$가 완전 고유 사상이면\  Lemma
\ref{duality-lemma-compare-with-pullback-flat-proper-noetherian}에 의해 (e)는 참이다.
\item 열린 덮개\  $X = \bigcup U_i$가 존재하고 (d)가\  $U_i \to Y$에 대해 참이면
(d)는\  $X \to Y$에 대해서도 참이다. (e)도 마찬가지이다. 이는\  $f^!$와\  $\mu_f$의
구성이 열린 부분스킴으로 넘어가는 것과 가환하기 때문이다.
\item $g : Y \to Z$가\  $\textit{FTS}_S$의 두 번째 완전 사상이고 (e)가\  $f$와\ 
$g$에 대해 성립하면,
$f^!g^!\mathcal{O}_Z =
Lf^*g^!\mathcal{O}_Z \otimes_{\mathcal{O}_X}^\mathbf{L} f^!\mathcal{O}_Y$
이고, Lemma \ref{duality-lemma-map-pullback-to-shriek-well-defined}의 가환도식에 의해
(e)는\  $g \circ f$에 대해서도 성립한다.
\item (d)와 (e)가\  $f$와\  $g$ 모두에 대해 성립하면, (d)와 (e)는\  $g \circ f$에
대해서도 성립한다. 실제로\  $f^!g^!\mathcal{O}_Z$는 (앞 항에 의해) 위로 유계이고\ 
$L(g \circ f)^*$는 유한 코호몰로지 차원을 가지며, 위에서 본 (e)로부터 (d)가
따른다.
\end{enumerate}
이 관찰들에 의해\  $X$가 아핀인 경우에 (d)와 (e)를 증명하면 충분하다. 몰입\ 
$X \to \mathbf{A}^n_Y$를 선택하자(Morphisms, Lemma
\ref{morphisms-lemma-quasi-affine-finite-type-over-S}). 이를\ 
$X \to U \to \mathbf{A}^n_Y \to Y$로 인수분해하자. 여기서\  $X \to U$는 닫힌
몰입이고\  $U \subset \mathbf{A}^n_Y$는 열린집합이다. More on Morphisms, Lemma
\ref{more-morphisms-lemma-perfect-permanence}에 의해\  $X \to U$가 완전 닫힌
몰입임에 유의하자. 따라서 완전 닫힌 몰입과 사영\ 
$\mathbf{A}^n_Y \to Y$에 대해 보조정리를 증명하면 충분하다.

\medskip\noindent
$f : X \to Y$를 완전 닫힌 몰입이라 하자. (d)가 성립한다는 것은 이미 알고 있다.
$K \in D^b_{\textit{Coh}}(\mathcal{O}_Y)$라 하자. 그러면\ 
$f^!K = R\SheafHom(\mathcal{O}_X, K)$
(Lemma \ref{duality-lemma-shriek-closed-immersion})이고\ 
$f_*f^!K = R\SheafHom_{\mathcal{O}_Y}(f_*\mathcal{O}_X, K)$이다. $f$가
완전하므로 복합체\  $f_*\mathcal{O}_X$는 완전하고, 따라서\ 
$R\SheafHom_{\mathcal{O}_Y}(f_*\mathcal{O}_X, K)$는 위로 유계이다. 이는 (d)가
성립함을 증명한다. 몇 가지 세부사항은 생략한다.

\medskip\noindent
$f : \mathbf{A}^n_Y \to Y$를 사영이라 하자. 그러면\  Lemma
\ref{duality-lemma-shriek-affine-line}을 반복하여 적용하면 (d)가 성립한다. 마지막으로
(e)는\  $\mathbf{P}^n_Y \to Y$에 대해 성립하고(평탄하고 고유),
$\mathbf{A}^n_Y \subset \mathbf{P}^n_Y$가 열린집합이므로 참이다.
\end{proof}
\begin{lemma}
\label{duality-lemma-flat-shriek-relatively-perfect}
Situation \ref{duality-situation-shriek}에서\  $f : X \to Y$를\  $\textit{FTS}_S$의 사상이라
하자. $f$가 평탄이면\  $f^!\mathcal{O}_Y$는\  $D(\mathcal{O}_X)$의\  $Y$-완전 대상이고,
$\mathcal{O}_X \to
R\SheafHom_{\mathcal{O}_X}(f^!\mathcal{O}_Y, f^!\mathcal{O}_Y)$
는 동형이다.
\end{lemma}

\begin{proof}
두 명제 모두\  $X$ 위에서 국소적이다. 따라서\  $X$와\  $Y$가 아핀이라고 가정해도
된다. 그러면\  Remark \ref{duality-remark-local-calculation-shriek}에 의해 보조정리는
대수 보조정리, 즉\  Dualizing Complexes, Lemma
\ref{dualizing-lemma-relative-dualizing-algebraic}로 바뀐다. (용어를 맞추려면\ 
Derived Categories of Schemes, Lemma
\ref{perfect-lemma-affine-locally-rel-perfect}를 사용한다.)
\end{proof}
\begin{lemma}
\label{duality-lemma-lci-shriek}
Situation \ref{duality-situation-shriek}에서\  $f : X \to Y$를\  $\textit{FTS}_S$의 사상이라
하자. $f : X \to Y$가 국소 완전 교차 사상이라고 가정하자. 그러면
\begin{enumerate}
\item $f^!\mathcal{O}_Y$는\  $D(\mathcal{O}_X)$의 가역 대상이고,
\item $f^!$는 완전 복합체를 완전 복합체로 보낸다.
\end{enumerate}
\end{lemma}

\begin{proof}
국소 완전 교차 사상이 완전하다는 것을 상기하자. More on Morphisms, Lemma
\ref{more-morphisms-lemma-lci-properties}를 보라. Lemma
\ref{duality-lemma-perfect-comparison-shriek}에 의해\  $f^!\mathcal{O}_Y$가\ 
$D(\mathcal{O}_X)$의 가역 대상임을 보이면 충분하다. 이 문제는\  $X$와\  $Y$
위에서 국소적이다. 따라서\  $X \to Y$가\  $X \to \mathbf{A}^n_Y \to Y$로
인수분해되며 첫 번째 화살표가\  Koszul 정칙 몰입이라고 가정해도 된다. More on
Morphisms, Section \ref{more-morphisms-section-lci}를 보라. 결과는\  Lemma
\ref{duality-lemma-shriek-affine-line}에 의해\  $\mathbf{A}^n_Y \to Y$에 대해 성립한다.
따라서\  $f$가\  Koszul 정칙 몰입일 때 보조정리를 증명하면 충분하다. 다시 국소적으로
작업하여\  $X = \Spec(A)$이고\  $Y = \Spec(B)$인 경우로 환원한다. 여기서 어떤 정칙열\ 
$f_1, \ldots, f_r \in B$에 대해\  $A = B/(f_1, \ldots, f_r)$이다(Noether 국소환에
대해서는\  Koszul 정칙과 정칙의 개념이 같다는 것을 사용한다. More on Algebra,
Lemma \ref{more-algebra-lemma-noetherian-finite-all-equivalent}를 보라). 따라서\ 
$X \to Y$는 합성
$$
X = X_r \to X_{r - 1} \to \ldots \to X_1 \to X_0 = Y
$$
이고, 각 화살표는 유효\  Cartier 약수의 포함이다. 이와 같이 유효\  Cartier 약수의
포함\  $i : D \to X$인 경우로 환원한다. 이 경우\  Lemma
\ref{duality-lemma-compute-for-effective-Cartier}에 의해\ 
$i^!\mathcal{O}_X = \mathcal{N}[1]$이고 증명이 끝난다.
\end{proof}






\section{상단 느낌표의 기저변환}
\label{duality-section-base-change-shriek}

\noindent
Situation \ref{duality-situation-shriek}에서
$$
\xymatrix{
X' \ar[r]_{g'} \ar[d]_{f'} & X \ar[d]^f \\
Y' \ar[r]^g & Y
}
$$
를\  $\textit{FTS}_S$의 카르테시안 도식이라 하자. 여기서\  $X$와\  $Y'$는\  $Y$ 위에서\ 
Tor 독립이다. 현재의 설정은 이 일반성에서 기저변환 사상\ 
$L(g')^* \circ f^! \to (f')^! \circ Lg^*$를 구성하기에 충분하지 않다. 그 이유는
일반적으로\  $\overline{X}$와\  $Y'$가\  $Y$ 위에서\  tor 독립이 되도록\  $Y$ 위에서\  X의
콤팩트화\  $j : X \to \overline{X}$를 선택할 수 없고, 따라서\  Section
\ref{duality-section-base-change-map}의 기저변환 사상 구성을 적용할 수 없기
때문이다\footnote{유도 대수기하학에 익숙한 독자는 이것이 ``진정한'' 문제가 아님을
알 것이다. 즉\  $Y$ 위에서\  $\overline{X}$와\  $Y'$의 유도 섬유곱을\ 
$\overline{X}'$로 놓으면, Lemma \ref{duality-lemma-base-change-shriek-flat}의 증명과
정확히 같이 논증하여 이 사상을 정의할 수 있다. 결국\  $X$와\  $Y'$의\  Tor 독립성은\ 
$X'$이 유도 스킴\  $\overline{X}'$의 열린 부분스킴이 됨을 보장한다.}.

\medskip\noindent
부분적인 해결책은\  Section \ref{duality-section-relative-dualizing-complexes}에서 찾을 수
있다. 즉 사상\  $f$가 평탄이면 상대 쌍대화 복합체라는 좋은 개념이 존재하고, Lemmas
\ref{duality-lemma-compactifyable-relative-dualizing}
\ref{duality-lemma-base-change-relative-dualizing},
and \ref{duality-lemma-perfect-comparison-shriek}를 사용하여 정준 기저변환 동형을 구성할 수
있다. 이를 사용할 필요가 생기면 이 장의 뒤쪽에서 정확한 명제와 증명을 추가하겠다.
\begin{lemma}
\label{duality-lemma-base-change-shriek-flat}
Situation \ref{duality-situation-shriek}에서
$$
\xymatrix{
X' \ar[r]_{g'} \ar[d]_{f'} & X \ar[d]^f \\
Y' \ar[r]^g & Y
}
$$
를\  $g$가 평탄인\  $\textit{FTS}_S$의 카르테시안 도식이라 하자. 그러면\ 
$D_\QCoh^+(\mathcal{O}_Y)$ 위에서 동형\ 
$L(g')^* \circ f^! \to (f')^! \circ Lg^*$가 존재한다.
\end{lemma}

\begin{proof}
실제로\  $g$가 평탄하므로, $Y$ 위에서\  $X$의 콤팩트화\ 
$j : X \to \overline{X}$를 어떻게 선택하든 스킴\  $\overline{X}$는\  $Y'$와\  Tor
독립이다. $j$의 기저변환을\  $j' : X' \to \overline{X}'$로, 사영을\ 
$\overline{g}' : \overline{X}' \to \overline{X}$로 나타내자. 기저변환 사상을 합성
$$
L(g')^* \circ f^! = L(g')^* \circ j^* \circ a =
(j')^* \circ L(\overline{g}')^* \circ a \longrightarrow
(j')^* \circ a' \circ Lg^* = (f')^! \circ Lg^*
$$
으로 정의한다. 여기서 가운데 화살표는 기저변환 사상
(\ref{duality-equation-base-change-map})이고, $a$와\  $a'$는 각각\ 
$\overline{X} \to Y$와\  $\overline{X}' \to Y'$에 대한\  Lemma
\ref{duality-lemma-twisted-inverse-image}의 밂의 오른쪽 수반이다. 이 구성은 콤팩트화의
선택과 무관하다(이 결과를 사용할 필요가 생기면 정확한 보조정리를 서술하고
증명하겠다).

\medskip\noindent
증명을 끝내려면 기저변환 사상\  $L(g')^* \circ a \to a' \circ Lg^*$가\ 
$D_\QCoh^+(\mathcal{O}_Y)$ 위에서 동형임을 보이면 충분하다. Lemma
\ref{duality-lemma-proper-noetherian}에 의해\  $a$와\  $a'$의 형성은\  $Y$와\  $Y'$의 아핀
열린집합으로 제한하는 것과 가환한다. 따라서\  Remark
\ref{duality-remark-check-over-affines}에 의해\  $Y$와\  $Y'$가 아핀이라고 가정해도 된다.
그러므로 결과는\  Lemma \ref{duality-lemma-more-base-change}에 의한다.
\end{proof}
\begin{lemma}
\label{duality-lemma-shriek-etale}
Situation \ref{duality-situation-shriek}에서\  $f : X \to Y$를\  $\textit{FTS}_S$의 \'etale
사상이라 하자. 그러면\  $D^+_\QCoh(\mathcal{O}_Y)$ 위의 함자들로서\ 
$f^! \cong f^*$이다.
\end{lemma}

\begin{proof}
\'etale 사상이 평탄이고\  syntomic이며 국소 완전 교차 사상이라는 사실을 사용할
것이다(Morphisms, Lemma \ref{morphisms-lemma-etale-syntomic}와\ 
\ref{morphisms-lemma-etale-flat}, 그리고\  More on Morphisms, Lemma
\ref{more-morphisms-lemma-flat-lci}). Lemma
\ref{duality-lemma-perfect-comparison-shriek}에 의해\ 
$f^!\mathcal{O}_Y = \mathcal{O}_X$임을 보이면 충분하다. Lemma
\ref{duality-lemma-lci-shriek}에 의해\  $f^!\mathcal{O}_Y$가 가역 가군임을 알고 있다.
가환도식
$$
\xymatrix{
X \times_Y X \ar[r]_{p_2} \ar[d]_{p_1} & X \ar[d]^f \\
X \ar[r]^f & Y
}
$$
과 대각사상\  $\Delta : X \to X \times_Y X$를 생각하자. $\Delta$는 열린
몰입이므로(Morphisms, Lemmas
\ref{morphisms-lemma-diagonal-unramified-morphism}와\ 
\ref{morphisms-lemma-etale-smooth-unramified}), Lemma
\ref{duality-lemma-shriek-open-immersion}에 의해\  $\Delta^! = \Delta^*$이다. Lemma
\ref{duality-lemma-upper-shriek-composition}에 의해\ 
$\Delta^! \circ p_1^! \circ f^! = f^!$이다. 위 도식에\  Lemma
\ref{duality-lemma-base-change-shriek-flat}을 적용하면\ 
$p_1^!\mathcal{O}_X = p_2^*f^!\mathcal{O}_Y$이다. 따라서
$$
f^!\mathcal{O}_Y = \Delta^!p_1^!f^!\mathcal{O}_Y =
\Delta^*(p_1^*f^!\mathcal{O}_Y \otimes p_1^!\mathcal{O}_X) =
\Delta^*(p_2^*f^!\mathcal{O}_Y \otimes p_1^*f^!\mathcal{O}_Y) =
(f^!\mathcal{O}_Y)^{\otimes 2}
$$
임을 얻는다. 두 번째 단계에서는\  Lemma
\ref{duality-lemma-perfect-comparison-shriek}을 다시 한 번 사용했다. 그러므로 원하는 대로\ 
$f^!\mathcal{O}_Y = \mathcal{O}_X$이다.
\end{proof}

\noindent
이 절의 나머지에서는 좋은 기저변환 사상 이론의 결과가 될, 증명하기 쉬운 몇 가지
결과를 서술한다.

\begin{lemma}[임시 기저변환]
\label{duality-lemma-base-change-locally}
Situation \ref{duality-situation-shriek}에서
$$
\xymatrix{
X' \ar[r]_{g'} \ar[d]_{f'} & X \ar[d]^f \\
Y' \ar[r]^g & Y
}
$$
를\  $\textit{FTS}_S$의 카르테시안 도식이라 하자.
$E \in D^+_\QCoh(\mathcal{O}_Y)$를\  $Lg^*E$가\  $D^+(\mathcal{O}_Y)$에 속하는
대상이라 하자. $f$가 평탄이면, $Y$, $Y'$, $X$의 아핀 열린집합들 안으로 사상되는
아핀 열린집합\  $U' \subset X'$에 대해\  $L(g')^*f^!E$와\  $(f')^!Lg^*E$를\ 
$D(\mathcal{O}_{U'})$의 대상으로 제한한 것들은 동형이다.
\end{lemma}

\begin{proof}
가정에 의해\  $X$, $Y$, $Y'$, $X'$가 아핀인 경우로 즉시 환원된다.
$Y = \Spec(R)$, $Y' = \Spec(R')$, $X = \Spec(A)$,
$X' = \Spec(A')$라 쓰자. 그러면\  $A' = A \otimes_R R'$이다.
$E$가\  $K \in D^+(R)$에 대응한다고 하자. 주어진 사상들을\ 
$\varphi : R \to A$와\  $\varphi' : R' \to A'$로 나타내면, Remark
\ref{duality-remark-local-calculation-shriek}에서\  $L(g')^*f^!E$와\  $(f')^!Lg^*E$가 각각\ 
$\varphi^!(K) \otimes_A^\mathbf{L} A'$와\ 
$(\varphi')^!(K \otimes_R^\mathbf{L} R')$에 대응함을 알 수 있다. 여기서\ 
$\varphi^!$와\  $(\varphi')^!$는\  Dualizing Complexes, Section
\ref{dualizing-section-relative-dualizing-complex-algebraic}의 함자들이다. 결과는\ 
Dualizing Complexes, Lemma \ref{dualizing-lemma-bc-flat}에서 따른다.
\end{proof}
\begin{lemma}
\label{duality-lemma-relative-dualizing-fibres}
Situation \ref{duality-situation-shriek}에서\  $f : X \to Y$를\  $\textit{FTS}_S$의 사상이라
하고, $f$가 평탄이라고 가정하자.
$D^b_{\textit{Coh}}(X)$에서\  $\omega_{X/Y}^\bullet = f^!\mathcal{O}_Y$로 놓자.
$y \in Y$라 하고\  $h : X_y \to X$를 사영이라 하자. 그러면\ 
$Lh^*\omega_{X/Y}^\bullet$은\  $X_y$ 위의 쌍대화 복합체이다.
\end{lemma}

\begin{proof}
Lemma \ref{duality-lemma-perfect-comparison-shriek}에 의해 복합체\ 
$\omega_{X/Y}^\bullet$은\  $D^b_{\textit{Coh}}$에 속한다. 쌍대화 복합체라는 것은
국소적 성질이다. 따라서\  Lemma \ref{duality-lemma-base-change-locally}에 의해\ 
$(X_y \to y)^!\mathcal{O}_y$가\  $X_y$ 위의 쌍대화 복합체임을 보이면 충분하다.
이는\  Lemma \ref{duality-lemma-shriek-dualizing}에서 따른다.
\end{proof}







\section{쌍대성 이론}
\label{duality-section-duality}

\noindent
이 절에서는 위의 매우 일반적인 결과들이 고정된\  Noether 기저 스킴 위의 유한형
분리 스킴들에 어떤 종류의 쌍대성 이론을 주는지 자세히 설명한다.

\medskip\noindent
Noether 스킴\  $X$ 위의 쌍대화 복합체는\  $D(\mathcal{O}_X)$의 대상으로서 아핀
국소적으로 대응하는 환들의 쌍대화 복합체를 준다는 것을 상기하자. Definition
\ref{duality-definition-dualizing-scheme}을 보라.

\medskip\noindent
Noether 스킴\  $S$가 주어졌을 때, $S$ 위에서 유한형이고 분리인 스킴들의 범주를\ 
$\textit{FTS}_S$로 나타내자. 그러면 다음이 성립한다.
\begin{enumerate}
\item 함자\  $f^!$들은\  $D_\QCoh^+$를\  $\textit{FTS}_S$ 위의 유사함자로 만들고,
\item $f : X \to Y$가\  $\textit{FTS}_S$의 고유 사상이면\  $f^!$는\ 
$Rf_* : D_\QCoh(\mathcal{O}_X) \to D_\QCoh(\mathcal{O}_Y)$의 오른쪽 수반을\ 
$D_\QCoh^+(\mathcal{O}_Y)$에 제한한 것이며, 모든\ 
$K \in D_{\textit{Coh}}^-(\mathcal{O}_X)$와\ 
$M \in D_\QCoh^+(\mathcal{O}_Y)$에 대해 정준 동형
$$
Rf_*R\SheafHom_{\mathcal{O}_X}(K, f^!M)
\to
R\SheafHom_{\mathcal{O}_Y}(Rf_*K, M)
$$
이 존재하고,
\item $\textit{FTS}_S$의 대상\  $X$가 쌍대화 복합체\ 
$\omega_X^\bullet$을 가지면 함자\ 
$D_X = R\SheafHom_{\mathcal{O}_X}(-, \omega_X^\bullet)$는\ 
$D_{\textit{Coh}}(\mathcal{O}_X)$의 대합을 정의하여\ 
$D_{\textit{Coh}}^+(\mathcal{O}_X)$와\ 
$D_{\textit{Coh}}^-(\mathcal{O}_X)$를 서로 바꾸고\ 
$D_{\textit{Coh}}^b(\mathcal{O}_X)$를 보존하며,
\item $f : X \to Y$가\  $\textit{FTS}_S$의 사상이고\ 
$\omega_Y^\bullet$이\  $Y$ 위의 쌍대화 복합체이면 다음이 성립한다.
\begin{enumerate}
\item $\omega_X^\bullet = f^!\omega_Y^\bullet$은\  $X$의 쌍대화 복합체이고,
\item $M \in D_{\textit{Coh}}^+(\mathcal{O}_Y)$에 대해 정준적으로\ 
$f^!M = D_X(Lf^*D_Y(M))$이며,
\item 또한\  $f$가 고유이면
$$
Rf_*R\SheafHom_{\mathcal{O}_X}(K, \omega_X^\bullet) =
R\SheafHom_{\mathcal{O}_Y}(Rf_*K, \omega_Y^\bullet)
$$
가\  $D^-_{\textit{Coh}}(\mathcal{O}_X)$의\  $K$에 대해 성립한다.
\end{enumerate}
\item $f : X \to Y$가\  $\textit{FTS}_S$의 닫힌 몰입이면\ 
$f^!(-) = R\SheafHom(\mathcal{O}_X, -)$이고,
\item $f : Y \to X$가\  $\textit{FTS}_S$의 유한 사상이면\ 
$f_*f^!(-) = R\SheafHom_{\mathcal{O}_X}(f_*\mathcal{O}_Y, -)$이고,
\item $f : X \to Y$가\  $\textit{FTS}_S$의 한 대상 안으로 들어가는 유효\  Cartier
약수의 포함이면\ 
$f^!(-) = Lf^*(-) \otimes_{\mathcal{O}_X} f^*\mathcal{O}_Y(X)[-1]$이고,
\item $f : X \to Y$가\  $\textit{FTS}_S$의 한 대상 안으로 들어가는 여차원\  $c$의\ 
Koszul 정칙 몰입이면\ 
$f^!(-) \cong Lf^*(-) \otimes_{\mathcal{O}_X} \wedge^c\mathcal{N}[-c]$이며,
\item $f : X \to Y$가\  $\textit{FTS}_S$에서 상대 차원\  $d$인 매끄러운 고유
사상이면\ 
$f^!(-) \cong Lf^*(-)  \otimes_{\mathcal{O}_X} \Omega^d_{X/Y}[d]$이다.
\end{enumerate}
이는\  Lemmas
\ref{duality-lemma-dualizing-schemes},
\ref{duality-lemma-iso-on-RSheafHom},
\ref{duality-lemma-twisted-inverse-image-closed},
\ref{duality-lemma-finite-twisted},
\ref{duality-lemma-sheaf-with-exact-support-effective-Cartier},
\ref{duality-lemma-regular-immersion},
\ref{duality-lemma-smooth-proper},
\ref{duality-lemma-upper-shriek-composition},
\ref{duality-lemma-pseudo-functor},
\ref{duality-lemma-shriek-closed-immersion},
\ref{duality-lemma-shriek-dualizing},
\ref{duality-lemma-shriek-via-duality}, 그리고\ 
\ref{duality-lemma-perfect-comparison-shriek}과\  Example
\ref{duality-example-iso-on-RSheafHom-noetherian}에서 따른다. 우리의 함자들은 궁극적으로
존재 정리(Derived Categories, Proposition \ref{derived-proposition-brown})를
호출하는 매우 추상적인 절차로 얻었다. 이는 일반적으로 함자\  $f^!$를 명시적으로
기술할 수 없다는 뜻이다. 때로는 이것이 문제가 될 수 있다. 하지만 실제로는
쌍대화 복합체의 존재와 쌍대성 동형을 아는 것만으로도\  $f^!$를 결정하기에 충분한
경우가 많다.





\section{쌍대화 복합체의 붙이기}
\label{duality-section-glue}

\noindent
이제\  Situation \ref{duality-situation-dualizing}에서와 같은 쌍\ 
$(S, \omega_S^\bullet)$가 주어졌을 때, 쌍대화 복합체를 붙여서\  $S$ 위의 모든
유한형 스킴에 적용되는 이론을 얻는다. 이는\  \cite[Remark on page 310]{RD}와
유사하다.

\begin{situation}
\label{duality-situation-dualizing}
여기서\  $S$는\  Noether 스킴이고\  $\omega_S^\bullet$은 쌍대화 복합체이다.
\end{situation}

\noindent
Situation \ref{duality-situation-dualizing}에서\  $X$를\  $S$ 위의 유한형 스킴이라 하자.
$\mathcal{U} : X = \bigcup_{i = 1, \ldots, n} U_i$를\  $\textit{FTS}_S$의
대상들로 이루어진\  $X$의 유한 열린 덮개라 하자. Situation
\ref{duality-situation-shriek}을 보라. 이는 사상\  $U_i \to S$가 분리라는 뜻일 뿐이다
(이미 유한형이기 때문이다). $S$ 위의 유한형 아핀 스킴은 모두\  Schemes, Lemma
\ref{schemes-lemma-compose-after-separated}에 의해\  $\textit{FTS}_S$의 대상이므로,
그러한 열린 덮개는 확실히 존재한다. 그러면 모든\ 
$i, j, k \in \{1, \ldots, n\}$에 대해 사상\ 
$p_i : U_i \to S$, $p_{ij} : U_i \cap U_j \to S$, 그리고\ 
$p_{ijk} : U_i \cap U_j \cap U_k \to S$는 분리이고, 이 스킴들은 각각\ 
$\textit{FTS}_S$의 대상이다. 이러한 열린 덮개로부터 다음을 얻는다.
\begin{enumerate}
\item $\omega_i^\bullet = p_i^!\omega_S^\bullet$은\  $U_i$ 위의 쌍대화
복합체이다. Section \ref{duality-section-duality}를 보라.
\item 모든\  $i, j$에 대해 정준 동형\ 
$\varphi_{ij} :
\omega_i^\bullet|_{U_i \cap U_j} \to \omega_j^\bullet|_{U_i \cap U_j}$가 있다.
\item
\label{duality-item-cocycle-glueing}
모든\  $i, j, k$에 대해\  $D(\mathcal{O}_{U_i \cap U_j \cap U_k})$에서
$$
\varphi_{ik}|_{U_i \cap U_j \cap U_k} =
\varphi_{jk}|_{U_i \cap U_j \cap U_k} \circ
\varphi_{ij}|_{U_i \cap U_j \cap U_k}
$$
가 성립한다.
\end{enumerate}
여기서 (2)에서는\  $(U_i \cap U_j \to U_i)^!$가 제한으로 주어진다는 사실
(Lemma \ref{duality-lemma-shriek-open-immersion})과\  Lemma
\ref{duality-lemma-upper-shriek-composition}에 의한 정준 동형
$$
(U_i \cap U_j \to U_i)^! \circ p_i^! = p_{ij}^! =
(U_i \cap U_j \to U_j)^! \circ p_j^!
$$
을 사용한다. (3)을 얻기 위해서는\  Lemma \ref{duality-lemma-pseudo-functor}에 의해 상단
느낌표 함자들이 유사함자를 이룬다는 사실을 사용한다.

\medskip\noindent
방금 설명한 상황에서
{\it $\omega_S^\bullet$과\  $\mathcal{U}$에 대해 정규화된 쌍대화 복합체}란\ 
$K \in D(\mathcal{O}_X)$와 동형\ 
$\alpha_i : K|_{U_i} \to \omega_i^\bullet$로 이루어진 쌍\  $(K, \alpha_i)$로서,
$\varphi_{ij}$가\ 
$\alpha_j|_{U_i \cap U_j} \circ \alpha_i^{-1}|_{U_i \cap U_j}$로 주어지는
것을 말한다. 스킴 위에서 쌍대화 복합체라는 것은 국소적 성질이므로,
$\omega_S^\bullet$과\  $\mathcal{U}$에 대해 정규화된 쌍대화 복합체는 실제로
쌍대화 복합체이다.
\begin{lemma}
\label{duality-lemma-good-dualizing-unique}
Situation \ref{duality-situation-dualizing}에서\  $X$를\  $S$ 위의 유한형 스킴이라 하고,
$\mathcal{U}$를\  $S$ 위에서 분리인 스킴들로 이루어진\  $X$의 유한 열린 덮개라
하자. $\omega_S^\bullet$과\  $\mathcal{U}$에 대해 정규화된 쌍대화 복합체가
존재하면, 그것은 유일한 동형을 제외하고 유일하다.
\end{lemma}

\begin{proof}
$(K, \alpha_i)$와\  $(K', \alpha_i')$가 두 개의 그러한 복합체라고 하자. 이때\ 
$L = R\SheafHom_{\mathcal{O}_X}(K, K')$를 생각한다. Lemma
\ref{duality-lemma-dualizing-unique-schemes}와 그 증명에 의해, 이는\ 
$D(\mathcal{O}_X)$의 가역 대상이다. $\alpha_i$와\  $\alpha'_i$를 사용하면 동형
$$
\alpha_i^t \otimes \alpha'_i :
L|_{U_i} \longrightarrow
R\SheafHom_{\mathcal{O}_X}(\omega_i^\bullet, \omega_i^\bullet) =
\mathcal{O}_{U_i}[0]
$$
을 얻는다. 이 사실만으로도\  $D(\mathcal{O}_X)$에서\  $L = H^0(L)[0]$임을 알 수
있다. 또한\  $H^0(L)$은\  $X$의 열린집합\  $U_i$들 위에 주어진 자명화를 갖는 가역
층이다. 마지막으로\ 
$\alpha_j|_{U_i \cap U_j} \circ \alpha_i^{-1}|_{U_i \cap U_j}$와\ 
$\alpha'_j|_{U_i \cap U_j} \circ (\alpha'_i)^{-1}|_{U_i \cap U_j}$가 모두\ 
$\varphi_{ij}$를 준다는 조건으로부터 전이 사상들이\  $1$임을 알 수 있고,
$H^0(L) = \mathcal{O}_X$인 동형을 얻는다.
\end{proof}
\begin{lemma}
\label{duality-lemma-good-dualizing-independence-covering}
Situation \ref{duality-situation-dualizing}에서\  $X$를\  $S$ 위의 유한형 스킴이라 하고,
$\mathcal{U}$와\  $\mathcal{V}$를\  $S$ 위에서 분리인 스킴들로 이루어진\  $X$의 두
유한 열린 덮개라 하자. $\omega_S^\bullet$과\  $\mathcal{U}$에 대해 정규화된
쌍대화 복합체가 존재하면, $\omega_S^\bullet$과\  $\mathcal{V}$에 대해 정규화된
쌍대화 복합체도 존재하며 이 두 복합체는 정준적으로 동형이다.
\end{lemma}

\begin{proof}
$\mathcal{U}$가 열린집합\  $U_1, \ldots, U_n$으로 주어지고\  $\mathcal{V}$가
열린집합\  $U_1, \ldots, U_{n + m}$으로 주어진 경우를 증명하면 충분하다.
실제로\  $m = 1$이라고 가정해도 되고 그렇게 하자. $\omega_S^\bullet$과\ 
$\mathcal{V}$에 대해 정규화된 쌍대화 복합체\  $(K, \alpha_i)$에서\ 
$\omega_S^\bullet$과\  $\mathcal{U}$에 대해 정규화된 쌍대화 복합체로 가는 것은\ 
$i = n + 1$에 대한\  $\alpha_i$를 잊으면 된다. 반대로\  $(K, \alpha_i)$를\ 
$\omega_S^\bullet$과\  $\mathcal{U}$에 대해 정규화된 쌍대화 복합체라 하자.
증명을 끝내려면 필요한 조건을 만족하는 사상\ 
$\alpha_{n + 1} : K|_{U_{n + 1}} \to \omega_{n + 1}^\bullet$을 구성해야 한다.
이를 위해\  $U_{n + 1} = \bigcup U_i \cap U_{n + 1}$이 열린 덮개임을 관찰한다.
$(K|_{U_{n + 1}}, \alpha_i|_{U_i \cap U_{n + 1}})$가\ 
$\omega_S^\bullet$과 덮개\  $U_{n + 1} = \bigcup U_i \cap U_{n + 1}$에 대해
정규화된 쌍대화 복합체임은 분명하다. 한편 조건
(\ref{duality-item-cocycle-glueing})에 의해 쌍\ 
$(\omega_{n + 1}^\bullet|_{U_{n + 1}}, \varphi_{n + 1i})$도\ 
$\omega_S^\bullet$과 덮개\ 
$U_{n + 1} = \bigcup U_i \cap U_{n + 1}$에 대해 정규화된 또 하나의 쌍대화
복합체이다. Lemma \ref{duality-lemma-good-dualizing-unique}에 의해, 주어진 국소 동형들과
호환되는 유일한 동형
$$
\alpha_{n + 1} : K|_{U_{n + 1}} \longrightarrow \omega_{n + 1}^\bullet
$$
을 얻는다. 이것이 필요한 성질을 만족한다는 것을 보이는 일은 유쾌한 연습문제이다.
\end{proof}
\begin{lemma}
\label{duality-lemma-existence-good-dualizing}
Situation \ref{duality-situation-dualizing}에서\  $X$를\  $S$ 위의 유한형 스킴이라 하고,
$\mathcal{U}$를\  $S$ 위에서 분리인 스킴들로 이루어진\  $X$의 유한 열린 덮개라
하자. 그러면\  $\omega_S^\bullet$과\  $\mathcal{U}$에 대해 정규화된 쌍대화
복합체가 존재한다.
\end{lemma}

\begin{proof}
$\mathcal{U} : X = \bigcup_{i = 1, \ldots, n} U_i$라 하자. $n$에 대한 귀납법으로
보조정리를 증명한다. $n = 1$인 기초 단계는 자명하다. $n > 1$이라 가정하자.
$X' = U_1 \cup \ldots \cup U_{n - 1}$로 놓고,
$(K', \{\alpha'_i\}_{i = 1, \ldots, n - 1})$를\  $\omega_S^\bullet$과\ 
$\mathcal{U}' : X' = \bigcup_{i = 1, \ldots, n - 1} U_i$에 대해 정규화된
쌍대화 복합체라 하자.
$(K'|_{X' \cap U_n}, \alpha'_i|_{U_i \cap U_n})$가\  $\omega_S^\bullet$과 덮개\ 
$X' \cap U_n = \bigcup_{i = 1, \ldots, n - 1} U_i \cap U_n$에 대해 정규화된
쌍대화 복합체임은 분명하다. 한편 조건 (\ref{duality-item-cocycle-glueing})에 의해 쌍\ 
$(\omega_n^\bullet|_{X' \cap U_n}, \varphi_{ni})$도\  $\omega_S^\bullet$과 덮개\ 
$X' \cap U_n = \bigcup_{i = 1, \ldots, n - 1} U_i \cap U_n$에 대해 정규화된
또 하나의 쌍대화 복합체이다. Lemma \ref{duality-lemma-good-dualizing-unique}에 의해,
주어진 국소 동형들과 호환되는 유일한 동형
$$
\epsilon : K'|_{X' \cap U_n} \longrightarrow \omega_i^\bullet|_{X' \cap U_n}
$$
을 얻는다. Cohomology, Lemma \ref{cohomology-lemma-glue}에 의해\ 
$K \in D(\mathcal{O}_X)$와 동형\  $\beta : K|_{X'} \to K'$ 및\ 
$\gamma : K|_{U_n} \to \omega_n^\bullet$을 얻으며, 이들은\ 
$\epsilon = \gamma|_{X'\cap U_n} \circ \beta|_{X' \cap U_n}^{-1}$을 만족한다.
이제
$$
\alpha_i = \alpha'_i \circ \beta|_{U_i}, i = 1, \ldots, n - 1,
\text{ and }
\alpha_n = \gamma
$$
로 정의한다. 여전히\  $\varphi_{ij}$가\ 
$\alpha_j|_{U_i \cap U_j} \circ \alpha_i^{-1}|_{U_i \cap U_j}$로 주어짐을
확인해야 한다. $i, j \leq n - 1$이면 이는\  $\alpha_i'$에 대한 해당 조건에서
따른다. $i = j = n$인 경우도 분명하다. $i < j = n$이면
$$
\alpha_n|_{U_i \cap U_n} \circ \alpha_i^{-1}|_{U_i \cap U_n} =
\gamma|_{U_i \cap U_n} \circ \beta^{-1}|_{U_i \cap U_n}
\circ (\alpha'_i)^{-1}|_{U_i \cap U_n} =
\epsilon|_{U_i \cap U_n} \circ (\alpha'_i)^{-1}|_{U_i \cap U_n}
$$
을 얻는다. 이는 바로\  $\epsilon$이 사상\  $\alpha_i'$와\  $\alpha_{ni}$에 호환되는
유일한 사상이기 때문에\  $\alpha_{in}$과 같다.
\end{proof}
\noindent
$(S, \omega_S^\bullet)$를\  Situation \ref{duality-situation-dualizing}에서와 같이 두자.
위 보조정리들의 요지는\  $S$ 위의 임의의 유한형 스킴\  $X$가 주어지면, 유일한
동형을 제외하고 유일하게 정해지는 쌍\  $(K, \alpha_U)$가 있다는 것이다. 이는
대상\  $K \in D(\mathcal{O}_X)$와, $S$ 위에서 분리인 모든 열린 부분스킴\ 
$U \subset X$에 대한 동형\ 
$\alpha_U : K|_U \to \omega_U^\bullet$로 이루어진다. 여기서\ 
$\omega_U^\bullet = (U \to S)^!\omega_S^\bullet$은\  $U$ 위의 쌍대화 복합체이다.
Section \ref{duality-section-duality}를 보라. 또한\ 
$\mathcal{U} : X = \bigcup U_i$가\  $S$ 위에서 분리인 열린집합들로 이루어진
유한 열린 덮개이면, $(K, \alpha_{U_i})$는\  $\omega_S^\bullet$과\ 
$\mathcal{U}$에 대해 정규화된 쌍대화 복합체이다. 실제로 유일한 동형을
제외한 유일성은\  Lemma \ref{duality-lemma-good-dualizing-unique}, 한 열린 덮개에 대한
존재성은\  Lemma \ref{duality-lemma-existence-good-dualizing}, 그리고 그때\  $K$가 모든
열린 덮개에 대해 작동한다는 사실은\  Lemma
\ref{duality-lemma-good-dualizing-independence-covering}이다.

\begin{definition}
\label{duality-definition-good-dualizing}
$S$를\  Noether 스킴이라 하고\  $\omega_S^\bullet$을\  $S$ 위의 쌍대화 복합체라
하자. $X$를\  $S$ 위의 유한형 스킴이라 하자. 위에서 구성한 복합체\  $K$를
{\it $\omega_S^\bullet$에 대해 정규화된 쌍대화 복합체}라 하며\ 
$\omega_X^\bullet$로 나타낸다.
\end{definition}

\noindent
용어가 시사하듯이, $\omega_S^\bullet$에 대해 정규화된 쌍대화 복합체는 단지\ 
$X$의 유도 범주의 대상이 아니라 위에서 설명한 국소 동형들도 함께 갖는다.
$X$가\  $S$ 위에서 분리이고\  $p : X \to S$가 구조 사상이면\ 
$\omega_X^\bullet = p^!\omega_S^\bullet$로 놓는 것과 이는 충돌하지 않는다.
더 일반적으로 다음 건전성 검사가 성립한다.

\begin{lemma}
\label{duality-lemma-good-over-both}
$(S, \omega_S^\bullet)$를\  Situation \ref{duality-situation-dualizing}에서와 같이 두자.
$f : X \to Y$를\  $S$ 위의 유한형 스킴들의 사상이라 하자.
$\omega_X^\bullet$과\  $\omega_Y^\bullet$을\  $\omega_S^\bullet$에 대해 정규화된
쌍대화 복합체라 하자. 그러면\  $\omega_X^\bullet$은\ 
$\omega_Y^\bullet$에 대해 정규화된 쌍대화 복합체이다.
\end{lemma}

\begin{proof}
이는 단지 기호를 정리하는 문제이다. 유한 아핀 열린 덮개\ 
$\mathcal{V} : Y = \bigcup V_j$를 택한다. 각\  $j$에 대해 유한 아핀 열린 덮개\ 
$f^{-1}(V_j) = U_{ji}$를 택하고\  $\mathcal{U} : X = \bigcup U_{ji}$로 놓는다.
스킴\  $V_j$와\  $U_{ji}$는\  $S$ 위에서 분리이므로,
$q_j : V_j \to S$, $p_{ji} : U_{ji} \to S$,
$f_{ji} : U_{ji} \to V_j$, 그리고\  $f_{ji}' : U_{ji} \to Y$에 대한 상단 느낌표
함자들을 갖는다. $(L, \beta_j)$를\  $\omega_S^\bullet$과\  $\mathcal{V}$에 대해
정규화된 쌍대화 복합체라 하고, $(K, \gamma_{ji})$를\ 
$\omega_S^\bullet$과\  $\mathcal{U}$에 대해 정규화된 쌍대화 복합체라 하자.
(달리 말하면\  $L = \omega_Y^\bullet$이고\  $K = \omega_X^\bullet$이다.) 다음과 같이
정의할 수 있다.
$$
\alpha_{ji} :
K|_{U_{ji}} \xrightarrow{\gamma_{ji}}
p_{ji}^!\omega_S^\bullet = f_{ji}^!q_j^!\omega_S^\bullet
\xrightarrow{f_{ji}^!\beta_j^{-1}} f_{ji}^!(L|_{V_j}) =
(f_{ji}')^!(L)
$$
증명에는\ 
$\alpha_{ji}|_{U_{ji} \cap U_{j'i'}}
\circ \alpha_{j'i'}^{-1}|_{U_{ji} \cap U_{j'i'}}$가 정준 동형\ 
$(f_{ji}')^!(L)|_{U_{ji} \cap U_{j'i'}} \to
(f_{j'i'}')^!(L)|_{U_{ji} \cap U_{j'i'}}$임을 보여야 한다. 이는 형식적이므로
세부 사항을 생략한다.
\end{proof}

\begin{lemma}
\label{duality-lemma-open-immersion-good-dualizing-complex}
$(S, \omega_S^\bullet)$를\  Situation \ref{duality-situation-dualizing}에서와 같이 두자.
$j : X \to Y$를\  $S$ 위의 유한형 스킴들의 열린 몰입이라 하자.
$\omega_X^\bullet$과\  $\omega_Y^\bullet$을\  $\omega_S^\bullet$에 대해 정규화된
쌍대화 복합체라 하자. 그러면 정준 동형\ 
$\omega_X^\bullet = \omega_Y^\bullet|_X$가 있다.
\end{lemma}

\begin{proof}
Definition \ref{duality-definition-good-dualizing} 바로 위에 주어진 정규화된 쌍대화
복합체의 구성에서 즉시 따른다.
\end{proof}
\begin{lemma}
\label{duality-lemma-proper-map-good-dualizing-complex}
$(S, \omega_S^\bullet)$를\  Situation \ref{duality-situation-dualizing}에서와 같이 두자.
$f : X \to Y$를\  $S$ 위의 유한형 스킴들의 고유 사상이라 하자.
$\omega_X^\bullet$과\  $\omega_Y^\bullet$을\  $\omega_S^\bullet$에 대해 정규화된
쌍대화 복합체라 하자. $a$를\  $f$에 대한\  Lemma
\ref{duality-lemma-twisted-inverse-image}의 오른쪽 수반이라 하자. 그러면 정준 동형\ 
$a(\omega_Y^\bullet) = \omega_X^\bullet$이 있다.
\end{lemma}

\begin{proof}
$p : X \to S$와\  $q : Y \to S$를 구조 사상이라 하자. $X$와\  $Y$가\  $S$ 위에서
분리이면, 이는\  $\omega_X^\bullet = p^!\omega_S^\bullet$,
$\omega_Y^\bullet = q^!\omega_S^\bullet$, $f^! = a$, 그리고\ 
$f^! \circ q^! = p^!$이라는 사실에서 따른다(Lemma
\ref{duality-lemma-upper-shriek-composition}). 일반적인 경우에는 먼저\  Lemma
\ref{duality-lemma-good-over-both}를 사용하여\  $Y = S$인 경우로 환원한다. 이 경우\  $X$와\ 
$Y$는\  $S$ 위에서 분리이고, 방금 결과를 보았다.
\end{proof}

\noindent
$(S, \omega_S^\bullet)$를\  Situation \ref{duality-situation-dualizing}에서와 같이 두자.
$S$ 위의 유한형 스킴\  $X$에 대해\  $\omega_X^\bullet$로\ 
$\omega_S^\bullet$에 대해 정규화된\  $X$의 쌍대화 복합체를 나타내자. Lemma
\ref{duality-lemma-dualizing-schemes}에서처럼\ 
$D_X(-) = R\SheafHom_{\mathcal{O}_X}(-, \omega_X^\bullet)$로 정의한다.
$f : X \to Y$를\  $S$ 위의 유한형 스킴들의 사상이라 하고
$$
f_{new}^! = D_X \circ Lf^* \circ D_Y :
D_{\textit{Coh}}^+(\mathcal{O}_Y)
\to
D_{\textit{Coh}}^+(\mathcal{O}_X)
$$
로 정의한다. $f : X \to Y$와\  $g : Y \to Z$가\  $S$ 위의 유한형 스킴들 사이의
합성 가능한 사상이면 다음과 같이 정의한다.
\begin{align*}
(g \circ f)^!_{new} & = D_X \circ L(g \circ f)^* \circ D_Z \\
& = D_X \circ Lf^* \circ Lg^* \circ D_Z \\
& \to D_X \circ Lf^* \circ D_Y \circ D_Y \circ Lg^* \circ D_Z \\
& = f^!_{new} \circ g^!_{new}
\end{align*}
여기서 화살표는\  Lemma \ref{duality-lemma-dualizing-schemes}에서 정의된다. 결과들을 다음
보조정리에 모은다.

\begin{lemma}
\label{duality-lemma-duality-bootstrap}
$(S, \omega_S^\bullet)$를\  Situation \ref{duality-situation-dualizing}에서와 같이 두자.
$S$ 위의 모든 유한형 스킴과 그 사상들에 대해\  $f^!_{new}$와\ 
$\omega_X^\bullet$을 위와 같이 정의하면 다음이 성립한다.
\begin{enumerate}
\item 함자\  $f^!_{new}$와 화살표\ 
$(g \circ f)^!_{new} \to f^!_{new} \circ g^!_{new}$는\ 
$D_{\textit{Coh}}^+$를\  $S$ 위의 유한형 스킴들의 범주에서 범주들의\ 
$2$-범주로 가는 유사함자로 만든다.
\item $\omega_X^\bullet = (X \to S)^!_{new} \omega_S^\bullet$이다.
\item 함자\  $D_X$는\  $D_{\textit{Coh}}(\mathcal{O}_X)$의 대합을 정의하여\ 
$D_{\textit{Coh}}^+(\mathcal{O}_X)$와\  $D_{\textit{Coh}}^-(\mathcal{O}_X)$를
서로 바꾸고\  $D_{\textit{Coh}}^b(\mathcal{O}_X)$를 보존한다.
\item $f : X \to Y$가\  $S$ 위의 유한형 스킴들의 사상이면\ 
$\omega_X^\bullet = f^!_{new}\omega_Y^\bullet$이다.
\item $M \in D_{\textit{Coh}}^+(\mathcal{O}_Y)$에 대해\ 
$f^!_{new}M = D_X(Lf^*D_Y(M))$이다.
\item 또한\  $f$가 고유이면, $f^!_{new}$는\ 
$Rf_* : D_\QCoh(\mathcal{O}_X) \to D_\QCoh(\mathcal{O}_Y)$의 오른쪽 수반을\ 
$D_{\textit{Coh}}^+(\mathcal{O}_Y)$에 제한한 것과 동형이고, 모든\ 
$K \in D^-_{\textit{Coh}}(\mathcal{O}_X)$와\ 
$M \in D_{\textit{Coh}}^+(\mathcal{O}_Y)$에 대해 정준 동형
$$
Rf_*R\SheafHom_{\mathcal{O}_X}(K, f^!_{new}M)
\to
R\SheafHom_{\mathcal{O}_Y}(Rf_*K, M)
$$
과
$$
Rf_*R\SheafHom_{\mathcal{O}_X}(K, \omega_X^\bullet) =
R\SheafHom_{\mathcal{O}_Y}(Rf_*K, \omega_Y^\bullet)
$$
이\  $K \in D^-_{\textit{Coh}}(\mathcal{O}_X)$에 대해 성립한다.
\end{enumerate}
$X$가\  $S$ 위에서 분리이면\  $\omega_X^\bullet$은 정준적으로\ 
$(X \to S)^!\omega_S^\bullet$과 동형이다. 또한\  $f$가\  $S$ 위에서 분리인
스킴들 사이의 사상이면, $D_{\textit{Coh}}^+$의\  $K$에 대해 정준 동형\footnote{이
동형들이\  $(g \circ f)^! \to f^! \circ g^!$ 및\ 
$(g \circ f)^!_{new} \to f^!_{new} \circ g^!_{new}$와 호환되는지는 확인하지
않았다. 나중에 필요하면 여기서 이를 확인할 것이다.}
$f_{new}^!K = f^!K$가 있다.
\end{lemma}

\begin{proof}
$f : X \to Y$, $g : Y \to Z$, $h : Z \to T$를\  $S$ 위의 유한형 스킴들의
사상이라 하자. 다음 도식이 가환임을 보여야 한다.
$$
\xymatrix{
(h \circ g \circ f)^!_{new} \ar[r] \ar[d] &
f^!_{new} \circ (h \circ g)^!_{new} \ar[d] \\
(g \circ f)^!_{new} \circ h^!_{new} \ar[r] &
f^!_{new} \circ g^!_{new} \circ h^!_{new}
}
$$
$\eta_Y : \text{id} \to D_Y^2$와\  $\eta_Z : \text{id} \to D_Z^2$를\  Lemma
\ref{duality-lemma-dualizing-schemes}의 정준 동형이라 하자. 그러면\  Categories, Lemma
\ref{categories-lemma-properties-2-cat-cats}를 사용한 계산(생략함)에 의해 두 화살표\ 
$(h \circ g \circ f)^!_{new} \to f^!_{new} \circ g^!_{new} \circ h^!_{new}$는
모두
$$
1 \star \eta_Y \star 1 \star \eta_Z \star 1 :
D_X \circ Lf^* \circ Lg^* \circ Lh^* \circ D_T
\longrightarrow
D_X \circ Lf^* \circ D_Y^2 \circ Lg^* \circ D_Z^2 \circ Lh^* \circ D_T
$$
로 주어진다. 이것으로 (1)을 증명했다. (2)는\  $(X \to S)^!_{new}$의 정의와\ 
$D_S(\omega_S^\bullet) = \mathcal{O}_S$라는 사실에서 즉시 따른다. (3)은\  Lemma
\ref{duality-lemma-dualizing-schemes}이다. (4)는 (2)와 같은 논증으로부터 따른다. (5)는\ 
$f^!_{new}$의 정의이다.

\medskip\noindent
(6)의 증명. $a$를\  $S$ 위의 유한형 스킴들의 고유 사상\  $f : X \to Y$에 대한\ 
Lemma \ref{duality-lemma-twisted-inverse-image}의 오른쪽 수반이라 하자. 문제는\  $X$나\ 
$Y$가\  $S$ 위에서 분리임을 알지 못하며(일반적으로 이는 참이 아니다), 따라서\ 
$S$ 위의\  $f$에\  Lemma \ref{duality-lemma-shriek-via-duality}를 즉시 적용할 수 없다는
것이다. 이를 피하기 위해\  Lemma
\ref{duality-lemma-proper-map-good-dualizing-complex}의 정준 식별\ 
$\omega_X^\bullet = a(\omega_Y^\bullet)$을 사용한다. 따라서 기저 스킴을\  $Y$로
두고\  $f : X \to Y$에\  Lemma \ref{duality-lemma-shriek-via-duality}를 적용하면,
$f^!_{new}$는\  $a$를\  $D_{\textit{Coh}}^+(\mathcal{O}_Y)$에 제한한 것이다.
표시된 등식들은\  Example \ref{duality-example-iso-on-RSheafHom-noetherian}에서 따른다.

\medskip\noindent
마지막 주장들은 정규화된 쌍대화 복합체의 구성과 이미 사용한\  Lemma
\ref{duality-lemma-shriek-via-duality}에서 따른다.
\end{proof}
\begin{remark}
\label{duality-remark-independent-omega-S}
$S$를 쌍대화 복합체를 갖는\  Noether 스킴이라 하자. $f : X \to Y$를\  $S$ 위의
유한형 스킴들의 사상이라 하자. 그러면 함자
$$
f_{new}^! : D^+_{Coh}(\mathcal{O}_Y) \to D^+_{Coh}(\mathcal{O}_X)
$$
는 쌍대화 복합체\  $\omega_S^\bullet$의 선택과 무관하며, 그 무관성은 정준
동형까지이다. 증명의 개요를 설명하자. 다른 쌍대화 복합체는 모두\ 
$\omega_S^\bullet \otimes_{\mathcal{O}_S}^\mathbf{L} \mathcal{L}$의 꼴이며,
여기서\  $\mathcal{L}$은\  $D(\mathcal{O}_S)$의 가역 대상이다. Lemma
\ref{duality-lemma-dualizing-unique-schemes}을 보라. 임의의 유한형 분리 사상\ 
$p : U \to S$에 대해\  Lemma \ref{duality-lemma-compare-with-pullback-perfect}에 의해\ 
$p^!(\omega_S^\bullet \otimes^\mathbf{L}_{\mathcal{O}_S} \mathcal{L}) =
p^!(\omega_S^\bullet) \otimes^\mathbf{L}_{\mathcal{O}_U} Lp^*\mathcal{L}$이다.
따라서\  $\omega_X^\bullet$과\  $\omega_Y^\bullet$이\  $\omega_S^\bullet$에 대해
정규화된 쌍대화 복합체이면,
$\omega_X^\bullet \otimes_{\mathcal{O}_X}^\mathbf{L} La^*\mathcal{L}$과\ 
$\omega_Y^\bullet \otimes_{\mathcal{O}_Y}^\mathbf{L} Lb^*\mathcal{L}$은\ 
$\omega_S^\bullet \otimes_{\mathcal{O}_S}^\mathbf{L} \mathcal{L}$에 대해 정규화된
쌍대화 복합체임을 알 수 있다(여기서\  $a : X \to S$와\  $b : Y \to S$는 구조
사상이다). 이제\  $K \in D^+_{Coh}(\mathcal{O}_Y)$에 대해 다음과 같으므로 결과가
따른다.
\begin{align*}
& R\SheafHom_{\mathcal{O}_X}(Lf^*R\SheafHom_{\mathcal{O}_Y}(K,
\omega_Y^\bullet \otimes_{\mathcal{O}_Y}^\mathbf{L} Lb^*\mathcal{L}),
\omega_X^\bullet \otimes_{\mathcal{O}_X}^\mathbf{L} La^*\mathcal{L}) \\
& = R\SheafHom_{\mathcal{O}_X}(Lf^*R(\SheafHom_{\mathcal{O}_Y}(K,
\omega_Y^\bullet) \otimes_{\mathcal{O}_Y}^\mathbf{L} Lb^*\mathcal{L}),
\omega_X^\bullet \otimes_{\mathcal{O}_X}^\mathbf{L} La^*\mathcal{L}) \\
& = R\SheafHom_{\mathcal{O}_X}(Lf^*R\SheafHom_{\mathcal{O}_Y}(K,
\omega_Y^\bullet) \otimes_{\mathcal{O}_X}^\mathbf{L} La^*\mathcal{L},
\omega_X^\bullet \otimes_{\mathcal{O}_X}^\mathbf{L} La^*\mathcal{L}) \\
& = R\SheafHom_{\mathcal{O}_X}(Lf^*R\SheafHom_{\mathcal{O}_Y}(K,
\omega_Y^\bullet), \omega_X^\bullet)
\end{align*}
마지막 등식은\  $La^*\mathcal{L}$이\  $D(\mathcal{O}_X)$에서 가역이기 때문이다.
\end{remark}

\begin{example}
\label{duality-example-trace-proper}
$S$를\  Noether 스킴이라 하고\  $\omega_S^\bullet$을 쌍대화 복합체라 하자.
$f : X \to Y$를\  $S$ 위의 유한형 스킴들의 고유 사상이라 하자.
$\omega_X^\bullet$과\  $\omega_Y^\bullet$을\  $\omega_S^\bullet$에 대해 정규화된
쌍대화 복합체라 하자. 이 상황에서\ 
$a(\omega_Y^\bullet) = \omega_X^\bullet$이고(Lemma
\ref{duality-lemma-proper-map-good-dualizing-complex}), 따라서 대각합 사상(Section
\ref{duality-section-trace})은 정준 화살표
$$
\text{Tr}_f : Rf_*\omega_X^\bullet \longrightarrow \omega_Y^\bullet
$$
이며, 이는 다음 동형들을 준다(Lemma \ref{duality-lemma-duality-bootstrap}).
$$
\Hom_X(L, \omega_X^\bullet) = \Hom_Y(Rf_*L, \omega_Y^\bullet)
$$
그리고
$$
Rf_*R\SheafHom_{\mathcal{O}_X}(L, \omega_X^\bullet) =
R\SheafHom_{\mathcal{O}_Y}(Rf_*L, \omega_Y^\bullet)
$$
여기서\  $L$은\  $D_\QCoh(\mathcal{O}_X)$의 대상이다.
\end{example}

\begin{remark}
\label{duality-remark-dualizing-finite}
$S$를\  Noether 스킴이라 하고\  $\omega_S^\bullet$을 쌍대화 복합체라 하자.
$f : X \to Y$를\  $S$ 위의 유한형 스킴들 사이의 유한 사상이라 하자.
$\omega_X^\bullet$과\  $\omega_Y^\bullet$을\  $\omega_S^\bullet$에 대해 정규화된
쌍대화 복합체라 하자. 그러면\  Lemmas \ref{duality-lemma-finite-twisted}와\ 
\ref{duality-lemma-proper-map-good-dualizing-complex}에 의해\ 
$D_\QCoh^+(f_*\mathcal{O}_X)$에서
$$
f_*\omega_X^\bullet = R\SheafHom(f_*\mathcal{O}_X, \omega_Y^\bullet)
$$
이다. Example \ref{duality-example-trace-proper}의 대각합 사상은
$$
\text{Tr}_f : Rf_*\omega_X^\bullet = f_*\omega_X^\bullet =
R\SheafHom(f_*\mathcal{O}_X, \omega_Y^\bullet) \longrightarrow
\omega_Y^\bullet
$$
이며, 흔히 ``$1$에서의 값매김''이라고 불린다.
\end{remark}

\begin{remark}
\label{duality-remark-relative-dualizing-complex-shriek}
$f : X \to Y$를\  Situation \ref{duality-situation-dualizing}에서와 같은 쌍\ 
$(S, \omega_S^\bullet)$ 위의 유한형 스킴들의 평탄 고유 사상이라 하자. 상대
쌍대화 복합체(Remark \ref{duality-remark-relative-dualizing-complex})는\ 
$\omega_{X/Y}^\bullet = a(\mathcal{O}_Y)$이다. Lemma
\ref{duality-lemma-proper-map-good-dualizing-complex}에 의해 다음 식의 첫 번째 정준
동형을 얻는다.
$$
\omega_X^\bullet = a(\omega_Y^\bullet) =
Lf^*\omega_Y^\bullet \otimes_{\mathcal{O}_X}^\mathbf{L} \omega_{X/Y}^\bullet
$$
이는\  $D(\mathcal{O}_X)$에서의 식이다. 두 번째 정준 동형은\  Remark
\ref{duality-remark-relative-dualizing-complex}의 논의에서 따른다.
\end{remark}



\section{차원 함수}
\label{duality-section-dimension-functions}

\noindent
한 스킴 위의 유한형 스킴으로 넘어갈 때 차원 함수가 어떻게 변하는지에 관한
정보가 조금 더 필요하다.

\begin{lemma}
\label{duality-lemma-good-dualizing-normalized}
$S$를\  Noether 스킴이라 하고\  $\omega_S^\bullet$을 쌍대화 복합체라 하자.
$X$를\  $S$ 위의 유한형 스킴이라 하고\  $\omega_X^\bullet$을\ 
$\omega_S^\bullet$에 대해 정규화된 쌍대화 복합체라 하자. $x \in X$가\  $S$의
닫힌점\  $s$ 위에 놓이는 닫힌점이면, $\omega_{S, s}^\bullet$이\ 
$\mathcal{O}_{S, s}$ 위의 정규화된 쌍대화 복합체일 때\ 
$\omega_{X, x}^\bullet$은\  $\mathcal{O}_{X, x}$ 위의 정규화된 쌍대화
복합체이다.
\end{lemma}

\begin{proof}
$X$를\  $x$의 아핀 근방으로 바꿀 수 있으므로, $f : X \to S$가 분리라고
가정할 수 있고 그렇게 하자. 그러면\ 
$\omega_X^\bullet = f^!\omega_S^\bullet$이다. 다음 복합체가 차수\  $0$에
집중됨을 보여야 한다.
$R\Hom_{\mathcal{O}_{X, x}}(\kappa(x), \omega_{X, x}^\bullet)$
$i_x : x \to X$를 포함 사상이라 하자. $x$가 닫힌점이므로 이는 닫힌 몰입이다.
따라서\  Lemma \ref{duality-lemma-shriek-closed-immersion}에 의해\ 
$R\Hom_{\mathcal{O}_{X, x}}(\kappa(x), \omega_{X, x}^\bullet)$은\ 
$i_x^!\omega_X^\bullet$을 나타낸다. 가환 도식
$$
\xymatrix{
x \ar[r]_{i_x} \ar[d]_\pi & X \ar[d]^f \\
s \ar[r]^{i_s} & S
}
$$
을 생각하자. Morphisms, Lemma
\ref{morphisms-lemma-closed-point-fibre-locally-finite-type}에 의해 확대\ 
$\kappa(x)/\kappa(s)$는 유한이고, 따라서\  $\pi$는 유한 사상이다. 그러므로
$$
i_x^!\omega_X^\bullet = i_x^! f^! \omega_S^\bullet =
\pi^! i_s^! \omega_S^\bullet
$$
이다. 따라서\  $\omega_{S, s}^\bullet$이\  $\mathcal{O}_{S, s}$ 위의 정규화된
쌍대화 복합체이면 위와 같은 논증에 의해\ 
$i_s^!\omega_S^\bullet = \kappa(s)[0]$이다. 또한
$$
R\pi_*(\pi^!(\kappa(s)[0])) =
R\SheafHom_{\mathcal{O}_s}(R\pi_*(\kappa(x)[0]), \kappa(s)[0]) =
\widetilde{\Hom_{\kappa(s)}(\kappa(x), \kappa(s))}
$$
이다. 첫 번째 등식은\  $L = \kappa(x)[0]$로 둔\  Example
\ref{duality-example-iso-on-RSheafHom-noetherian}에서 따르고, 두 번째 등식은\ 
$\pi_*$가 완전하기 때문에 성립한다. 따라서\  $\pi^!(\kappa(s)[0])$는 차수\ 
$0$에 집중되어 있고, 결론을 얻는다.
\end{proof}
\begin{lemma}
\label{duality-lemma-good-dualizing-dimension-function}
$S$를\  Noether 스킴이라 하고\  $\omega_S^\bullet$을 쌍대화 복합체라 하자.
$f : X \to S$를 유한형이라 하고\  $\omega_X^\bullet$을\ 
$\omega_S^\bullet$에 대해 정규화된 쌍대화 복합체라 하자. 모든\  $x \in X$에
대해
$$
\delta_X(x) - \delta_S(f(x)) = \text{trdeg}_{\kappa(f(x))}(\kappa(x))
$$
가 성립한다. 여기서\  $\delta_S$와\  $\delta_X$는 각각\ 
$\omega_S^\bullet$과\  $\omega_X^\bullet$의 차원 함수이다. Lemma
\ref{duality-lemma-dimension-function-scheme}을 보라.
\end{lemma}

\begin{proof}
$X$를\  $x$의 아핀 근방으로 바꿀 수 있다. 따라서\  $S$ 위의 콤팩트화\ 
$X \subset \overline{X}$가 존재한다고 가정할 수 있고 그렇게 하자. 이제\  $X$를\ 
$\overline{X}$로 바꾸어\  $X$가\  $S$ 위에서 고유라고 가정할 수 있다. 또한\  $X$를\ 
$X$의\  $x$를 포함하는 연결 성분으로 바꾸어\  $X$가 연결이라고 가정할 수 있다. 다음으로\ 
$\delta_X$와 함수\ 
$x \mapsto \delta_S(f(x)) + \text{trdeg}_{\kappa(f(x))}(\kappa(x))$는 모두\ 
$X$ 위의 차원 함수임을 상기하자. Morphisms, Lemma
\ref{morphisms-lemma-dimension-function-propagates}를 보라(또한\  Lemma
\ref{duality-lemma-dimension-function-scheme}에 의해\  $S$가 보편적\  catenary라는 사실을
사용한다). Topology, Lemma \ref{topology-lemma-dimension-function-unique}에 의해
그 차이는 국소 상수이고, $X$가 연결이므로 상수이다. 따라서\  $X$의 임의의 한
점에서 등식을 증명하면 충분하다. Properties, Lemma
\ref{properties-lemma-locally-Noetherian-closed-point}에 의해\  $X$에는 닫힌점\ 
$x$가 있다. $X \to S$가 고유이므로\  $x$의 상\  $s$는\  $S$에서 닫혀 있다.
따라서\  Lemma \ref{duality-lemma-good-dualizing-normalized}를 적용하여 결론을 얻는다.
\end{proof}
\begin{lemma}
\label{duality-lemma-shriek}
Situation \ref{duality-situation-shriek}에서\  $f : X \to Y$를\  $\textit{FTS}_S$의
사상이라 하자. $x \in X$이고 그 상이\  $y \in Y$라고 하자. 그러면
$$
H^i(f^!\mathcal{O}_Y)_x \not = 0
\Rightarrow - \dim_x(X_y) \leq i.
$$
\end{lemma}

\begin{proof}
명제는\  $X$ 위에서 국소적이므로\  $X$와\  $Y$가 아핀 스킴이라고 가정할 수 있다.
$X = \Spec(A)$와\  $Y = \Spec(R)$로 쓰자. 그러면\  $f^!\mathcal{O}_Y$는\ 
Dualizing Complexes, Section
\ref{dualizing-section-relative-dualizing-complexes-Noetherian}의 상대 쌍대화
복합체\  $\omega_{A/R}^\bullet$에 대응한다(Remark
\ref{duality-remark-local-calculation-shriek}). 따라서 보조정리는\  Dualizing
Complexes, Lemma
\ref{dualizing-lemma-relative-dualizing-trivial-vanishing}에서 따른다.
\end{proof}

\begin{lemma}
\label{duality-lemma-flat-shriek}
Situation \ref{duality-situation-shriek}에서\  $f : X \to Y$를\  $\textit{FTS}_S$의
사상이라 하자. $x \in X$이고 그 상이\  $y \in Y$라고 하자. $f$가 평탄이면
$$
H^i(f^!\mathcal{O}_Y)_x \not = 0
\Rightarrow - \dim_x(X_y) \leq i \leq 0.
$$
실제로\  $f$의 모든 올의 차원이\  $\leq d$이면,
$f^!\mathcal{O}_Y$는\  $D(X, f^{-1}\mathcal{O}_Y)$의 대상으로서\  $[-d, 0]$ 안의\ 
tor-진폭을 갖는다.
\end{lemma}

\begin{proof}
Lemma \ref{duality-lemma-shriek}의 증명과 정확히 같은 방식으로 논하면, 이는\  Dualizing
Complexes, Lemma \ref{dualizing-lemma-relative-dualizing-flat-vanishing}에서
따른다.
\end{proof}
\begin{lemma}
\label{duality-lemma-shriek-over-CM}
Situation \ref{duality-situation-shriek}에서\  $f : X \to Y$를\  $\textit{FTS}_S$의
사상이라 하자. $x \in X$이고 그 상이\  $y \in Y$라고 하자. 다음을 가정하자.
\begin{enumerate}
\item $\mathcal{O}_{Y, y}$는\  Cohen-Macaulay이다.
\item $x$를 포함하는\  $X$의 기약 성분의 임의의 일반점\  $\xi$에 대해\ 
$\text{trdeg}_{\kappa(f(\xi))}(\kappa(\xi)) \leq r$이다.
\end{enumerate}
그러면
$$
H^i(f^!\mathcal{O}_Y)_x \not = 0
\Rightarrow - r \leq i
$$
이고, 줄기\  $H^{-r}(f^!\mathcal{O}_Y)_x$는\  $\mathcal{O}_{X, x}$-가군으로서\ 
$(S_2)$이다.
\end{lemma}

\begin{proof}
$X$를\  $x$의 열린 근방으로 바꾼 뒤, $X$의 모든 기약 성분이\  $x$를 지난다고
가정할 수 있다. 그러면\  Lemma \ref{duality-lemma-shriek}의 증명과 정확히 같은 방식으로
논하여\  Dualizing Complexes, Lemma
\ref{dualizing-lemma-relative-dualizing-CM-vanishing}에서 결론을 얻는다.
\end{proof}

\begin{lemma}
\label{duality-lemma-flat-quasi-finite-shriek}
Situation \ref{duality-situation-shriek}에서\  $f : X \to Y$를\  $\textit{FTS}_S$의
사상이라 하자. $f$가 평탄하고 준유한이면
$$
f^!\mathcal{O}_Y = \omega_{X/Y}[0]
$$
이며, 여기서\  $\omega_{X/Y}$는\  $Y$ 위에서 평탄인 어떤 연접\ 
$\mathcal{O}_X$-가군이다.
\end{lemma}

\begin{proof}
Lemma \ref{duality-lemma-flat-shriek}와, Lemma \ref{duality-lemma-shriek-coherent}에 의해\ 
$f^!\mathcal{O}_Y$의 코호몰로지 층들이 연접이라는 사실의 귀결이다.
\end{proof}
\begin{lemma}
\label{duality-lemma-CM-shriek}
Situation \ref{duality-situation-shriek}에서\  $f : X \to Y$를\  $\textit{FTS}_S$의
사상이라 하자. $f$가\  Cohen-Macaulay이면(More on Morphisms, Definition
\ref{more-morphisms-definition-CM})
$$
f^!\mathcal{O}_Y = \omega_{X/Y}[d]
$$
이며, 여기서\  $\omega_{X/Y}$는\  $Y$ 위에서 평탄인 어떤 연접\ 
$\mathcal{O}_X$-가군이고\  $d$는\  $Y$ 위에서의\  $X$의 상대 차원을 주는\  $X$ 위의
국소 상수 함수이다.
\end{lemma}

\begin{proof}
상대 차원\  $d$는\  Morphisms, Lemma
\ref{morphisms-lemma-flat-finite-presentation-CM-fibres-relative-dimension}에 의해
잘 정의되고 국소 상수이다. $f^!\mathcal{O}_Y$의 코호몰로지 층들은\  Lemma
\ref{duality-lemma-shriek-coherent}에 의해 연접이다. $f^!\mathcal{O}_Y$의 다른
코호몰로지 층들이 영임을 보일 수 있다면, Lemma \ref{duality-lemma-flat-shriek}에서\ 
$\omega_{X/Y}$의 평탄성을 얻는다.

\medskip\noindent
문제는\  $X$ 위에서 국소적이므로\  $X$와\  $Y$가 아핀이고 사상의 상대 차원이\ 
$d$라고 가정할 수 있다. $d = 0$이면 결과는\  Lemma
\ref{duality-lemma-flat-quasi-finite-shriek}에서 직접 따른다. $d > 0$이면\  More on
Morphisms, Lemma
\ref{more-morphisms-lemma-flat-finite-presentation-characterize-CM}에 의해\  $g$가
준유한이고 평탄인 분해
$$
X \xrightarrow{g} \mathbf{A}^d_Y \xrightarrow{p} Y
$$
가 존재한다고 가정할 수 있다. 이때\  $f^! = g^! \circ p^!$이다. Lemma
\ref{duality-lemma-shriek-affine-line}에 의해\ 
$p^!\mathcal{O}_Y \cong \mathcal{O}_{\mathbf{A}^d_Y}[-d]$임을 알 수 있다.
$d = 0$인 경우를 적용하면 결론을 얻는다.
\end{proof}

\begin{remark}
\label{duality-remark-the-same-is-true}
$S$를 쌍대화 복합체\  $\omega_S^\bullet$이 주어진\  Noether 스킴이라 하자. 이때\ 
Lemmas \ref{duality-lemma-shriek}, \ref{duality-lemma-flat-shriek},
\ref{duality-lemma-flat-quasi-finite-shriek}, 그리고\  \ref{duality-lemma-CM-shriek}은\  $S$ 위의
유한형 스킴들의 임의의 사상\  $f : X \to Y$에 대해 성립하며, 다만\  $f^!$를\ 
$f_{new}^!$로 바꾸어야 한다. 각 경우 증명이 즉시 아핀인 경우로 환원되고,
그때\  Lemma \ref{duality-lemma-duality-bootstrap}에 의해\  $f^! = f_{new}^!$이므로 이는
분명하다.
\end{remark}



\section{쌍대화 가군}
\label{duality-section-dualizing-module}


\noindent
이 절은\  Dualizing Complexes, Section \ref{dualizing-section-dualizing-module}에
이어진다.

\medskip\noindent
$X$를\  Noether 스킴이라 하고\  $\omega_X^\bullet$을 쌍대화 복합체라 하자.
$H^n(\omega_X^\bullet)$이 영이 아닌 가장 작은 정수\  $n \in \mathbf{Z}$를
택하자. 달리 말하면\  $-n$은\  $\omega_X^\bullet$에 결부된 차원 함수의 최댓값이다
(Lemma \ref{duality-lemma-dimension-function-scheme}). 때로는\ 
$H^n(\omega_X^\bullet)$을\  $X$의 {\it 쌍대화 가군} 또는 {\it 쌍대화 층}이라
하며, 흔히\  $\omega_X$로 나타낸다. 앞으로 ``$\omega_X$를 쌍대화 가군이라
하자''라는 말은 위의 상황을 뜻한다.

\medskip\noindent
Noether 스킴\  $X$ 위의 쌍대화 가군\  $\omega_X$를 사용할 때에는 다음 사항에
주의해야 한다.
\begin{enumerate}
\item $X$에서\  $X$의 열린집합\  $U$로 넘어갈 때 정수\  $n$이 바뀔 수 있으며, 이 경우\ 
$\omega_X|_U = \omega_U$는 참이 아니다.
\item 쌍대화 복합체는 유일하지 않으며, 쌍대화 가군은 가역 가군을 텐서하는
것까지로만 유일하다.
\end{enumerate}
두 번째 문제는 흔히 중요하지 않다. 고정된 쌍대화 복합체\ 
$\omega_S^\bullet$이 주어진 기저변환\  $S$ 위의 유한형\  $X$를 다룰 것이고,
$\omega_X^\bullet$은\  $\omega_S^\bullet$에 대해 정규화된 쌍대화 복합체이기
때문이다. $X$가 등차원이면 첫 번째 문제는 생기지 않는다. 더 정확히 말해\ 
$\omega_X^\bullet$에 결부된 차원 함수(Lemma
\ref{duality-lemma-dimension-function-scheme})가\  $X$의 모든 일반점을 같은 정수로 보내면
그 문제는 생기지 않는다.
\begin{example}
\label{duality-example-proper-over-local}
$(A, \mathfrak m, \kappa)$가 국소\  Noether 환이고\  $S = \Spec(A)$이며,
$\omega_S^\bullet$이 정규화된 쌍대화 복합체\  $\omega_A^\bullet$에 대응한다고
하자. $f : X \to S$가\  $S$ 위에서 고유이고\ 
$\omega_X^\bullet = f^!\omega_S^\bullet$이면, 연접층
$$
\omega_X = H^{-\dim(X)}(\omega_X^\bullet)
$$
은 쌍대화 가군이고 흔히\  $X$의 쌍대화 가군이라 한다($S$와\ 
$\omega_S^\bullet$은 문맥에서 이해한다). 이것이 좋은 성질들을 가짐을 보게 될
것이다.
\end{example}

\begin{example}
\label{duality-example-equidimensional-over-field}
$X$가 체\  $k$ 위의 유한형 등차원 스킴이라고 하자. 이때\ 
$\omega_X^\bullet$을\  $k[0]$에 대해 정규화된 쌍대화 복합체로 택하고
$$
\omega_X = H^{-\dim(X)}(\omega_X^\bullet)
$$
을\  $X$의 쌍대화 가군이라 부르는 것이 관례이다. $X$가\  $k$ 위에서 분리이면,
Lemma \ref{duality-lemma-duality-bootstrap}에 의해\ 
$\omega_X^\bullet = f^!\mathcal{O}_{\Spec(k)}$이며, 여기서\ 
$f : X \to \Spec(k)$는 구조 사상이다. $X$가\  $k$ 위에서 고유이면 이는\  Example
\ref{duality-example-proper-over-local}의 특수한 경우이다.
\end{example}
\begin{lemma}
\label{duality-lemma-dualizing-module}
$X$를 연결\  Noether 스킴이라 하고\  $\omega_X$를\  $X$ 위의 쌍대화 가군이라 하자.
$\omega_X$의 지지집합은 임의의 차원 함수에 관해 최대 차원을 갖는 기약 성분들의
합집합이고, $\omega_X$는 성질\  $(S_2)$를 갖는 연접\  $\mathcal{O}_X$-가군이다.
\end{lemma}

\begin{proof}
위에서 설명한 규약에 의해, $\omega_X$가 가장 왼쪽의 영이 아닌 코호몰로지
층인 쌍대화 복합체\  $\omega_X^\bullet$이 존재한다. $X$가 연결이므로 임의의 두
차원 함수의 차이는 상수이다(Topology, Lemma
\ref{topology-lemma-dimension-function-unique}). 따라서\ 
$\omega_X^\bullet$에 결부된 차원 함수를 사용할 수 있다(Lemma
\ref{duality-lemma-dimension-function-scheme}). 이러한 관찰 아래 보조정리는 이제\ 
Dualizing Complexes, Lemma \ref{dualizing-lemma-depth-dualizing-module}와
정의들, 특히\  Cohomology of Schemes, Definition
\ref{coherent-definition-depth}에서 따른다.
\end{proof}
\begin{lemma}
\label{duality-lemma-vanishing-good-dualizing}
$X/A$, $\omega_X^\bullet$, 그리고\  $\omega_X$를\  Example
\ref{duality-example-proper-over-local}에서와 같이 두자. 그러면 다음이 성립한다.
\begin{enumerate}
\item $H^i(\omega_X^\bullet) \not = 0 \Rightarrow
i \in \{-\dim(X), \ldots, 0\}$이다.
\item $H^i(\omega_X^\bullet)$의 지지집합의 차원은\  $-i$ 이하이다.
\item $\text{Supp}(\omega_X)$는 차원이\  $\dim(X)$인 성분들의 합집합이다.
\item $\omega_X$는 성질\  $(S_2)$를 갖는다.
\end{enumerate}
\end{lemma}

\begin{proof}
$\delta_X$와\  $\delta_S$를\  Lemma
\ref{duality-lemma-good-dualizing-dimension-function}에서처럼\  $\omega_X^\bullet$과\ 
$\omega_S^\bullet$에 결부된 차원 함수라 하자. $X$가\  $A$ 위에서 고유이므로,
$X$의 모든 닫힌 부분스킴은 닫힌점\  $s \in S$로 사상되는 닫힌점\  $x$를 포함하며\ 
$\delta_X(x) = \delta_S(s) = 0$이다. 따라서 임의의 점\  $\xi \in X$에 대해\ 
$\delta_X(\xi) = \dim(\overline{\{\xi\}})$이다. 그러므로 보조정리의 각 주장은\ 
$\Spec(\mathcal{O}_{X, x})$ 위에서 일어나는 일을 보아 확인할 수 있고, 이 경우
결과는\  Dualizing Complexes, Lemmas
\ref{dualizing-lemma-sitting-in-degrees}와\ 
\ref{dualizing-lemma-depth-dualizing-module}에서 따른다. 일부 세부 사항은
생략한다. 마지막 두 주장은\  Lemma \ref{duality-lemma-dualizing-module}에서도 유도할 수
있다.
\end{proof}
\begin{lemma}
\label{duality-lemma-dualizing-module-proper-over-A}
쌍대화 가군\  $\omega_X$를 갖는\  $X/A$를\  Example
\ref{duality-example-proper-over-local}에서와 같이 두자. $d = \dim(X_s)$를 닫힌 올의
차원이라 하자. $\dim(X) = d + \dim(A)$이면, 쌍대화 가군\  $\omega_X$는 연접\ 
$\mathcal{O}_X$-가군들의 범주 위의 함자
$$
\mathcal{F} \longmapsto \Hom_A(H^d(X, \mathcal{F}), \omega_A)
$$
를 표현한다.
\end{lemma}

\begin{proof}
다음이 성립한다.
\begin{align*}
\Hom_X(\mathcal{F}, \omega_X)
& =
\Ext^{-\dim(X)}_X(\mathcal{F}, \omega_X^\bullet) \\
& =
\Hom_X(\mathcal{F}[\dim(X)], \omega_X^\bullet) \\
& =
\Hom_X(\mathcal{F}[\dim(X)], f^!(\omega_A^\bullet)) \\
& =
\Hom_S(Rf_*\mathcal{F}[\dim(X)], \omega_A^\bullet) \\
& =
\Hom_A(H^d(X, \mathcal{F}), \omega_A)
\end{align*}
첫 번째 등식은\  $i < -\dim(X)$에 대해\  $H^i(\omega_X^\bullet) = 0$이기 때문이다.
Lemma \ref{duality-lemma-vanishing-good-dualizing}와\  Derived Categories, Lemma
\ref{derived-lemma-negative-exts}를 보라. 두 번째 등식은\  Ext 군의 정의에서
따른다. 세 번째 등식은 우리가\  $\omega_X^\bullet$을 택한 방식이다. 네 번째
등식은\  $f^!$가\  $f$에 대한\  Lemma \ref{duality-lemma-twisted-inverse-image}의 오른쪽
수반이기 때문에 성립한다. Section \ref{duality-section-duality}를 보라. 마지막 등식은\ 
$i > d$에 대해\  $R^if_*\mathcal{F}$가 영이고(Cohomology of Schemes, Lemma
\ref{coherent-lemma-higher-direct-images-zero-above-dimension-fibre}),
$j < -\dim(A)$에 대해\  $H^j(\omega_A^\bullet)$이 영이기 때문에 성립한다.
\end{proof}







\section{Cohen-Macaulay 스킴}
\label{duality-section-CM}

\noindent
이 절은\  Dualizing Complexes, Section \ref{dualizing-section-CM}에 이어진다.
Cohen-Macaulay 스킴에서는 쌍대성이 특히 단순한 형태를 취한다.

\begin{lemma}
\label{duality-lemma-dualizing-module-CM-scheme}
$X$를 쌍대화 복합체\  $\omega_X^\bullet$을 갖는 국소\  Noether 스킴이라 하자.
\begin{enumerate}
\item $X$가\  Cohen-Macaulay이다\  $\Leftrightarrow$ $\omega_X^\bullet$이 국소적으로
영이 아닌 코호몰로지 층을 유일하게 갖는다.
\item $\mathcal{O}_{X, x}$가\  Cohen-Macaulay이다\  $\Leftrightarrow$
$\omega_{X, x}^\bullet$이 영이 아닌 코호몰로지를 유일하게 갖는다.
\item $U = \{x \in X \mid \mathcal{O}_{X, x}\text{ is Cohen-Macaulay}\}$는
열린집합이고\  Cohen-Macaulay이다.
\end{enumerate}
$X$가 연결이고\  Cohen-Macaulay이면, 어떤 정수\  $n$과\  Cohen-Macaulay인 연접\ 
$\mathcal{O}_X$-가군\  $\omega_X$가 존재하여\ 
$\omega_X^\bullet = \omega_X[-n]$이다.
\end{lemma}
\begin{proof}
정의와\  Dualizing Complexes, Lemma
\ref{dualizing-lemma-dualizing-localize}에 의해 모든\  $x \in X$에 대해 복합체\ 
$\omega_{X, x}^\bullet$은\  $\mathcal{O}_{X, x}$ 위의 쌍대화 복합체이다.
Dualizing Complexes, Lemma \ref{dualizing-lemma-apply-CM}에 의해 (2)가 성립한다.

\medskip\noindent
(3)을 보이기 위해\  $\mathcal{O}_{X, x}$가\  Cohen-Macaulay라고 가정하자.
$H^{n_{x}}(\omega_{X, x}^\bullet)$이 영이 아닌 유일한 정수를\  $n_x$라 하자.
$x$의 아핀 근방\  $V \subset X$에 대해\  $\omega_X^\bullet|_V$는\ 
$D^b_{\textit{Coh}}(\mathcal{O}_V)$에 속하므로, 영이 아닌 연접 가군\ 
$H^i(\omega_X^\bullet)|_V$는 유한 개뿐이다. 따라서\  $V$를 축소하여\ 
$H^{n_x}$만 영이 아니라고 가정할 수 있다. Modules, Lemma
\ref{modules-lemma-finite-type-stalk-zero}를 보라. 이로써 모든\  $v \in V$에 대해\ 
$\mathcal{O}_{X, v}$가\  Cohen-Macaulay임을 알 수 있다. 이는\  $U$가 열린집합이며\ 
Cohen-Macaulay 스킴임을 동시에 증명한다.

\medskip\noindent
(1)의 증명. 함의\  $\Leftarrow$는 (2)에서 따른다. 함의\  $\Rightarrow$는 앞 문단의
논의에서 따른다. 거기서\  $\mathcal{O}_{X, x}$가\  Cohen-Macaulay이면\  $x$의 한
근방에서 복합체\  $\omega_X^\bullet$이 영이 아닌 코호몰로지 층을 하나만 가짐을
보였다.

\medskip\noindent
$X$가 연결이고\  Cohen-Macaulay라고 가정하자. 위의 논의에 의해 사상\ 
$x \mapsto n_x$는 국소 상수이다. $X$가 연결이므로 이는 어떤\  $n$과 같은
상수이다. $\omega_X = H^n(\omega_X^\bullet)$로 두면 보조정리가 성립한다.
실제로\  Dualizing Complexes, Lemma \ref{dualizing-lemma-apply-CM}에 의해\ 
$\omega_X$는\  Cohen-Macaulay이다(Cohomology of Schemes, Definition
\ref{coherent-definition-Cohen-Macaulay}도 보라).
\end{proof}
\begin{lemma}
\label{duality-lemma-has-dualizing-module-CM-scheme}
$X$를 국소\  Noether 스킴이라 하자. $\omega_X[0]$이\  $X$ 위의 쌍대화 복합체가
되는 연접층\  $\omega_X$가 존재하면, $X$는\  Cohen-Macaulay 스킴이다.
\end{lemma}

\begin{proof}
이는\  Dualizing Complexes, Lemma
\ref{dualizing-lemma-has-dualizing-module-CM}와 우리의 정의들에서 즉시 따른다.
\end{proof}
\begin{lemma}
\label{duality-lemma-affine-flat-Noetherian-CM}
Situation \ref{duality-situation-shriek}에서\  $f : X \to Y$를\  $\textit{FTS}_S$의 사상이라
하자. $x \in X$라 하자. $f$가 평탄이면 다음은 동치이다.
\begin{enumerate}
\item $f$는\  $x$에서\  Cohen-Macaulay이다.
\item $f^!\mathcal{O}_Y$는\  $x$의 한 근방에서 영이 아닌 코호몰로지 층을
유일하게 갖는다.
\end{enumerate}
\end{lemma}

\begin{proof}
한 방향은\  Lemma \ref{duality-lemma-CM-shriek}에서 따른다. 역을 증명하기 위해\ 
$f^!\mathcal{O}_Y$가 영이 아닌 코호몰로지 층을 유일하게 갖는다고 가정할 수
있다. $y = f(x)$라 하자. $\xi_1, \ldots, \xi_n \in X_y$를\  $x$로 특수화되는
올\  $X_y$의 일반점들이라 하자. $d_1, \ldots, d_n$을\  $X_y$의 대응하는 기약
성분들의 차원이라 하자. 사상\  $f : X \to Y$는\  More on Morphisms, Lemma
\ref{more-morphisms-lemma-flat-finite-presentation-CM-open}에 의해\  $\eta_i$에서\ 
Cohen-Macaulay이다. 따라서\  Lemma \ref{duality-lemma-CM-shriek}에 의해\ 
$d_1 = \ldots = d_n$이다. 그 공통값을\  $d$라 하면\  $d = \dim_x(X_y)$이다.
$X$를 축소하여 모든 올의 차원이\  $d$ 이하라고 가정할 수 있다(Morphisms,
Lemma \ref{morphisms-lemma-openness-bounded-dimension-fibres}). 그러면 유일하게
영이 아닌 코호몰로지 층\  $\omega = H^{-d}(f^!\mathcal{O}_Y)$는\  Lemma
\ref{duality-lemma-flat-shriek}에 의해\  $Y$ 위에서 평탄이다. 따라서\ 
$h : X_y \to X$를 표준 사상이라 하면, Derived Categories of Schemes, Lemma
\ref{perfect-lemma-tor-independence-and-tor-amplitude}에 의해\ 
$Lh^*(f^!\mathcal{O}_Y) = Lh^*(\omega[d]) = (h^*\omega)[d]$이다.
그러므로\  Lemma \ref{duality-lemma-relative-dualizing-fibres}에 의해\  $h^*\omega[d]$는\ 
$X_y$의 쌍대화 복합체이다. 따라서\  Lemma
\ref{duality-lemma-dualizing-module-CM-scheme}에 의해\  $X_y$는\  Cohen-Macaulay이다.
이는 원하는 대로\  $f$가\  $x$에서\  Cohen-Macaulay임을 증명한다.
\end{proof}
\begin{remark}
\label{duality-remark-CM-morphism-compare-dualizing}
Situation \ref{duality-situation-shriek}에서\  $f : X \to Y$를\  $\textit{FTS}_S$의 사상이라
하자. $f$가 상대 차원\  $d$의\  Cohen-Macaulay 사상이라고 가정하자.
$\omega_{X/Y} = H^{-d}(f^!\mathcal{O}_Y)$를\  $f^!\mathcal{O}_Y$의 유일하게 영이
아닌 코호몰로지 층이라 하자. Lemma \ref{duality-lemma-CM-shriek}를 보라. 그러면\ 
$K \in D^+_\QCoh(\mathcal{O}_Y)$에 대해 정준 동형
$$
f^!K = Lf^*K \otimes_{\mathcal{O}_X}^\mathbf{L} \omega_{X/Y}[d]
$$
이 존재한다. Lemma \ref{duality-lemma-perfect-comparison-shriek}를 보라. 특히\  $S$가
쌍대화 복합체\  $\omega_S^\bullet$을 갖고,
$\omega_Y^\bullet = (Y \to S)^!\omega_S^\bullet$이며,
$\omega_X^\bullet = (X \to S)^!\omega_S^\bullet$이면
$$
\omega_X^\bullet =
Lf^*\omega_Y^\bullet \otimes_{\mathcal{O}_X}^\mathbf{L} \omega_{X/Y}[d]
$$
이다. 더 나아가\  $X$와\  $Y$가 연결이고\  Cohen-Macaulay이며\  $\omega_Y$와\ 
$\omega_X$가 각각\  $\omega_Y^\bullet$과\  $\omega_X^\bullet$의 유일하게 영이
아닌 코호몰로지 층을 나타낸다면
$$
\omega_X = f^*\omega_Y \otimes_{\mathcal{O}_X} \omega_{X/Y}.
$$
이다. $S$ 위의 임의의 유한형 스킴\  $X$와\  $Y$에 대해서도(즉, 반드시\  $S$ 위에서
분리일 필요는 없다) Section \ref{duality-section-glue}에서처럼\  $\omega_S^\bullet$에
대해 정규화된 쌍대화 복합체를 사용하면 비슷한 결과가 성립한다.
\end{remark}




\section{Gorenstein 스킴}
\label{duality-section-gorenstein}

\noindent
이 절은\  Dualizing Complexes, Section
\ref{dualizing-section-gorenstein}에 이어진다.

\begin{definition}
\label{duality-definition-gorenstein}
$X$를 스킴이라 하자. $X$가 국소\  Noether이고 모든\  $x \in X$에 대해\ 
$\mathcal{O}_{X, x}$가\  Gorenstein이면\  $X$를 {\it Gorenstein}이라 한다.
\end{definition}

\noindent
이 정의가 의미 있는 까닭은\  Noether 환이\  Gorenstein이라는 것이 그 모든 국소환이\ 
Gorenstein이라는 것과 동치이기 때문이다. Dualizing Complexes, Definition
\ref{dualizing-definition-gorenstein}을 보라.

\begin{lemma}
\label{duality-lemma-gorenstein-CM}
Gorenstein 스킴은\  Cohen-Macaulay이다.
\end{lemma}

\begin{proof}
아핀 국소적으로 보면 이는 대수에서의 대응하는 결과인\  Dualizing Complexes,
Lemma \ref{dualizing-lemma-gorenstein-CM}에서 따른다.
\end{proof}

\begin{lemma}
\label{duality-lemma-regular-gorenstein}
정칙 스킴은\  Gorenstein이다.
\end{lemma}

\begin{proof}
아핀 국소적으로 보면 이는 대수에서의 대응하는 결과인\  Dualizing Complexes,
Lemma \ref{dualizing-lemma-regular-gorenstein}에서 따른다.
\end{proof}
\begin{lemma}
\label{duality-lemma-gorenstein}
$X$를 국소\  Noether 스킴이라 하자.
\begin{enumerate}
\item $X$가 쌍대화 복합체\  $\omega_X^\bullet$을 가지면 다음이 성립한다.
\begin{enumerate}
\item $X$가\  Gorenstein이다\  $\Leftrightarrow$ $\omega_X^\bullet$은\ 
$D(\mathcal{O}_X)$의 가역 대상이다.
\item $\mathcal{O}_{X, x}$가\  Gorenstein이다\  $\Leftrightarrow$
$\omega_{X, x}^\bullet$은\  $D(\mathcal{O}_{X, x})$의 가역 대상이다.
\item $U = \{x \in X \mid \mathcal{O}_{X, x}\text{ is Gorenstein}\}$은 열린\ 
Gorenstein 부분스킴이다.
\end{enumerate}
\item $X$가\  Gorenstein이면, $X$가 쌍대화 복합체를 갖는 것은\ 
$\mathcal{O}_X[0]$이 쌍대화 복합체인 것과 동치이다.
\end{enumerate}
\end{lemma}

\begin{proof}
아핀 국소적으로 보면 이는 대수에서의 대응하는 결과인\  Dualizing Complexes,
Lemma \ref{dualizing-lemma-gorenstein}에서 따른다.
\end{proof}
\begin{lemma}
\label{duality-lemma-gorenstein-lci}
$f : Y \to X$가 국소 완전교차 사상이고\  $X$가\  Gorenstein 스킴이면, $Y$도\ 
Gorenstein이다.
\end{lemma}

\begin{proof}
More on Morphisms, Lemma \ref{more-morphisms-lemma-affine-lci}에 의해 환 사상에
관한 대응하는 명제를 증명하면 충분하다. 이것은\  Dualizing Complexes, Lemma
\ref{dualizing-lemma-gorenstein-lci}이다.
\end{proof}

\begin{lemma}
\label{duality-lemma-gorenstein-local-syntomic}
성질\  $\mathcal{P}(S) =$``$S$ is Gorenstein''은\  syntomic 위상에 관해 국소적이다.
\end{lemma}

\begin{proof}
$\{S_i \to S\}$를\  syntomic 피복이라 하자. 스킴\  $S$가 국소\  Noether인 것은 각\ 
$S_i$가\  Noether인 것과 동치이다. Descent, Lemma
\ref{descent-lemma-Noetherian-local-fppf}를 보라. 따라서 이제\  $S$와\  $S_i$가
국소\  Noether라고 가정할 수 있다. $S$가\  Gorenstein이면\  Lemma
\ref{duality-lemma-gorenstein-lci}에 의해 각\  $S_i$가\  Gorenstein이다. 역으로 각\  $S_i$가\ 
Gorenstein이면, 각 점\  $s \in S$에 대해\  $s$로 가는 어떤\  $i$와\  $t \in S_i$를
택할 수 있다. 그러면\  $\mathcal{O}_{S, s} \to \mathcal{O}_{S_i, t}$는\ 
Gorenstein인\  $\mathcal{O}_{S_i, t}$를 공역으로 갖는 평탄 국소환 준동형이다.
따라서\  Dualizing Complexes, Lemma
\ref{dualizing-lemma-flat-under-gorenstein}에 의해\  $\mathcal{O}_{S, s}$는\ 
Gorenstein이다.
\end{proof}





\section{Gorenstein 사상}
\label{duality-section-gorenstein-morphisms}

\noindent
이 절은 일련의 절 가운데 하나이다. 정규 사상, 정칙 사상, Cohen-Macaulay 사상에
대응하는 절들은\  More on Morphisms, Sections
\ref{more-morphisms-section-normal},
\ref{more-morphisms-section-regular}, 그리고\ 
\ref{more-morphisms-section-CM}에서 찾을 수 있다.

\medskip\noindent
다음 보조정리는 기하적으로\  Gorenstein인 스킴을 따로 정의하는 것은 의미가 없음을
말한다. 그런 스킴은\  Gorenstein 스킴과 같기 때문이다.

\begin{lemma}
\label{duality-lemma-gorenstein-base-change}
$X$를 체\  $k$ 위의 국소\  Noether 스킴이라 하자. $k'/k$를 유한 생성 체 확대라
하자. $x \in X$를 한 점이라 하고\  $x' \in X_{k'}$를\  $x$ 위에 놓이는 점이라
하자. 그러면
$$
\mathcal{O}_{X, x}\text{ is Gorenstein}
\Leftrightarrow
\mathcal{O}_{X_{k'}, x'}\text{ is Gorenstein}
$$
이다. $X$가\  $k$ 위에서 국소 유한형이면 임의의 체 확대\  $k'/k$에 대해서도 같은
결론이 성립한다.
\end{lemma}

\begin{proof}
두 경우 모두 환 사상\  $\mathcal{O}_{X, x} \to \mathcal{O}_{X_{k'}, x'}$은\ 
Noether 국소환들의 충실 평탄 국소 준동형이다. 따라서\ 
$\mathcal{O}_{X_{k'}, x'}$이\  Gorenstein이면\  Dualizing Complexes, Lemma
\ref{dualizing-lemma-flat-under-gorenstein}에 의해\  $\mathcal{O}_{X, x}$도\ 
Gorenstein이다. 위로 올라가는 방향에도\  Dualizing Complexes, Lemma
\ref{dualizing-lemma-flat-under-gorenstein}을 사용한다. 그러므로
$$
\mathcal{O}_{X_{k'}, x'}/\mathfrak m_x \mathcal{O}_{X_{k'}, x'} =
\kappa(x) \otimes_k k'
$$
이\  Gorenstein임을 보이면 된다. 첫 번째 경우에는\  $k \to k'$가 유한 생성이고,
두 번째 경우에는\  $k \to \kappa(x)$가 유한 생성임에 유의하자. 따라서 이는\ 
Gorenstein에 대해 성질 (A)가 성립한다는 사실에서 따른다. Dualizing Complexes,
Lemma \ref{dualizing-lemma-formal-fibres-gorenstein}을 보라.
\end{proof}
\noindent
위의 보조정리는 다음이\  Gorenstein 사상의 올바른 정의임을 보장한다.

\begin{definition}
\label{duality-definition-gorenstein-morphism}
$f : X \to Y$를 스킴의 사상이라 하자. 모든 올\  $X_y$가 국소\  Noether 스킴이라고
가정하자.
\begin{enumerate}
\item $x \in X$이고\  $y = f(x)$라 하자. $f$가\  $x$에서 평탄이고 스킴\  $X_y$의\ 
$x$에서의 국소환이\  Gorenstein이면, $f$가 {\it $x$에서\  Gorenstein}이라고 한다.
\item $f$가\  $X$의 모든 점에서\  Gorenstein이면\  $f$를 {\it Gorenstein 사상}이라
한다.
\end{enumerate}
\end{definition}

\noindent
이를 바꾸어 말하면 다음과 같다.

\begin{lemma}
\label{duality-lemma-gorenstein-morphism}
$f : X \to Y$를 스킴의 사상이라 하자. $f$의 모든 올이 국소\  Noether라고
가정하자. 다음은 동치이다.
\begin{enumerate}
\item $f$는\  Gorenstein이다.
\item $f$는 평탄이고 그 올들은\  Gorenstein 스킴이다.
\end{enumerate}
\end{lemma}

\begin{proof}
이는 정의들에서 직접 따른다.
\end{proof}
\begin{lemma}
\label{duality-lemma-gorenstein-CM-morphism}
Gorenstein 사상은\  Cohen-Macaulay이다.
\end{lemma}

\begin{proof}
Lemma \ref{duality-lemma-gorenstein-CM}와 정의들에서 따른다.
\end{proof}

\begin{lemma}
\label{duality-lemma-lci-gorenstein}
Syntomic 사상은\  Gorenstein이다. 동치로, 평탄 국소 완전교차 사상은\ 
Gorenstein이다.
\end{lemma}

\begin{proof}
Syntomic 사상은 평탄이고 그 올들은 체 위의 국소 완전교차임을 상기하자.
Morphisms, Lemma \ref{morphisms-lemma-syntomic-flat-fibres}를 보라. 체 위의 국소
완전교차는\  Lemma \ref{duality-lemma-gorenstein-lci}에 의해\  Gorenstein 스킴이므로 결론을
얻는다. 성질 ``syntomic''과 ``평탄 국소 완전교차 사상''은\  More on Morphisms,
Lemma \ref{more-morphisms-lemma-flat-lci}에 의해 동치이다.
\end{proof}
\begin{lemma}
\label{duality-lemma-composition-gorenstein}
$f : X \to Y$와\  $g : Y \to Z$를 사상이라 하자. $f$, $g$, $g \circ f$의 올\ 
$X_y$, $Y_z$, $X_z$가 국소\  Noether라고 가정하자.
\begin{enumerate}
\item $f$가\  $x$에서\  Gorenstein이고\  $g$가\  $f(x)$에서\  Gorenstein이면,
$g \circ f$는\  $x$에서\  Gorenstein이다.
\item $f$와\  $g$가\  Gorenstein이면\  $g \circ f$도\  Gorenstein이다.
\item $g \circ f$가\  $x$에서\  Gorenstein이고\  $f$가\  $x$에서 평탄이면, $f$는\ 
$x$에서\  Gorenstein이고\  $g$는\  $f(x)$에서\  Gorenstein이다.
\item $g \circ f$가\  Gorenstein이고\  $f$가 평탄이면, $f$는\  Gorenstein이고\ 
$g$는\  $f$의 상에 속하는 모든 점에서\  Gorenstein이다.
\end{enumerate}
\end{lemma}

\begin{proof}
이를 대수로 번역하면\  Dualizing Complexes, Lemma
\ref{dualizing-lemma-flat-under-gorenstein}에서 따른다.
\end{proof}
\begin{lemma}
\label{duality-lemma-flat-morphism-from-gorenstein-scheme}
\begin{slogan}
평탄 올뭉치의 전체공간이\  Gorenstein이면 기저와 올들도 그러하다.
\end{slogan}
$f : X \to Y$를 국소\  Noether 스킴들의 평탄 사상이라 하자. $X$가\ 
Gorenstein이면\  $f$는\  Gorenstein이고, 모든\  $x \in X$에 대해\ 
$\mathcal{O}_{Y, f(x)}$는\  Gorenstein이다.
\end{lemma}

\begin{proof}
이를 대수로 번역하면\  Dualizing Complexes, Lemma
\ref{dualizing-lemma-flat-under-gorenstein}에서 따른다.
\end{proof}
\begin{lemma}
\label{duality-lemma-base-change-gorenstein}
$f : X \to Y$를 스킴의 사상이라 하자. 모든 올\  $X_y$가 국소\  Noether 스킴이라고
가정하자. $Y' \to Y$가 국소 유한형이라고 하자. $f' : X' \to Y'$를\  $f$의
기저변환이라 하자. $x' \in X'$를 상이\  $x \in X$인 점이라 하자.
\begin{enumerate}
\item $f$가\  $x$에서\  Gorenstein이면\  $f' : X' \to Y'$는\  $x'$에서\ 
Gorenstein이다.
\item $f$가\  $x$에서 평탄이고\  $f'$가\  $x'$에서\  Gorenstein이면, $f$는\  $x$에서\ 
Gorenstein이다.
\item $Y' \to Y$가\  $f'(x')$에서 평탄이고\  $f'$가\  $x'$에서\  Gorenstein이면,
$f$는\  $x$에서\  Gorenstein이다.
\end{enumerate}
\end{lemma}

\begin{proof}
$Y' \to Y$에 관한 가정으로부터\  $y' \in Y'$가\  $y \in Y$로 갈 때 체 확대\ 
$\kappa(y')/\kappa(y)$가 유한 생성임에 유의하자. 따라서 모든 올\ 
$X'_{y'} = (X_y)_{\kappa(y')}$도 국소\  Noether이다. Varieties, Lemma
\ref{varieties-lemma-locally-Noetherian-base-change}를 보라. 그러므로 보조정리의
진술은 의미가 있다. $y' = f'(x')$와\  $y = f(x)$로 두자. 그러면 다음과 같은
국소환의 가환 도식을 얻는다.
$$
\xymatrix{
\mathcal{O}_{X', x'} & \mathcal{O}_{X, x} \ar[l] \\
\mathcal{O}_{Y', y'} \ar[u] & \mathcal{O}_{Y, y} \ar[l] \ar[u]
}
$$
여기서 왼쪽 위 모서리는 오른쪽 아래 모서리 위에서 오른쪽 위 모서리와 왼쪽 아래
모서리를 텐서곱한 것의 국소화이다.

\medskip\noindent
$f$가\  $x$에서\  Gorenstein이라고 가정하자.
$\mathcal{O}_{Y, y} \to \mathcal{O}_{X, x}$의 평탄성은\ 
$\mathcal{O}_{Y', y'} \to \mathcal{O}_{X', x'}$의 평탄성을 함의한다. Algebra,
Lemma \ref{algebra-lemma-base-change-flat-up-down}을 보라.
$\mathcal{O}_{X, x}/\mathfrak m_y\mathcal{O}_{X, x}$가\  Gorenstein이라는 사실은\ 
$\mathcal{O}_{X', x'}/\mathfrak m_{y'}\mathcal{O}_{X', x'}$가\  Gorenstein임을
함의한다. Lemma \ref{duality-lemma-gorenstein-base-change}를 보라. 따라서\  $f'$는\ 
$x'$에서\  Gorenstein이다.

\medskip\noindent
$f$가\  $x$에서 평탄이고\  $f'$가\  $x'$에서\  Gorenstein이라고 가정하자.
$\mathcal{O}_{X', x'}/\mathfrak m_{y'}\mathcal{O}_{X', x'}$가\  Gorenstein이라는
사실은\  $\mathcal{O}_{X, x}/\mathfrak m_y\mathcal{O}_{X, x}$가\  Gorenstein임을
함의한다. Lemma \ref{duality-lemma-gorenstein-base-change}를 보라. 따라서\  $f$는\ 
$x$에서\  Gorenstein이다.

\medskip\noindent
$Y' \to Y$가\  $y'$에서 평탄이고\  $f'$가\  $x'$에서\  Gorenstein이라고 가정하자.
$\mathcal{O}_{Y', y'} \to \mathcal{O}_{X', x'}$와\ 
$\mathcal{O}_{Y, y} \to \mathcal{O}_{Y', y'}$의 평탄성은\ 
$\mathcal{O}_{Y, y} \to \mathcal{O}_{X, x}$의 평탄성을 함의한다. Algebra,
Lemma \ref{algebra-lemma-base-change-flat-up-down}을 보라.
$\mathcal{O}_{X', x'}/\mathfrak m_{y'}\mathcal{O}_{X', x'}$가\  Gorenstein이라는
사실은\  $\mathcal{O}_{X, x}/\mathfrak m_y\mathcal{O}_{X, x}$가\  Gorenstein임을
함의한다. Lemma \ref{duality-lemma-gorenstein-base-change}를 보라. 따라서\  $f$는\ 
$x$에서\  Gorenstein이다.
\end{proof}
\begin{lemma}
\label{duality-lemma-flat-lft-base-change-gorenstein}
$f : X \to Y$를 평탄 국소 유한형인 스킴의 사상이라 하자. 그러면 집합\ 
$\{x \in X \mid f\text{ is Gorenstein at }x\}$의 형성은 임의의 기저변환과
가환한다.
\end{lemma}

\begin{proof}
가정에 의해\  $f$의 모든 올은 체 위에서 국소 유한형이고, 따라서 국소\  Noether이다.
임의의 기저변환에 대해서도 마찬가지이다. 그러므로 진술은 의미가 있다. 올을
살펴보면 다음 문제로 환원된다. $X$를 체\  $k$ 위의 국소 유한형 스킴이라 하고,
$K/k$를 체 확대라 하며, $x_K \in X_K$를 상이\  $x \in X$인 점이라 하자. 문제는\ 
$\mathcal{O}_{X_K, x_K}$가\  Gorenstein인 것과\  $\mathcal{O}_{X, x}$가\ 
Gorenstein인 것이 동치임을 보이는 것이다.

\medskip\noindent
이 문제는 다음과 같은 약간의 대수로 풀 수 있다. $x$를 포함하는 아핀 열린집합\ 
$\Spec(A) \subset X$를 택하자. $x$가\  $\mathfrak p \subset A$에 대응한다고 하자.
$A_K = A \otimes_k K$라 두면\  $\Spec(A_K) \subset X_K$가\  $x_K$를 포함한다.
$x_K$가\  $\mathfrak p_K \subset A_K$에 대응한다고 하자. $\omega_A^\bullet$을\ 
$A$의 쌍대화 복합체라 하자. Dualizing Complexes, Lemma
\ref{dualizing-lemma-base-change-dualizing-over-field}에 의해\ 
$\omega_A^\bullet \otimes_A A_K$는\  $A_K$의 쌍대화 복합체이다. 이제 결론을 얻는다.
실제로\  $A_\mathfrak p \to (A_K)_{\mathfrak p_K}$는\  Noether 환들의 평탄 국소
준동형이므로, $(\omega_A^\bullet)_\mathfrak p$가\  $D(A_\mathfrak p)$의 가역
대상인 것은\ 
$(\omega_A^\bullet)_\mathfrak p \otimes_{A_\mathfrak p} (A_K)_{\mathfrak p_K}$가\ 
$D((A_K)_{\mathfrak p_K})$의 가역 대상인 것과 동치이다. 몇 가지 세부 사항은
생략한다. 힌트: 코호몰로지 가군을 살펴보라.
\end{proof}
\begin{lemma}
\label{duality-lemma-affine-flat-Noetherian-gorenstein}
Situation \ref{duality-situation-shriek}에서\  $f : X \to Y$를\  $\textit{FTS}_S$의 사상이라
하자. $x \in X$라 하자. $f$가 평탄이면 다음은 동치이다.
\begin{enumerate}
\item $f$는\  $x$에서\  Gorenstein이다.
\item $f^!\mathcal{O}_Y$는\  $x$의 한 근방에서 가역 대상과 동형이다.
\end{enumerate}
특히\  $f$가\  Gorenstein인 점들의 집합은\  $X$에서 열린집합이다.
\end{lemma}

\begin{proof}
$\omega^\bullet = f^!\mathcal{O}_Y$로 두자. Lemma
\ref{duality-lemma-relative-dualizing-fibres}에 의해 이는 연접 코호몰로지 층을 갖는
유계 복합체이고, 올\  $X_y$로의 그 유도 제한\  $Lh^*\omega^\bullet$은\  $X_y$ 위의
쌍대화 복합체이다. 점의 포함 사상을\  $i : x \to X_y$라 쓰자. 그러면 다음은
동치이다.
\begin{enumerate}
\item $f$는\  $x$에서\  Gorenstein이다.
\item $\mathcal{O}_{X_y, x}$는\  Gorenstein이다.
\item $Lh^*\omega^\bullet$은\  $x$의 한 근방에서 가역이다.
\item $Li^* Lh^* \omega^\bullet$은\  $\kappa(x)$ 위에서 차원이\  $1$인 영이 아닌
코호몰로지를 정확히 하나 갖는다.
\item $L(h \circ i)^* \omega^\bullet$은\  $\kappa(x)$ 위에서 차원이\  $1$인 영이
아닌 코호몰로지를 정확히 하나 갖는다.
\item $\omega^\bullet$은\  $x$의 한 근방에서 가역이다.
\end{enumerate}
(1)과 (2)의 동치는 정의에 의한다($f$가 평탄이므로). (2)와 (3)의 동치는\ 
Lemma \ref{duality-lemma-gorenstein}에서 따른다. (3)과 (4)의 동치는\  More on Algebra,
Lemma \ref{more-algebra-lemma-lift-bounded-pseudo-coherent-to-perfect}에서 따른다.
(4)와 (5)의 동치는\  $Li^* Lh^* = L(h \circ i)^*$이기 때문이다. (5)와 (6)의
동치는\  More on Algebra, Lemma
\ref{more-algebra-lemma-lift-bounded-pseudo-coherent-to-perfect}에서 따른다.
따라서 보조정리는 분명하다.
\end{proof}
\begin{lemma}
\label{duality-lemma-flat-finite-presentation-characterize-gorenstein}
$f : X \to S$를 평탄 국소 유한 표시인 스킴의 사상이라 하자. $x \in X$이고 그
상이\  $s \in S$라고 하자. $d = \dim_x(X_s)$로 두자. 다음은 동치이다.
\begin{enumerate}
\item $f$는\  $x$에서\  Gorenstein이다.
\item $x$의 열린 근방\  $U \subset X$와\  $S$ 위의 국소 준유한 사상\ 
$U \to \mathbf{A}^d_S$가 존재하여, 그 사상은\  $x$에서\  Gorenstein이다.
\item $x$의 열린 근방\  $U \subset X$와\  $S$ 위의 국소 준유한\  Gorenstein 사상\ 
$U \to \mathbf{A}^d_S$가 존재한다.
\item $x$의 열린 근방\  $U \subset X$의 임의의\  $S$-사상\ 
$g : U \to \mathbf{A}^d_S$에 대해 다음이 성립한다.
$g$가\  $x$에서 준유한이다\  $\Rightarrow$ $g$는\  $x$에서\  Gorenstein이다.
\end{enumerate}
특히\  $f$가\  Gorenstein인 점들의 집합은\  $X$에서 열린집합이다.
\end{lemma}

\begin{proof}
$x \in U$인 아핀 열린집합\  $U = \Spec(A) \subset X$와\  $f(U) \subset V$인\ 
$V = \Spec(R) \subset S$를 택하자. 그러면\  $R \to A$는 유한 표시 평탄 환
사상이다. $\mathfrak p \subset A$를\  $x$에 대응하는 소 아이디얼이라 하자.
$A$를 주 국소화로 바꾸어 준유한 사상\  $R[x_1, \ldots, x_d]  \to A$가 존재한다고
가정할 수 있다. Algebra, Lemma
\ref{algebra-lemma-quasi-finite-over-polynomial-algebra}를 보라. 따라서\  $x$의 열린
근방\  $U \subset X$와 국소\footnote{$S$가 준분리이면\  $g$는 준유한이다.} 준유한
사상\  $g : U \to \mathbf{A}^d_S$로 이루어진 쌍\  $(U, g)$가 적어도 하나 존재한다.

\medskip\noindent
이를 전제로 보조정리는 다음 대수 문제로 번역된다(번역은 생략한다). $R \to A$가
평탄 유한 표시이고, $\mathfrak p \subset A$가 소 아이디얼이며,
$\varphi : R[x_1, \ldots, x_d] \to A$가\  $\mathfrak p$에서 준유한이라고 하자.
다음은 동치이다.
\begin{enumerate}
\item[(a)] $\Spec(\varphi)$는\  $\mathfrak p$에서\  Gorenstein이다.
\item[(b)] $\Spec(A) \to \Spec(R)$는\  $\mathfrak p$에서\  Gorenstein이다.
\item[(c)] $\Spec(A) \to \Spec(R)$는\  $\mathfrak p$의 한 열린 근방에서\ 
Gorenstein이다.
\end{enumerate}
각 경우\  $R[x_1, \ldots, x_n] \to A$는\  $\mathfrak p$에서 평탄하므로 평탄성의
열림성(Algebra, Theorem \ref{algebra-theorem-openness-flatness})에 의해\ 
$R[x_1, \ldots, x_n] \to A$가 평탄이라고 가정할 수 있다($A$를 적당한 주
국소화로 바꾼다). Algebra, Lemma
\ref{algebra-lemma-flat-finite-presentation-limit-flat}에 의해\ 
$R_0 \subset R$과\  $R_0[x_1, \ldots, x_n] \to A_0$가 존재하여, $R_0$는\ 
$\mathbf{Z}$ 위의 유한형이고\  $R_0 \to A_0$는 유한형이며\ 
$R_0[x_1, \ldots, x_n] \to A_0$는 평탄이다. 평탄 유한형 사상이\  Gorenstein인
점들의 집합은\  Lemma \ref{duality-lemma-base-change-gorenstein}에 의해 기저변환과
가환함에 유의하자. 이로써\  $R$이\  Noether인 경우로 환원된다.

\medskip\noindent
따라서\  $X$와\  $S$가 아핀이고\  $f$가 다음과 같이 분해된다고 가정할 수 있다.
$$
X \xrightarrow{g} \mathbf{A}^n_S \xrightarrow{p} S
$$
여기서\  $g$는 평탄 준유한이고\  $S$는\  Noether이다. 그러면\  $X$와\ 
$\mathbf{A}^n_S$는\  $S$ 위에서 분리이고,
$$
f^!\mathcal{O}_S = g^!p^!\mathcal{O}_S = g^!\mathcal{O}_{\mathbf{A}^n_S}[n]
$$
이다. 이는 상단 느낌표 함자의 알려진 성질들(Lemmas
\ref{duality-lemma-upper-shriek-composition}와\  \ref{duality-lemma-shriek-affine-line})에 의한다.
그러므로\  Lemma \ref{duality-lemma-affine-flat-Noetherian-gorenstein}에 의해 (a), (b),
(c)는 동치이다.
\end{proof}
\begin{lemma}
\label{duality-lemma-gorenstein-local-source-and-target}
성질\  $\mathcal{P}(f)=$``$f$의 올들이 국소\  Noether이고\  $f$가\  Gorenstein이다''는
공역 위의\  fppf 위상에 관해 국소적이고 정의역 위의\  syntomic 위상에 관해
국소적이다.
\end{lemma}

\begin{proof}
다음과 같이 두자.
$\mathcal{P}(f) =
\mathcal{P}_1(f) \wedge \mathcal{P}_2(f)$
여기서\  $\mathcal{P}_1(f)=$``$f$ is flat''이고,
$\mathcal{P}_2(f)=$``the fibres of $f$ are locally Noetherian
and Gorenstein''이다. $\mathcal{P}_1$이 정의역과 공역 위의\  fppf 위상에 관해
국소적임은 알려져 있다. Descent, Lemmas
\ref{descent-lemma-descending-property-flat}와\ 
\ref{descent-lemma-flat-fpqc-local-source}를 보라. 따라서\  $\mathcal{P}_2$를
다루면 된다.

\medskip\noindent
$f : X \to Y$를 스킴의 사상이라 하자. $\{\varphi_i : Y_i \to Y\}_{i \in I}$를\ 
$Y$의\  fppf 피복이라 하자. $f_i : X_i \to Y_i$를\  $\varphi_i$에 의한\  $f$의
기저변환이라 쓰자. $i \in I$이고\  $y_i \in Y_i$인 점을 택하자.
$y = \varphi_i(y_i)$로 두자. 다음에 유의하자.
$$
X_{i, y_i} = \Spec(\kappa(y_i)) \times_{\Spec(\kappa(y))} X_y.
$$
또한\  $\kappa(y_i)/\kappa(y)$는 유한 생성 체 확대이다. 따라서\  $X_y$가 국소\ 
Noether이면\  $X_{i, y_i}$도 국소\  Noether이다. Varieties, Lemma
\ref{varieties-lemma-locally-Noetherian-base-change}를 보라. 그리고\  $X_y$가
추가로\  Gorenstein이면\  Lemma \ref{duality-lemma-gorenstein-base-change}에 의해\ 
$X_{i, y_i}$도\  Gorenstein이다. 따라서\  $\mathcal{P}_2$는 공역 위에서\  fppf
국소적이다.

\medskip\noindent
$\{X_i \to X\}$를\  $X$의\  syntomic 피복이라 하자. $y \in Y$라 하자. 이 경우\ 
$\{X_{i, y} \to X_y\}$는 올의\  syntomic 피복이다. 따라서\  $\mathcal{P}_2$가
정의역 위의\  syntomic 위상에 관해 국소라는 것은\  Lemma
\ref{duality-lemma-gorenstein-local-syntomic}에서 따른다.
\end{proof}



\section{쌍대화 복합체에 관한 보충}
\label{duality-section-more-dualizing}

\noindent
다른 곳에 잘 들어맞지 않는 몇 가지 보조정리를 모은다.

\begin{lemma}
\label{duality-lemma-descent-ascent}
$f : X \to Y$를 국소\  Noether 스킴의 사상이라 하자. 다음 가운데 하나를
가정하자.
\begin{enumerate}
\item $f$는\  syntomic 전사이다.
\item $f$는 평탄 국소 완전교차 전사 사상이다.
\item $f$는 유한형\  Gorenstein 전사 사상이다.
\end{enumerate}
그러면\  $K \in D_\QCoh(\mathcal{O}_Y)$가\  $Y$ 위의 쌍대화 복합체인 것은\ 
$Lf^*K$가\  $X$ 위의 쌍대화 복합체인 것과 동치이다.
\end{lemma}

\begin{proof}
아핀 열린집합을 취하고\  Derived Categories of Schemes, Lemma
\ref{perfect-lemma-affine-compare-bounded}를 사용하면 이는\  Dualizing Complexes,
Lemma \ref{dualizing-lemma-descent-ascent}로 번역된다.
\end{proof}






\section{체 위의 고유 스킴에 대한 쌍대성}
\label{duality-section-duality-proper-over-field}

\noindent
이 절에서는 위의 쌍대화 복합체와 쌍대성에 관한 매우 일반적인 내용을 체 위의
고유 스킴에 적용하여 그 귀결들을 전개한다.

\begin{lemma}
\label{duality-lemma-duality-proper-over-field}
$X$를 체\  $k$ 위의 고유 스킴이라 하자. 다음 성질들을 갖는 쌍대화 복합체\ 
$\omega_X^\bullet$이 존재한다.
\begin{enumerate}
\item $H^i(\omega_X^\bullet)$은\  $i \in [-\dim(X), 0]$일 때만 영이 아닐 수 있다.
\item $\omega_X = H^{-\dim(X)}(\omega_X^\bullet)$는 지지집합이 차원\ 
$\dim(X)$의 기약 성분들인 연접\  $(S_2)$-가군이다.
\item $H^i(\omega_X^\bullet)$의 지지집합의 차원은\  $-i$ 이하이다.
\item 닫힌점\  $x \in X$에 대해 가군\ 
$H^i(\omega_{X, x}^\bullet) \oplus \ldots \oplus H^0(\omega_{X, x}^\bullet)$이
영이 아닌 것은\  $\text{depth}(\mathcal{O}_{X, x}) \leq -i$인 것과 동치이다.
\item $K \in D_\QCoh(\mathcal{O}_X)$에 대해 이동 및 완전 삼각형과 양립하는
함자적 동형\footnote{이 성질은\  Yoneda 보조정리에 의해\ 
$D_\QCoh(\mathcal{O}_X)$에서\  $\omega_X^\bullet$을 유일한 동형까지 특징짓는다.
$\omega_X^\bullet$은\  $D^b_{\textit{Coh}}(\mathcal{O}_X)$에 속하므로 실제로는\ 
$K \in D^b_{\textit{Coh}}(\mathcal{O}_X)$만 생각해도 충분하다.}
$$
\Ext^i_X(K, \omega_X^\bullet) = \Hom_k(H^{-i}(X, K), k)
$$
이 존재한다.
\item $X$ 위의 준연접\  $\mathcal{F}$에 대해 함자적 동형\ 
$\Hom(\mathcal{F}, \omega_X) = \Hom_k(H^{\dim(X)}(X, \mathcal{F}), k)$이
존재한다.
\item $X \to \Spec(k)$가 상대 차원\  $d$의 매끄러운 사상이면\ 
$\omega_X^\bullet \cong \wedge^d\Omega_{X/k}[d]$이고\ 
$\omega_X \cong \wedge^d\Omega_{X/k}$이다.
\end{enumerate}
\end{lemma}
\begin{proof}
$f : X \to \Spec(k)$를 구조 사상이라 쓰자. $a$를 이 사상의 추이 함자의 오른쪽
수반이라 하자. Lemma \ref{duality-lemma-twisted-inverse-image}를 보라. 상대 쌍대화 복합체
$$
\omega_X^\bullet = a(\mathcal{O}_{\Spec(k)})
$$
를 생각하자. Remark \ref{duality-remark-relative-dualizing-complex}와 비교하라. $f$가
고유이므로 정의에 의해\ 
$f^!(\mathcal{O}_{\Spec(k)}) = a(\mathcal{O}_{\Spec(k)})$이다. Section
\ref{duality-section-upper-shriek}를 보라. Lemma \ref{duality-lemma-shriek-dualizing}를 적용하면\ 
$\omega_X^\bullet$이 쌍대화 복합체임을 알 수 있다. 또한\  $\omega_X^\bullet$과\ 
$\omega_X$가\  Example \ref{duality-example-proper-over-local} 및\  Example
\ref{duality-example-equidimensional-over-field}에서와 같음을 알 수 있다.

\medskip\noindent
(1), (2), (3)은\  Lemma \ref{duality-lemma-vanishing-good-dualizing}에서 따른다.

\medskip\noindent
닫힌점\  $x \in X$에 대해\  $\omega_{X, x}^\bullet$은\ 
$\mathcal{O}_{X, x}$ 위의 정규화된 쌍대화 복합체이다. Lemma
\ref{duality-lemma-good-dualizing-normalized}를 보라. 그러면 (4)는\  Dualizing Complexes,
Lemma \ref{dualizing-lemma-depth-in-terms-dualizing-complex}에서 따른다.

\medskip\noindent
(5)는 구성에 의해 성립한다. 실제로\  $a$는\ 
$Rf_* : D_\QCoh(\mathcal{O}_X) \to D(\mathcal{O}_{\Spec(k)}) = D(k)$의 오른쪽
수반이고, 이 함자를\  $K \mapsto R\Gamma(X, K)$와 동일시할 수 있다. 또한\ 
$k$-가군의 유도 범주\  $D(k)$는 등급화\  $k$-벡터공간의 범주와 같다는 사실을
사용한다.

\medskip\noindent
(6)은 연접\  $\mathcal{F}$에 대해서는\  Lemma
\ref{duality-lemma-dualizing-module-proper-over-A}에서 따르고, 일반적으로는 (5)에서\ 
$K = \mathcal{F}[0]$과\  $i = -\dim(X)$로 놓아 풀어 쓰면 따른다.

\medskip\noindent
(7)은\  Lemma \ref{duality-lemma-smooth-proper}에서 따른다.
\end{proof}
\begin{remark}
\label{duality-remark-duality-proper-over-field}
$k$, $X$, $\omega_X^\bullet$을\  Lemma \ref{duality-lemma-duality-proper-over-field}에서와
같이 두자. 복합체\  $\omega_X^\bullet$의 항등 사상은 (5)의 함자적 동형을 통해
사상
$$
t : H^0(X, \omega_X^\bullet) \longrightarrow k
$$
에 대응한다. $D_\QCoh(\mathcal{O}_X)$의 임의의\  $K$에 대해 (5)의\ 
$\Hom(K, \omega_X^\bullet)$와\  $H^0(X, K)^\vee$의 동일시는 쌍
$$
\Hom_X(K, \omega_X^\bullet) \times H^0(X, K) \longrightarrow k,\quad
(\alpha, \beta) \longmapsto t(\alpha(\beta))
$$
에 대응한다. 이는 (5)의 동형들이 함자적이라는 사실에서 따른다. 마찬가지로 임의의\ 
$i \in \mathbf{Z}$에 대해 쌍
$$
\Ext^i_X(K, \omega_X^\bullet) \times H^{-i}(X, K) \longrightarrow k,\quad
(\alpha, \beta) \longmapsto t(\alpha(\beta))
$$
을 얻는다. 여기서는\  $\alpha$를 사상\  $K[-i] \to \omega_X^\bullet$으로,
$\beta$를\  $H^0(X, K[-i])$의 원소로 보아\  $\alpha(\beta)$를 정의한다. $K$가
일반적이면 이 쌍이 한쪽에서만 비퇴화임을 안다는 점에 유의하자. 즉 이 쌍은\ 
$\Hom_X(K, \omega_X^\bullet)$, 각각\  $\Ext^i_X(K, \omega_X^\bullet)$를\ 
$H^0(X, K)$, 각각\  $H^{-i}(X, K)$의\  $k$-선형 쌍대와 동형으로 만들지만 일반적으로
그 역은 성립하지 않는다. $K$가\  $D^b_{\textit{Coh}}(\mathcal{O}_X)$에 속하면\ 
$\Hom_X(K, \omega_X^\bullet)$, $\Ext_X(K, \omega_X^\bullet)$,
$H^0(X, K)$, $H^i(X, K)$는 유한 차원\  $k$-벡터공간이다(Derived Categories of
Schemes, Lemmas \ref{perfect-lemma-coherent-internal-hom}과\ 
\ref{perfect-lemma-direct-image-coherent-bdd-below}). 이 경우 쌍들은 통상적 의미에서
완전하다.
\end{remark}
\begin{remark}
\label{duality-remark-coherent-duality-proper-over-field}
Remark \ref{duality-remark-duality-proper-over-field}의 논의를 계속하며 같은 기호\  $k$, $X$,
$\omega_X^\bullet$, $t$를 사용한다. $\mathcal{F}$가 연접\  $\mathcal{O}_X$-가군이면
유한 차원\  $k$-벡터공간의 완전쌍
$$
\langle -, - \rangle :
\Ext^i_X(\mathcal{F}, \omega_X^\bullet) \times H^{-i}(X,\mathcal{F})
\longrightarrow k,\quad
(\alpha, \beta) \longmapsto t(\alpha(\beta))
$$
을 얻는다. 이 쌍들은 다음의 분명한 함자성을 만족한다. 연접\ 
$\mathcal{O}_X$-가군들의 준동형\  $\varphi : \mathcal{F} \to \mathcal{G}$가 주어지면
$$
\langle \alpha \circ \varphi, \beta \rangle =
\langle \alpha, \varphi(\beta) \rangle
$$
이다. 여기서\  $\alpha \in \Ext^i_X(\mathcal{G}, \omega_X^\bullet)$이고\ 
$\beta \in H^{-i}(X, \mathcal{F})$이다. 달리 말하면\  $\varphi$가 유도하는\ 
$k$-선형 사상\  $\Ext^i_X(\mathcal{G}, \omega_X^\bullet) \to
\Ext^i_X(\mathcal{F}, \omega_X^\bullet)$은 이 쌍들을 통해\  $\varphi$가 유도하는\ 
$k$-선형 사상\  $H^{-i}(X, \mathcal{F}) \to H^{-i}(X, \mathcal{G})$의\ 
$k$-선형 쌍대가 된다. 이 방식으로 서술하면\  $\varphi$가 준연접\ 
$\mathcal{O}_X$-가군들의 준동형일 때에도 성립한다.
\end{remark}
\begin{lemma}
\label{duality-lemma-duality-proper-over-field-perfect}
$k$, $X$, $\omega_X^\bullet$을\  Lemma \ref{duality-lemma-duality-proper-over-field}에서와
같이 두자. $t : H^0(X, \omega_X^\bullet) \to k$를\  Remark
\ref{duality-remark-duality-proper-over-field}에서와 같이 두자.
$E \in D(\mathcal{O}_X)$가 완전이라고 하자. 그러면 모든\  $i$에 대해 쌍
$$
H^i(X, \omega_X^\bullet \otimes_{\mathcal{O}_X}^\mathbf{L} E^\vee)
\times
H^{-i}(X, E) \longrightarrow k, \quad
(\xi, \eta) \longmapsto
t((1_{\omega_X^\bullet} \otimes \epsilon)(\xi \cup \eta))
$$
은 완전하다. 여기서\  $\cup$는\  Cohomology, Section
\ref{cohomology-section-cup-product}의 컵곱이고,
$\epsilon : E^\vee \otimes_{\mathcal{O}_X}^\mathbf{L} E \to \mathcal{O}_X$는\ 
Cohomology, Example \ref{cohomology-example-dual-derived}에서와 같다.
\end{lemma}

\begin{proof}
$E$를\  $E[-i]$로 바꾸면\  $i = 0$인 경우로 환원된다. Cohomology, Lemma
\ref{cohomology-lemma-ext-composition-is-cup}에 의해 이 쌍은\  Remark
\ref{duality-remark-duality-proper-over-field}에서 논의한 쌍과 같음을 알 수 있다.
따라서 그 주석의 논의로부터 결론이 따른다.
\end{proof}
\begin{lemma}
\label{duality-lemma-duality-proper-over-field-CM}
$X$를 체\  $k$ 위의 고유\  Cohen-Macaulay 스킴이라 하고, 차원\  $d$의
등차원이라고 하자. Lemma \ref{duality-lemma-duality-proper-over-field}의 가군\ 
$\omega_X$는 다음 성질들을 갖는다.
\begin{enumerate}
\item $\omega_X$는\  $X$ 위의 쌍대화 가군이다(Section
\ref{duality-section-dualizing-module}).
\item $\omega_X$는 지지집합이\  $X$인\  Cohen-Macaulay 연접 가군이다.
\item $K \in D_\QCoh(X)$에 대해 이동 및 완전 삼각형과 양립하는 함자적 동형\ 
$\Ext^i_X(K, \omega_X[d]) = \Hom_k(H^{-i}(X, K), k)$이 존재한다.
\item $X$ 위의 준연접\  $\mathcal{F}$에 대해 함자적 동형\ 
$\Ext^{d - i}(\mathcal{F}, \omega_X) = \Hom_k(H^i(X, \mathcal{F}), k)$이
존재한다.
\end{enumerate}
\end{lemma}

\begin{proof}
Lemma \ref{duality-lemma-duality-proper-over-field}에서\  $\omega_X$가 쌍대화 가군임은
분명하다(쌍대화 복합체의 가장 왼쪽에 있는 영이 아닌 코호몰로지 층이기 때문이다).
$\omega_X^\bullet = \omega_X[d]$이고\  $X$가\  Cohen-Macualay이므로\  Lemma
\ref{duality-lemma-dualizing-module-CM-scheme}에 의해\  $\omega_X$는\  Cohen-Macaulay이다.
나머지 명제들은 이 사실과\  Lemma \ref{duality-lemma-duality-proper-over-field}의 대응하는
명제들을 결합하면 따른다.
\end{proof}
\begin{remark}
\label{duality-remark-rework-duality-locally-free-CM}
$X$를 체\  $k$ 위의 고유\  Cohen-Macaulay 스킴이라 하고, 차원\  $d$의
등차원이라고 하자. $\omega_X^\bullet$과\  $\omega_X$를\  Lemma
\ref{duality-lemma-duality-proper-over-field}에서와 같이 두자. Lemma
\ref{duality-lemma-duality-proper-over-field-CM}에 의해\ 
$\omega_X^\bullet = \omega_X[d]$이다. $t : H^d(X, \omega_X) \to k$를\  Remark
\ref{duality-remark-duality-proper-over-field}의 사상이라 하자. $\mathcal{E}$를 쌍대\ 
$\mathcal{E}^\vee$를 갖는 유한 국소 자유\  $\mathcal{O}_X$-가군이라 하자. 그러면
완전쌍
$$
H^i(X, \omega_X \otimes_{\mathcal{O}_X} \mathcal{E}^\vee)
\times
H^{d - i}(X, \mathcal{E})
\longrightarrow
k,\quad
(\xi, \eta) \longmapsto t(1 \otimes \epsilon)(\xi \cup \eta))
$$
을 얻는다. 여기서\  $\cup$는 컵곱이고\ 
$\epsilon : \mathcal{E}^\vee \otimes_{\mathcal{O}_X} \mathcal{E}
\to \mathcal{O}_X$는 평가 사상이다. 이는\  Lemma
\ref{duality-lemma-duality-proper-over-field-perfect}의 특수한 경우이다.
\end{remark}
\noindent
쌍대화 복합체에 대한 건전성 점검은 다음과 같다.

\begin{lemma}
\label{duality-lemma-sanity-check-duality}
$X$를 체\  $k$ 위의 고유 스킴이라 하자. $\omega_X^\bullet$과\  $\omega_X$를\  Lemma
\ref{duality-lemma-duality-proper-over-field}에서와 같이 두자.
\begin{enumerate}
\item 어떤 체\  $k'$에 대해\  $X \to \Spec(k)$가\ 
$X \to \Spec(k') \to \Spec(k)$로 분해되면, $\omega_X^\bullet$과\  $\omega_X$는
사상\  $X \to \Spec(k')$에 대해\  Lemma
\ref{duality-lemma-duality-proper-over-field}에서와 같다.
\item $K/k$가 체 확대이면, $\omega_X^\bullet$과\  $\omega_X$를 기저변환\  $X_K$로
당긴 것은 사상\  $X_K \to \Spec(K)$에 대해\  Lemma
\ref{duality-lemma-duality-proper-over-field}에서와 같다.
\end{enumerate}
\end{lemma}

\begin{proof}
$f : X \to \Spec(k)$를 구조 사상이라 쓰고, $f' : X \to \Spec(k')$를 주어진
분해라 쓰자. Lemma \ref{duality-lemma-duality-proper-over-field}의 증명에서는\ 
$\omega_X^\bullet = a(\mathcal{O}_{\Spec(k)})$로 택했다. 여기서\  $a$는\  $f$에 대한\ 
Lemma \ref{duality-lemma-twisted-inverse-image}의 오른쪽 수반이다. 따라서\ 
$a(\mathcal{O}_{\Spec(k)}) \cong a'(\mathcal{O}_{\Spec(k)})$임을 보여야 한다.
여기서\  $a'$는\  $f'$에 대한\  Lemma \ref{duality-lemma-twisted-inverse-image}의 오른쪽
수반이다. $k' \subset H^0(X, \mathcal{O}_X)$이므로\  $k'/k$는 유한 확대이다
(Cohomology of Schemes, Lemma
\ref{coherent-lemma-proper-over-affine-cohomology-finite}). 수반의 유일성에 의해\ 
$a = a' \circ b$이다. 여기서\  $b$는\  $g : \Spec(k') \to \Spec(k)$에 대한\  Lemma
\ref{duality-lemma-twisted-inverse-image}의 오른쪽 수반이다. 달리 말하면\ 
$f^! = (f')^! \circ g^!$이다. 따라서\ 
$\Hom_k(k', k) \cong k'$가\  $k'$-가군으로서 성립함을 보이면 충분하다. Example
\ref{duality-example-affine-twisted-inverse-image}를 보라. 이는 양변이 같은 차원, 즉 차원\ 
$1$의\  $k'$-벡터공간이기 때문에 성립한다.

\medskip\noindent
(2)의 증명. 이는\  Lemma \ref{duality-lemma-more-base-change}에 의해\  $a$에 대한 기저변환이
성립하기 때문이다. Remark \ref{duality-remark-relative-dualizing-complex}의 논의를 보라.
\end{proof}







\section{상대 쌍대화 복합체}
\label{duality-section-relative-dualizing-complexes}

\noindent
유한 표시 평탄 고유 사상에 대해서는 강직한 상대 쌍대화 복합체가 있다. Remark
\ref{duality-remark-relative-dualizing-complex}와\  Lemma \ref{duality-lemma-van-den-bergh}를 보라.
Noether 스킴들의 유한형 분리 사상\  $f : X \to Y$에 대해서는\ 
$f^!\mathcal{O}_Y$를 생각할 수 있다. 이 절에서는 스킴들 사이의 평탄 국소 유한
표시 사상(반드시 준분리이거나 준콤팩트일 필요가 없고, 스킴들도 반드시 국소\ 
Noether일 필요가 없다)에 대한 상대 쌍대화 복합체를 정의한다. 그러한 복합체가
존재하고 유일한 동형까지 유일하며 위에서 언급한 경우들과 일치함을 보인다. 이 절을
읽기 전에\  Dualizing Complexes, Section
\ref{dualizing-section-relative-dualizing-complexes}를 읽기 바란다.

\begin{definition}
\label{duality-definition-relative-dualizing-complex}
$X \to S$를 평탄 국소 유한 표시인 스킴의 사상이라 하자. 대각 사상\ 
$\Delta_{X/S} : X \to X \times_S X$가 닫힌 몰입\  $\Delta : X \to W$를 통해
분해되는 임의의 열린집합\  $W \subset X \times_S X$를 택하자. {\it 상대 쌍대화
복합체}란 대상\  $K \in D(\mathcal{O}_X)$와 사상
$$
\xi : \Delta_*\mathcal{O}_X \longrightarrow L\text{pr}_1^*K|_W
$$
으로 이루어진 쌍\  $(K, \xi)$로서, 이 사상이\  $D(\mathcal{O}_W)$에 속하고 다음을
만족하는 것이다.
\begin{enumerate}
\item $K$는\  $S$-완전이다(Derived Categories of Schemes, Definition
\ref{perfect-definition-relatively-perfect}).
\item $\xi$는\  $\Delta_*\mathcal{O}_X$와\ 
$R\SheafHom_{\mathcal{O}_W}(
\Delta_*\mathcal{O}_X, L\text{pr}_1^*K|_W)$의 동형을 정한다.
\end{enumerate}
\end{definition}
\noindent
Lemma \ref{duality-lemma-sheaf-with-exact-support-ext}에 의해 조건 (2)는\ 
$D(\mathcal{O}_X)$의 동형
$$
\mathcal{O}_X \longrightarrow
R\SheafHom(\mathcal{O}_X, L\text{pr}_1^*K|_W)
$$
으로서\  $\Delta$에 의한 추이가\  $\xi$와 일치하는 동형의 존재와 동치이다.
$R\SheafHom(\mathcal{O}_X, L\text{pr}_1^*K|_W)$는 열린집합\  $W$의 선택과 무관하므로
쌍\  $(K, \xi)$들의 범주도 그러하다. $X \to S$가 분리이면\ 
$W = X \times_S X$로 택할 수 있다. 다음 보조정리를 사용하여 많은 논증을 환의
경우로 환원할 것이다.
\begin{lemma}
\label{duality-lemma-relative-dualizing-complex-algebra}
$X \to S$를 평탄 국소 유한 표시인 스킴의 사상이라 하고\  $(K, \xi)$를 상대
쌍대화 복합체라 하자. 수평 화살표들이 열린 몰입인 임의의 가환 도식
$$
\xymatrix{
\Spec(A) \ar[d] \ar[r] & X \ar[d] \\
\Spec(R) \ar[r] & S
}
$$
에 대해\  $K$를\  $\Spec(A)$에 제한한 것은\  Derived Categories of Schemes, Lemma
\ref{perfect-lemma-affine-compare-bounded}를 통해\  Dualizing Complexes, Definition
\ref{dualizing-definition-relative-dualizing-complex}의 의미에서\  $R \to A$의
상대 쌍대화 복합체에 대응한다.
\end{lemma}

\begin{proof}
$R\SheafHom$의 형성은 열린집합으로의 제한과 가환하므로\ 
$X = \Spec(A)$와\  $S = \Spec(R)$라고 가정해도 된다. $D(\mathcal{O}_X)$의 상대
완전 대상들은 유사연접이고, 따라서\  $D_\QCoh(\mathcal{O}_X)$에 속함에 유의하자
(Derived Categories of Schemes, Lemma \ref{perfect-lemma-pseudo-coherent}).
그러므로 진술은 의미가 있다. $\Delta_*$, $L\text{pr}_1^*$,
$R\SheafHom$을 취하는 것은\  Derived Categories of Schemes, Lemmas
\ref{perfect-lemma-quasi-coherence-pushforward},
\ref{perfect-lemma-quasi-coherence-pullback},
\ref{perfect-lemma-quasi-coherence-internal-hom}에 의해 대수 쪽에서 일어나는 일과
양립함에 유의하자. 마지막 보조정리를 위해서는\  $L\text{pr}_1^*K$가\  $S$-완전이고
(따라서 아래로 유계이며), $\Delta_*\mathcal{O}_X$가\ 
$D(\mathcal{O}_W)$의 유사연접 대상임을 사용한다. 이를 대수로 번역하면\  $A$가\ 
$A \otimes_R A$-가군으로서 유사연접이라는 뜻이고, 이는\  More on Algebra, Lemma
\ref{more-algebra-lemma-more-relative-pseudo-coherent-is-moot}를\ 
$R \to A \otimes_R A \to A$에 적용하면 따른다. 따라서\  Dualizing Complexes,
Definition \ref{dualizing-definition-relative-dualizing-complex}의 조건들을 정확히
회복한다.
\end{proof}
\begin{lemma}
\label{duality-lemma-relative-dualizing-RHom}
$X \to S$를 평탄 국소 유한 표시인 스킴의 사상이라 하고\  $(K, \xi)$를 상대
쌍대화 복합체라 하자. 그러면\ 
$\mathcal{O}_X \to R\SheafHom_{\mathcal{O}_X}(K, K)$는 동형이다.
\end{lemma}

\begin{proof}
아핀 국소적으로 보면\  Lemma \ref{duality-lemma-relative-dualizing-complex-algebra}에 의해
대수의 경우로 환원되고, 이는\  Dualizing Complexes, Lemma
\ref{dualizing-lemma-relative-dualizing-RHom}이다.
\end{proof}
\begin{lemma}
\label{duality-lemma-uniqueness-relative-dualizing}
$X \to S$를 평탄 국소 유한 표시인 스킴의 사상이라 하자. $(K, \xi)$와\ 
$(L, \eta)$가\  $X/S$ 위의 두 상대 쌍대화 복합체이면, $\xi$를\  $\eta$로 보내는
유일한 동형\  $K \to L$이 존재한다.
\end{lemma}

\begin{proof}
$U \subset X$를\  $S$의 한 아핀 열린집합으로 가는 아핀 열린집합이라 하자. Lemma
\ref{duality-lemma-relative-dualizing-complex-algebra}와\  Dualizing Complexes, Lemma
\ref{dualizing-lemma-uniqueness-relative-dualizing}에 의해 동형\ 
$K|_U \to L|_U$가 존재한다. 독자는 그 보조정리의 논증을 스킴의 경우에 다시
사용하여 이 경우의 증명을 얻을 수 있다. 여기서는 대신 붙이기 논증을 사용한다.

\medskip\noindent
동형\  $\alpha : K \to L$이 주어졌다고 하자. 그러면 어떤 가역 단면\ 
$u \in H^0(W, \Delta_*\mathcal{O}_X) = H^0(X, \mathcal{O}_X)$에 대해\ 
$\alpha(\xi) = u \eta$이다. 실제로 가정에 의해\  $\eta$와\  $\alpha(\xi)$는 모두
가역\  $\Delta_*\mathcal{O}_X$-가군의 생성원이다. 따라서\  $\alpha$를\ 
$u^{-1}\alpha$로 바꾸면\  $\alpha(\xi) = \eta$이다. Lemma
\ref{duality-lemma-relative-dualizing-RHom}에 의해\  $K$의 자기동형군은\ 
$H^0(X, \mathcal{O}_X^*)$이므로 이러한\  $\alpha$는 많아야 하나이다.

\medskip\noindent
$\mathcal{B}$를\  $S$의 한 아핀 열린집합으로 가는\  $X$의 아핀 열린집합들의 모임이라
하자. 각\  $U \in \mathcal{B}$에 대해 앞 두 문단의 논의로부터\  $\xi$를\  $\eta$로
보내는 유일한 동형\  $\alpha_U : K|_U \to L|_U$가 존재한다. Lemma
\ref{duality-lemma-relative-dualizing-RHom}에 의해\  $X$의 임의의 열린집합\  $U$와\  $i < 0$에
대해\  $\text{Ext}^i(K|_U, K|_U) = 0$임에 유의하자. Cohomology, Lemma
\ref{cohomology-lemma-vanishing-and-glueing}를\  $\text{id} : X \to X$에 적용하면,
모든\  $U \in \mathcal{B}$에 대해\  $\alpha_U$와 일치하는 유일한 사상\ 
$\alpha : K \to L$을 얻는다. 이는 국소적으로 참이므로\  $\alpha$는\  $\xi$를\ 
$\eta$로 보낸다.
\end{proof}
\begin{lemma}
\label{duality-lemma-existence-relative-dualizing}
$X \to S$를 평탄 국소 유한 표시인 스킴의 사상이라 하자. 상대 쌍대화 복합체\ 
$(K, \xi)$가 존재한다.
\end{lemma}

\begin{proof}
$\mathcal{B}$를\  $S$의 한 아핀 열린집합으로 가는\  $X$의 아핀 열린집합들의 모임이라
하자. 각\  $U$에 대해\  $S$ 위의\  $U$를 위한 상대 쌍대화 복합체\  $(K_U, \xi_U)$가
존재한다. 실제로\  $U \to X \to S$가\  $V$를 통과하도록 하는 아핀 열린집합\  $V \subset S$를
택하자. $U = \Spec(A)$와\  $V = \Spec(R)$로 쓰자. Dualizing Complexes, Lemma
\ref{dualizing-lemma-base-change-relative-dualizing}에 의해\  $R \to A$를 위한 상대
쌍대화 복합체\  $K_A \in D(A)$가 존재한다. Lemma
\ref{duality-lemma-relative-dualizing-complex-algebra}의 증명을 거꾸로 따라가면, 이는\ 
$V$-완전 대상\  $K_U \in D(\mathcal{O}_U)$와\  $D(\mathcal{O}_{U \times_V U})$의 사상
$$
\xi : \Delta_*\mathcal{O}_U \to L\text{pr}_1^*K_U
$$
을 정한다. $V$-완전인 것은\  $S$-완전인 것과 같고\ 
$U \times_V U = U \times_S U$이므로\  $(K_U, \xi_U)$는 원하는 바와 같다.
\medskip\noindent
$U' \subset U \subset X$이고\  $U', U \in \mathcal{B}$이면\  Lemma
\ref{duality-lemma-uniqueness-relative-dualizing}에 의해\  $D(\mathcal{O}_{U'})$에 유일한
동형\  $\rho_{U'}^U : K_U|_{U'} \to K_{U'}$가 존재하여\ 
$\xi_U|_{U' \times_S U'}$를\  $\xi_{U'}$로 보낸다. 여기서는 상대 쌍대화 복합체를
열린집합에 제한한 것이 자명하게 상대 쌍대화 복합체임을 사용했다. 유일성으로부터\ 
$\mathcal{B}$의\  $U'' \subset U' \subset U$에 대해\ 
$\rho^U_{U''} = \rho^V_{U''} \circ \rho ^U_{U'}|_{U''}$임이 보장된다. Lemma
\ref{duality-lemma-relative-dualizing-RHom}을\  $U/S$와\  $K_U$에 적용하면, 모든\ 
$U \in \mathcal{B}$와\  $i < 0$에 대해\  $\text{Ext}^i(K_U, K_U) = 0$임에 유의하자.
따라서\  BBD 붙이기 보조정리(Cohomology, Theorem
\ref{cohomology-theorem-glueing-bbd-general})에 의해 유일한 해가 존재한다. 즉 대상\ 
$K \in D(\mathcal{O}_X)$와 동형\  $\rho_U : K|_U \to K_U$가 존재하여 모든\ 
$U' \subset U$ 및\  $U, U' \in \mathcal{B}$에 대해\ 
$\rho^U_{U'} \circ \rho_U|_{U'} = \rho_{U'}$이다.
\medskip\noindent
증명을 마치려면\  $D(\mathcal{O}_W)$의 사상
$$
\xi : \Delta_*\mathcal{O}_X \longrightarrow L\text{pr}_1^*K|_W
$$
으로서\  $\Delta_*\mathcal{O}_X$에서\ 
$R\SheafHom_{\mathcal{O}_W}(\Delta_*\mathcal{O}_X, L\text{pr}_1^*K|_W)$로의
동형을 유도하는 사상을 구성해야 한다. $W$를 바꾸어도 되므로\ 
$W = \bigcup_{U \in \mathcal{B}} U \times_S U$로 택한다. $\rho_U$를 사용하면
동형
$$
R\SheafHom_{\mathcal{O}_W}(
\Delta_*\mathcal{O}_X, L\text{pr}_1^*K|_W)|_{U \times_S U}
\xrightarrow{\rho_U}
R\SheafHom_{\mathcal{O}_{U \times_S U}}(
\Delta_*\mathcal{O}_U, L\text{pr}_1^*K_U)
$$
을 얻는다. $W$는 열린집합\  $U \times_S U$들로 덮이므로\ 
$R\SheafHom_{\mathcal{O}_W}(\Delta_*\mathcal{O}_X, L\text{pr}_1^*K|_W)$의
코호몰로지 층은 차수\  $0$을 제외하면 영이다. 또한 동형
$$
H^0\left(U \times_S U, R\SheafHom_{\mathcal{O}_W}(\Delta_*\mathcal{O}_X,
L\text{pr}_1^*K|_W)\right)
\xrightarrow{\rho_U}
H^0\left((R\SheafHom_{\mathcal{O}_{U \times_S U}}(
\Delta_*\mathcal{O}_U, L\text{pr}_1^*K_U)\right)
$$
을 얻는다. 좌변에서 이 사상에 의해\  $\xi_U$로 가는 원소를\  $\tau_U$라 하자.
$U' \subset U \subset X$이고\  $U', U \in \mathcal{B}$일 때\ 
$\rho^U_{U'}$, $\xi_U$, $\xi_{U'}$, $\rho_U$, $\rho_{U'}$ 사이의 양립성은\ 
$\tau_U|_{U' \times_S U'} = \tau_{U'}$를 함의한다. 따라서 제$0$ 코호몰로지 층\ 
$H^0(R\SheafHom_{\mathcal{O}_W}(\Delta_*\mathcal{O}_X, L\text{pr}_1^*K|_W))$의
대역 단면\  $\tau$를 얻는다.
$R\SheafHom_{\mathcal{O}_W}(\Delta_*\mathcal{O}_X, L\text{pr}_1^*K|_W)$의
그 밖의 코호몰로지 층들은 영이므로 이 대역 단면\ 
$\tau$가 원하는 사상\  $\xi$를 정한다. $\xi$를\  $U \times_S U$에 제한한 것은\ 
$\xi_U$를 주므로\  Definition \ref{duality-definition-relative-dualizing-complex}의 마지막
조건을 만족한다.
\end{proof}
\begin{lemma}
\label{duality-lemma-base-change-relative-dualizing}
스킴들의\  Cartesian 정사각형
$$
\xymatrix{
X' \ar[d]_{f'} \ar[r]_{g'} & X \ar[d]^f \\
S' \ar[r]^g & S
}
$$
을 생각하자. $X \to S$가 평탄 국소 유한 표시라고 가정하자. $(K, \xi)$를\ 
$f$의 상대 쌍대화 복합체라 하자. $K' = L(g')^*K$로 두자. $\xi'$를\  $\xi$의
유도 기저변환이라 하자(증명을 보라). 그러면\  $(K', \xi')$는\  $f'$의 상대 쌍대화
복합체이다.
\end{lemma}

\begin{proof}
Cartesian 정사각형
$$
\xymatrix{
X' \ar[d]_{\Delta_{X'/S'}} \ar[r] & X \ar[d]^{\Delta_{X/S}} \\
X' \times_{S'} X' \ar[r]^{g' \times g'} & X \times_S X
}
$$
을 생각하자. $\Delta_{X/S}$가 닫힌 몰입\  $\Delta : X \to W$를 통해 분해되는
열린집합\  $W \subset X \times_S X$를 택하자. $\Delta_{X'/S'}$가 닫힌 몰입\ 
$\Delta' : X \to W'$를 통해 분해되고\  $(g' \times g')(W') \subset W$인 열린집합\ 
$W' \subset X' \times_{S'} X'$를 택하자. 유도된 사상도 계속\ 
$g' \times g' : W' \to W$로 쓰자. 다음이 성립한다.
$$
L(g' \times g')^*\Delta_*\mathcal{O}_X =
\Delta'_*\mathcal{O}_{X'}
\quad\text{and}\quad
L(g' \times g')^*L\text{pr}_1^*K|_W =
L\text{pr}_1^*K'|_{W'}
$$
첫 번째 등식은\  $X$와\  $X' \times_{S'} X'$가\  $X \times_S X$ 위에서\  tor-독립이기
때문에 성립한다(예를 들어\  More on Morphisms, Lemma
\ref{more-morphisms-lemma-case-of-tor-independence}를 보라). 두 번째 등식은 유도
당김의 추이성(Cohomology, Lemma
\ref{cohomology-lemma-derived-pullback-composition})에서 따른다. 따라서\ 
$\xi' = L(g' \times g')^*\xi$를 사상
$$
\xi' : \Delta'_*\mathcal{O}_{X'} \longrightarrow L\text{pr}_1^*K'|_{W'}
$$
으로 볼 수 있다. 이제 보조정리의 증명은 간단하다. 먼저\  Derived Categories of
Schemes, Lemma \ref{perfect-lemma-base-change-relatively-perfect}에 의해\  $K'$은\ 
$S'$-완전이다. $\xi'$가\  $\Delta'_*\mathcal{O}_{X'}$에서\ 
$R\SheafHom_{\mathcal{O}_{W'}}(
\Delta'_*\mathcal{O}_{X'}, L\text{pr}_1^*K'|_{W'})$로의 동형을 유도함을 확인할
때에는 아핀 국소적으로 작업해도 된다. Lemma
\ref{duality-lemma-relative-dualizing-complex-algebra}에 의해 대수의 대응하는 명제로
환원되고, 이는\  Dualizing Complexes, Lemma
\ref{dualizing-lemma-base-change-relative-dualizing}에서 증명된다.
\end{proof}
\begin{lemma}
\label{duality-lemma-flat-proper-relative-dualizing}
$S$를 준콤팩트 준분리 스킴이라 하자. $f : X \to S$를 유한 표시 평탄 고유
사상이라 하자. Remark \ref{duality-remark-relative-dualizing-complex}의 상대 쌍대화
복합체\  $\omega_{X/S}^\bullet$과 (\ref{duality-equation-pre-rigid})은\  Definition
\ref{duality-definition-relative-dualizing-complex}의 의미에서 상대 쌍대화 복합체이다.
\end{lemma}

\begin{proof}
Lemma \ref{duality-lemma-properties-relative-dualizing}에서\  $\omega_{X/S}^\bullet$이\ 
$S$-완전임을 증명했다. $c$를 대각 사상\  $\Delta : X \to X \times_S X$에 대한\ 
Lemma \ref{duality-lemma-twisted-inverse-image}의 오른쪽 수반이라 하자. 그러면
(\ref{duality-equation-pre-rigid})에\  $\Delta_*$를 적용하여 동형
$$
\Delta_*\mathcal{O}_X \to
\Delta_*(c(L\text{pr}_1^*\omega_{X/S}^\bullet)) =
R\SheafHom_{\mathcal{O}_{X \times_S X}}(
\Delta_*\mathcal{O}_X, L\text{pr}_1^*\omega_{X/S}^\bullet)
$$
을 얻는다. 등식은\  Lemmas \ref{duality-lemma-twisted-inverse-image-closed}와\ 
\ref{duality-lemma-sheaf-with-exact-support-ext}에 의해 성립한다. 이로써 증명이 끝난다.
\end{proof}
\begin{remark}
\label{duality-remark-relative-dualizing-complex-bis}
$X \to S$를 유한 표시 평탄 고유인 스킴의 사상이라 하자. Lemma
\ref{duality-lemma-existence-relative-dualizing}에 의해\  Definition
\ref{duality-definition-relative-dualizing-complex}의 의미에서 상대 쌍대화 복합체\ 
$(\omega_{X/S}^\bullet, \xi)$가 존재한다. $S'$가 준콤팩트 준분리인 임의의 사상\ 
$g : S' \to S$를 생각하자(예를 들어\  $S$의 아핀 열린집합). Lemma
\ref{duality-lemma-base-change-relative-dualizing}에 의해\ 
$(L(g')^*\omega_{X/S}^\bullet, L(g')^*\xi)$는\  Definition
\ref{duality-definition-relative-dualizing-complex}의 의미에서 기저변환\ 
$f' : X' \to S'$의 상대 쌍대화 복합체이다. $\omega_{X'/S'}^\bullet$을\  Remark
\ref{duality-remark-relative-dualizing-complex}의 의미에서\  $X' \to S'$의 상대 쌍대화
복합체라 하자. Lemmas \ref{duality-lemma-flat-proper-relative-dualizing}와\ 
\ref{duality-lemma-uniqueness-relative-dualizing}를 결합하면
(\ref{duality-equation-pre-rigid}) 및\  $L(g')^*\xi$와 양립하는 유일한 동형
$$
\omega_{X'/S'}^\bullet \longrightarrow L(g')^*\omega_{X/S}^\bullet
$$
이 존재함을 알 수 있다. 이 동형들은\  $S$ 위의 준콤팩트 준분리 스킴들 사이의
사상들 및\  Lemma \ref{duality-lemma-proper-flat-base-change}의 기저변환 동형들과 양립한다
(이 양립성이 필요해지면 여기에서 주의 깊게 서술하고 증명할 것이다).
\end{remark}
\begin{lemma}
\label{duality-lemma-compactifyable-relative-dualizing}
Situation \ref{duality-situation-shriek}에서\  $f : X \to Y$를\  $\textit{FTS}_S$의 사상이라
하자. $f$가 평탄이면\  $f^!\mathcal{O}_Y$는\  Definition
\ref{duality-definition-relative-dualizing-complex}의 의미에서\  $Y$ 위의\  $X$를 위한 상대
쌍대화 복합체의 첫 번째 성분이다.
\end{lemma}

\begin{proof}
Lemma \ref{duality-lemma-flat-shriek-relatively-perfect}에 의해\  $f^!\mathcal{O}_Y$는\ 
$Y$-완전이다. $f$가 분리이므로 대각 사상\ 
$\Delta : X \to X \times_Y X$는 닫힌 몰입이고\ 
$\Delta_*\Delta^!(-) =
R\SheafHom_{\mathcal{O}_{X \times_Y X}}(\mathcal{O}_X, -)$이다. Lemmas
\ref{duality-lemma-twisted-inverse-image-closed}와\ 
\ref{duality-lemma-sheaf-with-exact-support-ext}를 보라.
따라서 증명을 마치려면\ 
$\Delta^!(L\text{pr}_1^*f^!(\mathcal{O}_Y)) \cong \mathcal{O}_X$임을 보이면 충분하다.
여기서\  $\text{pr}_1 : X \times_Y X \to X$는 첫 번째 사영이다. 다음이 성립한다.
$$
\mathcal{O}_X = \Delta^! \text{pr}_1^!\mathcal{O}_X =
\Delta^! \text{pr}_1^! L\text{pr}_2^*\mathcal{O}_Y =
\Delta^!(L\text{pr}_1^* f^!\mathcal{O}_Y)
$$
여기서\  $\text{pr}_2 : X \times_Y X \to X$는 두 번째 사영이고, Lemma
\ref{duality-lemma-base-change-shriek-flat}의 기저변환 동형\ 
$\text{pr}_1^! \circ L\text{pr}_2^* = L\text{pr}_1^* \circ f^!$을 사용했다.
\end{proof}
\begin{lemma}
\label{duality-lemma-relative-dualizing-composition}
$f : Y \to X$와\  $X \to S$를 평탄 유한 표시인 스킴의 사상들이라 하자.
$(K, \xi)$와\  $(M, \eta)$를 각각\  $X \to S$와\  $Y \to X$를 위한 상대 쌍대화
복합체라 하자. $E = M \otimes_{\mathcal{O}_Y}^\mathbf{L} Lf^*K$로 두자. 그러면
적절한\  $\zeta$에 대해\  $(E, \zeta)$는\  $Y \to S$를 위한 상대 쌍대화 복합체이다.
\end{lemma}

\begin{proof}
Lemma \ref{duality-lemma-relative-dualizing-complex-algebra}와 이 보조정리의 대수적 판본
(Dualizing Complexes, Lemma
\ref{dualizing-lemma-relative-dualizing-composition})을 사용하면\  $E$가 아핀
국소적으로 상대 쌍대화 복합체의 첫 번째 성분임을 알 수 있다. 특히\  $E$가\ 
$S$-완전임도 알 수 있다. 이는 아핀 국소적으로 확인할 수 있기 때문이다. Derived
Categories of Schemes, Lemma
\ref{perfect-lemma-affine-locally-rel-perfect}를 보라.

\medskip\noindent
먼저 사상들\  $X \to S$와\  $Y \to X$가 분리인 경우에\  $\zeta$의 존재를 증명하자.
이 경우\  $\Delta_{X/S}$, $\Delta_{Y/X}$, $\Delta_{Y/S}$는 닫힌 몰입이다. 다음
그림을 생각하자.
$$
\xymatrix{
& & Y \ar@{=}[r] & Y \ar[d]^f \\
Y \ar[r]_{\Delta_{Y/X}} &
Y \times_X Y \ar[d]_m \ar[r]_\delta \ar[ru]_q &
Y \times_S Y \ar[d]^{f \times f} \ar[ru]_p & X\\
& X \ar[r]^{\Delta_{X/S}} & X \times_S X \ar[ru]_r
}
$$
여기서\  $p$, $q$, $r$은 첫 번째 사영들이다. Lemma
\ref{duality-lemma-sheaf-with-exact-support-internal-home}에 의해
$$
R\SheafHom_{\mathcal{O}_{Y \times_S Y}}(
\Delta_{Y/S, *}\mathcal{O}_Y, Lp^*E) =
R\delta_*\left(R\SheafHom_{\mathcal{O}_{Y \times_X Y}}(
\Delta_{Y/X, *}\mathcal{O}_Y,
R\SheafHom(\mathcal{O}_{Y \times_X Y}, Lp^*E))\right)
$$
이다. Lemma \ref{duality-lemma-sheaf-with-exact-support-tensor}에 의해
$$
R\SheafHom(\mathcal{O}_{Y \times_X Y}, Lp^*E) =
R\SheafHom(\mathcal{O}_{Y \times_X Y}, L(f \times f)^*Lr^*K)
\otimes_{\mathcal{O}_{Y \times_S Y}}^\mathbf{L} Lq^*M
$$
이다. Lemma \ref{duality-lemma-flat-bc-sheaf-with-exact-support}에 의해
$$
R\SheafHom(\mathcal{O}_{Y \times_X Y}, L(f \times f)^*Lr^*K) =
Lm^*R\SheafHom(\mathcal{O}_X, Lr^*K)
$$
이다. 마지막 식은\  $\xi$를 통해\ 
$Lm^*\mathcal{O}_X = \mathcal{O}_{Y \times_X Y}$와 동형이다. 따라서 앞의 식은\ 
$Lq^*M$과 동형이고, 따라서
$$
R\SheafHom_{\mathcal{O}_{Y \times_S Y}}(
\Delta_{Y/S, *}\mathcal{O}_Y, Lp^*E) =
R\delta_*\left(R\SheafHom_{\mathcal{O}_{Y \times_X Y}}(
\Delta_{Y/X, *}\mathcal{O}_Y, Lq^*M)\right)
$$
이다. 괄호 안의 대상은\  $\eta$를 통해\ 
$\Delta_{Y/X, *}*\mathcal{O}_X$와 동형이다. 이로써 분리인 경우의 증명이 끝난다.

\medskip\noindent
일반적인 경우에는\  $\Delta_{X/S}$가 닫힌 몰입\ 
$\Delta : X \to W$를 통해 분해되는 열린집합\  $W \subset X \times_S X$를 택하고,
$\Delta_{Y/X}$가 닫힌 몰입\  $\Delta' : Y \to V$를 통해 분해되는 열린집합\ 
$V \subset Y \times_X Y$를 택한다. 마지막으로\  $Y \times_X Y$와의 교집합이\ 
$V$이고\  $W$로 사상하는 열린집합\  $W' \subset Y \times_S Y$를 택한다. 그러면
다음 그림을 생각한다.
$$
\xymatrix{
& & Y \ar@{=}[r] & Y \ar[d]^f \\
Y \ar[r]_{\Delta'} &
V \ar[d]_m \ar[r]_\delta \ar[ru]_q &
W' \ar[d]^{f \times f} \ar[ru]_p & X\\
& X \ar[r]^\Delta & W \ar[ru]_r
}
$$
그리고 앞과 정확히 같은 논증을 사용한다.
\end{proof}





\section{lci 사상의 기본류}
\label{duality-section-fundamental-class}

\noindent
이 절에서는\  Section \ref{duality-section-examples}에서 수행한 계산들을 사용한다. 따라서
우리의 결과는 그 절에서와 같은 종류의 비유일성을 지닌다.
\begin{lemma}
\label{duality-lemma-determinant}
$X$를 국소환 달린 공간이라 하자.
$$
\mathcal{E}_1 \xrightarrow{\alpha} \mathcal{E}_0 \to \mathcal{F} \to 0
$$
을\  $\mathcal{O}_X$-가군들의 짧은 완전열이라 하자. $\mathcal{E}_1$과\ 
$\mathcal{E}_0$가 각각 계수\  $r_1, r_0$인 국소 자유 가군이라고 가정하자. 그러면
표준 사상
$$
\wedge^{r_0 - r_1}\mathcal{F}
\longrightarrow
\wedge^{r_1}(\mathcal{E}_1^\vee) \otimes \wedge^{r_0}\mathcal{E}_0
$$
이 존재하며, 이 사상이\  $x \in X$에서의 줄기 위에서 동형인 것은\  $x$의 어떤 열린
근방에서\  $\mathcal{F}$가 계수\  $r_0 - r_1$인 국소 자유 가군인 것과 동치이다.
\end{lemma}

\begin{proof}
$r_1 > r_0$이면 관례에 따라\  $\wedge^{r_0 - r_1}\mathcal{F} = 0$이고 유일한 사상은
동형일 수 없다. 따라서\  $r = r_0 - r_1 \geq 0$이라 가정해도 된다. 다음 공식으로
사상을 정의한다.
$$
s_1 \wedge \ldots \wedge s_r \mapsto
t_1^\vee \wedge \ldots \wedge t_{r_1}^\vee \otimes
\alpha(t_1) \wedge \ldots \wedge \alpha(t_{r_1}) \wedge
\tilde s_1 \wedge \ldots \wedge \tilde s_r
$$
여기서\  $t_1, \ldots, t_{r_1}$은\  $\mathcal{E}_1$의 국소 기저이고,
$t_1^\vee, \ldots, t_{r_1}^\vee$는 이에 대응하는\  $\mathcal{E}_1^\vee$의 쌍대
기저이며, $s'_i$는\  $s_i$를\  $\mathcal{E}_0$의 단면으로 국소적으로 올린 것이다.
이 정의가 잘 정의됨의 증명은 생략한다.

\medskip\noindent
$\mathcal{F}$가 계수\  $r$인 국소 자유 가군이면 이 사상이 동형임은 곧바로 확인된다.
역으로 이 사상이\  $x$에서의 줄기 위에서 동형이라고 가정하자. 그러면\ 
$\wedge^r\mathcal{F}_x$는 가역이다. 따라서\  $\mathcal{F}_x$는 많아야\  $r$개의
원소로 생성된다. 이는\  $\alpha$를\  $\mathfrak m_x$로 나눈 계수가\  $r$일 때에만
가능하다. 즉\  $\alpha$는\  $\mathfrak m_x$로 나눈 뒤 최대 계수를 갖는다. 따라서\ 
$\alpha$는\  $x$의 한 근방에서 최대 계수를 가지며, 원하는 대로\  $\mathcal{F}$는 그
근방에서 계수\  $r$인 국소 자유 가군이다.
\end{proof}
\begin{lemma}
\label{duality-lemma-fundamental-class-lci}
$Y$를 뇌터 스킴이라 하자. $f : X \to Y$를, 몰입\  $X \to P$와 고유 매끄러운 사상\ 
$P \to Y$의 합성으로 분해되는 국소 완전 교차 사상이라 하자. $r$을\  $X$ 위의 국소
상수 함수로서\ 
$\omega_{X/Y} = H^{-r}(f^!\mathcal{O}_Y)$가\  $f^!\mathcal{O}_Y$의 유일한 영 아닌
코호몰로지 층이 되게 하는 함수라 하자. Lemma \ref{duality-lemma-lci-shriek}를 보라. 그러면
사상
$$
\wedge^r\Omega_{X/Y} \longrightarrow \omega_{X/Y}
$$
이 존재하고, 이 사상이 한 점\  $x$에서의 줄기 위에서 동형인 것은\  $f$가\  $x$에서
매끄러운 것과 동치이다.
\end{lemma}

\begin{proof}
가정에 의해\  $X$는\  $Y$ 위에서 콤팩트화 가능하므로\  $f^!$가 정의된다. Section
\ref{duality-section-upper-shriek}를 보라. $X \to P$가 닫힌 몰입\  $i : X \to W$를 통해
분해되는 열린 부분스킴\  $j : W \to P$를 택하자. 또한\  $g : P \to Y$가 주어진
사상이면\  $f^! = i^! \circ j^! \circ g^!$이다. Lemma
\ref{duality-lemma-smooth-proper}에 의해\ 
$g^!\mathcal{O}_Y = \wedge^d\Omega_{P/Y}[d]$이다. 여기서\  $d$는\  $Y$ 위의\  $P$의
상대차원을 주는 국소 상수 함수이다. 또한\  $j^! = j^*$이고,
$i^!\mathcal{O}_W = \wedge^c\mathcal{N}[-c]$이다. 여기서\  $c$는\  $W$ 안의\  $X$의
여차원($X$ 위의 국소 상수 함수)이고, $\mathcal{N}$은\  Koszul-정칙 몰입\  $i$의
법선층이다. Lemma \ref{duality-lemma-regular-immersion}을 보라. 이를 종합하면
$$
f^!\mathcal{O}_Y =
\left(\wedge^c\mathcal{N} \otimes_{\mathcal{O}_X}
\wedge^d\Omega_{P/Y}|_X\right)[d - c]
$$
를 얻는다. 여기서\  Lemma \ref{duality-lemma-perfect-comparison-shriek}도 사용했다. 따라서\ 
$X$ 위의 국소 상수 함수로서\  $r = d|_X - c$이다. $X \to P$의 여법선층은\ 
$\mathcal{I}/\mathcal{I}^2$이다. 여기서\  $\mathcal{I} \subset \mathcal{O}_W$는\ 
$i$의 아이디얼층이다. Morphisms, Section
\ref{morphisms-section-conormal-sheaf}를 보라. Morphisms, Lemma
\ref{morphisms-lemma-differentials-relative-immersion}의 표준 완전열
$$
\mathcal{I}/\mathcal{I}^2 \to
\Omega_{P/Y}|_X \to \Omega_{X/Y} \to 0
$$
을 생각하자. Lemma \ref{duality-lemma-determinant}를 적용하여 원하는 사상을 얻는다.

\medskip\noindent
$f$가\  $x$에서 매끄러우면\  Lemma \ref{duality-lemma-determinant}와\  $\Omega_{X/Y}$가\ 
$x$에서 계수\  $r$인 국소 자유 가군이라는 사실로부터 이 사상은 동형이다. 역으로
우리의 사상이\  $x$에서의 줄기 위에서 동형이라고 가정하자. 그러면 보조정리에 의해\ 
$X$를\  $x$의 한 열린 근방으로 바꾼 뒤\  $\Omega_{X/Y}$는 계수\  $r$인 자유 가군이다.
한편 다음과 같이 가정해도 된다.
$X = \Spec(A)$이고\  $Y = \Spec(R)$이며\ 
$A = R[x_1, \ldots, x_n]/(f_1, \ldots, f_m)$이고\ 
$f_1, \ldots, f_m$은\  Koszul 정칙열이다(이는 국소 완전 교차 사상의 정의에서
따른다). 분명히\  $r = n - m$이다. 따라서\  $x$의 잉여체에서 편미분 행렬\ 
$\partial f_j/\partial x_i$의 계수는\  $m$이다. 그러므로 변수의 순서를 바꾸어\ 
$(\partial f_j/\partial x_i)_{1 \leq i, j \leq m}$의 행렬식이\  $x$의 한 열린
근방에서 가역이라고 가정해도 된다. 따라서 이 점에서\  $R \to A$는 매끄럽다. 예를
들어\  Algebra, Example \ref{algebra-example-make-standard-smooth}를 보라.
\end{proof}
\begin{lemma}
\label{duality-lemma-fundamental-class-almost-lci}
$f : X \to Y$를 스킴의 사상이라 하고\  $r \geq 0$이라 하자. 다음을 가정하자.
\begin{enumerate}
\item $Y$는\  Cohen--Macaulay이다(Properties, Definition
\ref{properties-definition-Cohen-Macaulay}).
\item $f$는\  $X \to P \to Y$로 분해되며, 첫 번째 사상은 몰입이고 두 번째 사상은
매끄럽고 고유하다.
\item $x \in X$이고\  $\dim(\mathcal{O}_{X, x}) \leq 1$이면\  $f$는\  $x$에서\ 
Koszul이다(More on Morphisms, Definition
\ref{more-morphisms-definition-lci}).
\item $\xi$가\  $X$의 한 기약 성분의 일반점이면\ 
$\text{trdeg}_{\kappa(f(\xi))} \kappa(\xi) = r$이다.
\end{enumerate}
그러면\  $\omega_{X/Y} = H^{-r}(f^!\mathcal{O}_Y)$로 둘 때 사상
$$
\wedge^r\Omega_{X/Y} \longrightarrow \omega_{X/Y}
$$
이 존재하며, $f$가 매끄러운 자취에서 동형이다.
\end{lemma}

\begin{proof}
$U \subset X$를\  $f$가 국소 완전 교차 사상인 열린 부분스킴이라 하자. 가정 (4)에
의해 모든 일반점에서\  $f$의 상대차원이\  $r$이므로\  Lemma
\ref{duality-lemma-fundamental-class-lci}의 국소 상수 함수는 상수값\  $r$을 가진다. 따라서
사상
$$
\wedge^r\Omega_{X/Y}|_U = \wedge^r \Omega_{U/Y}
\longrightarrow
\omega_{U/Y} = \omega_{X/Y}|_U
$$
을 얻으며, 이는\  $f$의 매끄러운 점들에서 동형이다(매끄러운 사상은 국소 완전 교차
사상이므로 이 자취는\  $U$에 포함된다). Lemma \ref{duality-lemma-shriek-over-CM}과\  $Y$가\ 
Cohen--Macaulay라는 가정에 의해 가군\  $\omega_{X/Y}$는\  $(S_2)$이다. 조건 (3)에
의해\  $U$는 여차원\  $1$인 모든 점을 포함하므로\  Divisors, Lemma
\ref{divisors-lemma-depth-2-hartog}를 사용하면\ 
$j_*\omega_{U/Y} = \omega_{X/Y}$를 얻는다. 따라서\  $U$ 위의 사상은\  $X$로
연장되고 증명이 끝난다.
\end{proof}





\section{연접 가군의 영에 의한 연장}
\label{duality-section-extension-by-zero}

\noindent
이 절과 이어지는 몇 절의 내용은\  \cite{RD}의\  Deligne 부록에서 찾을 수 있다.

\medskip\noindent
이 절에서\  $j : U \to X$는 뇌터 스킴들의 열린 몰입이다. 다음과 같이 구성되는\ 
$D^b_{\textit{Coh}}(\mathcal{O}_X)$의 역계\  $(K_n)$을 생각할 것이다.
$\mathcal{F}^\bullet$을 연접\  $\mathcal{O}_X$-가군들의 유계 복합체라 하고,
$\mathcal{I} \subset \mathcal{O}_X$를\  $V(\mathcal{I}) = X \setminus U$인 준연접
아이디얼층이라 하자. 그러면
$$
K_n = \mathcal{I}^n\mathcal{F}^\bullet
$$
로 둘 수 있다. 더 정확히 말하면\  $K_n$은 차수\  $q$의 항이\ 
$\mathcal{F}^q$의 연접 부분가군\  $\mathcal{I}^n\mathcal{F}^q$인 복합체로 표현되는\ 
$D^b_{\textit{Coh}}(\mathcal{O}_X)$의 대상이다. 사상들\ 
$\ldots \to K_3 \to K_2 \to K_1$은\  $U$에 제한하면 동형을 유도함에 유의하자.
이러한 계를 {\it Deligne 계}라 부르자.
\begin{lemma}
\label{duality-lemma-lift-map}
$j : U \to X$를 뇌터 스킴들의 열린 몰입이라 하자. $(K_n)$을\  Deligne 계라 하고,
상수계\  $(K_n|_U)$의 값을\  $K \in D^b_{\textit{Coh}}(\mathcal{O}_U)$로 쓰자.
$L$을\  $D^b_{\textit{Coh}}(\mathcal{O}_X)$의 대상이라 하자. 그러면\ 
$\colim \Hom_X(K_n, L) = \Hom_U(K, L|_U)$이다.
\end{lemma}

\begin{proof}
$L \to M \to N \to L[1]$을\  $D^b_{\textit{Coh}}(\mathcal{O}_X)$의 식별 삼각형이라
하자. 그러면 가환 그림
\begingroup\small
$$
\xymatrix{
\ldots \ar[r] &
\colim \Hom_X(K_n, L) \ar[r] \ar[d] &
\colim \Hom_X(K_n, M) \ar[r] \ar[d] &
\colim \Hom_X(K_n, N) \ar[r] \ar[d] &
\ldots \\
\ldots \ar[r] &
\Hom_U(K, L|_U) \ar[r] &
\Hom_U(K, M|_U) \ar[r] &
\Hom_U(K, N|_U) \ar[r] &
\ldots
}
$$
\endgroup
을 얻는다. Derived Categories, Lemma
\ref{derived-lemma-representable-homological}과\  Algebra, Lemma
\ref{algebra-lemma-directed-colimit-exact}에 의해 두 행은 완전하다. 따라서 보조정리의
명제가\  $N[-1]$, $L$, $N$, $L[1]$에 대해 성립하면\  5-보조정리에 의해\  $M$에
대해서도 성립한다. 그러므로\  $L$의 표준 절단에 대한 식별 삼각형들(Derived
Categories, Remark \ref{derived-remark-truncation-distinguished-triangle} 참조)을
사용하여\  $L$이 하나의 영 아닌 코호몰로지 층만 갖는 경우로 환원한다.

\medskip\noindent
연접\  $\mathcal{O}_X$-가군들의 유계 복합체\  $\mathcal{F}^\bullet$과\ 
$X \setminus U$를 잘라내는 준연접 아이디얼\ 
$\mathcal{I} \subset \mathcal{O}_X$를 택하여, $K_n$이\ 
$\mathcal{I}^n\mathcal{F}^\bullet$로 표현되게 하자. “우둔한” 절단을 사용하면 항별로
분열되는 양립 가능한 복합체들의 짧은 완전열
$$
0 \to \sigma_{\geq a + 1} \mathcal{I}^n\mathcal{F}^\bullet \to
\mathcal{I}^n\mathcal{F}^\bullet \to
\sigma_{\leq a} \mathcal{I}^n\mathcal{F}^\bullet \to 0
$$
을 얻고, 이는 다시\  $D^b_{\textit{Coh}}(\mathcal{O}_X)$의 식별 삼각형들의 양립
가능한 계에 대응한다. 위와 같이 논증하여\  $\mathcal{F}^\bullet$이 하나의 영 아닌
항만 갖는 경우로 환원한다. 이로써 다음 문단에서 다루는 경우로 환원된다.

\medskip\noindent
연접\  $\mathcal{O}_X$-가군\  $\mathcal{F}$와 연접\  $\mathcal{O}_X$-가군\ 
$\mathcal{G}$가 주어졌을 때, 모든\  $i \geq 0$에 대해 표준 사상
$$
\colim \Ext^i_X(\mathcal{I}^n\mathcal{F}, \mathcal{G})
\longrightarrow
\Ext^i_U(\mathcal{F}|_U, \mathcal{G}|_U)
$$
이 동형임을 보여야 한다. $i = 0$인 경우는\  Cohomology of Schemes, Lemma
\ref{coherent-lemma-homs-over-open}이다. 이제\  $i > 0$이라 가정하자.

\medskip\noindent
단사성. $\xi \in \Ext^i_X(\mathcal{I}^n\mathcal{F}, \mathcal{G})$를\  $U$에 제한하면
영이 되는 원소라 하자. 어떤\  $m \geq n$이 존재하여\  $\xi$를\ 
$\mathcal{I}^m\mathcal{F} = \mathcal{I}^{m - n}\mathcal{I}^n\mathcal{F}$에 제한한
것이 영임을 보여야 한다. $\mathcal{F}$를\  $\mathcal{I}^n\mathcal{F}$로 바꾸어\ 
$n = 0$이라 가정해도 된다. 즉\  $U$에 제한하면 영이 되는\ 
$\xi \in \Ext^i_X(\mathcal{F}, \mathcal{G})$가 주어진다. Derived Categories of
Schemes, Proposition \ref{perfect-proposition-DCoh}에 의해\ 
$D^b_{\textit{Coh}}(\mathcal{O}_X) = D^b(\textit{Coh}(\mathcal{O}_X))$이다. 따라서
연접\  $\mathcal{O}_X$-가군들의 아벨 범주에서\  $\Ext$ 군을 계산할 수 있다. 이로부터
전사\  $\alpha : \mathcal{F}'' \to \mathcal{F}$로서\  $\xi \circ \alpha = 0$인 것이
존재한다(여기에서\  $i > 0$임을 사용한다). $\mathcal{F}' = \Ker(\alpha)$로 두어
짧은 완전열
$$
0 \to \mathcal{F}' \to \mathcal{F}'' \to \mathcal{F} \to 0
$$
을 얻는다. 따라서\  $\xi$는 어떤 원소\ 
$\xi' \in \Ext^{i - 1}_X(\mathcal{F}', \mathcal{G})$의 상이며, 이를\  $U$에 제한한
것은\ 
$\Ext^{i - 1}_U(\mathcal{F}''|_U, \mathcal{G}|_U) \to
\Ext^{i - 1}_U(\mathcal{F}'|_U, \mathcal{G}|_U)$의 상에 속한다. Artin--Rees에
의해 역계\  $(\mathcal{I}^n\mathcal{F}')$와\ 
$(\mathcal{I}^n \mathcal{F}'' \cap \mathcal{F}')$는\  pro-동형이다. Cohomology of
Schemes, Lemma \ref{coherent-lemma-Artin-Rees}를 보라. 짧은 완전열들의 양립 가능한 계
$$
0 \to
\mathcal{F}' \cap \mathcal{I}^n\mathcal{F}'' \to
\mathcal{I}^n\mathcal{F}'' \to
\mathcal{I}^n\mathcal{F} \to 0
$$
이 있으므로 가환 그림
$$
\xymatrix{
\colim \Ext^{i - 1}_X(\mathcal{I}^n\mathcal{F}'', \mathcal{G})
\ar[r] \ar[d] &
\colim \Ext^{i - 1}_X(\mathcal{F}' \cap \mathcal{I}^n\mathcal{F}'', \mathcal{G})
\ar[r] \ar[d] &
\colim \Ext^i_X(\mathcal{I}^n\mathcal{F}, \mathcal{G})
\ar[d] \\
\Ext^{i - 1}_U(\mathcal{F}''|_U, \mathcal{G}|_U) \ar[r] &
\Ext^{i - 1}_U(\mathcal{F}'|_U, \mathcal{G}|_U) \ar[r] &
\Ext^{i - 1}_U(\mathcal{F}|_U, \mathcal{G}|_U)
}
$$
을 얻으며 두 행은 완전하다. $i$에 대한 귀납법과 위 역계에 관한 관찰에 의해 왼쪽
두 수직 화살표가 동형임을 안다. 이제\  $\xi$는 오른쪽 위 군의 원소를 주며, 이는
가운데 위 군의\  $\xi'$의 상이다. 다시 이것은 왼쪽 아래 군의 어떤 원소에서 오는
가운데 아래 군의 원소로 사상된다. 따라서 원하는 대로 어떤\  $n$에 대해\  $\xi$는\ 
$\Ext^i_X(\mathcal{I}^n\mathcal{F}, \mathcal{G})$에서 영으로 사상된다.

\medskip\noindent
전사성. $\xi \in \Ext^i_U(\mathcal{F}|_U, \mathcal{G}|_U)$라 하자. 위와 같이\ 
$i > 0$임을 사용하면 연접\  $\mathcal{O}_U$-가군들의 전사\ 
$\mathcal{H} \to \mathcal{F}|_U$로서\  $\xi$가\ 
$\Ext^i_U(\mathcal{H}, \mathcal{G}|_U)$에서 영으로 사상되는 것을 찾을 수 있다.
Properties, Lemma \ref{properties-lemma-extend-finite-presentation}에 의해, $U$에
제한하면\  $\mathcal{H} \to \mathcal{F}|_U$가 되는 연접\  $\mathcal{O}_X$-가군들의
사상\  $\varphi : \mathcal{F}'' \to \mathcal{F}$를 찾을 수 있다. 이 보조정리는\ 
$\varphi$가 전사임을 보장하지 않지만 이는 문제가 되지 않는다(조금 더 작업하면
전사인\  $\varphi$를 택할 수 있다). $\mathcal{F}' = \Ker(\varphi)$로 쓰자. 짧은
완전열
$$
0 \to \mathcal{F}'|_U \to \mathcal{F}''|_U \to \mathcal{F}|_U \to 0
$$
은\  $\xi$가\ 
$\Ext^{i - 1}_U(\mathcal{F}'|_U, \mathcal{G}|_U)$의\  $\xi'$의 상임을 보여준다.
$i$에 대한 귀납법으로 어떤\  $n$을 택하여\  $\xi'$가 어떤\  $\xi'_n$의 상이고 그 원소가\ 
$\Ext^{i - 1}_X(\mathcal{I}^n\mathcal{F}', \mathcal{G})$에 속하도록 할 수 있다.
Artin--Rees에 의해 어떤\  $m \geq n$에
대해\ 
$\mathcal{F}' \cap \mathcal{I}^m\mathcal{F}'' \subset
\mathcal{I}^n\mathcal{F}'$이다. 짧은 완전열
$$
0 \to \mathcal{F}' \cap \mathcal{I}^m\mathcal{F}'' \to
\mathcal{I}^m\mathcal{F}'' \to \mathcal{I}^m\Im(\varphi) \to 0
$$
을 사용하면\  $\xi'_n$의 상은\ 
$\Ext^{i - 1}_X(\mathcal{F}' \cap \mathcal{I}^m\mathcal{F}'', \mathcal{G})$에서
경계 사상에 의해\  $\Ext^i_X(\mathcal{I}^m\Im(\varphi), \mathcal{G})$의 원소\ 
$\xi_m$으로 사상되고, 이 원소는\  $\xi$로 사상된다. $\Im(\varphi)$와\ 
$\mathcal{F}$는\  $U$ 위에서 일치하므로\ 
$\mathcal{F}/\mathcal{I}^m\Im(\varphi)$는\  $X \setminus U$에 지지된다. 따라서
어떤\  $l \geq m$에 대해\ 
$\mathcal{I}^l\mathcal{F} \subset \mathcal{I}^m\Im(\varphi)$이다. Cohomology of
Schemes, Lemma \ref{coherent-lemma-power-ideal-kills-sheaf}를 보라. $\xi_m$의 상을\ 
$\Ext^i_X(\mathcal{I}^l\mathcal{F}, \mathcal{G})$에서 취하면 증명이 끝난다.
\end{proof}
\begin{lemma}
\label{duality-lemma-lift-map-plus}
Lemma \ref{duality-lemma-lift-map}의 결과는\  $L \in D^+_{\textit{Coh}}(\mathcal{O}_X)$에
대해서도 성립한다.
\end{lemma}

\begin{proof}
실제로\  $(K_n)$이\  Deligne 계이면 어떤\  $b \in \mathbf{Z}$가 존재하여\  $i > b$일 때\ 
$H^i(K_n) = 0$이다. 그러면\ 
$\Hom(K_n, L) = \Hom(K_n, \tau_{\leq b}L)$이고\ 
$\Hom(K, L) = \Hom(K, \tau_{\leq b}L)$이다. 따라서\  $\tau_{\leq b}L$에 대한
보조정리의 결과를 사용하면 끝난다.
\end{proof}
\begin{lemma}
\label{duality-lemma-extension-by-zero}
$j : U \to X$를 뇌터 스킴들의 열린 몰입이라 하자.
\begin{enumerate}
\item $(K_n)$과\  $(L_n)$을\  Deligne 계라 하자. 상수계\  $(K_n|_U)$와\  $(L_n|_U)$의
값을 각각\  $K$와\  $L$이라 하자. $D(\mathcal{O}_U)$의 사상\ 
$\alpha : K \to L$이 주어지면, $U$에 제한했을 때\  $\alpha$가 되는\ 
$D^b_{\textit{Coh}}(\mathcal{O}_X)$의\  pro-계의 유일한 사상\ 
$(K_n) \to (L_n)$이 존재한다.
\item $K \in D^b_{\textit{Coh}}(\mathcal{O}_U)$가 주어지면, $(K_n|_U)$가 값\ 
$K$인 상수계가 되는\  Deligne 계\  $(K_n)$이 존재한다.
\item (2)의\  $D^b_{\textit{Coh}}(\mathcal{O}_X)$의\  pro-대상\  $(K_n)$은
(pro-대상으로서) 유일한 동형 아래 유일하다.
\end{enumerate}
\end{lemma}

\begin{proof}
(1)은\  Lemma \ref{duality-lemma-lift-map}과\  pro-계 사이의 사상이 그 계들이 공표현하는
함자 사이의 사상과 같다는 사실의 직접적인 귀결이다. Categories, Remark
\ref{categories-remark-pro-category-copresheaves}를 보라.

\medskip\noindent
$K$가 (2)에서와 같다고 하자. $U$에 제한하면\  $K$와 동형인\ 
$K' \in D^b_{\textit{Coh}}(\mathcal{O}_X)$를 택할 수 있다. Derived Categories of
Schemes, Lemma \ref{perfect-lemma-lift-coherent}를 보라. Derived Categories of
Schemes, Proposition \ref{perfect-proposition-DCoh}에 의해\  $K'$를 연접\ 
$\mathcal{O}_X$-가군들의 유계 복합체\  $\mathcal{F}^\bullet$로 표현할 수 있다.
소멸 자취가\  $X \setminus U$인 준연접 아이디얼층\ 
$\mathcal{I} \subset \mathcal{O}_X$를 택하자(예를 들어\  $X \setminus U$ 위의
축약 유도 부분스킴 구조에 대응하는\  $\mathcal{I}$를 택한다). 그러면 이 절의
도입부에서처럼\  $K_n$을 복합체\  $\mathcal{I}^n\mathcal{F}^\bullet$로 표현되는
대상으로 둘 수 있다.

\medskip\noindent
(3)은 (1)과 (2)에서 즉시 따른다.
\end{proof}
\begin{lemma}
\label{duality-lemma-extension-by-zero-triangle}
$j : U \to X$를 뇌터 스킴들의 열린 몰입이라 하자.
$$
K \to L \to M \to K[1]
$$
을\  $D^b_{\textit{Coh}}(\mathcal{O}_U)$의 식별 삼각형이라 하자. 그러면\ 
$D^b_{\textit{Coh}}(\mathcal{O}_X)$의 식별 삼각형들의 역계
$$
K_n \to L_n \to M_n \to K_n[1]
$$
이 존재하여\  $(K_n)$, $(L_n)$, $(M_n)$은\  Deligne 계이고, 이 식별 삼각형들을\ 
$U$에 제한한 것은 처음의 식별 삼각형과 동형이다.
\end{lemma}

\begin{proof}
$(K_n)$을\  Lemma \ref{duality-lemma-extension-by-zero}의 (2)에서와 같이 택하자. $U$에
제한하면\  $L$이 되는\  $D^b_{\textit{Coh}}(\mathcal{O}_X)$의 대상\  $L'$을 택하자
(보조정리에 의해 가능하다). Lemma \ref{duality-lemma-lift-map}에 의해 어떤\  $n$과\ 
$X$ 위의 사상\  $K_n \to L'$을 택하여, 이를\  $U$에 제한하면 주어진 화살표\ 
$K \to L$이 되게 할 수 있다. 따라서\  $D^b_{\textit{Coh}}(\mathcal{O}_X)$의 사상\ 
$K' \to L'$로서\  $U$에 제한하면 주어진 화살표\  $K \to L$이 되는 것이 존재한다.

\medskip\noindent
Derived Categories of Schemes, Proposition \ref{perfect-proposition-DCoh}에 의해,
$K' \to L'$을 표현하는 연접\  $\mathcal{O}_X$-가군들의 유계 복합체의 사상\ 
$\alpha^\bullet : \mathcal{F}^\bullet \to
\mathcal{G}^\bullet$을 택할 수 있다. 소멸 자취가\  $X \setminus U$인 준연접
아이디얼층\  $\mathcal{I} \subset \mathcal{O}_X$를 택하고\ 
$K_n = \mathcal{I}^n\mathcal{F}^\bullet$ 및\ 
$L_n = \mathcal{I}^n\mathcal{G}^\bullet$로 둔다. $\alpha^\bullet$은 복합체의 사상\ 
$\alpha_n^\bullet : \mathcal{I}^n\mathcal{F}^\bullet \to
\mathcal{I}^n\mathcal{G}^\bullet$을 유도한다. Derived Categories, Section
\ref{derived-section-cones}의 원뿔 구성에서 다음이 분명하다.
$$
C(\alpha_n)^\bullet = \mathcal{I}^nC(\alpha^\bullet)
$$
그러므로\  $M_n = C(\alpha_n)^\bullet$로 둘 수 있다. 즉 식별 삼각형들의 양립
가능한 계(Derived Categories, Section
\ref{derived-section-canonical-delta-functor}의 논의 참조)
$$
K_n \to L_n \to M_n \to K_n[1]
$$
이 있으며, 이를\  $U$에 제한하면 공리\  TR3과\  Derived Categories, Lemma
\ref{derived-lemma-third-isomorphism-triangle}에 의해 처음의 식별 삼각형과 동형이다.
\end{proof}
\begin{remark}
\label{duality-remark-extension-by-zero}
$j : U \to X$를 뇌터 스킴들의 열린 몰입이라 하자.
$K \in D^b_{\textit{Coh}}(\mathcal{O}_U)$를, $U$에 제한하면 값\  $K$가 되는\  Deligne
계로 보내면 함자
$$
Rj_! :
D^b_{\textit{Coh}}(\mathcal{O}_U)
\longrightarrow
\text{Pro-}D^b_{\textit{Coh}}(\mathcal{O}_X)
$$
가 정해진다. 이 함자는\  Lemma \ref{duality-lemma-extension-by-zero-triangle}에 의해
“완전”하고, Lemma \ref{duality-lemma-lift-map}에 의해 함자\ 
$j^* : D^b_{\textit{Coh}}(\mathcal{O}_X) \to
D^b_{\textit{Coh}}(\mathcal{O}_U)$의 “왼쪽 수반”이다.
\end{remark}
\begin{remark}
\label{duality-remark-extension-by-zero-linear-pro-system}
$(A_n)$과\  $(B_n)$을 한 범주\  $\mathcal{C}$의 역계라 하자. $(A_n)$에서\  $(B_n)$으로의
선형\  pro-사상이란 어떤 고정된 정수\  $c, d \geq 1$에 대해 모든\  $n \geq 1$에
대하여 정의된 사상들의 양립 가능한 족\  $\varphi_n : A_{cn + d} \to B_n$으로
주어지는 것이라 하자. 어떤\  $c'' \geq \max(c, c')$와\ 
$d'' \geq \max(d, d')$가 존재하여 모든\  $n$에 대해 두 유도 사상\ 
$A_{c'' n + d''} \to B_n$이 같으면\ 
$(\varphi_n : A_{cn + d} \to B_n)$과\ 
$(\psi_n : A_{c'n + d'} \to B_n)$이 같은 사상을 정한다고 하자. $U$ 위의 값이
주어진\  Deligne 계\  $(K_n)$은 선형\  pro-동형 아래 잘 정의될 듯하다. 이것이 필요해지면
여기에 주의 깊게 정식화하고 증명할 것이다.
\end{remark}
\begin{lemma}
\label{duality-lemma-deligne-system-2-out-of-3}
$j : U \to X$를 뇌터 스킴들의 열린 몰입이라 하자.
$$
K_n \to L_n \to M_n \to K_n[1]
$$
을\  $D^b_{\textit{Coh}}(\mathcal{O}_X)$의 식별 삼각형들의 역계라 하자. $(K_n)$과\ 
$(M_n)$이\  Deligne 계와\  pro-동형이면\  $(L_n)$도 그러하다.
\end{lemma}

\begin{proof}
계\  $(K_n|_U)$와\  $(M_n|_U)$은 상수계와\  pro-동형이므로 본질적으로 상수이다. 그
값을\  $K$와\  $M$이라 쓰자. Derived Categories, Lemma
\ref{derived-lemma-essentially-constant-2-out-of-3}에 의해 역계\  $L_n|_U$도
본질적으로 상수이다. 그 값을\  $L$이라 쓰자.
$N \in D^b_{\textit{Coh}}(\mathcal{O}_X)$이라 하고 다음 가환 그림을 생각하자.
\begingroup\small
$$
\xymatrix{
\ldots \ar[r] &
\colim \Hom_X(M_n, N) \ar[r] \ar[d] &
\colim \Hom_X(L_n, N) \ar[r] \ar[d] &
\colim \Hom_X(K_n, N) \ar[r] \ar[d] &
\ldots \\
\ldots \ar[r] &
\Hom_U(M, N|_U) \ar[r] &
\Hom_U(L, N|_U) \ar[r] &
\Hom_U(K, N|_U) \ar[r] &
\ldots
}
$$
\endgroup
Lemma \ref{duality-lemma-lift-map}과 동형인\  ind-계들은 같은 쌍대극한을 갖는다는 사실에 의해,
가운데 것의 오른쪽 두 수직 화살표와 왼쪽 두 수직 화살표는 동형이다. 5-보조정리에
의해 가운데 수직 화살표도 동형이다. 이제\  $(L'_n)$을\  $U$에 제한했을 때 상수값\ 
$L$을 갖는\  Deligne 계라 하자(Lemma \ref{duality-lemma-extension-by-zero}에 의해 존재한다).
그러면\ 
$\colim \Hom_X(L'_n, N) = \Hom_U(L, N|_U)$도 성립한다. 따라서\  Categories,
Remark \ref{categories-remark-pro-category-copresheaves}에 의해\  pro-계\  $(L_n)$과\ 
$(L'_n)$은\  pro-동형이다.
\end{proof}
\begin{lemma}
\label{duality-lemma-consequence-Artin-Rees-bis}
$X$를 뇌터 스킴이라 하자. $\mathcal{I} \subset \mathcal{O}_X$를 준연접
아이디얼층이라 하고, $\mathcal{F}^\bullet$을 연접\  $\mathcal{O}_X$-가군들의
복합체라 하며, $p \in \mathbf{Z}$라 하자.
$\mathcal{H} = H^p(\mathcal{F}^\bullet)$ 및\ 
$\mathcal{H}_n = H^p(\mathcal{I}^n\mathcal{F}^\bullet)$로 두자. 그러면 표준\ 
$\mathcal{O}_X$-가군 사상들
$$
\ldots \to \mathcal{H}_3 \to \mathcal{H}_2 \to \mathcal{H}_1 \to \mathcal{H}
$$
이 있다. 어떤\  $c > 0$이 존재하여\  $n \geq c$일 때\ 
$\mathcal{H}_n \to \mathcal{H}$의 상은\  $\mathcal{I}^{n - c}\mathcal{H}$에 포함되고,
표준\  $\mathcal{O}_X$-가군 사상\ 
$\mathcal{I}^n\mathcal{H} \to \mathcal{H}_{n - c}$가 존재하며, 합성들
$$
\mathcal{I}^n \mathcal{H} \to \mathcal{H}_{n - c} \to
\mathcal{I}^{n - 2c}\mathcal{H}
\quad\text{and}\quad
\mathcal{H}_n \to \mathcal{I}^{n - c}\mathcal{H} \to \mathcal{H}_{n - 2c}
$$
은 표준 사상들이다. 특히 역계\  $(\mathcal{H}_n)$과\ 
$(\mathcal{I}^n\mathcal{H})$는\  $\textit{Mod}(\mathcal{O}_X)$의\  pro-대상으로서
동형이다.
\end{lemma}

\begin{proof}
$X$가 아핀이면 대수로 옮긴 이 명제는\  More on Algebra, Lemma
\ref{more-algebra-lemma-consequence-Artin-Rees-bis}이다. 일반적인 경우에는 그
보조정리의 증명에서 대수의\  Artin--Rees에 대한 참조를\  Cohomology of Schemes,
Lemma \ref{coherent-lemma-Artin-Rees}에 대한 참조로 바꾸고 정확히 같은 방식으로
논증한다. 세부사항은 생략한다.
\end{proof}
\begin{lemma}
\label{duality-lemma-characterize-extension-by-zero-algebra}
$j : U \to X$를 뇌터 스킴들의 열린 몰입이라 하자. $a \leq b$를 정수라 하고,
$(K_n)$을\  $D^b_{\textit{Coh}}(\mathcal{O}_X)$의 역계로서\  $i \not \in [a, b]$이면\ 
$H^i(K_n) = 0$인 것이라 하자. 다음은 동치이다.
\begin{enumerate}
\item $(K_n)$은\  Deligne 계와\  pro-동형이다.
\item 모든\  $p \in \mathbf{Z}$에 대해 연접\  $\mathcal{O}_X$-가군\  $\mathcal{F}$가
존재하여\  pro-계\  $(H^p(K_n))$과\  $(\mathcal{I}^n\mathcal{F})$가\  pro-동형이다.
\end{enumerate}
\end{lemma}

\begin{proof}
(1)을 가정하자. (2)를 증명할 때\  $(K_n)$이\  Deligne 계라고 가정해도 된다. 정의에
의해 연접\  $\mathcal{O}_X$-가군들의 유계 복합체\  $\mathcal{F}^\bullet$과\ 
$X \setminus U$를 잘라내는 준연접 아이디얼층을 택하여\  $K_n$이\ 
$\mathcal{I}^n\mathcal{F}^\bullet$로 표현되게 할 수 있다. 따라서 결과는\  Lemma
\ref{duality-lemma-consequence-Artin-Rees-bis}에서 따른다.

\medskip\noindent
(2)를 가정하자. $b - a$에 대한 귀납법으로\  $(K_n)$이 (1)에서와 같음을 증명한다.
$a = b$이면 본질적으로 가정에 의해 (1)이 성립한다. $a < b$이면 식별 삼각형들의
양립 가능한 계
$$
\tau_{\leq a}K_n \to K_n \to \tau_{\geq a + 1}K_n \to (\tau_{\leq a}K_n)[1]
$$
을 생각한다. Derived Categories, Remark
\ref{derived-remark-truncation-distinguished-triangle}을 보라. $b - a$에 대한 귀납법에
의해\  $\tau_{\leq a}K_n$과\  $\tau_{\geq a + 1}K_n$은\  Deligne 계와\  pro-동형이다.
Lemma \ref{duality-lemma-deligne-system-2-out-of-3}로 결론을 얻는다.
\end{proof}
\begin{lemma}
\label{duality-lemma-characterize-extension-by-zero}
$j : U \to X$를 뇌터 스킴들의 열린 몰입이라 하고, $(K_n)$을\ 
$D^b_{\textit{Coh}}(\mathcal{O}_X)$의 역계라 하자.
$X = W_1 \cup \ldots \cup W_r$을 열린 덮개라 하자. 다음은 동치이다.
\begin{enumerate}
\item $(K_n)$은\  Deligne 계와\  pro-동형이다.
\item 각\  $i$에 대해 제한\  $(K_n|_{W_i})$는 열린 몰입\ 
$U \cap W_i \to W_i$에 관한\  Deligne 계와\  pro-동형이다.
\end{enumerate}
\end{lemma}

\begin{proof}
$r$에 대해 귀납한다. $r = 1$이면 결과는 분명하다. $r > 1$이라 가정하고\ 
$V = W_1 \cup \ldots \cup W_{r - 1}$로 두자. 귀납법에 의해\  $(K_n|_V)$는\  Deligne
계이다. 이로써 다음 문단의 논의로 환원된다.

\medskip\noindent
$X = V \cup W$가 열린 덮개이고\  $(K_n|_W)$와\  $(K_n|_V)$가\  Deligne 계와\ 
pro-동형이라고 가정하자. $(K_n)$이\  Deligne 계와\  pro-동형임을 보여야 한다.
$(K_n|_{V \cap W})$는\  Deligne 계와\  pro-동형임에 유의하자(Deligne 계의 구성에서
열린집합으로의 제한이 이를 보존함이 즉시 따른다). 특히\  pro-계\ 
$(K_n|_{U \cap V})$, $(K_n|_{U \cap W})$, $(K_n|_{U \cap V \cap W})$는 본질적으로
상수이다. Cohomology, Lemma \ref{cohomology-lemma-exact-sequence-j-star}의 식별
삼각형들과\  Derived Categories, Lemma
\ref{derived-lemma-essentially-constant-2-out-of-3}에서\  $(K_n|_U)$가 본질적으로
상수임이 따른다. 이 계의 값을\  $K \in D^b_{\textit{Coh}}(\mathcal{O}_U)$로 쓰자.
$L$을\  $D^b_{\textit{Coh}}(\mathcal{O}_X)$의 대상이라 하고 다음 그림을 생각하자.
\begingroup\footnotesize
$$
\xymatrix{
\colim \Ext^{-1}(K_n|_V, L|_V) \oplus
\colim \Ext^{-1}(K_n|_W, L|_W) \ar[r] \ar[d] &
\Ext^{-1}(K|_{U \cap V}, L|_{U \cap V}) \oplus
\Ext^{-1}(K|_{U \cap W}, L|_{U \cap W}) \ar[d] \\
\colim \Ext^{-1}(K_n|_{V \cap W}, L|_{V \cap W}) \ar[r] \ar[d] &
\Ext^{-1}(K|_{U \cap V \cap W}, L|_{U \cap V \cap W}) \ar[d] \\
\colim \Hom(K_n, L) \ar[d] \ar[r] &
\Hom(K|_U, L|_U) \ar[d] \\
\colim \Hom(K_n|_V, L|_V) \oplus \colim \Hom(K_n|_W, L|_W) \ar[r] \ar[d] &
\Hom(K|_{U \cap V}, L|_{U \cap V}) \oplus
\Hom(K|_{U \cap W}, L|_{U \cap W}) \ar[d] \\
\colim \Hom(K_n|_{V \cap W}, L|_{V \cap W}) \ar[r] &
\Hom(K|_{U \cap V \cap W}, L|_{U \cap V \cap W})
}
$$
\endgroup
Cohomology, Lemma \ref{cohomology-lemma-mayer-vietoris-hom}과 여과 쌍대극한이
완전하다는 사실에 의해 두 수직 열은 완전하다. Lemma \ref{duality-lemma-lift-map}과\ 
pro-동형인 계들이 같은 쌍대극한을 갖는다는 사실에 의해 가운데 것을 제외한 모든
수평 화살표는 동형이다. 따라서\  5-보조정리에 의해 가운데 것도 동형이다. 이로부터\ 
$(K_n)$이\  $K$를 위한\  Deligne 계와\  pro-동형임이 따른다. 실제로\  $(K'_n)$을\  $U$에
제한하면 상수값\  $K$를 갖는\  Deligne 계라 하면(Lemma
\ref{duality-lemma-extension-by-zero}에 의해 존재한다),
$\colim \Hom_X(K'_n, L) = \Hom_U(K, L|_U)$도 성립한다. 따라서\  Categories,
Remark \ref{categories-remark-pro-category-copresheaves}에 의해\  pro-계\  $(K_n)$과\ 
$(K'_n)$은\  pro-동형이다.
\end{proof}
\begin{lemma}
\label{duality-lemma-tensoring-Deligne-system}
$j : U \to X$를 뇌터 스킴들의 열린 몰입이라 하고,
$\mathcal{I} \subset \mathcal{O}_X$를\  $V(\mathcal{I}) = X \setminus U$인 준연접
아이디얼층이라 하자. $K$를\  $D^b_{\textit{Coh}}(\mathcal{O}_X)$의 대상이라 하자.
그러면
$$
K \otimes_{\mathcal{O}_X}^\mathbf{L} \mathcal{I}^n
$$
은\  $U$ 위에서 상수값\  $K|_U$를 갖는\  Deligne 계와\  pro-동형이다.
\end{lemma}

\begin{proof}
Lemma \ref{duality-lemma-characterize-extension-by-zero}에 의해 문제는\  $X$ 위에서
국소적이다. 따라서\  $X$가 한 뇌터 환의 스펙트럼이라고 가정해도 된다. 이 경우
명제는 대수적 판본인\  More on Algebra, Lemma
\ref{more-algebra-lemma-tensoring-Deligne-system}에서 따른다.
\end{proof}






\section{콤팩트 지지 코호몰로지를 위한 준비}
\label{duality-section-preliminaries-compactly-supported}

\noindent
Situation \ref{duality-situation-shriek}에서\  $f : X \to Y$를 범주\  $\textit{FTS}_S$의
사상이라 하자. 앞 절의 구성들을 사용하여 함자
$$
Rf_! :
D^b_{\textit{Coh}}(\mathcal{O}_X)
\longrightarrow
\text{Pro-}D^b_{\textit{Coh}}(\mathcal{O}_Y)
$$
를 구성할 것이다. $f$가 열린 몰입이면 이는\  Remark
\ref{duality-remark-extension-by-zero}의 함자로 환원되고, 일반적으로는\  $f$의 콤팩트화를
사용하여 구성된다. 이 작업에 앞서 구성이 잘 정의됨을 증명하기 위해 다음 보조정리들이
필요하다.

\begin{lemma}
\label{duality-lemma-well-defined-pre}
$f : X \to Y$를 뇌터 스킴들의 고유 사상이라 하자. $V \subset Y$를 열린
부분스킴이라 하고\  $U = f^{-1}(V)$로 두자. 그림은 다음과 같다.
$$
\xymatrix{
U \ar[r]_j \ar[d]_g & X \ar[d]^f \\
V \ar[r]^{j'} & Y
}
$$
그러면 함자\ 
$D^b_{\textit{Coh}}(\mathcal{O}_U) \to
\text{Pro-}D^b_{\textit{Coh}}(\mathcal{O}_Y)$의 표준 동형\ 
$Rj'_! \circ Rg_* \to Rf_* \circ Rj_!$이 있다. 여기서\  $Rj_!$와\  $Rj'_!$은\ 
Remark \ref{duality-remark-extension-by-zero}에서와 같다.
\end{lemma}

\begin{proof}[첫 번째 증명]
$K$를\  $D^b_{\textit{Coh}}(\mathcal{O}_U)$의 대상이라 하자. $(K_n)$을\  $U \to X$를
위한\  Deligne 계로서\  $U$에 제한하면 상수값\  $K$를 갖는 것이라 하자. 물론 이는\ 
$(K_n)$이\  $\text{Pro-}D^b_{\textit{Coh}}(\mathcal{O}_X)$에서\  $Rj_!K$를 표현한다는
뜻이다. $Rj'_!Rg_*K$와\  $Rf_*Rj_!K$를\  $V$에 제한하면 모두 상수값\  $Rg_*K$인\ 
pro-대상이 됨에 유의하자. 첫 번째 것에는 자명하고, 두 번째 것은\ 
$(Rf_*K_n)|_V = Rg_*(K_n|_U) = Rg_*K$라는 사실에서 따른다. Lemma
\ref{duality-lemma-extension-by-zero}의\  Deligne 계의 유일성에 의해\  $(Rf_*K_n)$이\  Deligne
계와\  pro-동형임을 보이면 충분하다. 인용한 보조정리는 얻어진 동형이 함자적임도
보여준다.

\medskip\noindent
$(Rf_*K_n)$이\  Deligne 계와\  pro-동형임의 증명. 먼저 문제는\  $K$에 대응하는\  Deligne
계\  $(K_n)$의 선택과 무관함에 유의하자(앞서 말한 유일성에 의한다). Lemmas
\ref{duality-lemma-extension-by-zero-triangle}와\ 
\ref{duality-lemma-deligne-system-2-out-of-3}에 의해,
$D^b_{\textit{Coh}}(\mathcal{O}_U)$의 식별 삼각형
$$
K \to L \to M \to K[1]
$$
이 있고 결과가\  $K$와\  $M$에 대해 성립하면\  $L$에 대해서도 성립한다. 표준 절단의
식별 삼각형들(Derived Categories, Remark
\ref{derived-remark-truncation-distinguished-triangle})을 사용하여 다음 문단의 문제로
환원한다.

\medskip\noindent
$\mathcal{F}$를 연접\  $\mathcal{O}_X$-가군이라 하자.
$\mathcal{J} \subset \mathcal{O}_Y$를\  $Y \setminus V$를 잘라내는 준연접
아이디얼층이라 하자. $\mathcal{J}^n\mathcal{F}$로\ 
$f^*\mathcal{J}^n \otimes \mathcal{F} \to \mathcal{F}$의 상을 나타내자.
$(Rf_*(\mathcal{J}^n\mathcal{F}))$가\  Deligne 계임을 보여야 한다. Lemma
\ref{duality-lemma-characterize-extension-by-zero}에 의해 문제는\  $Y$ 위에서 국소적이다.
따라서\  $Y = \Spec(A)$가 아핀이고\  $\mathcal{J}$가 아이디얼\  $I \subset A$에
대응한다고 가정해도 된다. Lemma
\ref{duality-lemma-characterize-extension-by-zero-algebra}에 의해, 어떤 유한\  $A$-가군\ 
$M$에 대해 코호몰로지 가군의 역계\  $(H^p(X, I^n\mathcal{F}))$가 역계\  $(I^n M)$과\ 
pro-동형임을 보이면 충분하다. 이는\  Cohomology of Schemes, Lemma
\ref{coherent-lemma-cohomology-powers-ideal-application}에서 증명된다.
\end{proof}

\begin{proof}[두 번째 증명]
$K$를\  $D^b_{\textit{Coh}}(\mathcal{O}_U)$의 대상이라 하고\  $L$을\ 
$D^b_{\textit{Coh}}(\mathcal{O}_Y)$의 대상이라 하자. $K$와\  $L$에 함자적인 전단사
$$
\Hom_{\text{Pro-}D^b_{\textit{Coh}}(\mathcal{O}_Y)}(Rj'_!Rg_*K, L)
\longrightarrow
\Hom_{\text{Pro-}D^b_{\textit{Coh}}(\mathcal{O}_Y)}(Rf_*Rj_!K, L)
$$
를 구성할 것이다. $K$를 고정하면\  Categories, Remark
\ref{categories-remark-pro-category-copresheaves}에 의해 이는\  pro-대상의 동형\ 
$Rf_*Rj_!K \to Rj'_!Rg_*K$를 정하고, $K$를 변화시키면 함자의 동형이 정해진다.
실제로 이 동형을 만들기 위해 다음 함자적인 등식들의 열을 사용한다.
\begin{align*}
\Hom_{\text{Pro-}D^b_{\textit{Coh}}(\mathcal{O}_Y)}(Rj'_!Rg_*K, L)
& =
\Hom_V(Rg_*K, L|_V) \\
& =
\Hom_U(K, g^!(L|_V)) \\
& =
\Hom_U(K, f^!L|_U)) \\
& =
\Hom_{\text{Pro-}D^b_{\textit{Coh}}(\mathcal{O}_X)}(Rj_!K, f^!L) \\
& =
\Hom_{\text{Pro-}D^b_{\textit{Coh}}(\mathcal{O}_Y)}(Rf_*Rj_!K, L)
\end{align*}
첫 번째 등식은\  Lemma \ref{duality-lemma-lift-map}에 의해 참이다. 두 번째 등식은\  $g$가\ 
$f$를\  $V$로 기저변환한 고유 사상이므로, 구성상\  $g^!$가 밀어앞으로 보내기의 오른쪽
수반이라는 사실에서 따른다. Section \ref{duality-section-upper-shriek}를 보라. 세 번째
등식은\  Lemma \ref{duality-lemma-restrict-before-or-after}에 의해\ 
$g^!(L|_V) = f^!L|_U$이므로 성립한다. Lemma \ref{duality-lemma-shriek-coherent}에 의해\ 
$f^!L$은\  $D^+_{\textit{Coh}}(\mathcal{O}_X)$에 속하므로 네 번째 등식은\  Lemma
\ref{duality-lemma-lift-map-plus}에서 따른다. 다섯 번째 등식도\  $f$가 고유이어서\  $f^!$가\ 
$Rf_*$의 오른쪽 수반이므로 성립한다.
\end{proof}
\begin{lemma}
\label{duality-lemma-well-defined}
$j : U \to X$를 뇌터 스킴들의 열린 몰입이라 하자. $j' : U \to X'$를\  $X$ 위의\ 
$U$의 콤팩트화라 하고(증명을 보라), 구조 사상을\  $f : X' \to X$로 쓰자. 그러면
함자\  $D^b_{\textit{Coh}}(\mathcal{O}_U) \to
\text{Pro-}D^b_{\textit{Coh}}(\mathcal{O}_X)$의 표준 동형\ 
$Rj_! \to Rf_* \circ R(j')_!$이 있다. 여기서\  $Rj_!$와\  $Rj'_!$은\  Remark
\ref{duality-remark-extension-by-zero}에서와 같다.
\end{lemma}

\begin{proof}
$X'$가\  $X$ 위의\  $U$의 콤팩트화라는 것은 정확히\  $f : X' \to X$가 고유이고,
$j'$가 열린 몰입이며, $j = f \circ j'$라는 뜻이다. More on Flatness, Section
\ref{flat-section-compactify}를 보라. $j'(U) = f^{-1}(j(U))$이면 보조정리는\  Lemma
\ref{duality-lemma-well-defined-pre}에서 즉시 따른다. $j'(U) \not = f^{-1}(j(U))$이면\ 
$X'' \subset X'$로\  $j' : U \to X'$의 스킴론적 폐포를 나타내고,
$j'' : U \to X''$로 이에 대응하는 열린 몰입을 나타내자. 그림은 다음과 같다.
$$
\xymatrix{
& & X'' \ar[d]^{f'} \\
& & X' \ar[d]^f \\
U \ar[rr]^j \ar[rru]^{j'} \ar[rruu]^{j''} & & X
}
$$
More on Flatness, Lemma \ref{flat-lemma-compactifications-cofiltered}의 (c)와 위 논의에
의해 동형들\ 
$Rf'_* \circ Rj''_! = Rj'_!$ 및\  $R(f \circ f')_* \circ Rj''_! = Rj_!$이 있다.
$R(f \circ f')_* = Rf_* \circ Rf'_*$이므로 결론을 얻는다.
\end{proof}
\begin{remark}
\label{duality-remark-covariance-open-j-lower-shriek}
$X \supset U \supset U'$를 한 뇌터 스킴\  $X$의 열린 부분스킴들이라 하자. 포함
사상들을\  $j : U \to X$와\  $j' : U' \to X$로 쓰자. 다음 표준 사상이 존재한다고
주장한다.
$$
Rj'_!(K|_{U'}) \longrightarrow Rj_!K
$$
이는\  $D^b_{\textit{Coh}}(\mathcal{O}_U)$의\  $K$에 대해 함자적이다. 실제로\  Lemma
\ref{duality-lemma-lift-map}에 의해\  $D^b_{\textit{Coh}}(\mathcal{O}_X)$의 임의의\  $L$에
대하여 사상
\begin{align*}
\Hom_{\text{Pro-}D^b_{\textit{Coh}}(\mathcal{O}_X)}(Rj_!K, L)
& =
\Hom_U(K, L|_U) \\
& \to
\Hom_{U'}(K|_{U'}, L|_{U'}) \\
& =
\Hom_{\text{Pro-}D^b_{\textit{Coh}}(\mathcal{O}_X)}(Rj'_!(K|_{U'}), L)
\end{align*}
이 있으며, 이는\  $L$과\  $K'$에 함자적이다. $L$에 대한 함자성과\  Categories,
Remark \ref{categories-remark-pro-category-copresheaves}에 의해 표준 사상\ 
$Rj'_!(K|_{U'}) \to Rj_!K$를 얻고, 위 화살표의\  $K$에 대한 함자성에 의해 이 사상도\ 
$K$에 대해 함자적이다.

\medskip\noindent
이 화살표를 명시적으로 구성해 보자. $\mathcal{F}^\bullet$을 연접\ 
$\mathcal{O}_X$-가군들의 유계 복합체로서\  $U$에 제한하면 유도범주에서\  $K$를
표현하는 것이라 하자. Lemma \ref{duality-lemma-extension-by-zero}의 증명에서 이러한
복합체가 항상 존재함을 보았다. $\mathcal{I}$, resp.\ $\mathcal{I}'$를 각각\  $X$
위의 준연접 아이디얼층으로서\ 
$V(\mathcal{I}) = X \setminus U$, resp.\ $V(\mathcal{I}') = X \setminus U'$인
것이라 하자. $\mathcal{I}$를\  $\mathcal{I} + \mathcal{I}'$으로 바꾸어\ 
$\mathcal{I}' \subset \mathcal{I}$이라 가정해도 된다. 구성상\  $Rj_!K$, resp.\
$Rj'_!(K|_{U'})$는\  $D^b_{\textit{Coh}}(\mathcal{O}_X)$의 역계\  $(K_n)$,
resp.\ $(K'_n)$으로 표현되며
$$
K_n = \mathcal{I}^n\mathcal{F}^\bullet
\quad\text{resp.}\quad
K'_n = (\mathcal{I}')^n\mathcal{F}^\bullet
$$
이다. 위에서 구성한 사상은 포함\ 
$(\mathcal{I}')^n \subset \mathcal{I}^n$에서 오는 사상\  $K'_n \to K_n$들로
주어짐이 분명하다.
\end{remark}





\section{연접 가군의 콤팩트 지지 코호몰로지}
\label{duality-section-compactly-supported}

\noindent
Situation \ref{duality-situation-shriek}에서\  $\textit{FTS}_S$의 사상\  $f : X \to Y$가
주어졌을 때 함자
$$
Rf_! : D^b_{\textit{Coh}}(\mathcal{O}_X)
\longrightarrow
\text{Pro-}D^b_{\textit{Coh}}(\mathcal{O}_Y)
$$
를 정의하자. 실제로\  More on Flatness, Theorem \ref{flat-theorem-nagata}와\  Lemma
\ref{flat-lemma-compactifyable}에 의해 가능한\  $Y$ 위의 콤팩트화\ 
$j : X \to \overline{X}$를 택한다. 구조 사상을\ 
$\overline{f} : \overline{X} \to Y$로 쓰고,
$K \in D^b_{\textit{Coh}}(\mathcal{O}_X)$에 대해
$$
Rf_!K  = R\overline{f}_* Rj_! K
$$
로 둔다. 여기서\  $Rj_!$은\  Remark \ref{duality-remark-extension-by-zero}에서와 같다.

\begin{lemma}
\label{duality-lemma-lower-shriek-well-defined}
함자\  $Rf_!$은 동형을 제외하면 콤팩트화의 선택과 무관하다.
\end{lemma}

\noindent
사실 함자\  $Rf_!$은\  $f^!$의 “왼쪽 수반”으로 특징지어질 것이며, 이는 유일한 동형을
제외하고 이 함자를 결정한다.

\begin{proof}
More on Flatness, Theorem \ref{flat-theorem-nagata}와\  Lemmas
\ref{flat-lemma-compactifications-cofiltered},
\ref{flat-lemma-compactifyable}에 의해\  $Y$ 위의\  $X$의 콤팩트화들의 범주는
쌍대여과된다. 콤팩트화의 각 선택
$$
j : X \to \overline{X},\quad \overline{f} : \overline{X} \to Y
$$
에 위 구성은 함자\  $R\overline{f}_* \circ Rj_!$을 대응시킨다. $Y$ 위의 콤팩트화들\ 
$j_i : X \to \overline{X}_i$ 사이의 사상\ 
$g : \overline{X}_1 \to \overline{X}_2$가 주어졌다고 하자. 그러면 동형
$$
R\overline{f}_{2, *} \circ Rj_{2, !} =
R\overline{f}_{2, *} \circ Rg_* \circ j_{1, !} =
R\overline{f}_{1, *} \circ Rj_{1, !}
$$
을 얻는다. 첫 번째 등식에서\  Lemma \ref{duality-lemma-well-defined}를 사용했다. 이로써 우리
함자가 동형을 제외하면 콤팩트화의 선택과 무관함을 알 수 있다.
\end{proof}
\begin{proposition}
\label{duality-proposition-duality-compactly-supported}
Situation \ref{duality-situation-shriek}에서\  $f : X \to Y$를\  $\textit{FTS}_S$의 사상이라
하자. 그러면 함자\  $Rf_!$와\  $f^!$는 다음 의미에서 수반이다. 모든\ 
$K \in D^b_{\textit{Coh}}(\mathcal{O}_X)$ 및\ 
$L \in D^+_{\textit{Coh}}(\mathcal{O}_Y)$에 대해
$$
\Hom_X(K, f^!L) =
\Hom_{\text{Pro-}D^+_{\textit{Coh}}(\mathcal{O}_Y)}(Rf_!K, L)
$$
이며, 이는\  $K$와\  $L$에 이중함자적이다.
\end{proposition}

\begin{proof}
$Y$ 위의 콤팩트화\  $j : X \to \overline{X}$를 택하고 구조 사상을\ 
$\overline{f} : \overline{X} \to Y$로 쓰자. 그러면
\begin{align*}
\Hom_X(K, f^!L)
& =
\Hom_X(K, j^*\overline{f}{}^!L) \\
& =
\Hom_{\text{Pro-}D^+_{\textit{Coh}}(\mathcal{O}_{\overline{X}})}
(Rj_!K, \overline{f}{}^!L) \\
& =
\Hom_{\text{Pro-}D^+_{\textit{Coh}}(\mathcal{O}_Y)}(Rf_*Rj_!K, L) \\
& =
\Hom_{\text{Pro-}D^+_{\textit{Coh}}(\mathcal{O}_Y)}(Rf_!K, L)
\end{align*}
이다. 첫 번째 등식은\  Section \ref{duality-section-upper-shriek}의\  $f^!$ 구성에서 즉시
따른다. Lemma \ref{duality-lemma-shriek-coherent}에 의해\  $\overline{f}{}^!L$은\ 
$D^+_{\textit{Coh}}(\mathcal{O}_{\overline{X}})$에 속하므로 두 번째 등식은\  Lemma
\ref{duality-lemma-lift-map-plus}에서 따른다. $\overline{f}$가 고유이므로 구성상 함자\ 
$\overline{f}{}^!$는 밀어앞으로 보내기의 오른쪽 수반이다. 이로부터 세 번째 등식을
얻는다. 네 번째 등식은\  $Rf_! = Rf_* Rj_!$이기 때문에 성립한다.
\end{proof}
\begin{lemma}
\label{duality-lemma-compactly-supported-triangle}
Situation \ref{duality-situation-shriek}에서\  $f : X \to Y$를\  $\textit{FTS}_S$의 사상이라
하자.
$$
K \to L \to M \to K[1]
$$
을\  $D^b_{\textit{Coh}}(\mathcal{O}_X)$의 식별 삼각형이라 하자. 그러면\ 
$D^b_{\textit{Coh}}(\mathcal{O}_Y)$의 식별 삼각형들의 역계
$$
K_n \to L_n \to M_n \to K_n[1]
$$
이 존재하여\  pro-계\  $(K_n)$, $(L_n)$, $(M_n)$은 각각\  $Rf_!K$, $Rf_!L$, $Rf_!M$을
준다.
\end{lemma}

\begin{proof}
$Y$ 위의 콤팩트화\  $j : X \to \overline{X}$를 택하고 구조 사상을\ 
$\overline{f} : \overline{X} \to Y$로 쓰자. 열린 몰입\  $j$와 주어진 식별 삼각형에
대응하여\  Lemma \ref{duality-lemma-extension-by-zero-triangle}에서와 같은\ 
$D^b_{\textit{Coh}}(\mathcal{O}_{\overline{X}})$의 식별 삼각형들의 역계
$$
\overline{K}_n \to \overline{L}_n \to \overline{M}_n \to \overline{K}_n[1]
$$
을 택한다. $K_n = R\overline{f}_*\overline{K}_n$로 두고\  $L_n$과\  $M_n$도
마찬가지로 둔다. $Rf_!$의 정의에 의해 이것으로 충분하다.
\end{proof}
\begin{remark}
\label{duality-remark-compose-inverse-systems}
$\mathcal{C}$를 범주라 하자. $\mathcal{C}$의\  pro-대상 범주에서 역계들의 역계
$$
\ldots \xrightarrow{\alpha_4} (M_{3, n}) \xrightarrow{\alpha_3} (M_{2, n})
\xrightarrow{\alpha_2} (M_{1, n})
$$
이 주어졌다고 하자. 다시 말해 화살표\  $\alpha_i$들은\  pro-대상의 사상이다.
Categories, Example \ref{categories-example-pro-morphism-inverse-systems}에 의해 각\ 
$\alpha_i$를 쌍\  $(m_i, a_i)$으로 표현할 수 있다. 여기서\ 
$m_i : \mathbf{N} \to \mathbf{N}$은 증가함수이고\ 
$a_{i, n} : M_{i, m_i(n)} \to M_{i - 1, n}$은\  $\mathcal{C}$의 사상으로서 그림들
$$
\xymatrix{
\ldots \ar[r] &
M_{i, m_i(3)} \ar[d]^{a_{i, 3}} \ar[r] &
M_{i, m_i(2)} \ar[d]^{a_{i, 2}} \ar[r] &
M_{i, m_i(1)} \ar[d]^{a_{i, 1}} \\
\ldots \ar[r] &
M_{i - 1, 3} \ar[r] &
M_{i - 1, 2} \ar[r] &
M_{i - 1, 1}
}
$$
을 가환하게 한다. $m_i(n)$을\  $\max(n, m_i(n))$으로 바꾸고 그에 맞추어 사상\ 
$a_i(n)$을 조정하면(인용한 예에서와 같이) $m_i(n) \geq n$이라 가정해도 된다.
이 상황에서 역계
$$
\ldots \to
M_{4, m_4(m_3(m_2(4)))} \to
M_{3, m_3(m_2(3))} \to
M_{2, m_2(2)} \to
M_{1, 1}
$$
을 생각하자. 그 일반항은
$$
M_k = M_{k, m_k(m_{k - 1}(\ldots (m_2(k))\ldots))}
$$
이다. $\mathcal{C}$의 임의의 대상\  $N$에 대해
$$
\colim_i \colim_n \Mor_\mathcal{C}(M_{i, n}, N) =
\colim_k \Mor_\mathcal{C}(M_k, N) 
$$
이다. 세부사항은 생략한다. 다시 말해 역계\  $(M_k)$는 다음 성질을 갖는다.
$$
\colim_i \Mor_{\text{Pro-}\mathcal{C}}((M_{i, n}), N) =
\Mor_{\text{Pro-}\mathcal{C}}((M_k), N)
$$
Categories, Remark \ref{categories-remark-pro-category-copresheaves}의 논의에 의해 이
성질은\  pro-동형을 제외하고 역계\  $(M_k)$를 결정한다. 이런 방식으로\ 
$\text{Pro-}\mathcal{C}$의 어떤 역계들을 가산 지표 범주를 갖는\  pro-대상으로
바꿀 수 있다.
\end{remark}
\begin{remark}
\label{duality-remark-composition-lower-shriek}
Situation \ref{duality-situation-shriek}에서\  $f : X \to Y$와\  $g : Y \to Z$를\ 
$\textit{FTS}_S$의 합성 가능한 사상들이라 하자. 합성
$$
Rg_! \circ Rf_! :
D^b_{\textit{Coh}}(\mathcal{O}_X)
\longrightarrow
\text{Pro-}D^b_{\textit{Coh}}(\mathcal{O}_Z)
$$
을 정의하자. 실제로\  $Rf_!$의 구성에 의해\ 
$D^b_{\textit{Coh}}(\mathcal{O}_X)$의\  $K$에 대한 출력\  $Rf_!K$는\ 
$D^b_{\textit{Coh}}(\mathcal{O}_Y)$의 역계\  $(M_n)$의\  pro-동형류이다. $Rg_!$도
마찬가지로 구성되므로
$$
\ldots \to Rg_!M_3 \to Rg_!M_2 \to Rg_!M_1
$$
은\  $\text{Pro-}D^b_{\textit{Coh}}(\mathcal{O}_Y)$의 역계이다. Remark
\ref{duality-remark-compose-inverse-systems}의 논의에 의해\ 
$D^b_{\textit{Coh}}(\mathcal{O}_Z)$의 역계 중 다음을 만족하는 유일한\ 
pro-동형류가 있으며, 이를\  $Rg_! Rf_! K$로 나타낸다.
$$
\Hom_{\text{Pro-}D^b_{\textit{Coh}}(\mathcal{O}_Z)}(Rg_!Rf_!K, L) =
\colim_n \Hom_{\text{Pro-}D^b_{\textit{Coh}}(\mathcal{O}_Z)}(Rg_!M_n, L)
$$
다음 보조정리에서 즉시 따를 것이므로 이 구성이\  $K$에 대해 함자적임을 보이는 데
필요한 논의는 생략한다.
\end{remark}
\begin{lemma}
\label{duality-lemma-composition-lower-shriek}
Situation \ref{duality-situation-shriek}에서\  $f : X \to Y$와\  $g : Y \to Z$를\ 
$\textit{FTS}_S$의 합성 가능한 사상들이라 하자. Remark
\ref{duality-remark-composition-lower-shriek}의 표기 아래\ 
$Rg_! \circ Rf_! = R(g \circ f)_!$이다.
\end{lemma}

\begin{proof}
Categories, Remark \ref{categories-remark-pro-category-copresheaves}의 논의에 의해\ 
$D^b_{\textit{Coh}}(\mathcal{O}_Z)$의\  $L$로 가는\  $\Hom$을 계산할 때 같은 답을
얻음을 보이면 충분하다. Remark \ref{duality-remark-composition-lower-shriek}의 표기를
사용하여 다음과 같이 계산한다.
\begin{align*}
\Hom_Z(Rg_!Rf_!K, L)
& =
\colim_n \Hom_Z(Rg_!M_n, L) \\
& =
\colim_n \Hom_Y(M_n, g^!L) \\
& =
\Hom_Y(Rf_!K, g^!L) \\
& =
\Hom_X(K, f^!g^!L) \\
& =
\Hom_X(K, (g \circ f)^!L) \\
& =
\Hom_Z(R(g \circ f)_!K, L)
\end{align*}
첫 번째 등식은\  $Rg_!Rf_!K$의 정의이다. 두 번째 등식은\  $g$에 대한\  Proposition
\ref{duality-proposition-duality-compactly-supported}이다. 세 번째 등식은\  $Rf_!K$가\ 
$(M_n)$으로 주어진다는 사실이다. 네 번째 등식은\  $f$에 대한\  Proposition
\ref{duality-proposition-duality-compactly-supported}이다. 다섯 번째 등식은\  Lemma
\ref{duality-lemma-upper-shriek-composition}이다. 여섯 번째 등식은\  $g \circ f$에 대한\ 
Proposition \ref{duality-proposition-duality-compactly-supported}이다.
\end{proof}
\begin{remark}
\label{duality-remark-covariance-open-lower-shriek}
Situation \ref{duality-situation-shriek}에서\  $f : X \to Y$를\  $\textit{FTS}_S$의 사상이라
하고\  $U \subset X$를 열린집합이라 하자. $g = f|_U : U \to Y$로 두자. 그러면
$$
Rg_!(K|_U) \longrightarrow Rf_!K
$$
라는 표준 사상이 있으며, 이는\  $D^b_{\textit{Coh}}(\mathcal{O}_X)$의\  $K$에 대해
함자적이고 적어도 세 가지 방식으로 정의할 수 있다.
\begin{enumerate}
\item 포함 사상을\  $i : U \to X$로 쓰자. Lemma
\ref{duality-lemma-composition-lower-shriek}에 의해\  $Rg_! = Rf_! \circ Ri_!$이고, Remark
\ref{duality-remark-covariance-open-j-lower-shriek}의 특수한 경우인 사상\ 
$Ri_!(K|_U) \to K$에\  $Rf_!$을 적용할 수 있다.
\item $Y$ 위의\  $X$의 콤팩트화\  $j : X \to \overline{X}$를 택하고 구조 사상을\ 
$\overline{f} : \overline{X} \to Y$로 쓰자. $j' = j \circ i : U \to \overline{X}$로
두자. $Rf_! = R\overline{f}_* \circ Rj_!$과\ 
$Rg_! = R\overline{f}_* \circ Rj'_!$을 사용하고, Remark
\ref{duality-remark-covariance-open-j-lower-shriek}의 사상\ 
$Rj'_!(K|_U) \to Rj_!K$에\  $R\overline{f}_*$를 적용할 수 있다.
\item 다음을 사용할 수 있다.
\begin{align*}
\Hom_{\text{Pro-}D^b_{\textit{Coh}}(\mathcal{O}_Y)}(Rf_!K, L)
& =
\Hom_X(K, f^!L) \\
& \to
\Hom_U(K|_U, f^!L|_U) \\
& =
\Hom_U(K|_U, g^!L) \\
& =
\Hom_{\text{Pro-}D^b_{\textit{Coh}}(\mathcal{O}_Y)}(Rg_!(K|_U), L)
\end{align*}
이는\  $L$과\  $K$에 대해 함자적이다. 여기서는\  Proposition
\ref{duality-proposition-duality-compactly-supported}을 두 번 사용했고, 상단 느낌표
함자의 구성에서 따르는\  $g^! = i^* \circ f^!$를 사용했다. $L$에 대한 함자성과\ 
Categories, Remark \ref{categories-remark-pro-category-copresheaves}에 의해\ 
$\text{Pro-}D^b_{\textit{Coh}}(\mathcal{O}_Y)$의 표준 사상\ 
$Rg_!(K|_U) \to Rf_!K$를 얻고, 위 화살표의\  $K$에 대한 함자성에 의해 이 사상도\ 
$K$에 대해 함자적이다.
\end{enumerate}
이 세 구성은 모두 같은 화살표를 준다. 세부사항은 생략한다.
\end{remark}
\begin{remark}
\label{duality-remark-covariance-etale-lower-shriek}
Remark \ref{duality-remark-covariance-open-lower-shriek}에서 주어진 콤팩트 지지
코호몰로지의 공변성을 \'etale 사상으로 일반화하자. 즉\  Situation
\ref{duality-situation-shriek}에서\  $\textit{FTS}_S$의 가환 그림
$$
\xymatrix{
U \ar[rr]_h \ar[rd]_g & & X \ar[ld]^f \\
& Y
}
$$
이 주어지고\  $h$가 \'etale이라고 하자. 그러면 표준 사상
$$
Rg_!(h^*K) \longrightarrow Rf_!K
$$
이 있으며, 이는\  $D^b_{\textit{Coh}}(\mathcal{O}_X)$의\  $K$에 대해 함자적이다.
다음 사상들의 열로 이 변환을 정의한다.
\begin{align*}
\Hom_{\text{Pro-}D^b_{\textit{Coh}}(\mathcal{O}_Y)}(Rf_!K, L)
& =
\Hom_X(K, f^!L) \\
& \to
\Hom_U(h^*K, h^*(f^!L)) \\
& =
\Hom_U(h^*K, h^!f^!L) \\
& =
\Hom_U(h^*K, g^!L) \\
& =
\Hom_{\text{Pro-}D^b_{\textit{Coh}}(\mathcal{O}_Y)}(Rg_!(h^*K), L)
\end{align*}
이는\  $L$과\  $K$에 대해 함자적이다. 여기서\  Proposition
\ref{duality-proposition-duality-compactly-supported}을 두 번 사용했고, Lemma
\ref{duality-lemma-shriek-etale}의 등식\  $h^* = h^!$과\  Lemma
\ref{duality-lemma-upper-shriek-composition}의 등식\  $h^! \circ f^! = g^!$를 사용했다.
$L$에 대한 함자성과\  Categories, Remark
\ref{categories-remark-pro-category-copresheaves}에 의해\ 
$\text{Pro-}D^b_{\textit{Coh}}(\mathcal{O}_Y)$의 표준 사상\ 
$Rg_!(h^*K) \to Rf_!K$를 얻고, 위 화살표의\  $K$에 대한 함자성에 의해 이 사상도\ 
$K$에 대해 함자적이다.
\end{remark}
\begin{remark}
\label{duality-remark-covariance-lower-shriek}
Remarks \ref{duality-remark-covariance-open-lower-shriek}와\ 
\ref{duality-remark-covariance-etale-lower-shriek}에서 콤팩트 지지 코호몰로지의 구성이 열린
몰입과 \'etale 사상에 대해 공변임을 보았다. 사실 올바른 일반성은 다음과 같다.
$\textit{FTS}_S$의 가환 그림
$$
\xymatrix{
U \ar[rr]_h \ar[rd]_g & & X \ar[ld]^f \\
& Y
}
$$
이 주어지고\  $h$가 평탄 준유한이면 표준 변환
$$
Rg_! \circ h^* \longrightarrow Rf_!
$$
이 존재한다. Remark \ref{duality-remark-covariance-etale-lower-shriek}에서와 같이 이 사상은\ 
$D^+_{\textit{Coh}}(\mathcal{O}_X)$ 위의 함자 변환\  $h^* \to h^!$을 사용하여
구성할 수 있다. 다음을 상기하자.
$h^!K = h^*K \otimes \omega_{U/X}$이며, 여기서\ 
$\omega_{U/X} = h^!\mathcal{O}_X$는 평탄 준유한 사상\  $h$의 상대 쌍대화 층이다
(Lemmas \ref{duality-lemma-perfect-comparison-shriek}와\ 
\ref{duality-lemma-flat-quasi-finite-shriek}를 보라). $\omega_{U/X}$는\  Discriminants,
Remark \ref{discriminant-remark-relative-dualizing-for-quasi-finite}에서 구성될 상대
쌍대화 가군과 같으며, 이는\  Discriminants, Lemma
\ref{discriminant-lemma-compare-dualizing}에 따른다. 따라서\  Discriminants, Remark
\ref{discriminant-remark-relative-dualizing-for-flat-quasi-finite}에서 구성될 대각합
원소\  $\tau_{U/X} : \mathcal{O}_U \to \omega_{U/X}$를 사용하여 우리의 변환을 정의할
수 있다. 이것이 필요해지면 여기에 정확히 정식화하고 증명할 것이다.
\end{remark}






\section{콤팩트 지지 코호몰로지의 쌍대성}
\label{duality-section-duality-compactly-supported}

\noindent
$k$를 체라 하자. $U$를\  $k$ 위의 유한형 분리 스킴이라 하자.
$K$를\  $D^b_{\textit{Coh}}(\mathcal{O}_U)$의 대상이라 하자. $K$의 콤팩트 지지
코호몰로지\  $H^i_c(U, K)$를 다음과 같이 정의하자.
$k$ 위의 고유 스킴으로 가는 열린 몰입\  $j : U \to X$와, $U$로 제한하면 값이\ 
$K$인 상수 계가 되는\  $j : U \to X$의\  Deligne 계\  $(K_n)$을 택한다. 이제
$$
H^i_c(U, K) = \lim H^i(X, K_n)
$$
으로 둔다. 이 공간에는 역극한 위상을 주어 위상\  $k$-벡터 공간으로 본다
(More on Algebra, Section \ref{more-algebra-section-topological-ring}을 보라).
여기에는 몇 가지 언급할 점이 있다.

\medskip\noindent
첫째, 이 정의는\  $X$와\  $(K_n)$의 선택에 무관하다. 실제로\ 
$p : U \to \Spec(k)$가 구조 사상이면\  $Rp_!K = (R\Gamma(X, K_n))$가\ 
$D(k)$에서\  pro-동형까지 잘 정의됨을 이미 알고 있으므로, $H^i_c(U, K)$를
정의하는 극한도 그러하다.

\medskip\noindent
둘째, 다음 식을 쓰는 편이 더 자연스러워 보일 수 있다.
$$
H^i(R\lim R\Gamma(X, K_n)) = R\Gamma(X, R\lim K_n)
$$
그러나 이것도 같은 답을 준다. $k$-벡터 공간\  $H^j(X, K_n)$들이 유한 차원이므로
이 역계들은\  Mittag-Leffler 조건을 만족하며, 따라서\  Cohomology, Lemma
\ref{cohomology-lemma-RGamma-commutes-with-Rlim}의\  $R^1\lim$ 항들은 사라진다.

\medskip\noindent
$U' \subset U$가 열린 부분스킴이면 표준 사상
$$
H^i_c(U', K|_{U'}) \longrightarrow H^i_c(U, K)
$$
이 있고, 이는\  $D^b_{\textit{Coh}}(\mathcal{O}_U)$의\  $K$에 대해 함자적이다.
예를 들어\  Remark \ref{duality-remark-covariance-open-lower-shriek}을 보라.
사실\  Remark \ref{duality-remark-covariance-etale-lower-shriek}을 사용하면, 더 일반적으로\ 
$k$ 위의 유한형 분리 스킴 사이의 \'etale 사상\  $U' \to U$에 대해서도 이러한
사상이 존재함을 알 수 있다.

\medskip\noindent
$V$가\  $k$-벡터 공간이면\  $\Hom_k(V, k)$에 다음과 같이 위상을 준다.
$V = \bigcup V_i$를 유한 차원\  $k$-벡터 부분공간들의 여과 합집합으로 쓰고\ 
$\Hom_k(V, k) = \lim \Hom_k(V_i, k)$에 역극한 위상을 사용한다.
$\dim_k V < \infty$이면\  $\Hom_k(V, k)$의 위상은 이산 위상이다. 더 일반적으로,
$V = \colim_n V_n$가 유한 차원 벡터 공간들의 유향 여극한으로 주어지면\ 
$\Hom_k(V, k) = \lim \Hom_k(V_n, k)$가 위상 벡터 공간으로서 성립한다.
\begin{lemma}
\label{duality-lemma-duality-compact-support}
$k$가 체일 때\  $p : U \to \Spec(k)$가 유한형 분리 사상이라 하자.
$\omega_{U/k}^\bullet = p^!\mathcal{O}_{\Spec(k)}$라 하자. 위상\  $k$-벡터 공간의
표준 동형
$$
\Hom_k(H^i(U, K), k) =
H^{-i}_c(U, R\SheafHom_{\mathcal{O}_U}(K, \omega_{U/k}^\bullet))
$$
이 있으며, 이는\  $D^b_{\textit{Coh}}(\mathcal{O}_U)$의\  $K$에 대해 함자적이다.
\end{lemma}
\begin{proof}
$k$ 위의 콤팩트화\  $j : U \to X$를 택한다.
$V(\mathcal{I}) = X \setminus U$인 준연접 아이디얼 층\ 
$\mathcal{I} \subset \mathcal{O}_X$를 택하자.
Derived Categories of Schemes, Proposition \ref{perfect-proposition-DCoh}에 의해\ 
$K = M|_U$인\  $M \in D^b_{\textit{Coh}}(\mathcal{O}_X)$를 택할 수 있다.
Lemma \ref{duality-lemma-lift-map}에 의해
$$
H^i(U, K) =
\Ext^i_U(\mathcal{O}_U, M|_U) =
\colim \Ext^i_X(\mathcal{I}^n, M) =
\colim H^i(X, R\SheafHom_{\mathcal{O}_X}(\mathcal{I}^n, M))
$$
이다. $\mathcal{I}^n$은 연접\  $\mathcal{O}_X$-가군이므로\ 
$\mathcal{I}^n$은\  $D^-_{\textit{Coh}}(\mathcal{O}_X)$에 속하고, 따라서\ 
Derived Categories of Schemes, Lemma
\ref{perfect-lemma-coherent-internal-hom}에 의해\ 
$R\SheafHom_{\mathcal{O}_X}(\mathcal{I}^n, M)$은\ 
$D^+_{\textit{Coh}}(\mathcal{O}_X)$에 속한다.

\medskip\noindent
$q : X \to \Spec(k)$가 구조 사상일 때\ 
$\omega_{X/k}^\bullet = q^!\mathcal{O}_{\Spec(k)}$라 하자.
Section \ref{duality-section-duality-proper-over-field}을 보라. 그러면
\begin{align*}
\Hom_k(
&
H^i(X, R\SheafHom_{\mathcal{O}_X}(\mathcal{I}^n, M)), k) \\
& =
\Ext^{-i}_X(R\SheafHom_{\mathcal{O}_X}(\mathcal{I}^n, M),
\omega_{X/k}^\bullet) \\
& =
H^{-i}(X, R\SheafHom_{\mathcal{O}_X}(R\SheafHom_{\mathcal{O}_X}(
\mathcal{I}^n, M), \omega_{X/k}^\bullet))
\end{align*}
임을\  Lemma \ref{duality-lemma-duality-proper-over-field}에서 얻는다.
Lemma \ref{duality-lemma-internal-hom-evaluate-isom}의 (1)에 의해 표준 사상
$$
R\SheafHom_{\mathcal{O}_X}(M, \omega_{X/k}^\bullet)
\otimes_{\mathcal{O}_X}^\mathbf{L} \mathcal{I}^n
\longrightarrow
R\SheafHom_{\mathcal{O}_X}(R\SheafHom_{\mathcal{O}_X}(
\mathcal{I}^n, M), \omega_{X/k}^\bullet)
$$
은 동형이다. $p^!$는\  $q^!$와\  $U$로의 제한을 합성하여 구성되므로\ 
$\omega^\bullet_{U/k} = \omega^\bullet_{X/k}|_U$임에 유의하자.
따라서\  $R\SheafHom_{\mathcal{O}_X}(M, \omega_{X/k}^\bullet)$은\ 
$D^b_{\textit{Coh}}(\mathcal{O}_X)$의 대상이며, $U$로 제한하면\ 
$R\SheafHom_{\mathcal{O}_U}(K, \omega_{U/k}^\bullet)$가 된다.
그러므로\  Lemma \ref{duality-lemma-tensoring-Deligne-system}에 의해
$$
\lim H^{-i}(X, R\SheafHom_{\mathcal{O}_X}(M, \omega_{X/k}^\bullet)
\otimes_{\mathcal{O}_X}^\mathbf{L} \mathcal{I}^n)
$$
이 등식 오른쪽 항을 나타내는 한 모형임을 얻는다.
이와 같이 얻은 모든 동형을 합성하면 보조정리의 동형을 얻는다.
\end{proof}
\begin{lemma}
\label{duality-lemma-duality-compact-support-restrict-open}
Lemma \ref{duality-lemma-duality-compact-support}의 표기를 사용하고\ 
$U' \subset U$가 열린 부분스킴이라 하자. 그러면 그림
$$
\xymatrix{
\Hom_k(H^i(U, K), k) \ar[rr] & &
H^{-i}_c(U, R\SheafHom_{\mathcal{O}_U}(K, \omega_{U/k}^\bullet)) \\
\Hom_k(H^i(U', K|_{U'}), k) \ar[rr] \ar[u] & &
H^{-i}_c(U', R\SheafHom_{\mathcal{O}_{U'}}(K, \omega_{U'/k}^\bullet)) \ar[u]
}
$$
은 가환한다. 여기서 수평 화살표들은\  Lemma
\ref{duality-lemma-duality-compact-support}의 동형들이고, 왼쪽 수직 화살표는 제한 사상\ 
$H^i(U, K) \to H^i(U', K|_{U'})$의 쌍대 사상이며, 오른쪽 수직 화살표는\ 
Remark \ref{duality-remark-covariance-open-lower-shriek}의 사상이다
(보조정리 앞의 논의를 보라).
\end{lemma}
\begin{proof}
독자에게 이 증명은 건너뛰기를 강력히 권한다. Lemma
\ref{duality-lemma-duality-compact-support}의 증명에서처럼\  $X$와\  $M$을 택한다.
$R\SheafHom$과\  $\otimes^\mathbf{L}$에서 아래첨자\  $\mathcal{O}_X$를 생략하겠다.
다음과 같이 쓴다.
$$
H^i(U, K) = \colim H^i(X, R\SheafHom(\mathcal{I}^n, M))
$$
그리고
$$
H^i(U', K|_{U'}) = \colim H^i(X, R\SheafHom((\mathcal{I}')^n, M))
$$
로 쓴다. 이는\  Lemma \ref{duality-lemma-duality-compact-support}의 증명과 같되,
Remark \ref{duality-remark-covariance-open-j-lower-shriek}의 논의처럼\ 
$\mathcal{I}' \subset \mathcal{I}$를 택하여 사상\ 
$H^i(U, K) \to H^i(U', K|_{U'})$가 사상들\ 
$(\mathcal{I}')^n \to \mathcal{I}^n$으로부터 유도되게 한 것이다.
마찬가지로
$$
H^i_c(U, R\SheafHom(K, \omega_{U/k}^\bullet)) =
\lim H^i(X,  R\SheafHom(M, \omega_{X/k}^\bullet)
\otimes^\mathbf{L} \mathcal{I}^n)
$$
및
$$
H^i_c(U', R\SheafHom(K|_{U'}, \omega_{U'/k}^\bullet)) =
\lim H^i(X,  R\SheafHom(M, \omega_{X/k}^\bullet)
\otimes^\mathbf{L} (\mathcal{I}')^n)
$$
로 쓰면, 화살표\ 
$H^i_c(U', R\SheafHom(K|_{U'}, \omega_{U'/k}^\bullet))
\to H^i_c(U, R\SheafHom(K, \omega_{U/k}^\bullet))$도 마찬가지로 사상들\ 
$(\mathcal{I}')^n \to \mathcal{I}^n$에서 유도된다. 그림들
$$
\xymatrix{
R\SheafHom(M, \omega_{X/k}^\bullet)
\otimes^\mathbf{L} \mathcal{I}^n
\ar[rr] & &
R\SheafHom(R\SheafHom(\mathcal{I}^n, M), \omega_{X/k}^\bullet) \\
R\SheafHom(M, \omega_{X/k}^\bullet) \otimes^\mathbf{L} (\mathcal{I}')^n
\ar[rr] \ar[u] & &
R\SheafHom(R\SheafHom((\mathcal{I}')^n, M), \omega_{X/k}^\bullet) \ar[u]
}
$$
은 가환한다. Cohomology, Lemma
\ref{cohomology-lemma-internal-hom-evaluate}에서 수평 화살표의 구성이 세 변수 모두에
대해 함자적이기 때문이다. 따라서 마지막으로 그림들
$$
\xymatrix{
\Hom_k(H^i(X, R\SheafHom(\mathcal{I}^n, M)), k) \ar[r] &
H^{-i}(X, R\SheafHom(R\SheafHom(
\mathcal{I}^n, M), \omega_{X/k}^\bullet)) \\
\Hom_k(H^i(X, R\SheafHom((\mathcal{I}')^n, M)), k) \ar[r] \ar[u] &
H^{-i}(X, R\SheafHom(R\SheafHom(
(\mathcal{I}')^n, M), \omega_{X/k}^\bullet)) \ar[u]
}
$$
이 가환한다는 주장만 남는다. 이는 쌍대성 동형
$$
\Hom_k(H^i(X, L), k) = \Ext^{-i}_X(L, \omega_{X/k}^\bullet) =
H^{-i}(X, R\SheafHom(L, \omega_{X/k}^\bullet))
$$
이\  $D_\QCoh(\mathcal{O}_X)$의\  $L$에 대해 함자적이므로 참이다.
\end{proof}
\begin{lemma}
\label{duality-lemma-h0-compactly-supported}
$X$를 체\  $k$ 위의 고유 스킴이라 하자.
$i < 0$에 대해\  $H^i(K) = 0$인\ 
$K \in D^b_{\textit{Coh}}(\mathcal{O}_X)$를 택하고\ 
$\mathcal{F} = H^0(K)$라 하자. $Z \subset X$가 닫혀 있고 그 여집합이\ 
$U = X \setminus U$라 하자. 그러면
$$
H^0_c(U, K|_U) \subset H^0(X, \mathcal{F})
$$
은\  $Z$의 어떤 열린 근방에서 사라지는\  $\mathcal{F}$의 대역 단면들로 주어진다.
\end{lemma}

\begin{proof}
Remark \ref{duality-remark-covariance-open-lower-shriek}의 사상\ 
$H^0_c(U, K|_U) \to H^0_X(X, K) = H^0(X, K) = H^0(X, \mathcal{F})$을
생각하자. 이를 연구하기 위해\  $i < 0$이면\  $\mathcal{F}^i = 0$인 유계 복합체\ 
$\mathcal{F}^\bullet$로\  $K$를 나타낸다. 그러면 정의에 의해
$$
H^0_c(U, K|_U) = \lim H^0(X, \mathcal{I}^n\mathcal{F}^\bullet)
= \lim \Ker(
H^0(X, \mathcal{I}^n\mathcal{F}^0) \to
H^0(X, \mathcal{I}^n\mathcal{F}^1))
$$
이다. Artin--Rees 정리(Cohomology of Schemes, Lemma
\ref{coherent-lemma-Artin-Rees})에 의해 이는\ 
$\lim H^0(X, \mathcal{I}^n\mathcal{F})$와 같다.
따라서 화살표\  $H^0_c(U, K|_U) \to H^0(X, \mathcal{F})$는 단사이고, 그 상은
임의의\  $n$에 대해 부분층\  $\mathcal{I}^n\mathcal{F}$에 속하는\ 
$\mathcal{F}$의 대역 단면들로 이루어진다. 이들이\  $Z$의 어떤 근방에서 사라지는
단면들이라는 특성화는\  $\mathcal{F}$의 줄기들에\  Krull의 교집합 정리
(Algebra, Lemma \ref{algebra-lemma-intersect-powers-ideal-module-zero})를 적용하여
얻는다. 함수의 경우에는\  Algebra, Remark
\ref{algebra-remark-intersection-powers-ideal}의 논의를 보라.
\end{proof}



\section{Lichtenbaum의 정리}
\label{duality-section-lichtenbaum}

\noindent
아래 정리는\  Lichtenbaum이 추측했고\  Grothendieck이 증명했다
(\cite{Hartshorne-local-cohomology}를 보라). 이 정리에 대한\  Kleiman의 매우 훌륭한
증명이\  \cite{Kleiman-Lichtenbaum}에 있다. 지지를 가진 코호몰로지의 경우에 대한
일반화는\  \cite{Lyubeznik-Lichtenbaum}에서 찾을 수 있다.
논증에서 가장 흥미로운 부분은 다음 보조정리의 증명에 들어 있다.
\begin{lemma}
\label{duality-lemma-lichtenbaum}
$U$를 대수다양체라 하고\  $\mathcal{F}$를 연접\  $\mathcal{O}_U$-가군이라 하자.
$H^d(U, \mathcal{F})$가 영이 아니면\  $\dim(U) \geq d$이고, 등호가 성립하면\ 
$U$는 고유하다.
\end{lemma}

\begin{proof}
Grothendieck의 소멸 결과인\  Cohomology, Proposition
\ref{cohomology-proposition-vanishing-Noetherian}에 의해\ 
$\dim(U) \geq d$를 얻는다.
$\dim(U) = d$라 하자. $U$가\  $X$에서 조밀한 콤팩트화\  $U \to X$를 택한다.
(이는\  More on Flatness, Theorem \ref{flat-theorem-nagata}와\ 
Lemma \ref{flat-lemma-compactifyable}에 의해 가능하다.)
$X$를 그 축약으로 바꾸면\  $X$는 차원\  $d$인 고유 대수다양체이고,
$U$가 고유일 필요충분조건이\  $U = X$임을 알 수 있다.
$Z = X \setminus U$라 하자. $Z$가 공집합이 아니면\ 
$H^d(U, \mathcal{F})$가 영임을 보이겠다.

\medskip\noindent
Properties, Lemma \ref{properties-lemma-lift-finite-presentation}에 따라\ 
$U$로 제한하면\  $\mathcal{F}$가 되는 연접\  $\mathcal{O}_X$-가군\ 
$\mathcal{G}$를 택한다. Section \ref{duality-section-duality-proper-over-field}에서처럼\ 
$\omega_X^\bullet$가\  $X$의 쌍대화 복합체를 나타내게 하고\ 
$\omega_U^\bullet = \omega_X^\bullet|_U$라 하자.
Lemma \ref{duality-lemma-duality-compact-support}에 의해\ 
$H^d(U, \mathcal{F})$는
$$
H^{-d}_c(U, R\SheafHom_{\mathcal{O}_U}(\mathcal{F}, \omega_U^\bullet))
$$
의 쌍대 공간이다. Lemma \ref{duality-lemma-duality-proper-over-field}에 의해\ 
$\omega_X^\bullet$의 코호몰로지 층들은 차수\  $< -d$에서 사라지고,
$H^{-d}(\omega_X^\bullet) = \omega_X$는\  $(S_2)$이고 지지가\  $X$인 연접\ 
$\mathcal{O}_X$-가군이다. 특히\  Divisors, Lemma
\ref{divisors-lemma-torsion-free-finite-noetherian-domain}에 의해\ 
$\omega_X$는 꼬임이 없다. 따라서 코호몰로지 층
$$
H^{-d}(R\SheafHom_{\mathcal{O}_X}(\mathcal{G}, \omega_X^\bullet)) =
\SheafHom(\mathcal{G}, \omega_X)
$$
도\  Divisors, Lemma \ref{divisors-lemma-hom-into-torsion-free}에 의해
꼬임이 없다. 그러므로 이 층에는\  $X$의 어떤 공집합이 아닌 열린집합에서 사라지는
영이 아닌 단면이 없다(그러한 단면은 꼬임 단면일 것이다).
따라서\  Lemma \ref{duality-lemma-h0-compactly-supported}에서\ 
$H^{-d}_c(U, R\SheafHom_{\mathcal{O}_U}(\mathcal{F}, \omega_U^\bullet))$가
영임이 따르고, 이에 따라 원하는 대로\  $H^d(U, \mathcal{F})$도 영이다.
\end{proof}
\begin{theorem}
\label{duality-theorem-lichtenbaum}
$X$를 체\  $k$ 위의 공집합이 아닌 유한형 분리 스킴이라 하자.
$d = \dim(X)$라 하자. 다음은 동치이다.
\begin{enumerate}
\item $X$ 위의 모든 연접\  $\mathcal{O}_X$-가군\  $\mathcal{F}$에 대해\ 
$H^d(X, \mathcal{F}) = 0$이다.
\item $X$ 위의 모든 준연접\  $\mathcal{O}_X$-가군\  $\mathcal{F}$에 대해\ 
$H^d(X, \mathcal{F}) = 0$이다.
\item 차원이\  $d$인 기약 성분\  $X' \subset X$ 중\  $k$ 위에서 고유한 것은 없다.
\end{enumerate}
\end{theorem}

\begin{proof}
차원이\  $d$이고 고유한 기약 성분\  $X' \subset X$가 존재한다고 하자
(이를 정역적 닫힌 부분스킴으로 본다).
$\omega_{X'}$를\  $k$ 위의\  $X'$의 쌍대화 가군이라 하자.
Lemma \ref{duality-lemma-duality-proper-over-field}를 보라. 그러면 이 보조정리에 의해\ 
$H^d(X', \omega_{X'})$는\  $H^0(X', \mathcal{O}_{X'})$의 쌍대 공간이므로 영이
아니다. 따라서\ 
$H^d(X, \omega_{X'}) = H^d(X', \omega_{X'})$가 영이 아니며, (1)이 성립하지
않음을 얻는다. 이로써 (1)은 (3)을 함의한다.

\medskip\noindent
이제 (3)이 (1)을 함의함을 증명하자.
$\mathcal{F}$를\  $H^d(X, \mathcal{F})$가 영이 아닌 연접\ 
$\mathcal{O}_X$-가군이라 하자. Cohomology of Schemes, Lemma
\ref{coherent-lemma-coherent-filter}에서와 같은 여과
$$
0 = \mathcal{F}_0 \subset \mathcal{F}_1 \subset
\ldots \subset \mathcal{F}_m = \mathcal{F}
$$
를 택한다. 완전열
$$
H^d(X, \mathcal{F}_i) \to H^d(X, \mathcal{F}_{i + 1}) \to
H^d(X, \mathcal{F}_{i + 1}/\mathcal{F}_i)
$$
들을 얻는다. 따라서 어떤\  $i \in \{1, \ldots, m\}$에 대해\ 
$H^d(X, \mathcal{F}_{i + 1}/\mathcal{F}_i)$가 영이 아님을 알 수 있다.
여과의 선택에 의해 이는 정역적 닫힌 부분스킴\  $Z \subset X$와 영이 아닌
연접 아이디얼 층\  $\mathcal{I} \subset \mathcal{O}_Z$가 존재하여\ 
$H^d(Z, \mathcal{I})$가 영이 아님을 뜻한다.
Lemma \ref{duality-lemma-lichtenbaum}에 의해\  $\dim(Z) = d$이고\  $Z$가\  $k$ 위에서
고유임을 얻는데, 이는 (3)에 모순이다. 따라서 (3)은 (1)을 함의한다.

\medskip\noindent
끝으로 임의의\  Noether 스킴\  $X$에 대해 (1)과 (2)가 동치임을 보이자.
(2)가 (1)을 함의함은 자명하다. 반대로 (1)을 가정하고\ 
$\mathcal{F}$를 준연접\  $\mathcal{O}_X$-가군이라 하자.
Properties, Lemma \ref{properties-lemma-quasi-coherent-colimit-finite-type}에 의해\ 
$\mathcal{F} = \colim \mathcal{F}_i$를 연접 부분가군들의 여과 여극한으로 쓸 수
있다. 그러면\  Cohomology, Lemma
\ref{cohomology-lemma-quasi-separated-cohomology-colimit}에 의해\ 
$H^d(X, \mathcal{F}) = \colim H^d(X, \mathcal{F}_i) = 0$이다.
따라서 (2)가 성립한다.
\end{proof}
